% Options for packages loaded elsewhere
\PassOptionsToPackage{unicode}{hyperref}
\PassOptionsToPackage{hyphens}{url}
\PassOptionsToPackage{dvipsnames,svgnames,x11names}{xcolor}
\PassOptionsToPackage{space}{xeCJK}
%
\documentclass[
  11pt,
  oneside]{book}
\usepackage{amsmath,amssymb}
\usepackage{iftex}
\ifPDFTeX
  \usepackage[T1]{fontenc}
  \usepackage[utf8]{inputenc}
  \usepackage{textcomp} % provide euro and other symbols
\else % if luatex or xetex
  \usepackage{unicode-math} % this also loads fontspec
  \defaultfontfeatures{Scale=MatchLowercase}
  \defaultfontfeatures[\rmfamily]{Ligatures=TeX,Scale=1}
\fi
\usepackage{lmodern}
\ifPDFTeX\else
  % xetex/luatex font selection
  \setmainfont[BoldFont=DejaVuSerif-Bold.ttf,ItalicFont=DejaVuSerif-Italic.ttf,BoldItalicFont=DejaVuSerif-BoldItalic.ttf]{DejaVuSerif.ttf}
  \setmonofont[BoldFont=DejaVuSansMono-Bold.ttf,ItalicFont=DejaVuSansMono-Oblique.ttf,BoldItalicFont=DejaVuSansMono-BoldOblique.ttf]{DejaVuSansMono.ttf}
  \ifXeTeX
    \usepackage{xeCJK}
    \setCJKmainfont[BoldFont=NotoSerifCJKsc-Bold.otf,AutoFakeSlant=0.15]{NotoSerifCJKsc-Regular.otf}
          \fi
  \ifLuaTeX
    \usepackage[]{luatexja-fontspec}
    \setmainjfont[BoldFont=NotoSerifCJKsc-Bold.otf,AutoFakeSlant=0.15]{NotoSerifCJKsc-Regular.otf}
  \fi
\fi
% Use upquote if available, for straight quotes in verbatim environments
\IfFileExists{upquote.sty}{\usepackage{upquote}}{}
\IfFileExists{microtype.sty}{% use microtype if available
  \usepackage[]{microtype}
  \UseMicrotypeSet[protrusion]{basicmath} % disable protrusion for tt fonts
}{}
\makeatletter
\@ifundefined{KOMAClassName}{% if non-KOMA class
  \IfFileExists{parskip.sty}{%
    \usepackage{parskip}
  }{% else
    \setlength{\parindent}{0pt}
    \setlength{\parskip}{6pt plus 2pt minus 1pt}}
}{% if KOMA class
  \KOMAoptions{parskip=half}}
\makeatother
\usepackage{xcolor}
\usepackage[a4paper,left=24mm,right=24mm,top=25mm,bottom=25mm]{geometry}
\usepackage{longtable,booktabs,array}
\usepackage{calc} % for calculating minipage widths
% Correct order of tables after \paragraph or \subparagraph
\usepackage{etoolbox}
\makeatletter
\patchcmd\longtable{\par}{\if@noskipsec\mbox{}\fi\par}{}{}
\makeatother
% Allow footnotes in longtable head/foot
\IfFileExists{footnotehyper.sty}{\usepackage{footnotehyper}}{\usepackage{footnote}}
\makesavenoteenv{longtable}
\setlength{\emergencystretch}{3em} % prevent overfull lines
\providecommand{\tightlist}{%
  \setlength{\itemsep}{0pt}\setlength{\parskip}{0pt}}
\setcounter{secnumdepth}{-\maxdimen} % remove section numbering
% Shared print conventions; the two books keep separate prose and pagination.
\usepackage{fvextra}
\usepackage{xurl}
\usepackage{etoolbox}
\usepackage{fancyhdr}
\setCJKmonofont{NotoSansMonoCJKsc-Regular.otf}[AutoFakeBold=1.5]
\setlength{\emergencystretch}{3em}
\setlength{\parindent}{1.8em}
\setlength{\parskip}{0.32em}
\linespread{1.12}
\renewcommand{\contentsname}{目录}
\renewcommand{\figurename}{图}
\renewcommand{\tablename}{表}
% Reserve enough room for three-digit page numbers in the mathematics contents.
\makeatletter
\renewcommand{\@pnumwidth}{2.4em}
\renewcommand{\@tocrmarg}{3.4em}
\makeatother
\DefineVerbatimEnvironment{verbatim}{Verbatim}{breaklines,breakanywhere,fontsize=\small}
\AtBeginEnvironment{longtable}{\small\setlength{\tabcolsep}{3pt}\renewcommand{\arraystretch}{1.15}}
\pagestyle{fancy}
\fancyhf{}
\fancyfoot[C]{\thepage}
\renewcommand{\headrulewidth}{0pt}
\fancypagestyle{plain}{\fancyhf{}\fancyfoot[C]{\thepage}\renewcommand{\headrulewidth}{0pt}}
\Urlmuskip=0mu plus 1mu
\widowpenalty=10000
\clubpenalty=10000
\ifLuaTeX
  \usepackage{selnolig}  % disable illegal ligatures
\fi
\IfFileExists{bookmark.sty}{\usepackage{bookmark}}{\usepackage{hyperref}}
\IfFileExists{xurl.sty}{\usepackage{xurl}}{} % add URL line breaks if available
\urlstyle{same}
\hypersetup{
  pdftitle={UEOT Core 完整数学理论},
  colorlinks=true,
  linkcolor={MidnightBlue},
  filecolor={Maroon},
  citecolor={Blue},
  urlcolor={MidnightBlue},
  pdfcreator={LaTeX via pandoc}}

\title{UEOT Core 完整数学理论}
\usepackage{etoolbox}
\makeatletter
\providecommand{\subtitle}[1]{% add subtitle to \maketitle
  \apptocmd{\@title}{\par {\large #1 \par}}{}{}
}
\makeatother
\subtitle{从原始对象---演化思想到预测、实现、持续、组合、目的与学习的统一架构}
\author{}
\date{2026-10-30}

\begin{document}
\frontmatter
\maketitle

{
\hypersetup{linkcolor=MidnightBlue}
\setcounter{tocdepth}{1}
\tableofcontents
}
\mainmatter
\hypertarget{ux5168ux6587ux5bfcux822a}{%
\section{全文导航}\label{ux5168ux6587ux5bfcux822a}}

\begin{itemize}
\item
  \protect\hyperlink{ch-0}{0. 文档定位与原始思想的数学化}
\item
  \protect\hyperlink{ch-1}{1. 基础约定与可复用数学工具}
\item
  \protect\hyperlink{ch-2}{2. 世界过程、实验协议与时间}
\item
  \protect\hyperlink{ch-3}{3. 因子、表示协变与对象候选}
\item
  \protect\hyperlink{ch-4}{4. 规范预测身份：存在、最小性、细化与更新}
\item
  \protect\hyperlink{ch-5}{5. 预测信息、记忆复杂度与可压缩性}
\item
  \protect\hyperlink{ch-6}{6. 结构化开放、环境接口与因果边界}
\item
  \protect\hyperlink{ch-7}{7. 动态闭合、一般过程与多尺度一致性}
\item
  \protect\hyperlink{ch-8}{8. 持续性：不变结构、可行性与相对身份}
\item
  \protect\hyperlink{ch-9}{9. 恢复、扰动与规范恢复势}
\item
  \protect\hyperlink{ch-10}{10. 准平稳、谱与持续时间窗}
\item
  \protect\hyperlink{ch-11}{11. 历史身份、运输与发展定律}
\item
  \protect\hyperlink{ch-12}{12. 对象证书与因果完整性 Ω}
\item
  \protect\hyperlink{ch-13}{13. 预测身份的实现、冗余与毁损对偶}
\item
  \protect\hyperlink{ch-14}{14. 分辨率、粗粒化与微观实现的不确定区间}
\item
  \protect\hyperlink{ch-15}{15. 组合、重叠与父对象的生成}
\item
  \protect\hyperlink{ch-16}{16. 有限数据认证：响应、载体与协议}
\item
  \protect\hyperlink{ch-17}{17. 逆对象、逆动力学与逆目的的可识别性}
\item
  \protect\hyperlink{ch-18}{18. 可行未来、价值与目的桥接}
\item
  \protect\hyperlink{ch-19}{19. 最优控制与 GOD}
\item
  \protect\hyperlink{ch-20}{20. 对象商如何保持决策与价值}
\item
  \protect\hyperlink{ch-21}{21. GOA：长期结构的存在与稳定性}
\item
  \protect\hyperlink{ch-22}{22. 路径信息代价、可达作用与变分实现}
\item
  \protect\hyperlink{ch-23}{23. Π/Φ、约束几何与双驱动的严格版本}
\item
  \protect\hyperlink{ch-24}{24. 多层级价值、可积性与协调机制}
\item
  \protect\hyperlink{ch-25}{25. 繁殖、遗传与正算子演化}
\item
  \protect\hyperlink{ch-26}{26. 个体性转变与内生目的桥}
\item
  \protect\hyperlink{ch-27}{27. 反身结构、信念学习与能动性}
\item
  \protect\hyperlink{ch-28}{28. 完整接口、算法与组合契约}
\item
  \protect\hyperlink{ch-29}{29. 应用方法：从模型到可检验证书}
\item
  \protect\hyperlink{ch-30}{30. 完整算例与独立可执行检验}
\item
  \protect\hyperlink{ch-31}{31. 统一装配：核心能共同推出什么}
\item
  \protect\hyperlink{ch-32}{32. 有类型的统一：八类层级与四个生成机制}
\item
  \protect\hyperlink{ch-33}{33. 关系缺陷、定量校准与规范协变}
\item
  \protect\hyperlink{ch-34}{34. 操作对象、目的与条件生命周期的完整接口}
\item
  \protect\hyperlink{ch-35}{35. 同一有限源上的预测、载体、生存、修复与最优控制}
\item
  \protect\hyperlink{ch-36}{36. 实际有限路径、后验分解与部署覆盖}
\item
  \protect\hyperlink{ch-37}{37. 部分观察下的有限支持生存：双向定理与可执行检验}
\item
  \protect\hyperlink{ch-38}{38. 从注册采样到数据选择后的真实联合风险}
\item
  \protect\hyperlink{ch-39}{39. 执行、故障、复制与目的的同源构造}
\item
  \protect\hyperlink{appendix-a}{附录 A}
\item
  \protect\hyperlink{appendix-b}{附录 B}
\item
  \protect\hyperlink{appendix-c}{附录 C}
\item
  \protect\hyperlink{appendix-d}{附录 D}
\item
  \protect\hyperlink{appendix-e}{附录 E}
\item
  \protect\hyperlink{appendix-f}{附录 F. 维护版源码对应与证据层次}
\end{itemize}

\hypertarget{ux6587ux6863ux5b9aux4f4dux4e0eux539fux59cbux601dux60f3ux7684ux6570ux5b66ux5316}{%
\chapter{0. 文档定位与原始思想的数学化}\label{ux6587ux6863ux5b9aux4f4dux4e0eux539fux59cbux601dux60f3ux7684ux6570ux5b66ux5316}}

本文件是 UEOT（Universal Evolutionary Object Theory）的完整通用数学主文。它从原始论文提出的问题出发：对象为何能在材料和状态变化中持续存在；多个过程如何组成一个个体；对象通过什么接口与环境交互；功能、目标、学习与演化怎样在同一架构中获得严格表达。

本文将前期 Mathematical Foundations I--IV、统一数学 v1.1 至 v2.1、通用实现接口及后续审计中仍有效的结果重构为一套连续正文。定义、证明、条件、反例、算法接口与应用流程均在本文件内陈述，不要求先阅读旧版或依赖``上一版定理''才能理解主链。本版不收录量子专属理论、粒子选择或场论章节。后续分支产生的数学，仅在其证明不需要这些专属结构时回归本核心。

这里的``完整''指\textbf{本版所声明数学架构与应用接口的完整陈述}。它不表示所有开放问题都已解决，也不表示从几个哲学句子已推出自然界的全部规律。测度扩张、可测选择、谱论等基础数学采用明确标记的标准定理；本理论的主要连接结论均给出证明。

\hypertarget{ux539fux59cbux6982ux5ff5ux4e0eux672cux7248ux6570ux5b66ux5bf9ux8c61}{%
\section{0.1 原始概念与本版数学对象}\label{ux539fux59cbux6982ux5ff5ux4e0eux672cux7248ux6570ux5b66ux5bf9ux8c61}}

\begin{longtable}[]{@{}
  >{\raggedright\arraybackslash}p{(\columnwidth - 4\tabcolsep) * \real{0.3333}}
  >{\raggedright\arraybackslash}p{(\columnwidth - 4\tabcolsep) * \real{0.3333}}
  >{\raggedright\arraybackslash}p{(\columnwidth - 4\tabcolsep) * \real{0.3333}}@{}}
\toprule\noalign{}
\begin{minipage}[b]{\linewidth}\raggedright
原始问题或概念
\end{minipage} & \begin{minipage}[b]{\linewidth}\raggedright
本版数学实现
\end{minipage} & \begin{minipage}[b]{\linewidth}\raggedright
关键限定
\end{minipage} \\
\midrule\noalign{}
\endhead
\bottomrule\noalign{}
\endlastfoot
Being 是持续的 Becoming & 过程、历史分支、经验证的发展运输 & 持续不要求状态不变 \\
Ω-loop 与关系身份 & 预测---因果闭合、恢复/不变结构、机制毁损超图 & 几何环是特殊机制 \\
结构化开放、specific coupling & 内部因子与环境接口的条件核分解 & 资源交换不等于身份丧失 \\
Type I / Type II & 被动持续、内生维持与适应能力的证书剖面 & 不强制所有系统二分 \\
Π/Φ 双驱动 & 独立锚定的资源约束与信息约束；低维响应假设 & 数值分解和二维性须另检验 \\
GOD & 给定评价问题的最优动作对应 & 通常集合值；不等于即时梯度 \\
GOA & 最优闭环的长期集合、占据分布或条件生存结构 & 唯一固定点只属特例 \\
子对象---父对象 & 父层预测核心、切割不可约性、子层保真度 & 强相关不独自证明父对象 \\
自修改与自由 & 增广结构状态、策略扩展、目标修订能力 & 扩展有成本与失效条件 \\
意义 & 功能贡献、行为评价、规范桥接分别表达 & 不从稳定性直接推出义务 \\
\end{longtable}

原论文的思想在这里承担研究问题和架构选择的来源。数学定理承担条件性蕴涵的责任。两者联系紧密，但不能把哲学动机本身当作证明。

\hypertarget{ux9605ux8bfbux4e3bux7ebf}{%
\section{0.2 阅读主线}\label{ux9605ux8bfbux4e3bux7ebf}}

正文依次展开：

\begin{enumerate}
\def\labelenumi{\arabic{enumi}.}
\tightlist
\item
  第 1--7 章：一般过程、可测因子、规范预测身份、信息与接口、动态闭合。
\item
  第 8--12 章：持续与恢复、准平稳结构、历史发展、对象证书及 Ω。
\item
  第 13--17 章：实现、毁损、分辨率、组合、有限样本和逆问题。
\item
  第 18--24 章：可行未来、价值、控制、GOD/GOA、信息代价、双驱动与层级协调。
\item
  第 25--27 章：遗传、选择、个体性转变、内生目的与反身学习。
\item
  第 28--31 章：完整接口、应用方法、算例和跨模块装配定理。
\item
  第 32--38 章：类型化层级、缺陷校准、客体与科学端口、同源运行、有限支持生存及训练选择后的真实联合风险。
\item
  附录：符号、106 结果与完整来源索引、原数值程序、历史记录与维护约定。
\end{enumerate}

\textbf{依赖方向}是``世界与协议 → 候选表示 → 预测/因果/历史证书 → 实现与完整性''，以及在增加评价机制后``可行策略 → 价值 → GOD → GOA''。选择可以改变控制器与未来对象结构，但每一代的机制模型须先明确；不能用逻辑循环替代时间中的反馈。

\hypertarget{ux9648ux8ff0ux72b6ux6001}{%
\section{0.3 陈述状态}\label{ux9648ux8ff0ux72b6ux6001}}

\begin{itemize}
\tightlist
\item
  \textbf{D}：定义，固定一个数学对象或操作语义。
\item
  \textbf{P}：本文件给出证明的结果；不表示该结果为 UEOT 首创。
\item
  \textbf{K}：调用的标准定理，列明适用条件和承担的作用。
\item
  \textbf{H}：需要独立说明或检验的机制、领域或低维假设。
\item
  \textbf{E}：实际数据或指定实验能支持的证书。
\item
  \textbf{N}：反例或不可识别性结果。
\item
  \textbf{O}：开放问题，不进入已证依赖链。
\end{itemize}

每个主要正向结果使用稳定编号，例如 P-PRED-01。发现前提不足时修改该条及其依赖，不通过增加一个未经检验的``普适原则''保住结论。

\hypertarget{ux7ef4ux62a4ux7248-3.1ux81eaux8db3ux7406ux8bbaux4e0eux5f62ux5f0fux5316ux5b9eux73b0ux7684ux5bf9ux5e94}{%
\section{0.4 维护版 3.1：自足理论与形式化实现的对应}\label{ux7ef4ux62a4ux7248-3.1ux81eaux8db3ux7406ux8bbaux4e0eux5f62ux5f0fux5316ux5b9eux73b0ux7684ux5bf9ux5e94}}

维护版 3.1 完整承接第 0--31 章及附录中的 106 个稳定 P 编号，保留其原有陈述、证明、标准依赖、算法和完整数值代码。新增连接使用 M-32.1 等章内编号，既不重新定义已有 P-ID，也不通过增加一个宏大总定理取消领域前提。

本版新增第 32--38 章，把实际 Lean 中已经完成的层级分型、四生成元压缩、定量缺陷校准、操作对象与条件生命周期、同源控制与修复、有限支持可行性、数据选择后的联合风险重写为连续数学正文。各章同时给关键证明、必要性反例及精确 Lean 映射；读者不必先接受哲学解释才能检查这些命题。

有三种``完整''需要区分：本文内容自足；声明范围内的命题有证明；具体世界模型已被实验支持。第一项不自动推出后两项。106/106 源端点状态、P0--P12 工作包 CLOSED、C1--C7 的 package closure 也不是同一种事实。当前计数压缩是 11 个生成端点、89 个保留适配、6 个保留边界，见第 32 章。P12 的明确边界及科学接口的 real-world support 尚未确认，分别进入第 34 章正文。

附录 D 的 2371 项数值检查及其输出按 v3.0 历史记录保留。本轮新 Lean 检查是另一个证据集合，见第 38.8 节；不把旧代码输出冒称本轮重新运行。

\hypertarget{ux57faux7840ux7ea6ux5b9aux4e0eux53efux590dux7528ux6570ux5b66ux5de5ux5177}{%
\chapter{1. 基础约定与可复用数学工具}\label{ux57faux7840ux7ea6ux5b9aux4e0eux53efux590dux7528ux6570ux5b66ux5de5ux5177}}

\hypertarget{ux7a7aux95f4ux6838ux4e0eux7b26ux53f7}{%
\section{1.1 空间、核与符号}\label{ux7a7aux95f4ux6838ux4e0eux7b26ux53f7}}

随机变量所在空间默认标准 Borel。具有拓扑、极限或紧性结论时，再明确采用 Polish 空间、紧度量空间或光滑流形。标准 Borel 假设保证通常的正则条件概率存在；它不保证任意等价类商或任意选择映射自动可测。

概率核写为 \[
K:\mathcal X\longrightarrow\mathcal P(\mathcal Y),\qquad
K(x,B)=P(Y\in B\mid x).
\] 若 \(f:\mathcal Y\to\mathcal Z\) 可测，推送核为 \[
(f_\#K)(x,A)=K(x,f^{-1}A).
\] 概率分布采用行向量，状态函数采用列向量；因此分布演化为 \(\mu P\)，函数演化为 \(Pv\)。

用 \(H_t\) 表示历史，\(T\) 表示时域，\(U_t\) 表示环境接口，\(V\) 表示有限载体全集，\(\mathcal V\) 表示持续性评估域。折扣为 \(\beta\in(0,1)\)，吸收率为 \(\lambda\)，置信失效概率为 \(\alpha\)。对象尺度为 \(s\)，选择/评价层为 \(\ell\)。

\hypertarget{ux4fe1ux606fux8dddux79bbux4e0eux5355ux4f4d}{%
\section{1.2 信息、距离与单位}\label{ux4fe1ux606fux8dddux79bbux4e0eux5355ux4f4d}}

全文默认自然对数，熵与 KL 的单位为 nat；标明 bit 的算例使用底数 2。对数基改变时常数也须相应改变。 \[
D_{\rm TV}(P,Q)=\sup_A|P(A)-Q(A)|.
\] 有限空间上 \(D_{\rm TV}=\frac12\sum_y|p_y-q_y|\)。设 \(P\ll Q\)， \[
D_{\rm KL}(P\|Q)=\int\log\frac{dP}{dQ}\,dP;
\] 否则定义为 \(+\infty\)。\(0\log0=0\)。

\textbf{P-MET-01：核与读出的数据处理。} 对可测映射 \(f\) 和概率核 \(K\)， \[
D_{\rm TV}(f_\#P,f_\#Q)\le D_{\rm TV}(P,Q),\qquad
D_{\rm TV}(PK,QK)\le D_{\rm TV}(P,Q).
\] 若 \(f\) 是双可测双射，第一式取等号。

\textbf{证明。} 第一式将事件限制为逆像 \(f^{-1}A\)。第二式写成对 \(0\le K(x,A)\le1\) 的积分，使用有界函数形式的 TV 对偶式。双射情形应用 \(f^{-1}\) 得反向不等式。\(\blacksquare\)

KL 数据处理是 \textbf{K-INF-01}：共同随机通道处理不增加 KL；可由条件 KL 链式法则证明。对有限变量，该链式法则直接来自联合概率的乘法分解；标准 Borel 版本由正则条件核与相对熵分解得到。

\textbf{P-MET-02：跨度估计。} 对有界实函数 \(v\)，记 \(\operatorname{sp}(v)=\sup v-\inf v\)，则 \[
|Pv-Qv|\le\operatorname{sp}(v)D_{\rm TV}(P,Q).
\] \textbf{证明。} 常数的积分差为零。减去 \((\sup v+\inf v)/2\) 后，函数上确界范数为跨度的一半；再用 \(|\int g\,d(P-Q)|\le2\|g\|_\infty D_{\rm TV}(P,Q)\)。\(\blacksquare\)

\textbf{K-INF-02：Pinsker。} \[
D_{\rm TV}(P,Q)\le\sqrt{\tfrac12D_{\rm KL}(P\|Q)}.
\] 该常数对应自然对数。信息链式法则及该不等式是第 5、6、22 章的基础，而不是新的 UEOT 公设。

\hypertarget{ux4e09ux7c7bux8befux5deeux4e0dux80fdux6df7ux5408}{%
\section{1.3 三类误差不能混合}\label{ux4e09ux7c7bux8befux5deeux4e0dux80fdux6df7ux5408}}

\begin{itemize}
\tightlist
\item
  点态/最坏误差：对所有声明可达输入成立。
\item
  平均误差：在指定输入分布 \(\mu\) 下成立。
\item
  估计误差：以至少 \(1-\alpha\) 的采样概率成立。
\end{itemize}

平均误差不是最坏误差；训练分布误差不是部署分布误差；一次仿真观测到的小误差不是一个置信上界。第 16 章给出可合法转换它们的条件。

\hypertarget{ux53efux6d4bux5546ux7684ux4f7fux7528ux89c4ux5219}{%
\section{1.4 可测商的使用规则}\label{ux53efux6d4bux5546ux7684ux4f7fux7528ux89c4ux5219}}

不直接把``核相等的等价类集合''假定成良好可测空间。优先把状态定义为概率测度值随机元，或固定底概率空间内的子 \(\sigma\)-代数。需要截面、商流形、连续选择、弱极限时逐条增加条件。有限空间中的逐点结论与一般空间中的模零结论分开。

\hypertarget{ux540cux4e00ux5173ux7cfbux540cux4e00ux5ea6ux91cfux4e0eux540cux4e00ux8bc1ux636e}{%
\section{1.5 同一关系、同一度量与同一证据}\label{ux540cux4e00ux5173ux7cfbux540cux4e00ux5ea6ux91cfux4e0eux540cux4e00ux8bc1ux636e}}

同一零误差关系可以有许多定量表示。例如离散缺陷 \(d(x,y)=\mathbf1_{x\ne y}\) 与 \(d'(x,y)=100d(x,y)\) 具有相同的零集，但预算 \(d\le1\) 与 \(d'\le1\) 不等价。故有损商、发展管道和控制成本中的实数必须带上度量及标度。

第 33 章使用 \(a d\le d'\le b d\) 的校准来传递预算；只保留零关系时不能作这种数值传递。对概率还要区分训练测度 \(\mu\)、给定训练记录的未来条件律，以及该未来过程实际访问状态---动作对的占据律 \(\nu_n\)。它们通常位于不同空间，不能因都叫``分布''而代换。

Lean 对应为 \nolinkurl{MinCore.RelationDefectPresentation} 与 DefectCalibration，以及 \nolinkurl{UnifiedClosure.selectedTrainingDeployment} 和 coverageDeploymentRowLaw；完整命名与来源见附录 F。

\hypertarget{ux4e16ux754cux8fc7ux7a0bux5b9eux9a8cux534fux8baeux4e0eux65f6ux95f4}{%
\chapter{2. 世界过程、实验协议与时间}\label{ux4e16ux754cux8fc7ux7a0bux5b9eux9a8cux534fux8baeux4e0eux65f6ux95f4}}

\hypertarget{ux4e00ux822cux53d7ux63a7ux8fc7ux7a0b}{%
\section{2.1 一般受控过程}\label{ux4e00ux822cux53d7ux63a7ux8fc7ux7a0b}}

\textbf{D-WORLD。} 在时刻 \(t\)，历史 \[
H_t=(X_0,A_0,X_1,\ldots,A_{t-1},X_t)
\] 取值于历史空间 \(\mathcal H_t\)。环境、传感记录、内部结构等可作为 \(X_t\) 的分量。下一步由核 \[
q_t(dx'\mid h,a)
\] 给定，策略为因果核 \(\pi_t(da\mid h)\)。先给初始分布，再依次乘接策略与状态核。

\textbf{K-PROC-01：路径律存在。} 可数时间索引、标准 Borel 空间和可测概率核下，Ionescu--Tulcea 扩张给出唯一的相应路径概率律。它保证指定过程模型的概率一致性，不保证模型对应自然系统。

\textbf{P-PROC-01：历史状态化。} 任意上述历史依赖过程，在增广状态 \(Z_t=(t,H_t)\) 上都是时间齐次 Markov 控制过程。

\textbf{证明。} 给定 \((t,h)\) 与动作 \(a\)，抽样 \(x'\sim q_t(\cdot\mid h,a)\)，下一增广状态唯一为 \((t+1,h,a,x')\)。其分布不再依赖增广状态之前的记录。\(\blacksquare\)

这一构造可能无限增长，因此``能够 Markov 化''不等于``已经找到有用的有限记忆''。

\hypertarget{ux8fdeux7eedux65f6ux95f4}{%
\section{2.2 连续时间}\label{ux8fdeux7eedux65f6ux95f4}}

连续时间模块使用以下已明确选择的模型之一：

\begin{itemize}
\tightlist
\item
  连续半流 \(\phi_t\)；
\item
  强 Markov/Feller 半群 \(P_t\)；
\item
  有限状态生成元 \(L=(l_{xy})\)；
\item
  有适当存在唯一性、非爆炸与可积条件的扩散 \[
  dX_t=b(X_t,a_t)\,dt+\sigma(X_t,a_t)\,dW_t.
  \]
\end{itemize}

路径空间可采用连续路径 \(C([0,T],\mathcal X)\) 或适当的 càdlàg 路径空间。不能把任意带不可测选择的``历史集合''自动当成标准路径模型。

\hypertarget{ux534fux8baeux4e0eux5e72ux9884}{%
\section{2.3 协议与干预}\label{ux534fux8baeux4e0eux5e72ux9884}}

\textbf{D-PROT。} 协议包含观测映射、可实施动作、读出顺序、时间窗、停止规则、未来测试族及实施约束。协议 \(i\) 对历史 \(h\) 诱导记录核 \(K_i(h,\cdot)\)。

后续干预默认发生在已经形成的历史之后；这样不同协议共享历史参考分布 \(\mu\)。若实验改变了历史生成机制，则必须引入不同环境参考分布，采用共同结构核或支配测度，并重新说明比较的支撑。

观测条件概率 \(P(Y\mid A=a)\) 与干预响应 \(P(Y\mid do(a))\) 是不同输入。随机化试验、可验证的干预模型或充分因果识别假设才能连接两者。

\hypertarget{ux6709ux9650ux534fux8baeux5b8cux6574ux672aux6765ux4e0eux53efux6570ux6027}{%
\section{2.4 有限协议、完整未来与可数性}\label{ux6709ux9650ux534fux8baeux5b8cux6574ux672aux6765ux4e0eux53efux6570ux6027}}

数学上的``完整未来''是一个被声明的未来记录随机元，不是全宇宙所有事实。对无限未来，柱集构成可数决定类需依赖可数时间和可数生成可测结构。

本核心的规范预测定理采用有限或可数协议族。不可数动作空间可用于控制，但若要同时用所有动作定义规范状态，须另有共同点态结构核及可测性，或可数稠密族加连续性。该保护条件不能由未来空间可数生成来替代。

\hypertarget{ux6761ux4ef6ux90e8ux7f72ux7684ux751fux6210ux6b21ux5e8f}{%
\section{2.5 条件部署的生成次序}\label{ux6761ux4ef6ux90e8ux7f72ux7684ux751fux6210ux6b21ux5e8f}}

若先观察训练记录 \(\omega\sim\mu\)，再据此选择模型或策略 \(s(\omega)\)，应由声明的条件部署核 \(K_{s(\omega)}\) 生成未来 \(h\)。有限情形的联合质量是

\[J(\omega,h)=\mu(\omega)K_{s(\omega)}(h).\]

这是一个实际构造，其训练边缘等于 \(\mu\)，未来边缘是 \(\sum_\omega\mu(\omega)K_{s(\omega)}\)。选择可以依赖训练；相应未来风险并不要求选择与训练独立。需要验证的是给定训练记录后未来确实按这个核生成。

固定候选未来的边缘分布吻合，不足以推出数据选择后的条件吻合。第 38 章证明真实 joint 的 tower，再讨论错误条件律为何会破坏风险结论。这样``学习→控制→未来''由时间顺序实现，而不把要证明的风险等式存入一个证书字段。

\hypertarget{ux56e0ux5b50ux8868ux793aux534fux53d8ux4e0eux5bf9ux8c61ux5019ux9009}{%
\chapter{3. 因子、表示协变与对象候选}\label{ux56e0ux5b50ux8868ux793aux534fux53d8ux4e0eux5bf9ux8c61ux5019ux9009}}

\hypertarget{ux5d4cux5165ux56e0ux5b50}{%
\section{3.1 嵌入因子}\label{ux5d4cux5165ux56e0ux5b50}}

\textbf{D-FAC。} 对固定 \((\mathcal H,\mathscr H,\mu)\) 和可测统计量 \(M=f(H)\)，定义 \[
\mathscr A_f=\sigma(f)\subseteq\mathscr H.
\] 在同一底空间内模 \(\mu\)-零集识别，偏序为 \(\mathscr A_f\preceq\mathscr A_g\iff\mathscr A_f\subseteq\mathscr A_g\)。

在标准 Borel 值统计量下，若 \(f\) 对 \(\sigma(g)\) 可测，Doob--Dynkin 因子化给出 \(f=\varphi\circ g\) 的相应几乎处处版本。因而偏序表示``后者包含恢复前者所需的信息''。

同一统计量的可逆重编码不改变嵌入因子。两个独立均匀比特的因子虽然抽象同构，却是不同嵌入子代数。这一区分保留了对象实例，而不仅保留对象类型。

\hypertarget{ux8868ux793aux534fux53d8}{%
\section{3.2 表示协变}\label{ux8868ux793aux534fux53d8}}

\textbf{P-FAC-01。} 设 \(T:\mathcal X\to\mathcal X'\) 为双可测双射，同时运输 \[
P'(Tx,a,T B)=P(x,a,B),\qquad f'=f\circ T^{-1},
\] 以及观测、干预和奖励。则宏观路径律、预测充分性、精确动态闭合及控制价值保持。

\textbf{证明。} \(f'(Tx)=f(x)\)。一步推送相同；逐步条件化归纳得到所有有限路径记录律相同。路径律一致推出任何声明未来记录分布一致。若有奖励，逐项期望相同，取策略上确界得到价值相同。\(\blacksquare\)

这不表示任意有损传感器都保留对象；有损映射只服从数据处理。

\hypertarget{ux56e0ux5b50ux5750ux6807ux548cux8f7dux4f53}{%
\section{3.3 因子、坐标和载体}\label{ux56e0ux5b50ux5750ux6807ux548cux8f7dux4f53}}

例：\(Y_1=X_1\oplus X_2,\ Y_2=X_2\) 是可逆重编码。原变量 \(X_1\) 在新坐标下需要 \(Y_1,Y_2\) 才能恢复。最小坐标数量因重编码改变，预测关系却被正确运输。

反例的另一个方向：两块不同硬件上存放相同比特，可生成同一预测因子，却有不同物理失效机制。因此预测因子不唯一决定物理位置。第 13 章专门处理实现。

\hypertarget{ux56e0ux5b50ux8c31ux4e0eux5c3aux5ea6}{%
\section{3.4 因子谱与尺度}\label{ux56e0ux5b50ux8c31ux4e0eux5c3aux5ea6}}

对指定证书，收集所有通过的嵌入因子，按包含关系形成因子谱。独立过程 \(A,B\) 及联合过程 \((A,B)\) 都可能预测闭合。除非指定目标、压缩准则或组合条件，动力学不提供唯一对象分割。

表示最小性、载体极小性、物理不可分性和父对象不可约性分别属于不同问题。父对象可包含真实子对象；这不会与某个固定预测任务的最小充分状态相冲突。

\hypertarget{ux4eceux9884ux6d4bux56e0ux5b50ux5230ux7ec4ux7ec7ux5165ux53e3}{%
\section{3.5 从预测因子到组织入口}\label{ux4eceux9884ux6d4bux56e0ux5b50ux5230ux7ec4ux7ec7ux5165ux53e3}}

一个因子保存某组响应，并不意味着它保存组织的所有谓词。实际同源核心先把需要保留的 observation、organization、currentProgram、currentResource、safe 明确注册为输出，再对同一有限转移与奖励作稳定细化。商中这些属性因此可恢复，而不是因为商被命名为``对象''就自动存在。

\nolinkurl{OrganizationalCoreAssembly} 的 \nolinkurl{organizational_terminal_preserves_registered_interfaces} 与 \nolinkurl{organizational_terminal_is_coarsest} 分别证明属性保持和指定初始区分下的最粗性。其组织标签的物理含义仍需建模，最粗性也只相对于注册输出、动作、奖励和真实核。第 35 章将此入口接到实际源上的生存和安全控制。

\hypertarget{ux89c4ux8303ux9884ux6d4bux8eabux4efdux5b58ux5728ux6700ux5c0fux6027ux7ec6ux5316ux4e0eux66f4ux65b0}{%
\chapter{4. 规范预测身份：存在、最小性、细化与更新}\label{ux89c4ux8303ux9884ux6d4bux8eabux4efdux5b58ux5728ux6700ux5c0fux6027ux7ec6ux5316ux4e0eux66f4ux65b0}}

\hypertarget{ux6982ux7387ux6d4bux5ea6ux503cux7684ux89c4ux8303ux72b6ux6001}{%
\section{4.1 概率测度值的规范状态}\label{ux6982ux7387ux6d4bux5ea6ux503cux7684ux89c4ux8303ux72b6ux6001}}

设 \(H\sim\mu\)，协议 \(\mathcal I=\{i_1,i_2,\ldots\}\)，未来空间均为标准 Borel；对有限协议采用有限乘积。固定可测响应核 \(K_i(h,\cdot)\)，定义 \[
C_{\mathcal I}(h)=\bigl(K_{i_1}(h,\cdot),K_{i_2}(h,\cdot),\ldots\bigr).
\] \(\mathcal P(\mathcal Y)\) 使用评价 \(\sigma\)-代数；可数乘积仍是标准 Borel。称 \(C_{\mathcal I}\) 为规范干预预测状态。未来记录取完整声明路径时，称 FFIPS。

对象候选 \(f\) 的未来记为 \(Y_f^+\)，相应状态记为 \(C^{f}_{\mathcal I}\)。\textbf{对象只需保存其声明未来与接口所需的信息，不要求保存整个世界的未来。} 全局 FFIPS 是一个统一比较工具，不能直接与单个对象等同。

\hypertarget{ux6700ux5c0fux6027ux5b9aux7406}{%
\section{4.2 最小性定理}\label{ux6700ux5c0fux6027ux5b9aux7406}}

\textbf{P-PRED-01：共同协议的规范最小性。} 若统计量 \(S=s(H)\) 取值于标准 Borel，并对每个 \(i\) 存在核 \(L_i\) 满足 \[
K_i(H,\cdot)=L_i(S,\cdot)\quad\mu\text{-a.s.},
\] 则存在可测 \(\varphi\) 使 \[
C_{\mathcal I}=\varphi(S)\quad\mu\text{-a.s.}
\] 且 \(C_{\mathcal I}\) 本身对每个协议充分。因此 \(\sigma(C_{\mathcal I})\) 是共同响应所需的最小嵌入因子。

\textbf{证明。} 每个未来空间有可数决定概率测度的事件类。对全部协议及其决定事件的概率一集合取可数交，得到共同概率一集合。在其上 \[
C_{\mathcal I}=(L_{i_1}(S,\cdot),L_{i_2}(S,\cdot),\ldots),
\] 右侧是 \(S\) 的可测函数。每个 \(K_i\) 又是 \(C_{\mathcal I}\) 的坐标。对未来事件 \(B\)，塔式性质给 \[
E[K_i(H,B)\mid C_{\mathcal I}]=K_i(H,B),
\] 所以以规范状态为条件的未来核与完整历史的核相同。\(\blacksquare\)

若另一个规范最小统计量满足同一要求，两者生成同一嵌入因子。此结论针对给定协议、未来和参考分布，不保证自然世界存在唯一终极对象清单。

\hypertarget{ux4e0dux53efux6570ux7248ux672cux7684ux969cux788d}{%
\section{4.3 不可数版本的障碍}\label{ux4e0dux53efux6570ux7248ux672cux7684ux969cux788d}}

\textbf{N-PRED-01。} 令 \(H\) 在 \([0,1]\) 上均匀，所有协议的真实未来为常数零。对每个 \(i\in[0,1]\)，取一个条件律版本 \[
K_i(h,\{1\})=\mathbf1_{\{h=i\}}.
\] 每个版本只在一个零集上被改动，常数统计量对每个协议分别充分。但逐点收集所有坐标会区分每个 \(h\)。因此不能无条件把``分别几乎处处''改成``同时在共同集合上''。

可采用的扩展方式是：共同结构核在所有可达点明确定义；或存在可数稠密协议集，响应对协议连续，且连续性与版本均在一个共同概率一集合上成立。此时稠密集上的等式通过连续性推广。若使用不同协议对应的不同历史分布，还需显式说明共同支撑。

\hypertarget{ux76eeux6807ux5b8cux5907ux6027ux4e0eux7ec6ux5316}{%
\section{4.4 目标完备性与细化}\label{ux76eeux6807ux5b8cux5907ux6027ux4e0eux7ec6ux5316}}

\textbf{P-PRED-02：目标推送与协议细化。} 若 \(Y_1=T(Y_2)\)，则其规范预测核满足 \[
C_1=T_\#C_2.
\] 若 \(\mathcal I_1\subseteq\mathcal I_2\)，则 \[
\sigma(C_{\mathcal I_1})\subseteq\sigma(C_{\mathcal I_2}).
\] 对可数递增族， \[
\sigma(C_{\cup_n\mathcal I_n})=\bigvee_n\sigma(C_{\mathcal I_n}).
\]

\textbf{证明。} 目标变换对每个条件核作同一推送。小协议族的坐标是大协议族的投影；联合族的每个坐标在某一级出现，因此两侧生成同一 \(\sigma\)-代数。\(\blacksquare\)

这也给出有限未来窗 \(T_1<T_2\) 的单调精化。任意有限平台仍可能被下一项尚未测试的协议打破：可构造前 \(k\) 个测试对两历史完全相同，第 \(k+1\) 个才不同。

\hypertarget{ux5168ux672aux6765ux72b6ux6001ux7684ux9012ux5f52ux66f4ux65b0}{%
\section{4.5 全未来状态的递归更新}\label{ux5168ux672aux6765ux72b6ux6001ux7684ux9012ux5f52ux66f4ux65b0}}

\textbf{P-PRED-03：有限字母过程的递归更新。} 假设所有响应来自同一个因果过程；协议族对行动前缀和正概率观察后的允许延续封闭，并包含相应联合未来柱集。则预测等价在下一步更新中保持，存在 \[
C_{t+1}=\mathcal U_t(C_t,a_t,o_{t+1})
\] 的可测版本。

\textbf{证明。} 对下一观察 \(o\) 和任何允许延续记录事件 \(B\)，后续条件概率为 \[
P(B\mid h,a,o)
=\frac{P(o,B\mid h,a)}{P(o\mid h,a)}.
\] 正概率分母和分子都是当前规范状态保存的协议概率。若 \(h,h'\) 具有相同规范状态，更新后全部延续核相同。有限字母及可数核坐标使更新可测。零概率分支可任意扩展为数学核，但不因此获得经验认证。\(\blacksquare\)

只有单步边缘相同不足以保证递归充分。若协议不允许后续自适应操作，也不能把为该协议构造的状态直接用于更大的控制策略类。

\hypertarget{ux9012ux5f52ux5145ux5206ux6027ux7684ux6b63ux5f0fux5b9eux73b0ux8fb9ux754c}{%
\section{4.6 递归充分性的正式实现边界}\label{ux9012ux5f52ux5145ux5206ux6027ux7684ux6b63ux5f0fux5b9eux73b0ux8fb9ux754c}}

第 39.7、39.12 节已从声明的非空标准 Borel 潜状态、原 P/O 核、有限／可数协议及实际完整历史程序，构造后验响应递归、真实过去---未来联合律和模零最小预测消费者。原过去／filter 联合律只需观察可数生成；联合可测的 ambient 响应更新及以响应数组构造的 fresh／记忆预测核，还保留标准 Borel 观察、可数且单点可测动作的条件。继承控制记忆时保留响应与记忆，新未来程序以登记的空记忆开始时才取得相应 canonical-only 结论。未覆盖的模型仍须按本章条件及第 4.3 节的不可数版本保护处理；这些结果不从有限数据自动推断整个源。

P-PRED-03 的 Lean 源 PredictionUpdate 证明分子、分母已经通过当前状态因子分解时的 Bayes 更新因子化，并证明有限动作/观察及相应坐标条件下的联合可测性。它没有从任意数据中自动产生这些充分坐标。一般空间中的参数化解体见第 34.6 节的 C1。

有限受控源还可以抽取一个较弱但很有用的对象：给定粗历史的微观末状态支持。第 37 章从真实历史条件分布证明，该支持恰好等于按已记录动作和新观察递推的 traceSupport。此支持递推只表达哪些状态仍可能；完整后验权重仍由概率决定。正历史与零历史的区分不能省去。

\hypertarget{ux9884ux6d4bux4fe1ux606fux8bb0ux5fc6ux590dux6742ux5ea6ux4e0eux53efux538bux7f29ux6027}{%
\chapter{5. 预测信息、记忆复杂度与可压缩性}\label{ux9884ux6d4bux4fe1ux606fux8bb0ux5fc6ux590dux6742ux5ea6ux4e0eux53efux538bux7f29ux6027}}

\hypertarget{ux4fe1ux606fux4fddux7559ux6052ux7b49ux5f0f}{%
\section{5.1 信息保留恒等式}\label{ux4fe1ux606fux4fddux7559ux6052ux7b49ux5f0f}}

令 \(M=f(H)\)，暂固定一个实验分布，并假设有关互信息有限。

\textbf{P-INFO-01。} \[
I(H;Y)=I(M;Y)+I(H;Y\mid M).
\] 因此当 \(I(H;Y\mid M)\le\varepsilon\) 时， \[
I(M;Y)\ge I(H;Y)-\varepsilon.
\] 若 \(M\) 离散， \[
H(M)\ge I(H;Y)-\varepsilon.
\]

\textbf{证明。} 因 \(M\) 是 \(H\) 的函数，\(I(H,M;Y)=I(H;Y)\)。按另一顺序展开联合互信息即得；再用 \(I(M;Y)\le H(M)\)。\(\blacksquare\)

若 \(Y=H\) 是均匀 \(N\) 值变量的未来复制，则精确充分状态至少有 \(N\) 个值和 \(\log N\) 熵。任何对象发现算法都不能保证所有系统可非平凡压缩。

\hypertarget{ux4eceux4fe1ux606fux6b8bux5deeux5230ux9884ux6d4bux8befux5dee}{%
\section{5.2 从信息残差到预测误差}\label{ux4eceux4fe1ux606fux6b8bux5deeux5230ux9884ux6d4bux8befux5dee}}

\textbf{P-INFO-02。} 对标准 Borel 变量 \(H,Y,M,U\)，若 \(M,U\) 是历史的函数，令 \[
\varepsilon=I(H;Y\mid M,U).
\] 则 \[
E D_{\rm TV}\bigl(P(Y\mid H),P(Y\mid M,U)\bigr)
\le\sqrt{\varepsilon/2}.
\]

\textbf{证明。} 条件互信息等于两条件核 KL 的平均。逐条件用 Pinsker，再用平方根的 Jensen 不等式。\(\blacksquare\)

\textbf{推论。} 若两表示 \(Z_1,Z_2\) 的信息残差分别不超过 \(\varepsilon_1,\varepsilon_2\)，则 \[
E D_{\rm TV}(P(Y\mid Z_1),P(Y\mid Z_2))
\le\sqrt{\varepsilon_1/2}+\sqrt{\varepsilon_2/2}.
\] 这是预测律的一致性，不是两套坐标、物理载体或因果机制必相同。

\hypertarget{ux9884ux6d4bux7387ux5931ux771f}{%
\section{5.3 预测率失真}\label{ux9884ux6d4bux7387ux5931ux771f}}

\textbf{D-RD。} 允许编码器 \(P(M\mid H,U)\) 使用独立随机性，解码器为 \(Q(\cdot\mid M,U)\)。定义 \[
R_{\rm obj}(\varepsilon)=
\inf I(H;M\mid U)
\] 其中约束为 \[
E D_{\rm TV}\bigl(P(Y\mid H,U),Q(\cdot\mid M,U)\bigr)
\le\varepsilon.
\] 编码器不能读取尚未发生的 \(Y\)。该量衡量指定平均预测误差所需的记忆信息。

\textbf{P-INFO-03：零失真复杂度。} 若 \(C=P(Y\in\cdot\mid H,U)\) 是离散随机元且 \(H(C\mid U)<\infty\)，则 \[
R_{\rm obj}(0)=H(C\mid U).
\]

\textbf{证明。} 零平均 TV 意味着两个核几乎处处相同，故 \(C\) 可由 \((M,U)\) 恢复。于是 \[
I(H;M\mid U)\ge I(C;M\mid U)=H(C\mid U).
\] 取 \(M=C\)，解码器输出该核，恰好达到下界。证明对随机编码器也成立。\(\blacksquare\)

\(\varepsilon\) 增大时可行编码器集合增大，所以 \(R_{\rm obj}\) 非增。该定义是信息约束，并不单独给出计算时间、能耗或生物记忆容量。

\hypertarget{ux8eabux4efdux4fddux7559ux7684ux4fe1ux606fux4e0bux754c}{%
\section{5.4 身份保留的信息下界}\label{ux8eabux4efdux4fddux7559ux7684ux4fe1ux606fux4e0bux754c}}

\textbf{P-INFO-04：Fano 型身份下界。} 令 \(J\) 均匀取 \(K\ge2\) 个身份，未来记录 \(Y_T\) 可解码 \(\widehat J\)，错误率为 \(e\)。则 \[
I(J;Y_T)\ge\log K-h_2(e)-e\log(K-1).
\]

\textbf{证明。} 令错误指示 \(E=\mathbf1_{\{\widehat J\ne J\}}\)。给定记录与错误标记，正确时身份已知，错误时至多有 \(K-1\) 种可能。因此 \(H(J\mid Y_T)\le H(E)+e\log(K-1)\)，再从 \(\log K\) 中相减。\(\blacksquare\)

条件二元版本：若未来身份变量 \(B\in\{0,1\}\) 由 \((M,U)\) 以错误率不超过 \(\epsilon\le1/2\) 解码， \[
I(M;B\mid U)\ge H(B\mid U)-h_2(\epsilon).
\] 这里是可区分身份系综的信息，而非单个确定对象必须携带正熵的断言。

\hypertarget{ux9057ux6f0fux4e0e-packing-ux4e0bux754c}{%
\section{5.5 遗漏与 packing 下界}\label{ux9057ux6f0fux4e0e-packing-ux4e0bux754c}}

\textbf{P-INFO-05。} 若两历史被同一表示 \(R\) 合并，但某协议的未来核相距 \(\delta>0\)，则任意仅依赖 \(R\) 的解码核至少在其中一个历史上有 TV 误差 \(\delta/2\)。

\textbf{证明。} 两个误差的和至少为两真实核的距离。\(\blacksquare\)

若 \(N\) 个历史两两响应距离大于 \(2\varepsilon\)，一个统一误差不超过 \(\varepsilon\) 的确定性表示至少有 \(N\) 个取值；否则鸽巢原理与前述三角不等式矛盾。若输入分布恰为这 \(N\) 个历史的均匀分布，其表示熵至少为 \(\log N\)。

\hypertarget{ux4fe1ux606fux969cux788dux5728ux7edfux4e00ux7406ux8bbaux4e2dux7684ux4f4dux7f6e}{%
\section{5.6 信息障碍在统一理论中的位置}\label{ux4fe1ux606fux969cux788dux5728ux7edfux4e00ux7406ux8bbaux4e2dux7684ux4f4dux7f6e}}

InformationPacking 的 \nolinkurl{common_decoder_half_distance}、packing\_forces\_injective 和 packing\_ncard\_image 从真实概率律的 TV 三角关系给出本章的精确障碍。P-INFO-05 因而是压缩账本保留的边界，不能为了宣称所有端点都被生成元``生成''而消去它。

这些下界解释为何某些区别不能在指定精度下被合并。它们不决定计算时间、能耗或唯一信息驱动力；将 \(\Phi\) 选成 TV、Bayes 风险或记忆率之后，还需证明该量怎样进入同一策略评价和具体测量。第 23、33、34 章分别处理规范、定量校准和目的忠实性。

\hypertarget{ux7ed3ux6784ux5316ux5f00ux653eux73afux5883ux63a5ux53e3ux4e0eux56e0ux679cux8fb9ux754c}{%
\chapter{6. 结构化开放、环境接口与因果边界}\label{ux7ed3ux6784ux5316ux5f00ux653eux73afux5883ux63a5ux53e3ux4e0eux56e0ux679cux8fb9ux754c}}

\hypertarget{ux8054ux5408ux95edux5408ux5224ux636e}{%
\section{6.1 联合闭合判据}\label{ux8054ux5408ux95edux5408ux5224ux636e}}

令历史分解 \(H=(H^S,H^E)\)，内部状态 \(M=f(H^S)\)，接口 \(U=g(H^E)\)。对候选对象未来 \(Y_f^+\)，定义 \[
\varepsilon_{\rm joint}
=I(H^S,H^E;Y_f^+\mid M,U).
\]

\textbf{P-INT-01：结构化充分性的充要条件。} \[
Y_f^+\perp H\mid(M,U)
\iff C^f=\Psi(M,U)\quad\text{a.s.}
\] 对全部可数干预族采用共同版本时逐协议成立。

\textbf{证明。} 正向是 P-PRED-01 的最小性。反向由 \(E[C^f(B)\mid M,U]=C^f(B)\) 和塔式性质得到完整历史与表示条件下的未来概率相同。\(\blacksquare\)

\textbf{N-INT-01：分开的小残差不足。} 令独立均匀比特 \(X,E\)，\(Y=X\oplus E\)，\(M,U\) 都常数。则 \[
I(X;Y)=I(E;Y)=0,\qquad I(X,E;Y)=\log2.
\] 所以分别检验``内部剩余信息''和``环境剩余信息''，可能漏掉协同作用。

正确链式分解为 \[
I(X,E;Y\mid M,U)
=I(X;Y\mid M,U)+I(E;Y\mid X,M,U).
\] 第二项必须额外条件于 \(X\)。本版以联合残差为认证入口，两个未联合的指标只作诊断。

\hypertarget{ux6709ux9650ux7ed3ux6784ux6838ux7684ux6700ux5c0fux5185ux90e8ux73afux5883ux5206ux89e3}{%
\section{6.2 有限结构核的最小内部---环境分解}\label{ux6709ux9650ux7ed3ux6784ux6838ux7684ux6700ux5c0fux5185ux90e8ux73afux5883ux5206ux89e3}}

给定在全部允许 \((x,e,i)\) 上明确定义的有限响应核 \(K_i(x,e)\)。定义 \[
x\sim_Sx'\iff K_i(x,e)=K_i(x',e)\quad\forall e,i,
\] \[
e\sim_Ee'\iff K_i(x,e)=K_i(x,e')\quad\forall x,i.
\] 相应等价类为 \(M^*(x),U^*(e)\)。

\textbf{P-INT-02。} 存在良定义核 \[
K_i(x,e)=Q_i(M^*(x),U^*(e)).
\] 任何另一精确分离表示 \(K_i(x,e)=\widetilde Q_i(f(x),g(e))\)，均须分别细于 \(M^*,U^*\)。

\textbf{证明。} 先替换等价的内部点，再替换等价的环境点，核均不变，故商核良定义。若 \(f(x)=f(x')\)，则对每个环境和协议，两个核同为 \(\widetilde Q_i(f(x),g(e))\)，因此 \(x\sim_Sx'\)。环境方向相同。\(\blacksquare\)

若只有观测数据，零概率 \((x,e)\) 组合不能凭空得到结构核。上述点态全产品结论必须有全支撑或实验/机制定义；一般模零分解应采用 P-INT-01。

\hypertarget{ux56e0ux679cux63a5ux53e3}{%
\section{6.3 因果接口}\label{ux56e0ux679cux63a5ux53e3}}

将外部因素分成接口 \(U\) 和远端变量 \(E_{\rm far}\)。\textbf{H-INT} 声明，在允许的机制保持干预下， \[
P(Y_f^+\mid M,U,do(E_{\rm far}=e))
\] 不依赖 \(e\)。这是因果屏蔽证书的理想形式。实际数据须说明干预是否直接改变内部机制、是否改变接口、是否存在遗漏通道。

平均干预残差与最坏残差分别为 \[
\sup_i E_{\mu_i}D_{\rm TV}(K_i^{\rm full},K_i^{\rm obj}),
\qquad
\sup_i\operatorname*{ess\,sup}_{\mu_i}
D_{\rm TV}(K_i^{\rm full},K_i^{\rm obj}).
\] 后者更强。不能不注明语义而混用同一个误差符号。

\hypertarget{ux88abux52a8ux89c2ux6d4bux4e0dux80fdux552fux4e00ux786eux5b9aux56e0ux679cux65b9ux5411}{%
\section{6.4 被动观测不能唯一确定因果方向}\label{ux88abux52a8ux89c2ux6d4bux4e0dux80fdux552fux4e00ux786eux5b9aux56e0ux679cux65b9ux5411}}

取 \(0<|a|<1\)。模型一： \(X\sim N(0,1),\ Y=aX+\epsilon\)， \(\epsilon\sim N(0,1-a^2)\) 独立。 模型二交换 \(X,Y\) 的生成方向。两模型的联合高斯律相同，但 \[
E_1[Y\mid do(X=x)]=ax,\qquad
E_2[Y\mid do(X=x)]=0.
\] 任意数量的同类被动样本不能区分它们；选择性干预可以。这是对象因果边界的识别限制。

\hypertarget{ux6700ux5c0f-markov-ux8fb9ux754cux7684ux552fux4e00ux6027ux6761ux4ef6}{%
\section{6.5 最小 Markov 边界的唯一性条件}\label{ux6700ux5c0f-markov-ux8fb9ux754cux7684ux552fux4e00ux6027ux6761ux4ef6}}

相对于有限变量集 \(W\)，\(B\subseteq W\) 是目标 \(Y\) 的 blanket，若 \(Y\perp W\setminus B\mid B\)；极小 blanket 称边界。

\textbf{P-INT-03。} 若条件独立满足半图公理与 intersection，极小边界唯一。

\textbf{证明。} 假设 \(B,C\) 都是边界。记 \(A=B\setminus C\)，\(D=C\setminus B\)，\(F=B\cap C\)，\(R=W\setminus(B\cup C)\)。blanket 性给 \(Y\perp D\mid A,F\) 与 \(Y\perp A\mid D,F\)。 intersection 得 \(Y\perp(A,D)\mid F\)；再与 \(Y\perp R\mid A,D,F\) 用 contraction，得 \(F\) 也是 blanket。极小性迫使 \(F=B=C\)。\(\blacksquare\)

有限离散全产品支撑上的严格正联合分布保证 intersection：对固定 \(f\)，条件核同时只依赖 \(a\) 和只依赖 \(d\)，完整正支撑使其只能依赖 \(f\)。冗余副本 \(B=A\) 不满足此条件，可产生多个极小边界。观测边界唯一仍不等于因果父节点集合已识别。

\hypertarget{ux5206ux79bbux8868ux793aux4e0eux5f62ux6210binding-ux4e0dux80fdux4e92ux76f8ux66ffux4ee3}{%
\section{6.6 分离表示与形成---binding 不能互相替代}\label{ux5206ux79bbux8868ux793aux4e0eux5f62ux6210binding-ux4e0dux80fdux4e92ux76f8ux66ffux4ee3}}

P-INT-02 的两个输入商只证明响应核有良定义的双侧因子化。形成某个父对象还需要父未来、干预、载体与历史条件。C4 将 FormationRelation 与 \nolinkurl{BindingTransportCompatible} 分开，证明形成误差经 \(L\)-Lipschitz binding 放大并加上传输误差的 \(L\varepsilon_F+\varepsilon_B\) 界。

形成可交换不蕴含 binding 可交换，源码已有有限反例。因而``内外接口已闭合→父对象已生成''的叙述必须补齐这些独立条件。第 34.4 节给出完整组合命题及证明。

\hypertarget{ux52a8ux6001ux95edux5408ux4e00ux822cux8fc7ux7a0bux4e0eux591aux5c3aux5ea6ux4e00ux81f4ux6027}{%
\chapter{7. 动态闭合、一般过程与多尺度一致性}\label{ux52a8ux6001ux95edux5408ux4e00ux822cux8fc7ux7a0bux4e0eux591aux5c3aux5ea6ux4e00ux81f4ux6027}}

\hypertarget{ux7cbeux786e-markov-ux5546}{%
\section{7.1 精确 Markov 商}\label{ux7cbeux786e-markov-ux5546}}

\textbf{P-DYN-01：strong lumpability。} 设 \(X_t\) 为 Markov 链，\(f:\mathcal X\to\mathcal M\) 为可测满射，并已存在可测候选宏观核 \(\bar P\)。以下等价：

\begin{enumerate}
\def\labelenumi{\arabic{enumi}.}
\tightlist
\item
  对任意微观初始分布，\(f(X_t)\) 以同一个 \(\bar P\) 为 Markov 转移核。
\item
  对所有 \(x,B\)， \[
  P(x,f^{-1}B)=\bar P(f(x),B).
  \]
\end{enumerate}

\textbf{证明。} 从 2 出发，先对完整微观历史条件化，再对宏观历史取条件期望，得到宏观 Markov 性。从 1 出发，取微观初态 \(\delta_x\) 的第一步概率即得 2。\(\blacksquare\)

在有限状态空间，纤维内投影行相同便可直接定义 \(\bar P\)。一般空间仅有``纤维常数''时还需可测因子化条件。受控版本要求每个允许动作成立；单一被动策略下的闭合不够。

\hypertarget{ux8fdeux7eedux65f6ux95f4ux751fux6210ux5143ux5224ux636e}{%
\section{7.2 连续时间生成元判据}\label{ux8fdeux7eedux65f6ux95f4ux751fux6210ux5143ux5224ux636e}}

\textbf{P-DYN-02。} 有限 CTMC 生成元 \(L\) 与分块 \(B_1,\ldots,B_m\) 满足精确分块商，当且仅当 \[
\sum_{y\in B_j}l_{xy}
\] 对同一 \(B_i\) 中所有 \(x\) 不变。公共值构成宏观生成元 \(\bar L\)。

\textbf{证明。} 若半群商成立，对 \(e^{tL}\) 在零点求导得到该条件。反向，令 \(F\) 为分块指示矩阵，则条件为 \(LF=F\bar L\)；归纳有 \(L^nF=F\bar L^n\)，故 \(e^{tL}F=Fe^{t\bar L}\)。\(\blacksquare\)

\hypertarget{ux8fd1ux4f3cux6838ux4e0eux8defux5f84ux754c}{%
\section{7.3 近似核与路径界}\label{ux8fd1ux4f3cux6838ux4e0eux8defux5f84ux754c}}

若纤维内直径 \[
\epsilon(f)=\sup_{f(x)=f(x'),a}
D_{\rm TV}(f_\#P_{x,a},f_\#P_{x',a})
\] 小，且有可测截面 \(s(m)\)，取 \(\bar P_{m,a}=f_\#P_{s(m),a}\)，得到统一投影误差不超过 \(\epsilon(f)\)。有限空间总有截面；一般空间不能无条件假定。

\textbf{P-DYN-03：同策略路径误差。} 两套有限字母因果记录过程在相同已发生历史上的下一步核差不超过 \(\epsilon_t\in[0,1]\)，使用同一因果策略且初始记录相同，则 \[
D_{\rm TV}(P_{0:T},\widehat P_{0:T})
\le1-\prod_{t=0}^{T-1}(1-\epsilon_t)
\le\sum_{t=0}^{T-1}\epsilon_t.
\]

\textbf{证明。} 只要历史尚未分歧，动作可共同抽样，下一记录作最大耦合，条件相同概率至少 \(1-\epsilon_t\)。归纳得整条路径相同概率的乘积下界，耦合不等式给结论。\(\blacksquare\)

初始记录律已有误差 \(\epsilon_0^{\rm init}\) 时，可加上该误差。边缘时刻分布界是路径界的推送推论。任意单个滞后构造的近似核不自动组成半群。

\hypertarget{ux540cux4e00ux5e95ux8fc7ux7a0bux7684ux8de8ux5c3aux5ea6ux5546}{%
\section{7.4 同一底过程的跨尺度商}\label{ux540cux4e00ux5e95ux8fc7ux7a0bux7684ux8de8ux5c3aux5ea6ux5546}}

设 \(f_r=c_{sr}\circ f_s\)，两尺度对同一底核的误差分别为 \(\epsilon_r,\epsilon_s\)。

\textbf{P-DYN-04。} 在可达像 \(f_s(\mathcal X)\) 上， \[
\sup_{m,a}D_{\rm TV}
\bigl((c_{sr})_\#P_s(m,a),P_r(c_{sr}m,a)\bigr)
\le\epsilon_s+\epsilon_r.
\] 精确时得到 intertwining。

\textbf{证明。} 对 \(m=f_s(x)\)，插入共同底过程推送 \((f_r)_\#P(x,a)\)。第一项由数据处理不超过 \(\epsilon_s\)，第二项不超过 \(\epsilon_r\)。\(\blacksquare\)

若要在全宏观空间取上确界，须满射或控制未达点的扩展。这一结论不由载体集合的合法粗粒化自动推出。

\hypertarget{ux975e-markov-ux591aux65f6ux523bux63a5ux53e3}{%
\section{7.5 非 Markov 多时刻接口}\label{ux975e-markov-ux591aux65f6ux523bux63a5ux53e3}}

用 \(\mathbf T[\pi]\) 表示协议树 \(\pi\) 的完整记录律。若粗协议通过 \(J\) 实施于精细层，读出通过 \(C\) 转换，定义 \[
(\mathcal R\mathbf T)[\pi]=C_\#\mathbf T[J\pi].
\] 协议 lift 与读出推送具有结合律，因此两个逐级接口可复合。它只比较允许实施的协议树；不要求过程有有限记忆。

若精细层还有粗层无法表达的动作，过程商并不保证最优价值相同。价值保持必须另加策略完备性和奖励对应，见第 20、28 章。

\hypertarget{ux4eceux9010ux70b9ux8befux5deeux5230ux5b9eux9645ux5360ux636eux52a0ux6743ux8befux5dee}{%
\section{7.6 从逐点误差到实际占据加权误差}\label{ux4eceux9010ux70b9ux8befux5deeux5230ux5b9eux9645ux5360ux636eux52a0ux6743ux8befux5dee}}

统一的 \(\epsilon_t\) 可能很保守。若某行 \((x,a)\) 的模型误差为 \(e_t(x,a)\)，实际前缀与同一策略诱导占据律 \(\nu_t\)，则第 36 章从真实 PMF 延长证明

\[D_{\rm TV}(\text{实际粗路径},\text{参考路径})\le\sum_{t<T}E_{\nu_t}e_t.\]

若只在认证行集 \(C_t\) 上有 \(e_t\le\epsilon_t\in[0,1]\)，则

\[E_{\nu_t}e_t\le\epsilon_t+(1-\epsilon_t)\nu_t(C_t^c).\]

这不是用平均误差偷换统一误差：平均所用的 \(\nu_t\) 就是比较中的真实过程访问律，证明保留了它的来源。训练行分布不相同的情况下仍需 P-STAT-09 或新的覆盖结论。

\hypertarget{ux6301ux7eedux6027ux4e0dux53d8ux7ed3ux6784ux53efux884cux6027ux4e0eux76f8ux5bf9ux8eabux4efd}{%
\chapter{8. 持续性：不变结构、可行性与相对身份}\label{ux6301ux7eedux6027ux4e0dux53d8ux7ed3ux6784ux53efux884cux6027ux4e0eux76f8ux5bf9ux8eabux4efd}}

\hypertarget{ux8eabux4efdux5206ux652fux4e0eux6301ux7eedux57df}{%
\section{8.1 身份分支与持续域}\label{ux8eabux4efdux5206ux652fux4e0eux6301ux7eedux57df}}

可追踪的历史分支不必预先等同于某个固定状态集合。需要研究``在多久内仍符合给定结构要求''时，才明确选取评估域 \(\mathcal V\)、退出时间 \[
\tau_{\mathcal V}=\inf\{t\ge0:X_t\notin\mathcal V\}.
\] 动态身份也可使用增广状态 \((X_t,\Theta_t)\) 和其中的固定域，或一族明确的移动域。选择持续域属于模型与任务规定；最大化持续概率则是进一步的评价桥接。

\hypertarget{ux786eux5b9aux6027ux6301ux7eedux7684ux5fc5ux8981ux4e0dux53d8ux6838}{%
\section{8.2 确定性持续的必要不变核}\label{ux786eux5b9aux6027ux6301ux7eedux7684ux5fc5ux8981ux4e0dux53d8ux6838}}

\textbf{P-PER-01。} 连续半流 \(\phi_t\) 的轨道长期位于闭集 \(\mathcal V\)，且轨道预紧，则 \[
\omega(x)=\bigcap_{T>0}\overline{\{\phi_t(x):t\ge T\}}
\] 非空、紧、包含于 \(\mathcal V\)，并满足 \(\phi_s(\omega(x))=\omega(x)\)。

\textbf{证明。} 轨道闭包紧，递减的非空紧尾集有非空交。轨道位于闭域，交也位于该域。若 \(\phi_{t_n}(x)\to y\)，连续性给 \(\phi_s(y)\in\omega(x)\)。反向，取 \(t_n>s\)，从预紧序列 \(\phi_{t_n-s}(x)\) 抽收敛子列，其极限 \(z\in\omega(x)\) 满足 \(\phi_s(z)=y\)。\(\blacksquare\)

这不是邻域吸引性定理。有限时段持续也不推出紧不变核：平移 \(x_t=x_0+t\) 可在任意指定有限时间窗留在足够大的区间内，但无非空紧不变集。

\hypertarget{ux5bf9ux79f0ux7ea6ux5316}{%
\section{8.3 对称约化}\label{ux5bf9ux79f0ux7ea6ux5316}}

若群 \(G\) 的作用与动力学等变， \[
\phi_t(gx)=g\phi_t(x),
\] 并且轨道商是良好的 Hausdorff 可度量空间，则商动力学 \(\bar\phi_t([x])=[\phi_t(x)]\) 良定义。若商轨道预紧，可在商上用 P-PER-01。这样移动形态可保持相对身份。

没有等变性时，更换代表元可能改变商上的未来。不能只因``平移看起来无关''就擅自取商。

\hypertarget{ux968fux673aux6301ux7eedux4e0eux5360ux636eux6838}{%
\section{8.4 随机持续与占据核}\label{ux968fux673aux6301ux7eedux4e0eux5360ux636eux6838}}

\textbf{P-PER-02。} 设 \(P_t\) 为 Polish 空间上将 \(C_b\) 映入 \(C_b\) 的 Feller 半群。若 \[
\bar\mu_T=\frac1T\int_0^T\mu_0P_t\,dt
\] 在 \(T\ge1\) 上紧，则沿 \(T_n\to\infty\) 存在弱收敛子列；任一这样的极限 \(\nu\) 不变。若各时刻分布均支撑在闭集 \(\mathcal V\)，则 \(\nu(\mathcal V)=1\)。

\textbf{证明。} Prokhorov 给弱收敛子列。对 \(f\in C_b\)， \[
|(\bar\mu_TP_s-\bar\mu_T)f|
\le2s\|f\|_\infty/T.
\] Feller 性保证弱极限也可作用于 \(P_sf\)，故 \(\nu P_sf=\nu f\)。闭支撑由 Portmanteau 保持。\(\blacksquare\)

这里平均的是时刻分布，不是声称每条样本路径经验频率都收敛；后者需要遍历条件。

\hypertarget{ux6709ux9650ux72b6ux6001ux7684ux53efux884cux6838}{%
\section{8.5 有限状态的可行核}\label{ux6709ux9650ux72b6ux6001ux7684ux53efux884cux6838}}

给有限状态与动作、允许持续域 \(K_0=\mathcal V\)，递推 \[
K_{n+1}=\{x\in K_n:\exists a,\ P(K_n\mid x,a)=1\}.
\]

\textbf{P-PER-03。} 递推在有限步稳定于 \(K_\infty\)。它恰为存在策略使系统以概率一永不退出 \(\mathcal V\) 的状态集，并存在保持它的平稳确定性策略。

\textbf{证明。} 集合单调递减且有限。稳定后每点存在把下一状态留在同一集的动作，逐点选取即可归纳保证永不退出。反向，若某状态能永不退出，则第一步必须以概率一进入同样有此能力的状态；归纳它不可能在任何删除步骤被剔除。\(\blacksquare\)

该结论对有限模型的概率一持续成立。期望寿命长、有限时域高生存概率和此可行核是不同要求。

\hypertarget{ux8fdeux7eedux72b6ux6001ux7684ux5207ux5411ux6761ux4ef6}{%
\section{8.6 连续状态的切向条件}\label{ux8fdeux7eedux72b6ux6001ux7684ux5207ux5411ux6761ux4ef6}}

对闭域 \(\mathcal V\subseteq\mathbb R^d\)，定义 contingent cone \[
T_{\mathcal V}(x)=
\left\{v:\liminf_{h\downarrow0}\frac{d(x+hv,\mathcal V)}h=0\right\}.
\]

\textbf{P-PER-04：切向必要性。} 若一条在零点可微的可行轨迹从 \(x\) 出发并留在 \(\mathcal V\)，则 \(\dot x(0)\in T_{\mathcal V}(x)\)。

\textbf{证明。} \(x(h)=x+h\dot x(0)+o(h)\in\mathcal V\)，故距离比不超过 \(o(h)/h\)。\(\blacksquare\)

\textbf{K-VIA-01。} 对满足 Marchaud 条件的微分包含（上半连续、非空紧凸值、线性增长），闭集的可行性由 \(F(x)\cap T_{\mathcal V}(x)\ne\varnothing\) 的相应 viability theorem 刻画。对局部 Lipschitz、非爆炸单值 ODE，可采用 Nagumo 正不变性判据。原始理论中``约束在边界维持身份''的思想在此有精确接口；可行性与对所有控制都不退出须分开。

\hypertarget{ux6709ux9650ux6e90ux7684ux5b9eux9645ux8defux5f84ux4e0eux90e8ux5206ux89c2ux5bdfux7684ux6269ux5c55}{%
\section{8.7 有限源的实际路径与部分观察的扩展}\label{ux6709ux9650ux6e90ux7684ux5b9eux9645ux8defux5f84ux4e0eux90e8ux5206ux89c2ux5bdfux7684ux6269ux5c55}}

P-PER-03 的现有源 ViabilitySource.p\_per\_03 已构造真正的无限轨迹测度：对支撑在稳定核的初始 PMF，有平稳确定策略使所有时刻留在核中的事件概率为一。有限动作非空、源 PMF 和安全集是显式输入；稳定核可能为空。这里没有把``所有时刻安全''作为结论字段中的假设。

部分观察时，动作只能依赖粗历史，不可以对未观测微观状态分别选择。第 37 章因此在真实后验支持上要求一个共同认证动作，并对它的全部正概率后继支持作后向递推，证明有限时域策略存在的充要条件。它与本章完全观察、平稳无限时域核是不同定理，不能把有限支持结果直接取 \(T\to\infty\)。

\hypertarget{ux6062ux590dux6270ux52a8ux4e0eux89c4ux8303ux6062ux590dux52bf}{%
\chapter{9. 恢复、扰动与规范恢复势}\label{ux6062ux590dux6270ux52a8ux4e0eux89c4ux8303ux6062ux590dux52bf}}

\hypertarget{ux79bbux6563ux6062ux590dux7ba1ux9053}{%
\section{9.1 离散恢复管道}\label{ux79bbux6563ux6062ux590dux7ba1ux9053}}

\textbf{P-REC-01。} 设 \(R_n\ge0\)，\(ER_0<\infty\)，且 \[
E[R_{n+1}\mid\mathcal F_n]\le\kappa R_n+\eta,
\quad0\le\kappa<1,\quad\eta\ge0.
\] 则 \[
ER_n\le\kappa^nER_0+\eta\frac{1-\kappa^n}{1-\kappa},
\] 并对 \(a>0\) 有 \[
P(R_n\ge a)\le
\frac{\kappa^nER_0+\eta(1-\kappa^n)/(1-\kappa)}a.
\]

\textbf{证明。} 取全期望得到标量递推，归纳展开几何级数。尾界由 Markov 不等式。\(\blacksquare\)

持续噪声产生误差底 \(\eta/(1-\kappa)\)。例如 \(x_{n+1}=(x_n+1)/2,\ x_0=0\)，有 \(x_n=1-2^{-n}\)，不趋零。单次恢复相对``刚被扰动的点''收缩，不足以忽略下一次新增扰动。

\hypertarget{ux8fdeux7eedux968fux673aux6062ux590d}{%
\section{9.2 连续随机恢复}\label{ux8fdeux7eedux968fux673aux6062ux590d}}

\textbf{P-REC-02。} 设 \(W\ge c\,d(\cdot,\mathcal V)^2\)，属于生成元 \(\mathcal L\) 的适当定义域，并具备 Dynkin 公式、局部化极限与可积条件，使 \(m(t)=EW(X_t)\) 局部绝对连续。若 \[
\mathcal LW\le-aW+b,\qquad a,c>0,\ b\ge0,
\] 则 \[
EW(X_t)\le e^{-at}W(x)+\frac ba(1-e^{-at}),
\] 且除以 \(c\) 得均方距离界。

\textbf{证明。} 几乎处处 \(m'(t)=E\mathcal LW(X_t)\le-am(t)+b\)。乘 \(e^{at}\) 后积分即得。\(\blacksquare\)

\(\mathcal LW\) 的形式表达不自动保证可用 Dynkin；爆炸、不可积随机积分或定义域不符时必须先修复模型。

\hypertarget{ux6062ux590dux65f6ux95f4ux751fux6210ux89c4ux8303ux52bf}{%
\section{9.3 恢复时间生成规范势}\label{ux6062ux590dux65f6ux95f4ux751fux6210ux89c4ux8303ux52bf}}

令 \(\tau_A=\inf\{n\ge0:X_n\in A\}\)， \(V_A(x)=E_x\tau_A\in[0,\infty]\)，且在 \(A\) 上取零。

\textbf{P-REC-03。} 对 \(x\notin A\) 且 \(V_A(x)<\infty\)， \[
PV_A(x)-V_A(x)=-1.
\]

\textbf{证明。} 第一步分析与 Markov 性给 \(V_A(x)=1+E_xV_A(X_1)\)。有限性也同时保证右侧可积。\(\blacksquare\)

\textbf{P-REC-04：漂移推出恢复时间。} 若有限非负函数 \(V\) 在 \(A^c\) 满足 \(PV-V\le-c<0\)，则 \[
E_x\tau_A\le V(x)/c.
\]

\textbf{证明。} 对 \(\tau_A\wedge N\) 逐步条件期望， \[
E V(X_{\tau_A\wedge N})+cE(\tau_A\wedge N)\le V(x).
\] 舍弃非负第一项，再用单调收敛。\(\blacksquare\)

连续时间可在扩展生成元与停时 Dynkin 条件下写为 \(\mathcal LV_A=-1\)。有限恢复不要求严格距离收缩：\(x\mapsto\max(x-1,0)\) 在非负整数上必恢复，但 Lipschitz 常数为 1。

\hypertarget{ux3c9-ux7684ux673aux5236ux5256ux9762}{%
\section{9.4 Ω 的机制剖面}\label{ux3c9-ux7684ux673aux5236ux5256ux9762}}

至少区分恢复性、不变性、周期性、拓扑保护、准平稳性、内生控制可行性。\(x'=-x\) 在可缩空间上有稳定点，说明非零一阶同调不是普遍必要条件；中性圆周旋转不收缩却持续，说明严格收缩也不是必要条件。

组织闭合可以报告关闭过程后的退出风险差、寿命差或 hazard 比。比值须有正且有限的分母，hazard 须实际存在；有原子退出分布时应使用区间失败概率。一个机制图中的有向环不能独自保证稳定，环上增益也可能导致发散。

\hypertarget{ux6062ux590dux8fc7ux7a0bux6062ux590dux7a0bux5e8fux4e0eux540cux4e00ux6e90ux6838}{%
\section{9.5 恢复过程、恢复程序与同一源核}\label{ux6062ux590dux8fc7ux7a0bux6062ux590dux7a0bux5e8fux4e0eux540cux4e00ux6e90ux6838}}

恢复矩或返回时间的估计不自动提供可执行修复程序。实际 \nolinkurl{RepairProgramSafeControlBridge} 另给可信修复底座、对同一源转移核有效的程序编码，复用程序恢复定理，证明恢复后程序实施的核恰好等于已经认证的安全源动作核。

因此可先用生存固定点证明存在安全动作，再用程序桥证明该动作由恢复的程序实现。程序来源、解释器与编解码正确性仍必须说明；它们不能由 P-REC 的收缩或漂移不等式生成。连续分支的 RecoveryDynkin 也显式消费过程特定的 \nolinkurl{DynkinExpectationCertificate}，然后证明标量微分不等式的后果。

\hypertarget{ux51c6ux5e73ux7a33ux8c31ux4e0eux6301ux7eedux65f6ux95f4ux7a97}{%
\chapter{10. 准平稳、谱与持续时间窗}\label{ux51c6ux5e73ux7a33ux8c31ux4e0eux6301ux7eedux65f6ux95f4ux7a97}}

\hypertarget{killed-ux8fc7ux7a0bux4e0e-qsd}{%
\section{10.1 killed 过程与 QSD}\label{killed-ux8fc7ux7a0bux4e0e-qsd}}

在评估域 \(\mathcal V\) 上，定义次概率半群 \[
P_t^{\mathcal V}f(x)=E_x[f(X_t)\mathbf1_{\{t<\tau_{\mathcal V}\}}].
\] 概率 \(q\) 为 QSD，若存活条件下分布始终为 \(q\)。

\textbf{P-QSD-01：条件稳定化的逆向结论。} 若 \[
\mu_t^c=\frac{\mu P_t^{\mathcal V}}{\mu P_t^{\mathcal V}\mathbf1}
\longrightarrow q
\quad\text{in TV},
\] 对每个有限 \(s>0\) 有 \(qP_s^{\mathcal V}\mathbf1>0\)，且其生存函数在零点右连续，则 \(q\) 是 QSD，并有某 \(\lambda\ge0\) 满足 \[
qP_s^{\mathcal V}=e^{-\lambda s}q.
\]

\textbf{证明。} 半群性质给 \[
\mu_{t+s}^c(f)=
\frac{\mu_t^c(P_s^{\mathcal V}f)}
{\mu_t^c(P_s^{\mathcal V}\mathbf1)}.
\] TV 收敛允许对有界函数取极限，正分母得到 QSD 恒等式。令 \(c(s)=qP_s^{\mathcal V}\mathbf1\)，半群性质给 \(c(s+t)=c(s)c(t)\)。正值、右连续及 \(c(0)=1\) 使 \(c(s)=e^{-\lambda s}\)；次概率性使 \(\lambda\ge0\)。\(\blacksquare\)

最终必吸收时 \(\lambda>0\)，无吸收退化情形允许零。

\hypertarget{ux6709ux9650-perron-ux8868ux793a}{%
\section{10.2 有限 Perron 表示}\label{ux6709ux9650-perron-ux8868ux793a}}

\textbf{K-PF-01。} 对本原非负有限矩阵 \(Q\)，Perron--Frobenius 定理给唯一正主特征值 \(\rho>0\)、正左右特征向量及主模渐近分解。本原表示某个矩阵幂逐项严格正。

当 \(Q\) 为 killed 次随机核，\(\rho\le1\)。归一化 \[
qQ=\rho q,\quad Qh=\rho h,\quad q\mathbf1=1,\quad qh=1.
\]

\textbf{P-QSD-02。} \(q\) 是 QSD；Doob 变换 \[
Q^\uparrow_{ij}=\frac{Q_{ij}h_j}{\rho h_i}
\] 是随机核，且其不变概率为 \(\varpi_i=q_i h_i\)。

\textbf{证明。} \(qQ^n=\rho^nq\)，除以生存质量 \(\rho^n\) 即得 QSD。行和由 \(Qh=\rho h\) 为一；再计算 \[
\sum_iq_i h_i Q^\uparrow_{ij}
=\frac{h_j}{\rho}(qQ)_j=q_jh_j.
\] 总质量 \(qh=1\)。\(\blacksquare\)

从初态 \(\mu\) 出发，\(\mu Q^n\sim\rho^n(\mu h)q\)。\(q\) 与 \(\varpi\) 分别描述``活到此刻''与``条件永续过程''，一般不同。连续 CTMC 用 killed 生成元 \(K\)、\(qK=-\lambda q\)、\(Kh=-\lambda h\)，其变换为 \[
K^\uparrow=\operatorname{diag}(h)^{-1}K\operatorname{diag}(h)+\lambda I.
\]

\hypertarget{ux540cux521dux6001ux7684ux6301ux7eedux65f6ux95f4ux7a97}{%
\section{10.3 同初态的持续时间窗}\label{ux540cux521dux6001ux7684ux6301ux7eedux65f6ux95f4ux7a97}}

\textbf{P-QSD-03。} 对同一初态 \(\mu\)，假设在 \(t\ge t_0\) 上 \[
\|\mu_t^c-q\|_{\rm TV}\le Ce^{-\gamma t},\qquad
P_\mu(\tau>t)\ge ce^{-\lambda t},
\] 其中 \(C,c,\gamma>0,\lambda>0\)。给定 \(\epsilon>0\)、\(0<p\le c\)，若区间 \[
\left[\max\{t_0,0,\log(C/\epsilon)/\gamma\},
\ \log(c/p)/\lambda\right]
\] 非空，则其中同时满足条件稳定误差不超过 \(\epsilon\) 和生存率至少 \(p\)。

\textbf{证明。} 分别解两个指数不等式，再取交集。闭区间端点相等仍有一个时刻；要正长度须严格不等。\(\blacksquare\)

\(\lambda=0\) 时生存下界没有此有限上界；\(C=0\) 表示条件分布已精确稳定，混合下界直接取 \(t_0\)。不可把一个初态的混合界和另一初态的生存界拼接。

\hypertarget{ux53efux9006ux6269ux6563ux7684ux8c31ux7248ux672c}{%
\section{10.4 可逆扩散的谱版本}\label{ux53efux9006ux6269ux6563ux7684ux8c31ux7248ux672c}}

假设 killed 生成元 \(L_D\) 在有限测度 \(m(D)<\infty\) 的 \(L^2(m)\) 中自伴、紧 resolvent，正首特征函数 \(\phi_1\)， \[
-L_D\phi_n=\lambda_n\phi_n,\qquad
0<\lambda_1<\lambda_2\le\cdots.
\] 初始密度 \(g\in L^2(m)\)，\(a_1=\langle g,\phi_1\rangle>0\)。

\textbf{P-QSD-04。} 存在 \(C_g,c_g>0,t_0\)，使 \[
q(dx)=\frac{\phi_1(x)m(dx)}{\int\phi_1\,dm},\qquad
\|\mu_t^c-q\|_{\rm TV}\le C_ge^{-(\lambda_2-\lambda_1)t},
\] \[
P_g(\tau>t)\ge c_ge^{-\lambda_1t}\quad(t\ge t_0).
\]

\textbf{证明。} 谱展开为 \(P_t^Dg=e^{-\lambda_1t}a_1\phi_1+R_t\)，且 \(\|R_t\|_2\le e^{-\lambda_2t}\|g-a_1\phi_1\|_2\)。 Cauchy--Schwarz 将其转成 \(L^1\) 界。生存质量的主项系数 \(a_1\int\phi_1dm>0\)，充分大 \(t\) 时余项小于其一半。分子、分母同除 \(e^{-\lambda_1t}\)，比较条件密度与 \(q\)，得到所述速率。\(\blacksquare\)

一般非自伴核存在非正交模或 Jordan 项，不能未经控制直接套用同一个精确谱间隔常数。一般 QSD 存在与收敛条件可调用 Champagnat--Villemonais 的原始研究，本文只将其作为条件模块使用。\href{https://arxiv.org/abs/1404.1349}{原始论文}

\hypertarget{ux6b63ux7279ux5f81ux6570ux636eux4e0eux8c31ux5b58ux5728ux6027ux5206ux5c42}{%
\section{10.5 正特征数据与谱存在性分层}\label{ux6b63ux7279ux5f81ux6570ux636eux4e0eux8c31ux5b58ux5728ux6027ux5206ux5c42}}

M-PE-01 重用本章的归一化代数：给定正的左右特征数据，就能证明 QSD、Doob 随机核、平稳权重及随机读出的鞅性质。\nolinkurl{PositiveEigenstructure.qsdDoob_bundle} 不负责替任意矩阵产生 Perron 数据；P-EVO 的严格主谱极限仍消费 KPF01Certificate。

这解释为什么``统一正特征结构''与``所有长期结果被一个谱定理生成''不同。P-QSD-01 的条件极限、P-QSD-03 的同初态窗口、P-QSD-04 的可逆谱分析以及一般 GOA 稳定性仍各有适配。模型改变了支撑或失去谱隙时，应回到这些独立条件。

\hypertarget{ux5386ux53f2ux8eabux4efdux8fd0ux8f93ux4e0eux53d1ux5c55ux5b9aux5f8b}{%
\chapter{11. 历史身份、运输与发展定律}\label{ux5386ux53f2ux8eabux4efdux8fd0ux8f93ux4e0eux53d1ux5c55ux5b9aux5f8b}}

\hypertarget{ux6709ux5411ux5ef6ux7eed}{%
\section{11.1 有向延续}\label{ux6709ux5411ux5ef6ux7eed}}

\textbf{D-ID。} 一个对象出现 \(O_t^\alpha\) 是带时间和证书的实例。延续边 \[
\Gamma:O_t^\alpha\rightsquigarrow O_{t'}^\beta,\quad t<t'
\] 记录祖先关系、结构运输与可追踪机制。按时间复合得到有向历史路径。

在分裂 \(A\to B,A\to C\) 后，\(B,C\) 不因共享祖先而成为同时存在的同一实例。无向连通分量只表示谱系关联。融合也不能倒推出原有多个对象一直是同一个。

\hypertarget{ux51bbux7ed3ux6a21ux578bux7684-tcic}{%
\section{11.2 冻结模型的 TCIC}\label{ux51bbux7ed3ux6a21ux578bux7684-tcic}}

对不同时刻的冻结结构模型，用端核 \(K_t:\mathcal M_t\to\mathcal P(\mathcal M_t)\)； 运输 \(\Gamma_{s,t}:\mathcal M_s\to\mathcal M_t\) 满足 \(\Gamma_{r,t}=\Gamma_{s,t}\circ\Gamma_{r,s}\)。 精确交换为 \[
(\Gamma_{s,t})_\#K_s(m,\cdot)=K_t(\Gamma_{s,t}m,\cdot).
\]

\textbf{P-ID-01：运输缺损累计。} 若相邻端核交换误差至多 \(\epsilon_t\)，则 \[
\sup_mD_{\rm TV}\bigl((\Gamma_{0,n})_\#K_0(m),
K_n(\Gamma_{0,n}m)\bigr)\le\sum_{t<n}\epsilon_t.
\]

\textbf{证明。} 在第 \(n+1\) 步插入 \((\Gamma_{n,n+1})_\#K_n(\Gamma_{0,n}m)\)； 旧误差经推送不增加，新误差至多 \(\epsilon_n\)。归纳即得。\(\blacksquare\)

这些核是冻结模型的端核，不是真实的一步时间演化。若 \(P_t:\mathcal M_t\to\mathcal P(\mathcal M_{t+1})\)，两套实际历程的比较应写 \[
(\Gamma_{t+1})_\#P_t(m,\cdot)
\approx\widehat P_t(\Gamma_tm,\cdot).
\] 下一状态必须用下一时刻的运输映射。

\hypertarget{ux53d1ux5c55ux8fd0ux8f93-dtrc}{%
\section{11.3 发展运输 DTRC}\label{ux53d1ux5c55ux8fd0ux8f93-dtrc}}

结构描述子 \(\Theta_t\) 可包含预测因子、转移核、接口、载体族与机制证书。它位于事先声明的度量或伪度量空间；不同分量可保留向量误差。

\textbf{D-DTRC。} 发展模型 \[
\widehat\Theta_{t+1}=\Psi_t(\Theta_t,E_t,a_t)
\] 必须由机制、独立个体、早期数据或训练环境预先约束。测试残差 \[
r_t=d(\Theta_{t+1},\Psi_t(\Theta_t,E_t,a_t))
\] 在冻结模型下评价，并检查预测到的因果后果，而不只检查描述子拟合。

任意有限轨迹都能被时间查表模型零残差拟合，因此零训练残差不证明发展机制。

\textbf{P-ID-02：发展管道。} 在相同已声明外部输入下，参考轨迹 \(\bar\Theta_{t+1}=\Psi_t(\bar\Theta_t)\)。若实际模型残差至多 \(\epsilon_t\)，且 \(\Psi_t\) 是 \(L_t\)-Lipschitz，则 \[
e_{t+1}\le L_te_t+\epsilon_t,
\] \[
e_n\le\Bigl(\prod_{j<n}L_j\Bigr)e_0+
\sum_{k<n}\epsilon_k\prod_{j=k+1}^{n-1}L_j,
\quad e_t=d(\Theta_t,\bar\Theta_t).
\]

\textbf{证明。} 插入 \(\Psi_t(\Theta_t)\)，用三角不等式和 Lipschitz 条件，再展开递推。\(\blacksquare\)

若 \(L_t\le L<1,\epsilon_t\le\epsilon\)，得到 \(e_n\le L^ne_0+\epsilon(1-L^n)/(1-L)\)。 若实际和参考输入不同，应再加入输入敏感度误差；不能当作同一轨迹模型。

\hypertarget{ux5206ux88c2ux878dux5408ux4e0eux4fddux771fux5ea6}{%
\section{11.4 分裂、融合与保真度}\label{ux5206ux88c2ux878dux5408ux4e0eux4fddux771fux5ea6}}

延续记录保存 \[
\mathbf F_\Gamma=
(F_{\rm pred},F_{\rm causal},F_{\rm carrier},
F_{\rm control},F_{\rm heredity}),
\] 每项都有独立比较任务、归一化与置信范围。未研究的分量为空，不能填为一。

两个分支的剖面 \((1,0)\) 和 \((0,1)\) 没有自然的唯一百分比排序。若两轴同等地位，交换对称性迫使任何对称标量给相同分数；打破平局需要额外优先级。默认使用逐分量偏序。

\hypertarget{ux4fddux6301ux540cux4e00ux5bf9ux8c61ux4e0eux4ea7ux751fux5b50ux4ee3}{%
\section{11.5 保持同一对象与产生子代}\label{ux4fddux6301ux540cux4e00ux5bf9ux8c61ux4e0eux4ea7ux751fux5b50ux4ee3}}

LineageSemantics 中 SameObject 比较身份，OffspringOf 同时要求真实父子关系与不同身份，因此有 offspring\_not\_sameObject 和 \nolinkurl{sameObjectRepair_not_offspring}。结构相似、同祖先和修复同一对象分别是不同关系。

P12 的 \nolinkurl{homeostaticMarginals_preserve_lifecycleParent} 从 carrier 上每个组织对象均满足与同一 parent 的 SameObject 前提，推出各有限时刻被支持边缘的身份保持。这是一条条件身份桥，不是从边缘分布接近自动推出身份相同。第 33 章的缺陷界也不应直接解释成身份的百分比。

\hypertarget{ux5bf9ux8c61ux8bc1ux4e66ux4e0eux56e0ux679cux5b8cux6574ux6027-ux3c9}{%
\chapter{12. 对象证书与因果完整性 Ω}\label{ux5bf9ux8c61ux8bc1ux4e66ux4e0eux56e0ux679cux5b8cux6574ux6027-ux3c9}}

\hypertarget{ux5206ux5c42ux8bb0ux5f55}{%
\section{12.1 分层记录}\label{ux5206ux5c42ux8bb0ux5f55}}

\textbf{D-OBJ。} 局部操作对象证书至少包含：

\begin{enumerate}
\def\labelenumi{\arabic{enumi}.}
\tightlist
\item
  候选嵌入因子、声明的自身未来及环境接口；
\item
  预测闭合与必要的动态/递归闭合；
\item
  允许干预下的因果接口条件及其证据；
\item
  表示协变的运输规则；
\item
  协议、支撑、时域、误差和未决项。
\end{enumerate}

历史对象另含有向延续与 TCIC/DTRC。需要``可局域实现''``恢复''``自维持''``可遗传''或``具有功能目标''的主张时，分别加入相应模块。

\textbf{D-OBJ-STRONG。} 通过预先指定的预测、因果、历史和非平凡性条件的候选，称该研究上下文中的强操作对象。这是明确证书类的定义及其充分判据，不是所有自然对象的无条件充要定理。

非平凡性至少要说明保留了什么可检测区别，以及压缩或机制解释优于何种基线。全历史恒等表示不算自动发现了宏观对象，完全常数状态也不算自动发现了有内容的身份。

\hypertarget{ux7531ux8bc1ux4e66ux5bfcux51fa-ux3c9}{%
\section{12.2 由证书导出 Ω}\label{ux7531ux8bc1ux4e66ux5bfcux51fa-ux3c9}}

先固定基线对象证书与评估时域 \(T\)，再给有限机制集 \(\mathcal M\)。对移除集 \(D\subseteq\mathcal M\)，定义 \[
\chi_{O,T}(D)=1
\] 当且仅当指定移除操作后，分支仍通过\textbf{事先冻结}的证书。若移除有直接副作用，这些副作用属于干预语义，必须一并记录。

若失败单调， \[
D\subseteq D',\ \chi(D)=0\Rightarrow\chi(D')=0,
\] 定义 \[
\mathcal C_{\Omega,T}=\min_{\subseteq}\{D:\chi_{O,T}(D)=0\},
\qquad
\Omega_T(O)=(\mathcal M,\mathcal C_{\Omega,T}).
\]

\textbf{P-OMG-01：极小毁损重建。} \[
\chi_{O,T}(D)=0
\iff \exists C\in\mathcal C_{\Omega,T},\ C\subseteq D.
\] \textbf{证明。} 正向在有限失败子集中取极小元；反向由单调性。\(\blacksquare\)

若删掉抑制器反而恢复功能，失败不单调，应保存完整响应表或类型化机制模型。不得强制套用单一超图。

\hypertarget{ux5b8cux6574ux6027ux88d5ux91cf}{%
\section{12.3 完整性裕量}\label{ux5b8cux6574ux6027ux88d5ux91cf}}

在预先给定模型距离 \(d\) 与非空失败模型集 \(\mathcal F_O\) 下， \[
\Gamma_O(T)=\inf_{F\in\mathcal F_O}d(T,F).
\]

\textbf{P-OMG-02。} \[
|\Gamma_O(T)-\Gamma_O(S)|\le d(T,S).
\] \textbf{证明。} 对每个 \(F\)，\(d(T,F)\le d(T,S)+d(S,F)\)，取下确界，再交换 \(S,T\)。\(\blacksquare\)

若 \(\widehat\Gamma_O>\eta\) 且 \(d(T,\widehat T)\le\eta\)，则真实裕量为正。正裕量说明离该失败集有可控距离；``未落入失败集''在失败集非闭时不一定给正裕量。

\hypertarget{ux8bcaux65adux5256ux9762ux4e0eux539fux59cbux53c2ux6570}{%
\section{12.4 诊断剖面与原始参数}\label{ux8bcaux65adux5256ux9762ux4e0eux539fux59cbux53c2ux6570}}

\(L_c,L_o,\rho,\eta_{\rm pref}\) 保留为领域诊断入口：

\begin{itemize}
\tightlist
\item
  \(L_c\)：可由极小实现数、毁损集结构、有效模型复杂度等向量描述，避免任意单总分。
\item
  \(L_o\)：资源流量、接口带宽、外部支架依赖分别记录，不与身份退出率混用。
\item
  对非零漂移 \(F\) 和最优方向集 \(D^*\)，可用有符号对齐 \[
  A(F,D^*)=\sup_{d\in D^*,\,d\ne0}
  \frac{\langle F,d\rangle}{\|F\|\|d\|}.
  \] 反向为负，正交为零；单用正弦会混淆同向与反向。零漂移或无非零最优方向时记未定义或按任务另订约定。
\item
  驱动/调节能力比须同量纲、正分母，且有独立测量，不是普适心理指数。
\end{itemize}

\hypertarget{ux64cdux4f5cux5bf9ux8c61ux89d2ux8272ux4e0eux81eaux521bux751fux8fb9ux754c}{%
\section{12.5 操作对象角色与自创生边界}\label{ux64cdux4f5cux5bf9ux8c61ux89d2ux8272ux4e0eux81eaux521bux751fux8fb9ux754c}}

OperationalObject 分别记录预测身份、形成出处、相互作用区分、组成持续、物理可恢复性和语义纤维身份。某一项通过不能填补另一项缺失。特别是统计上恢复一个父候选，是关于认识与识别的输入，并非父对象已经被物理构造的证明。

第 34 章将 P0--P12 的终端写为有条件的对象生命周期。P12 最终审计为 PARTIAL / EXPLICIT BOUNDARY：已有组织构造、平均稳态、资源可行性与历史身份链，但输入仍包含程序来源与可信实现。工作包 CLOSED 不能被改写为无前提的生命起源或任意对象自生产定理。

\hypertarget{ux9884ux6d4bux8eabux4efdux7684ux5b9eux73b0ux5197ux4f59ux4e0eux6bc1ux635fux5bf9ux5076}{%
\chapter{13. 预测身份的实现、冗余与毁损对偶}\label{ux9884ux6d4bux8eabux4efdux7684ux5b9eux73b0ux5197ux4f59ux4e0eux6bc1ux635fux5bf9ux5076}}

\hypertarget{ux5355ux8c03ux8bbfux95eeux7cfbux7edf}{%
\section{13.1 单调访问系统}\label{ux5355ux8c03ux8bbfux95eeux7cfbux7edf}}

\textbf{D-CAR。} 给有限载体全集 \(V\)。对 \(S\subseteq V\)，读出 \(R_S(H)\) 满足限制一致性：若 \(S\subseteq T\)，则 \(R_S\) 是 \(R_T\) 的可测函数。

对固定规范身份 \(C\)，定义 \[
\mathcal A_0=\{S:\exists g_S,\ C=g_S(R_S)\},
\qquad
\mathcal M_0=\min_{\subseteq}\mathcal A_0.
\] 等式采用统一点态或同一参考测度下的共同模零语义。

\textbf{P-CAR-01。} \(\mathcal A_0\) 向上闭合，且在有限 \(V\) 上每个充分集包含一个极小充分集。

\textbf{证明。} 大读出能恢复小读出，再接原解码器。对充分集内有限多个充分子集取极小元。\(\blacksquare\)

\(\mathcal M_0\) 是无相互严格包含成员的集合族（clutter）。它可为空（整个装置都不充分），也可为 \(\{\varnothing\}\)（目标常数）。无限载体集可能没有极小元，本定理不覆盖。

\hypertarget{ux8bbfux95eeux64e6ux9664ux4e0e-blocker}{%
\section{13.2 访问擦除与 blocker}\label{ux8bbfux95eeux64e6ux9664ux4e0e-blocker}}

定义 \[
b_V(\mathcal M)=
\min_{\subseteq}\{D\subseteq V:\forall M\in\mathcal M,\ D\cap M\ne\varnothing\}.
\]

\textbf{P-CAR-02。} 在纯访问擦除中，删除 \(D\) 后仍能恢复 \(C\)，当且仅当存在 \(M\in\mathcal M_0\) 与 \(D\) 不交。因此极小访问毁损集为 \(b_V(\mathcal M_0)\)。

\textbf{证明。} 剩余访问集为 \(V\setminus D\)。其充分当且仅当包含极小充分集。否定这一命题恰为击中全部极小充分集。\(\blacksquare\)

\textbf{P-CAR-03：双 blocker。} \[
b_V(b_V(\mathcal M))=\mathcal M.
\]

\textbf{证明。} 任意 \(S\subseteq V\)，若包含某 \(M\in\mathcal M\)，则它击中每个击中 \(\mathcal M\) 的集，因而击中 \(b_V(\mathcal M)\)。若不包含任何 \(M\)，其补集 \(V\setminus S\) 与每个 \(M\) 相交，因而包含一个极小击中集 \(D\)，且 \(S\cap D=\varnothing\)。所以击中 \(b_V(\mathcal M)\) 的集合恰为 \(\mathcal M\) 的上闭包，取极小元得到结论。\(\blacksquare\)

退化约定为 \(b(\varnothing)=\{\varnothing\}\)、\(b(\{\varnothing\})=\varnothing\)，保持双对偶。该有限组合学结果是标准 blocker 理论的应用。

\hypertarget{ux8bbfux95eeux6bc1ux635fux4f55ux65f6ux662fux5b9eux9645ux5931ux6548}{%
\section{13.3 访问毁损何时是实际失效}\label{ux8bbfux95eeux6bc1ux635fux4f55ux65f6ux662fux5b9eux9645ux5931ux6548}}

\textbf{H-CAR-FRAME。} 将访问结论解释为时域内物理失效，至少要求：

\begin{enumerate}
\def\labelenumi{\arabic{enumi}.}
\tightlist
\item
  删除操作与载体集合一一对应；
\item
  剩余读出及其动力学按模型声明保持；
\item
  剩余读出充分性与该时域的目标证书等价；
\item
  修复、补偿、供能、控制和再生机制已纳入模型。
\end{enumerate}

``不允许修复''不能替代第 2 项。两份比特备份 \(a,b\) 中若 \(a\) 同时给 \(b\) 供电，访问毁损须擦除二者，实际删除 \(a\) 却能使二者在随后都失效。

因此存在两个相关但不同的超图：访问实现的 blocker 和真实机制干预的 \(\Omega_T\)。满足 frame 条件时才把二者对应。

\hypertarget{ux7f3aux635fux76f4ux5f84ux4e0eux89e3ux7801ux534aux5f84}{%
\section{13.4 缺损直径与解码半径}\label{ux7f3aux635fux76f4ux5f84ux4e0eux89e3ux7801ux534aux5f84}}

在有限历史和协议表中，定义 \[
e(S)=\max_{R_S(h)=R_S(h')}\max_i
D_{\rm TV}(p_{h,i},p_{h',i}).
\] 则 \(e(S)=0\) 当且仅当精确充分；增加载体使 \(e(S)\) 非增。

定义最坏最优解码半径 \[
r(S)=\inf_{\{q_{z,i}\}}\max_{h,i}
D_{\rm TV}(p_{h,i},q_{R_S(h),i}).
\]

\textbf{P-CAR-04。} \[
\tfrac12e(S)\le r(S)\le e(S).
\] \textbf{证明。} 对任何中心，同纤维两真实分布的距离不超过两个中心误差之和，给下界。每个读出纤维选一个代表历史的响应作为中心，给上界。\(\blacksquare\)

这把``同读出响应直径''和``最优解码误差''精确连接。二者不能直接共用同一数值；后续阈值恢复采用明确的 \(e(S)\)。

\hypertarget{ux4eceux6ce8ux518cux63a5ux53e3ux4fddux6301ux5230ux7269ux7406ux5b9eux73b0}{%
\section{13.5 从注册接口保持到物理实现}\label{ux4eceux6ce8ux518cux63a5ux53e3ux4fddux6301ux5230ux7269ux7406ux5b9eux73b0}}

同源核心能把需要保留的服务或程序属性放入有限模型输出，再通过最粗稳定商保持它们；这样载体读出可以明确说明恢复哪一个已注册目标。这个连接加强了预测---载体主线，却不取消 H-CAR-FRAME。

例如同样的两个冗余读出可由独立供电或共享单一电源实现；其访问 blocker 相同，实际毁损族不同。即便下游策略已被形式化证明维持源模型安全，模型是否遗漏共同电源仍是物理适配问题。对失效阈值的数值比较还需第 33 章的缺陷校准。

\hypertarget{ux5206ux8fa8ux7387ux7c97ux7c92ux5316ux4e0eux5faeux89c2ux5b9eux73b0ux7684ux4e0dux786eux5b9aux533aux95f4}{%
\chapter{14. 分辨率、粗粒化与微观实现的不确定区间}\label{ux5206ux8fa8ux7387ux7c97ux7c92ux5316ux4e0eux5faeux89c2ux5b9eux73b0ux7684ux4e0dux786eux5b9aux533aux95f4}}

\hypertarget{ux4e00ux822cux95edux5305ux7cfbux7edf}{%
\section{14.1 一般闭包系统}\label{ux4e00ux822cux95edux5305ux7cfbux7edf}}

\textbf{D-RES。} 有限载体全集 \(V\) 上的闭包系统 \(\mathcal L\subseteq2^V\) 包含 \(V\)，并对交封闭。定义 \[
\operatorname{cl}_{\mathcal L}(S)=
\bigcap\{A\in\mathcal L:S\subseteq A\}.
\] 粗尺度可表达族较小：\(\mathcal L_r\subseteq\mathcal L_s\)。

对固定向上闭合性质 \(\mathcal A\)，令 \(\mathcal M_r=\min(\mathcal A\cap\mathcal L_r)\)。

\textbf{P-RES-01：极小性质的粗粒化。} \[
\mathcal M_r=
\min_{\subseteq}\{\operatorname{cl}_r(M):M\in\mathcal M_s\}.
\]

\textbf{证明。} 每个粗极小充分集 \(A\) 在精细族中也充分，故包含一个精细极小充分集 \(M\)。其粗闭包仍包含 \(M\)、属于粗族且仍充分，又位于 \(A\) 内；粗极小性使它等于 \(A\)。反向，每个精细极小集的粗闭包充分，因而包含一个粗极小集。两族产生同一上闭包，极小元相同。\(\blacksquare\)

\textbf{P-RES-02：复合律。} 若 \(\mathcal L_r\subseteq\mathcal L_s\subseteq\mathcal L_t\)，则 \[
\operatorname{cl}_r\operatorname{cl}_s=\operatorname{cl}_r.
\] 在极小集合族上定义 \(\mathcal R_{r\leftarrow s}(\mathcal M)= \min\{\operatorname{cl}_r(M):M\in\mathcal M\}\)，则 \[
\mathcal R_{r\leftarrow t}
=\mathcal R_{r\leftarrow s}\mathcal R_{s\leftarrow t}.
\]

\textbf{证明。} \(S\subseteq\operatorname{cl}_sS\subseteq\operatorname{cl}_rS\)；再取粗闭包给等式。中间删去被包含的非极小元不会改变后续闭包生成的上闭包。\(\blacksquare\)

性质、目标、时域和底载体必须固定。改变这些量不是同一次纯分辨率变换。

\hypertarget{boolean-ux5220ux9664ux8bedux4e49}{%
\section{14.2 Boolean 删除语义}\label{boolean-ux5220ux9664ux8bedux4e49}}

一般闭包系统不保证 blocker 对偶。例： \[
\mathcal L=\{\varnothing,\{1\},\{1,2\}\},\quad
\mathcal M=\{\{1,2\}\}.
\] 在 \(\mathcal L\) 内限制击中集时，第一次 blocker 为 \(\{\{1\}\}\)，第二次仍相同，而不回到 \(\mathcal M\)。

因此完整访问删除采用有限原子 Boolean 代数。粗化通过满射 \(\pi:V\twoheadrightarrow W\) 表达，每个粗原子是一整条纤维。定义 \[
\pi_*\mathcal M=\min\{\pi(M):M\in\mathcal M\}.
\]

\textbf{P-RES-03：blocker 的尺度自然性。} \[
b_W(\pi_*\mathcal M)=\pi_*b_V(\mathcal M).
\]

\textbf{证明。} \(T\subseteq W\) 击中 \(\pi_*\mathcal M\) 当且仅当饱和原像 \(\pi^{-1}T\) 击中全部 \(M\)。这又当且仅当 \(\pi^{-1}T\) 包含某精细极小击中集 \(D\)，等价于 \(\pi(D)\subseteq T\)。所以粗击中集族正是精细极小击中集像的上闭包。取极小元即得。\(\blacksquare\)

这里粗删除指删除整个纤维；它不表示粗删除成本等于单个精细原子的成本。

\hypertarget{ux4e09ux4e2aux4f34ux968fux64cdux4f5c}{%
\section{14.3 三个伴随操作}\label{ux4e09ux4e2aux4f34ux968fux64cdux4f5c}}

对精细上闭族 \(\mathcal A_s\) 和粗上闭族 \(\mathcal A_r\)，定义 \[
\operatorname{Res}_\pi(\mathcal A_s)
=\{B\subseteq W:\pi^{-1}B\in\mathcal A_s\},
\] \[
\operatorname{Ext}_\pi^-(\mathcal A_r)
=\{S\subseteq V:\exists B\in\mathcal A_r,\ \pi^{-1}B\subseteq S\},
\] \[
\operatorname{Ext}_\pi^+(\mathcal A_r)
=\{S\subseteq V:\pi(S)\in\mathcal A_r\}.
\]

\textbf{P-RES-04：伴随三元组。} \[
\operatorname{Ext}^-_\pi\dashv\operatorname{Res}_\pi
\dashv\operatorname{Ext}^+_\pi.
\] 具体是 \[
\operatorname{Ext}^-_\pi(\mathcal A_r)\subseteq\mathcal A_s
\iff\mathcal A_r\subseteq\operatorname{Res}_\pi(\mathcal A_s),
\] \[
\operatorname{Res}_\pi(\mathcal A_s)\subseteq\mathcal A_r
\iff\mathcal A_s\subseteq\operatorname{Ext}^+_\pi(\mathcal A_r).
\]

\textbf{证明。} 第一式两边都表示每个粗充分集的完整原像精细充分。第二式正向：若 \(S\in\mathcal A_s\)，则其饱和扩张 \(\pi^{-1}\pi(S)\) 也充分，故 \(\pi(S)\in\mathcal A_r\)。反向取 \(S=\pi^{-1}B\)，利用满射的 \(\pi(\pi^{-1}B)=B\)。\(\blacksquare\)

\hypertarget{ux7cbeux786eux5faeux89c2ux4e0dux786eux5b9aux533aux95f4}{%
\section{14.4 精确微观不确定区间}\label{ux7cbeux786eux5faeux89c2ux4e0dux786eux5b9aux533aux95f4}}

给粗极小族 \(\mathcal D\)，定义 \[
J^-_\pi(\mathcal D)=\{\pi^{-1}D:D\in\mathcal D\},
\] \[
J^+_\pi(\mathcal D)=
\min\{S\subseteq V:\pi(S)\in\uparrow\mathcal D\}.
\]

\textbf{P-RES-05。} \[
\pi_*\mathcal C=\mathcal D
\iff
\uparrow J^-_\pi(\mathcal D)
\subseteq\uparrow\mathcal C
\subseteq\uparrow J^+_\pi(\mathcal D).
\]

\textbf{证明。} 若粗像为 \(\mathcal D\)，每个 \(D\) 是某精细极小边的像，所以完整原像含精细充分边；而每个精细充分边的像含粗极小边，得两包含。反向，左包含对每个 \(D\) 给出 \(F\in\mathcal C\) 且 \(F\subseteq\pi^{-1}D\)；右包含给某 \(D'\in\mathcal D\) 满足 \(D'\subseteq\pi(F)\subseteq D\)。极小族内无严格包含，故 \(D'=D=\pi(F)\)。同时所有精细边的像都包含粗极小边，所以粗极小像恰为 \(\mathcal D\)。\(\blacksquare\)

\textbf{P-RES-06：对象相关的唯一性。} 上述两端相同，当且仅当 \[
|\pi^{-1}(w)|=1
\quad\forall w\in\bigcup\mathcal D.
\]

\textbf{证明。} 活跃纤维均为单点时，满足粗充分性的最小截面就是完整原像，两端一致。若某活跃 \(w\) 有两个原子，取含 \(w\) 的粗极小边 \(D\)，每条纤维只选一个原子组成截面 \(S\)。它属于上端；却不包含 \(\pi^{-1}D\)，也不包含其他粗极小边的完整原像，因为其他边不能严格包含于 \(D\)。所以 \(S\) 不属下端。\(\blacksquare\)

粗尺度数据通常确定一族允许实现，而非唯一微观边界。非活跃纤维不影响这一对象特定结论。

\hypertarget{ux5206ux8fa8ux7387ux7684ux7c7bux578bux4e0eux5c3aux5ea6ux4f20ux8f93}{%
\section{14.5 分辨率的类型与尺度传输}\label{ux5206ux8fa8ux7387ux7684ux7c7bux578bux4e0eux5c3aux5ea6ux4f20ux8f93}}

本章的闭包、Boolean 删除和实现伴随属于 H0-RES。它与 H0-PQ 的预测商、H0-DYN 的核交换、H0-COMP 的父形成及 H0-HIST 的历史运输相关，却不是同一个算子。第 32 章列出八种类型及允许桥。

ObjectScaleCalculus 能复合尺度映射，并在分别给定身份、目的、GOD、GOA、修复、稳态保持证书后运输这些性质。当前 \nolinkurl{WilsonianRGInterpretation} 仅记录 couplingFlow，未提供它与某个实际尺度映射或源动力学的相容关系；其两个常值流不同的定理不能称作给定物理粗粒化的 Wilson 流非唯一性定理。

\hypertarget{ux7ec4ux5408ux91cdux53e0ux4e0eux7236ux5bf9ux8c61ux7684ux751fux6210}{%
\chapter{15. 组合、重叠与父对象的生成}\label{ux7ec4ux5408ux91cdux53e0ux4e0eux7236ux5bf9ux8c61ux7684ux751fux6210}}

\hypertarget{ux4e09ux79cdux5171ux540cux6027}{%
\section{15.1 三种共同性}\label{ux4e09ux79cdux5171ux540cux6027}}

多个系统具有共同外因、共同预测信息、共同目标，均不独自等于形成父对象。父对象必须有自己的规范未来、动态接口和历史证书；``组合''另研究其相对于声明子系统分割是否可约。

对子路径 \(\mathbf X_1,\ldots,\mathbf X_m\) 及外部接口记录 \(U\)，定义 \[
I_\pi=
E_U D_{\rm KL}\left(P_{\mathbf X_{1:m}\mid U}
\middle\|\bigotimes_{B\in\pi}P_{\mathbf X_B\mid U}\right).
\]

\textbf{P-COMP-01。} \(I_\pi=0\) 当且仅当该划分的路径块给定 \(U\) 后独立。有限个非平凡划分中， \[
\Sigma_{\rm path}=\min_\pi I_\pi>0
\] 当且仅当没有这种路径因子化。

\textbf{证明。} 每个条件 KL 非负，平均为零当且仅当几乎处处两个条件律相同。有限族的最小值为正当且仅当每项为正。\(\blacksquare\)

共同驱动例：\(X_1=U+\epsilon_1,\ X_2=U+\epsilon_2\)，噪声条件独立。裸相关可大，而条件整合为零。若 \(U\) 是内部结果的后果，条件化还可能制造 collider 偏差；外部接口的因果地位必须先说明。

\hypertarget{ux5e72ux9884ux4e0dux53efux7ea6ux6027}{%
\section{15.2 干预不可约性}\label{ux5e72ux9884ux4e0dux53efux7ea6ux6027}}

对每个非平凡划分 \(\pi\)，声明保持块内机制而改变跨块连接的可实施切割 \(\mathcal K_\pi\)。比较原记录律 \(P\) 与切割律 \(P^\pi_{\rm cut}\)。

可用总变差，也可用 \[
JS(P,Q)=\tfrac12D_{\rm KL}(P\|M)+\tfrac12D_{\rm KL}(Q\|M),
\quad M=(P+Q)/2.
\]

\textbf{P-COMP-02。} \(0\le JS(P,Q)\le\log2\)，且 \(JS=0\iff P=Q\)。 因此有限切割族的最小 JS 正，当且仅当每项切割都改变声明记录律。

\textbf{证明。} 非负性来自 KL。\(M\ge P/2,Q/2\) 给上界；零值迫使 \(P=M=Q\)。\(\blacksquare\)

检测不到切割效应不等于绝无微观耦合。要作后一解释，需额外的观察完备性和 cut-faithfulness：真实相关耦合被正确切断时必须在允许记录中显现。

\hypertarget{ux7ec4ux5408ux88d5ux91cf}{%
\section{15.3 组合裕量}\label{ux7ec4ux5408ux88d5ux91cf}}

令模型距离 \(d\)，且每个切割操作满足 \(d(\mathcal K_\pi T,\mathcal K_\pi S)\le L_\pi d(T,S)\)。 定义 \(\Gamma_{\rm comp}(T)=\min_\pi d(T,\mathcal K_\pi T)\)。

\textbf{P-COMP-03。} 若 \(L=\max_\pi L_\pi<\infty\)，则 \[
|\Gamma_{\rm comp}(T)-\Gamma_{\rm comp}(S)|
\le(1+L)d(T,S).
\]

\textbf{证明。} 每项距离差至多 \(d(T,S)+d(\mathcal K_\pi T,\mathcal K_\pi S)\)；取相同有限索引族的最小值保持界。\(\blacksquare\)

非扩张切割时常数为 2。不能对未经证明非扩张的``拟合后重新切割模型''直接使用常数 2。

\hypertarget{ux7236ux5c42ux65b0ux589eux9884ux6d4bux4fe1ux606f}{%
\section{15.4 父层新增预测信息}\label{ux7236ux5c42ux65b0ux589eux9884ux6d4bux4fe1ux606f}}

令子层规范核心 \(C_1,\ldots,C_m\)、父层核心 \(C_P\)、接口 \(U\)，均采用有限条件熵。定义 \[
\Delta_P=H(C_P\mid C_1,\ldots,C_m,U).
\]

\textbf{P-COMP-04。} 若父层表示 \(M_P\) 精确充分，则 \[
H(M_P\mid C_1,\ldots,C_m,U)\ge\Delta_P.
\]

\textbf{证明。} 最小性给 \(C_P=g(M_P,U)\)。条件于子核心与接口，联合熵 \(H(M_P,C_P\mid C_{1:m},U)=H(M_P\mid C_{1:m},U)\)，同时不小于 \(H(C_P\mid C_{1:m},U)\)。\(\blacksquare\)

一个一致的过程例：比特对的奇偶 \(S=X_1\oplus X_2\) 永久保持，每步重采样均匀 \(R\)，令 \[
X'_1=R,\qquad X'_2=R\oplus S.
\] 给定完整历史，每个单独子坐标的全部未来都是均匀独立序列，其单独预测核心为常数；联合未来的奇偶却精确保存 \(S\)。均匀初态下 \(\Delta_P=\log2\)。它展示联合预测信息，尚不单独证明两个子坐标都是完整子对象。

\hypertarget{ux5b50ux5bf9ux8c61ux8054ux76dfux4f5cux4e3aux7b2cux4e8cux5c42ux8f7dux4f53}{%
\section{15.5 子对象联盟作为第二层载体}\label{ux5b50ux5bf9ux8c61ux8054ux76dfux4f5cux4e3aux7b2cux4e8cux5c42ux8f7dux4f53}}

对已认证子对象集合 \([m]\)，定义能恢复父核心的联盟上闭族及其极小族 \(\mathcal M_P^{\rm child}\)。在单调逻辑接口擦除下，再用 blocker 获得极小联盟毁损。

交集 \(\bigcap_{M\in\mathcal M_P^{\rm child}}M\) 是出现在每个充分联盟中的必需子对象。若该交集为空，表示存在冗余组织方式，不表示父对象不存在。

\hypertarget{ux91cdux53e0ux51e0ux4f55ux7684ux89c4ux8303-boolean-ux5316}{%
\section{15.6 重叠几何的规范 Boolean 化}\label{ux91cdux53e0ux51e0ux4f55ux7684ux89c4ux8303-boolean-ux5316}}

给重叠区域 \(V_1,\ldots,V_m\subseteq V\)，成员签名 \[
s(x)=\{i:x\in V_i\},\quad A_J=\{x:s(x)=J\}.
\]

\textbf{P-COMP-05。} 非空 \(A_J\) 构成唯一最小包含全部 \(V_i\) 的 Boolean 代数的原子。

\textbf{证明。} 每点有唯一签名，故这些集合分割 \(V\)。又 \[
A_J=\bigcap_{i\in J}V_i\cap\bigcap_{i\notin J}V_i^c.
\] 每个包含生成区域的 Boolean 代数必须含这些原子；每个生成区域又是相应原子的并，故该代数最小且唯一。\(\blacksquare\)

只记录哪些交集非空的 nerve，一般不能恢复完整成员签名和物理删除语义。

对子联盟 \(S\)，令 \(V(S)=\bigcup_{i\in S}V_i\)。若联盟充分性恰由其并区域充分性定义，则

\textbf{P-COMP-06：物理载体到联盟的提升。} \[
\mathcal M_P^{\rm child}
=\min_{\subseteq}
\bigcup_{M\in\mathcal M_P^{\rm phys}}
\min_{\subseteq}\{S:M\subseteq V(S)\}.
\]

\textbf{证明。} 充分联盟的并区域含某物理极小载体，因而联盟含覆盖该载体的一个极小子联盟；反向任何这样的覆盖都充分。\(\blacksquare\)

逻辑关闭子接口不会自动删除共享物理区域。两类 blocker 的底集与操作含义不同。

\hypertarget{ux91cdux6784ux4e0eux751fux6210ux4e8bux4ef6}{%
\section{15.7 重构与生成事件}\label{ux91cdux6784ux4e0eux751fux6210ux4e8bux4ef6}}

若区域结构改变，先通过微观运输拉回到共同底空间，再取两组区域生成的共同 Boolean 细化，比较载体与机制。没有嵌套关系时不得伪造单向分辨率箭头。

\begin{itemize}
\tightlist
\item
  父对象出现：新父证书成立，子分支仍继续。
\item
  融合：新分支承接多个前驱，部分旧分支不再独立延续。
\item
  分裂：一个前驱产生多个同时独立的认证后继。
\item
  出生/死亡：相对于允许延续关系，缺少认证前驱/后继。
\item
  GOA 分岔：闭环长期结构改变，不等同于对象出生或融合。
\end{itemize}

\hypertarget{ux4fddux7559ux4e2aux4f53ux6027ux7684ux7ec4ux5408ux7a97ux53e3}{%
\section{15.8 保留个体性的组合窗口}\label{ux4fddux7559ux4e2aux4f53ux6027ux7684ux7ec4ux5408ux7a97ux53e3}}

\textbf{P-COMP-07。} 在紧耦合参数区间 \(K\) 上，整合量 \(D(k)\) 连续非减、子保真度 \(F(k)\) 连续非增。阈值可行集均非空，令 \[
k_D=\min\{k:D(k)\ge d_*\},\quad
k_F=\max\{k:F(k)\ge f_*\}.
\] 则联合可行集为 \([k_D,k_F]\)（若 \(k_D\le k_F\)），否则为空。

\textbf{证明。} 连续性保证边界达到，单调性使两个可行集分别为上区间和下区间。其交即所述区间。\(\blacksquare\)

单调性是具体模型条件；自然系统可能有多个分离窗口，不能因为``耦合更强''就预设整合一定更好。

\hypertarget{ux5f62ux6210bindingux4f20ux8f93ux7684ux771fux5b9eux7ec4ux5408ux4e49ux52a1}{%
\section{15.9 形成---binding---传输的真实组合义务}\label{ux5f62ux6210bindingux4f20ux8f93ux7684ux771fux5b9eux7ec4ux5408ux4e49ux52a1}}

设 \(f_0,f_1\) 分别形成两端候选，\(t_x\) 运输底层输入，\(B_0,B_1\) 实现 binding，\(t_p\) 运输父层。形成近似可交换与 binding 近似可交换是两项假设。若 \(B_1\) 对中间度量为 \(L\)-Lipschitz，则

\[d\bigl(B_1(f_1(t_xx)),t_p(B_0(f_0x))\bigr)\le L\epsilon_F+\epsilon_B.\]

证明插入 \(B_1(t_f(f_0x))\)，对第一段用 Lipschitz，第二段用 binding 传输界。C4 的 \nolinkurl{certified_formed_structural_continuation} 还要求两端实际形成关系及预算在 margin 内；它不把结构距离界等同于父 token 身份相等。完整类型见第 34.4 节。

\hypertarget{ux6709ux9650ux6570ux636eux8ba4ux8bc1ux54cdux5e94ux8f7dux4f53ux4e0eux534fux8bae}{%
\chapter{16. 有限数据认证：响应、载体与协议}\label{ux6709ux9650ux6570ux636eux8ba4ux8bc1ux54cdux5e94ux8f7dux4f53ux4e0eux534fux8bae}}

\hypertarget{ux6709ux9650ux54cdux5e94ux8868ux4e0eux540cux65f6ux7f6eux4fe1ux754c}{%
\section{16.1 有限响应表与同时置信界}\label{ux6709ux9650ux54cdux5e94ux8868ux4e0eux540cux65f6ux7f6eux4fe1ux754c}}

给 \(L\) 个历史---协议单元、每单元 \(N\) 个独立未来样本、共同结果字母表大小 \(K\)。独立性只需各单元内；单元间并合界不要求独立。真实和经验响应为 \(p_j,\widehat p_j\)。

\textbf{P-STAT-01：有限字母 TV 同时界。} \[
P\left(\max_{j\le L}D_{\rm TV}(p_j,\widehat p_j)>\eta\right)
\le L\,2^{K+1}e^{-2N\eta^2}.
\] 因此可取 \[
\eta_{N,K,L}(\alpha)=
\min\left\{1,\sqrt{\frac{(K+1)\log2+\log(L/\alpha)}{2N}}\right\}.
\]

\textbf{证明。} 对每个结果子集 \(A\)，经验质量是 Bernoulli 平均，Hoeffding 给双侧误差概率至多 \(2e^{-2N\eta^2}\)。对至多 \(2^K\) 个子集与 \(L\) 个单元并合。有限空间 TV 正是这些事件差的最大值。若公式超过 1，TV 的确定上界 1 直接有效。\(\blacksquare\)

Bernoulli 响应可直接取更紧的 \[
\eta_{N,L}(\alpha)=
\min\left\{1,\sqrt{\frac{\log(2L/\alpha)}{2N}}\right\}.
\] 各单元样本数不同时可分别分配 \(\alpha_j\)，最后取各单元有效半径的最大值。连续条件历史不能因单次出现就视为``已独立重复采样同一历史''。

\hypertarget{ux7f3aux635fux8befux5deeux4e0dux989dux5916ux4f9dux8d56ux8f7dux4f53ux5019ux9009ux6570ux91cf}{%
\section{16.2 缺损误差不额外依赖载体候选数量}\label{ux7f3aux635fux8befux5deeux4e0dux989dux5916ux4f9dux8d56ux8f7dux4f53ux5019ux9009ux6570ux91cf}}

\textbf{P-STAT-02。} 在同时响应误差事件上， \[
|\widehat e(S)-e(S)|\le2\eta
\quad\text{对所有载体 }S.
\]

\textbf{证明。} 每一对响应 TV 的估计误差至多两端误差之和。取相同读出纤维及协议索引的最大值，误差仍不超过 \(2\eta\)。该事件已覆盖底层全部响应，所以不必再对每个 \(S\) 分配一次概率。\(\blacksquare\)

这减少统计并合负担，但不消除枚举 \(2^{|V|}\) 个载体的计算成本。读出映射本身若也由数据估计，必须另计纤维误判风险。

\hypertarget{ux8fd1ux4f3cux65cfux4e0eux7cbeux786eux65cf}{%
\section{16.3 近似族与精确族}\label{ux8fd1ux4f3cux65cfux4e0eux7cbeux786eux65cf}}

定义 \[
\mathcal M_\delta=\min\{S:e(S)\le\delta\}.
\]

\textbf{P-STAT-03：固定容差族恢复。} 若 \[
\min_S|e(S)-\delta|>2\eta,
\] 则经验阈值 \(\widehat e(S)\le\delta\) 恢复 \(\mathcal M_\delta\)。

\textbf{证明。} 每个真缺损与阈值间的距离大于可能估计误差，全部分类不变，再取极小元。\(\blacksquare\)

\textbf{P-STAT-04：精确族恢复。} 若非零缺损间隔 \[
\Delta=\min_{S:e(S)>0}e(S)>0
\] 已由结构假设或独立证明给定，且 \[
2\eta<\tau<\Delta-2\eta,
\] 则经验阈值恢复 \(\mathcal M_0\)。

\textbf{证明。} 真零缺损估计至多 \(2\eta\)；真非零缺损估计至少 \(\Delta-2\eta\)。两类分离。\(\blacksquare\)

若无非零候选，全部充分的退化情况直接处理。条件 \(\Delta>4\eta\) 是存在此阈值的充分间隔要求。

反例：两响应参数 \(0.50,0.55\)，无载体缺损 \(0.05\)，读出历史后缺损零。取 \(\delta=0.1\)，即使 \(\eta=0\)，近似极小族也是 \(\{\varnothing\}\)，不等于精确极小族。

\hypertarget{ux4e09ux6001ux8ba4ux8bc1ux4e0eux7f6eux4fe1ux96c6ux5408}{%
\section{16.4 三态认证与置信集合}\label{ux4e09ux6001ux8ba4ux8bc1ux4e0eux7f6eux4fe1ux96c6ux5408}}

对缺损置信区间 \[
[\ell_S,u_S]=
[\max(0,\widehat e(S)-2\eta),\min(1,\widehat e(S)+2\eta)],
\] 若 \(u_S\le\delta\)，认证 \(\delta\)-充分；若 \(\ell_S>\delta\)，排除；否则未决。

对所有与置信区间和单调性相容的真缺损表，保留其极小载体族组成的结构置信集合。置信区间覆盖零不证明真实缺损恰等于零。精确结构结论来自采样界\textbf{加上}结构间隔，而非从有限噪声数据中无条件检测``等于零''。

\hypertarget{ux89c4ux8303ux9884ux6d4bux7c7bux6062ux590d}{%
\section{16.5 规范预测类恢复}\label{ux89c4ux8303ux9884ux6d4bux7c7bux6062ux590d}}

\textbf{P-STAT-05。} 对有限历史集合，若不同真实预测类在 \[
d(h,h')=\max_iD_{\rm TV}(p_{h,i},p_{h',i})
\] 下至少相距 \(\gamma>0\)，且 \(2\eta<\gamma/2\)，则以 \(\gamma/2\) 为经验距离阈值可恢复精确类。

\textbf{证明。} 同类经验距离至多 \(2\eta\)，异类至少 \(\gamma-2\eta\)。在此事件中阈值相邻关系恰为真实等价关系；无需靠不受控制的近邻传递来猜测类别。\(\blacksquare\)

未知间隔时，阈值关系可能不传递；应输出候选分割和未决对，不把单链聚类的连通性当成数学身份。

\hypertarget{ux8fdeux7eedux672aux6765ux7684-characteristic-ux5d4cux5165}{%
\section{16.6 连续未来的 characteristic 嵌入}\label{ux8fdeux7eedux672aux6765ux7684-characteristic-ux5d4cux5165}}

设可测正定核 \(k\) 的 RKHS 可分，\(k(y,y)\le1\)，特征映射可 Bochner 积分。定义 \[
\mu_P=E_Pk(Y,\cdot),\qquad
\operatorname{MMD}_k(P,Q)=\|\mu_P-\mu_Q\|.
\] \textbf{K-KER-01。} 在声明分布类上 characteristic 意味着 \(\mu_P=\mu_Q\iff P=Q\)。它给连续未来响应的忠实表示，不自动提供连续历史条件律的学习保证。\href{https://www.jmlr.org/papers/volume13/gretton12a/gretton12a.pdf}{原始研究}

\textbf{P-STAT-06：嵌入同时误差。} 对 \(L\) 个固定条件响应，各 \(N\) 个独立样本， \[
\max_j\|\widehat\mu_j-\mu_j\|
\le r_N(\alpha)=
\frac{1+\sqrt{2\log(L/\alpha)}}{\sqrt N}
\] 以至少 \(1-\alpha\) 概率成立。

\textbf{证明。} 独立中心化特征的交叉内积均值为零，故 \(E\|\widehat\mu-\mu\|^2\le1/N\)，再由 Jensen 得均值不超过 \(N^{-1/2}\)。替换一个样本使范数统计量变化至多 \(2/N\)。McDiarmid 给偏离均值 \(t\) 的概率至多 \(e^{-Nt^2/2}\)；对 \(L\) 个响应并合。\(\blacksquare\)

类间 MMD 间隔为 \(\gamma\) 且 \(r_N<\gamma/4\) 时，用 \(\gamma/2\) 阈值恢复类，证明与 P-STAT-05 相同。

\textbf{P-STAT-07：核运输协变。} 对双可测双射 \(T\)，同步定义 \(k_T(Tx,Ty)=k(x,y)\)，则 \[
\operatorname{MMD}_{k_T}(T_\#P,T_\#Q)
=\operatorname{MMD}_k(P,Q).
\] \textbf{证明。} MMD 平方中的三个核期望在变量替换后分别相同。\(\blacksquare\)

MMD 小一般不推出 TV 小；使用控制价值界时，可直接对 RKHS 范数有界的价值函数得到 \(|Pv-Qv|\le\|v\|_{\mathcal H_k}\operatorname{MMD}(P,Q)\)，或另证 TV 转换。不能用 characteristic 性质偷换误差尺度。

\hypertarget{ux534fux8baeux5e73ux53f0ux4e0eux91cdux590dux68c0ux9a8c}{%
\section{16.7 协议平台与重复检验}\label{ux534fux8baeux5e73ux53f0ux4e0eux91cdux590dux68c0ux9a8c}}

对预先声明的有限新增协议批，若所有应区分的响应有至少 \(\gamma\) 间隔，且所有嵌入误差小于 \(\gamma/4\)，则不会漏掉这一批可检测细化。错误概率由同时界控制。

这不覆盖未测协议。若反复查看数据直到``通过''，须用独立留出数据、有效置信序列，或预先分配各轮失效预算 \(\sum_k\alpha_k\le\alpha\)。每轮固定 \(95\%\) 不能自动组成整体 \(95\%\)。

\hypertarget{ux76f8ux5173ux8f68ux8ff9ux53d1ux73b0ux4e00ux81f4ux6027ux4e0eux6a21ux578bux9009ux62e9}{%
\section{16.8 相关轨迹、发现一致性与模型选择}\label{ux76f8ux5173ux8f68ux8ff9ux53d1ux73b0ux4e00ux81f4ux6027ux4e0eux6a21ux578bux9009ux62e9}}

\textbf{P-STAT-08：有限候选发现。} 对 \(m\) 个候选，有界损失 \(\ell_f\in[0,1]\)，独立验证样本数 \(N\)。令 \[
u=\sqrt{\frac{\log(2m/\alpha)}{2N}}.
\] 经验风险最小候选 \(\widehat f\) 满足 \[
R(\widehat f)\le\min_fR(f)+2u
\] 以至少 \(1-\alpha\) 概率成立。

\textbf{证明。} Hoeffding 与并合界给所有风险误差至多 \(u\)；插入经验最优值和真实最优候选的经验风险。\(\blacksquare\)

对平稳遍历序列，有限族的各样本平均可由遍历定理收敛，从而渐近一致；没有额外混合条件时不提供上述有限样本速率。

对绝对正则（\(\beta\)-mixing）过程，采用间隔 \(b\) 的 \(m\) 个样本块时，可调用耦合引理：与相同块边缘的独立副本联合律之 TV 差至多 \((m-1)\beta(b)\)。因此独立块集中界外再加此罚项。分块后按稀有历史条件筛样并不自动保留同一混合与样本数条件。

\hypertarget{ux90e8ux7f72ux5206ux5e03ux4e0eux8986ux76d6}{%
\section{16.9 部署分布与覆盖}\label{ux90e8ux7f72ux5206ux5e03ux4e0eux8986ux76d6}}

\textbf{P-STAT-09：平均误差运输。} 对 \(0\le e\le1\)，若部署分布 \(\nu\ll\mu\) 且 \(d\nu/d\mu\le C\)，则 \[
E_\nu e\le C E_\mu e.
\] 若 \(D_{\rm TV}(\nu,\mu)\le\delta\)，则 \[
E_\nu e\le E_\mu e+\delta.
\]

\textbf{证明。} 第一式对密度上界积分；第二式用 P-MET-02 及 \(\operatorname{sp}(e)\le1\)。\(\blacksquare\)

训练中罕见而控制策略经常访问的状态可能决定全部失败风险。因此对象认证的``有效支撑''必须与实际使用的协议、策略和环境对应。

\hypertarget{good-ux4e8bux4ef6ux6570ux636eux540eux9009ux62e9ux4e0eux5b9eux9645ux90e8ux7f72}{%
\section{16.10 Good 事件、数据后选择与实际部署}\label{good-ux4e8bux4ef6ux6570ux636eux540eux9009ux62e9ux4e0eux5b9eux9645ux90e8ux7f72}}

把注册行集固定在训练前，令 Good 为其所有真实行与经验行的 TV 距离不超过指定半径。旧同时校准的 Good 子集包含于这个规范事件，所以后者拥有相同坏训练质量上界。选择表示、参考模型或策略可以依赖训练记录，只要需要的拟合、覆盖和数值预算在同一 Good 上成立。

第 38 章直接构造训练---部署联合律，证明这种 Good-only 前提足以给整体风险 \(\alpha+(1-\alpha)\beta\)。不要求所有极罕见坏样本也满足小预算。只采样部分注册行仍不自动保证部署访问受覆盖：第 36 章保留真实 \(\nu_n(C_n^c)\)，第 37 章给出一类可以从实际支持可行性导出的零覆盖罚项。

\hypertarget{ux9006ux5bf9ux8c61ux9006ux52a8ux529bux5b66ux4e0eux9006ux76eeux7684ux7684ux53efux8bc6ux522bux6027}{%
\chapter{17. 逆对象、逆动力学与逆目的的可识别性}\label{ux9006ux5bf9ux8c61ux9006ux52a8ux529bux5b66ux4e0eux9006ux76eeux7684ux7684ux53efux8bc6ux522bux6027}}

\hypertarget{ux5b9eux9a8cux7b49ux4ef7ux7c7bux4e0eux6700ux4f18ux68c0ux9a8c}{%
\section{17.1 实验等价类与最优检验}\label{ux5b9eux9a8cux7b49ux4ef7ux7c7bux4e0eux6700ux4f18ux68c0ux9a8c}}

两个模型若对所有允许实验产生相同记录律，则任何只用这些记录的统计规则都不能区分它们。

\textbf{P-INV-01：二元检验下界。} 两模型等先验，数据律为 \(P_0,P_1\)，最小平均错误为 \[
R^*=\tfrac12(1-D_{\rm TV}(P_0,P_1)).
\]

\textbf{证明。} 选择模型一的区域为 \(A\)，则错误 \[
R(A)=\tfrac12[P_0(A)+P_1(A^c)]
=\tfrac12[1-(P_1(A)-P_0(A))].
\] 最优事件差为 TV。一般空间可取相对 \(P_0+P_1\) 的密度比较事件达到最优。\(\blacksquare\)

``预测准确''``边界识别''``动力学识别''``奖励识别''具有不同等价类，不能用其中一个成功代替全部成功。

\hypertarget{fisher-ux7684ux89c4ux8303ux96f6ux65b9ux5411}{%
\section{17.2 Fisher 的规范零方向}\label{fisher-ux7684ux89c4ux8303ux96f6ux65b9ux5411}}

正则统计模型 \(p_\theta\) 有 score \(s_\theta=\nabla\log p_\theta\)， \(\mathcal I=E_\theta s_\theta s_\theta^\top\)。

\textbf{P-INV-02。} 若沿光滑群轨道的似然不变，轨道切向量属于 \(\ker\mathcal I\)。

\textbf{证明。} 对任意轨道切向 \(v\)，似然方向导数为零，故 \(v^\top s_\theta=0\) 几乎处处，进而 \(\mathcal I v=0\)。\(\blacksquare\)

\textbf{P-INV-03。} 对条件独立、满足零均值 score 条件的实验， \[
\mathcal I_{\rm total}=\sum_e\mathcal I_e,\qquad
\ker\mathcal I_{\rm total}=\bigcap_e\ker\mathcal I_e.
\]

\textbf{证明。} 联合 score 为和，交叉项期望为零。半正定矩阵二次型之和为零当且仅当每项为零。\(\blacksquare\)

若 KL 在横截方向有二阶展开，且 Fisher 二次型下界为 \(m\|h\|^2\)，则充分小邻域中 \[
D_{\rm KL}(P_\theta\|P_{\theta+h})\ge(m/4)\|h\|^2.
\] 证明是将 \(o(\|h\|^2)\) 控制为不超过 \(m\|h\|^2/4\)。该局部二阶可分性不排除全局离散等价；Fisher 奇异也不总意味着高阶意义完全不可识别。

\hypertarget{ux53d1ux5c55ux5b9aux5f8bux7684ux8bbeux8ba1ux8bc6ux522b}{%
\section{17.3 发展定律的设计识别}\label{ux53d1ux5c55ux5b9aux5f8bux7684ux8bbeux8ba1ux8bc6ux522b}}

考虑预先指定线性参数族 \[
Y_t=\phi_t^\top\theta_*+\xi_t,\quad
G_N=\sum_{t=1}^N\phi_t\phi_t^\top.
\]

\textbf{P-INV-04。} 在无噪声、参数取整个 \(\mathbb R^d\) 的指定模型中，参数由该设计唯一识别，当且仅当 \(G_N\succ0\)。

\textbf{证明。} 写设计矩阵 \(\Phi\)，有 \(G_N=\Phi^\top\Phi\)。正定等价于 \(\ker\Phi=\{0\}\)。若奇异，取非零 \(v\in\ker\Phi\)，所有 \(\theta_*+cv\) 产生同样响应。\(\blacksquare\)

多个环境共享同一个参数时，合并 Gram 的零空间是各零空间之交。若各环境参数本来不同，直接合并不会创造共同定律。

\hypertarget{ux53efux9884ux6d4bux8bbeux8ba1ux7684-ols-ux754c}{%
\section{17.4 可预测设计的 OLS 界}\label{ux53efux9884ux6d4bux8bbeux8ba1ux7684-ols-ux754c}}

\textbf{P-INV-05。} 假设 \(\phi_t\) 对 \(\mathcal F_{t-1}\) 可测， \(\|\phi_t\|_\infty\le B\)，并且 \[
E[e^{\lambda\xi_t}\mid\mathcal F_{t-1}]
\le e^{\lambda^2\sigma^2/2}
\quad\forall\lambda\in\mathbb R.
\] 在 \(G_N\succeq N\kappa I,\ \kappa>0\) 的样本上，除概率至多 \(\alpha\) 的事件外， \[
\|\widehat\theta-\theta_*\|_2
\le\frac{\sigma B}{\kappa}
\sqrt{\frac{2d\log(2d/\alpha)}N}.
\]

\textbf{证明。} 每坐标 \(Z_j=\sum_t\phi_{tj}\xi_t\) 的条件矩母函数迭代给 \(E e^{\lambda Z_j}\le e^{N\lambda^2\sigma^2B^2/2}\)。 Chernoff 和坐标并合界得 \(\|Z\|_2\le\sigma B\sqrt{2Nd\log(2d/\alpha)}\)。 再用 \(\widehat\theta-\theta_*=G_N^{-1}Z\)。\(\blacksquare\)

若 Gram 下界以概率 \(1-\alpha_G\) 成立，总覆盖至少 \(1-\alpha-\alpha_G\)。这是固定时点保守界，标准自归一化方法能给更精细的序贯版本。\href{https://arxiv.org/abs/1102.2670}{原始研究}

反例：\(\phi_t=\xi_t\) 为独立均匀正负一，\(\theta_*=0,Y_t=\xi_t\)。噪声独立且次高斯、Gram 正定，但 OLS 恒为 1。缺少设计---噪声条件时，原来的消失误差界不成立。

若实际模型还有偏差 \(b_t\)，估计误差必须额外加 \[
\left\|G_N^{-1}\sum_t\phi_tb_t\right\|_2.
\] 测量误差、错误特征与非平稳参数不能塞进``均值零噪声''而不检验。

\hypertarget{ux4e3bux52a8ux5e72ux9884ux8bbeux8ba1}{%
\section{17.5 主动干预设计}\label{ux4e3bux52a8ux5e72ux9884ux8bbeux8ba1}}

候选模型两两构成待区分对，干预 \(i\) 覆盖某对，当它的真实响应在所需置信与样本预算下可区分。选取覆盖全部目标竞争对的干预集合，形成带成本的 test-cover 优化。

这比只检验一个被认定正确的边界更适合无真值场景。贪心覆盖是计算方法；最坏情况下所有干预都可能必需。需要估计的响应差本身也有不确定性，应在设计---验证数据之间隔离或显式采用序贯置信规则。

\hypertarget{ux53efux8bc6ux522bux5bf9ux8c61ux7a0bux5e8fux4e0eux6761ux4ef6ux90e8ux7f72ux7684ux4e0dux540cux5546}{%
\section{17.6 可识别对象、程序与条件部署的不同商}\label{ux53efux8bc6ux522bux5bf9ux8c61ux7a0bux5e8fux4e0eux6761ux4ef6ux90e8ux7f72ux7684ux4e0dux540cux5546}}

\nolinkurl{ProgramPurposeInformationBoundary} 对已给编码证明：存在完整左逆解码器当且仅当编码单射。编码本身是否能由某个物理种子产生，是不同问题；同一物理记录对应两个不同程序时不可能同时无误恢复它们。

机制观察若要识别目的，C6 要求同观测机制诱导相同策略排序。训练后使用某候选时，则要求给定训练记录的未来条件核与声明实际源吻合。任何固定候选的边缘拟合都不能省去这种条件同源性。应先判断当前证据要识别哪个商，再设计区分实验。

\hypertarget{ux53efux884cux672aux6765ux4ef7ux503cux4e0eux76eeux7684ux6865ux63a5}{%
\chapter{18. 可行未来、价值与目的桥接}\label{ux53efux884cux672aux6765ux4ef7ux503cux4e0eux76eeux7684ux6865ux63a5}}

\hypertarget{ux53efux884cux8defux5f84ux4e0eux53efux5b9eux65bdux7b56ux7565}{%
\section{18.1 可行路径与可实施策略}\label{ux53efux884cux8defux5f84ux4e0eux53efux5b9eux65bdux7b56ux7565}}

给定世界、对象分支、资源预算及行动约束，定义允许策略族 \(\mathfrak P_{\rm adm}\)，每个策略诱导路径律 \(P^\pi\)。可以声明硬约束 \(P^\pi(\Gamma_{\rm adm})=1\)，也可声明机会约束 \(P^\pi(\Gamma_{\rm adm})\ge p\)。二者不是同一问题。

仅列出满足状态约束的轨迹，不保证每条轨迹能由一个因果控制器实现。可行路径集、可达路径律集、可实施策略族分开记录。

\hypertarget{ux4ef7ux503cux7684ux7c7bux578b}{%
\section{18.2 价值的类型}\label{ux4ef7ux503cux7684ux7c7bux578b}}

\textbf{D-BRIDGE。} 评价桥接 \[
B^{s,\ell}=
(\mathfrak P_{\rm adm},J,\text{时域、预算、评价主体/层级})
\] 说明对象尺度 \(s\) 的行为按哪个评价层 \(\ell\) 排序。\(J\) 可为路径泛函的期望、风险函数、约束优先序或多目标向量。只有满足相应动态一致性时，才直接使用 Bellman 递推。

对确定路径，严格单调变换保持排序；但对随机路径，``先变换后取期望''一般改变风险偏好。保持所有期望效用比较的安全变换是正仿射变换，而非任意递增函数。

\hypertarget{ux5bf9ux8c61ux5e76ux4e0dux552fux4e00ux751fux6210ux76eeux6807}{%
\section{18.3 对象并不唯一生成目标}\label{ux5bf9ux8c61ux5e76ux4e0dux552fux4e00ux751fux6210ux76eeux6807}}

\textbf{N-TEL-01。} 在同一动力学和对象证书中，只要有两个可行动作 \(a,b\)，路径目标 \[
J_a=\mathbf1_{\{A_0=a\}},\qquad J_b=\mathbf1_{\{A_0=b\}}
\] 给相反的最优选择。因此对象存在不单独决定价值。

把观察到的策略定义为唯一高分策略，会使任何行为事后``最优''；它不可证伪。实际目的解释必须预先约束目标类或由独立机制识别，然后预测未参与拟合的环境和预算响应。

\hypertarget{ux4ef7ux503cux7b49ux4ef7ux4e0eux5956ux52b1ux5851ux5f62}{%
\section{18.4 价值等价与奖励塑形}\label{ux4ef7ux503cux7b49ux4ef7ux4e0eux5956ux52b1ux5851ux5f62}}

两个评价在指定策略类与初态范围上诱导相同最优策略集合，称操作目的等价。若动态实现与最优策略集合都保持，相应 GOD 与定义良好的 GOA 家族也保持。

\textbf{P-TEL-01：折扣势函数等价。} 有界势函数 \(\psi\)、常数 \(c\)、\(a>0\)，将一步随机奖励改为 \[
r'(x,a_t,x')=a\,r(x,a_t,x')+\beta\psi(x')-\psi(x)+c.
\] 则对每个策略 \[
V'^\pi(x)=aV^\pi(x)-\psi(x)+\frac c{1-\beta},
\] 最优策略不变。

\textbf{证明。} 前 \(T\) 步折扣势项望远镜相消为 \(-\psi(x_0)+\beta^T\psi(x_T)\)。有界性使最后项趋零，其余是正仿射变换。\(\blacksquare\)

有限时域需相应调整终端奖励，不能漏掉边界项。该等价给逆奖励识别的具体不唯一方向。

\hypertarget{ux76eeux7684ux5b9eux73b0ux8bc1ux4e66-tvrc}{%
\section{18.5 目的实现证书 TVRC}\label{ux76eeux7684ux5b9eux73b0ux8bc1ux4e66-tvrc}}

冻结桥接与机制参数，在留出环境比较：

\begin{itemize}
\tightlist
\item
  观察策略的价值后悔；
\item
  允许最优或近最优策略族所能产生的路径律；
\item
  独立预算变化引起的策略和影子价格响应；
\item
  模型在异常和失败情境下的预测。
\end{itemize}

路径残差应定义为 \[
\Delta_{\rm path}=
\inf_{\pi\in\mathfrak P^*_\epsilon}
D(P_{\rm obs},P^\pi),
\] 而非固定选择一条最优策略后要求所有真实行为都与它相同。多个等值最优策略可有不同路径和长期结构。该改写避免把合理的最优策略多样性误报为机制失败。

\hypertarget{ux6570ux503cux76eeux6807ux6392ux5e8fux548cux6700ux4f18ux96c6ux5408ux7684ux5f3aux5f31}{%
\section{18.6 数值目标、排序和最优集合的强弱}\label{ux6570ux503cux76eeux6807ux6392ux5e8fux548cux6700ux4f18ux96c6ux5408ux7684ux5f3aux5f31}}

同一个数值函数必有同一策略排序，同一排序必有同一最大化集合，反向通常不成立。TeleologicalGauge 保留这些不同等价关系；两个目的在策略排序上相同，也不一定落在预先指定的代数双驱动规范轨道。

C6 的 \nolinkurl{mechanismObservation_identifies_teleologicalClass} 只在机制---目的忠实性下识别排序等价类。它不提供奖励数值、Bellman 可实现性或生存可行性。这三项分别由评价校准、控制理论与生存理论承担。第 33.5 节给同一数值目标的完整规范刻画，第 34.2 节给排序层边界。

\hypertarget{ux6700ux4f18ux63a7ux5236ux4e0e-god}{%
\chapter{19. 最优控制与 GOD}\label{ux6700ux4f18ux63a7ux5236ux4e0e-god}}

\hypertarget{ux6709ux9650ux6298ux6263ux63a7ux5236ux7684ux5b8cux6574ux6784ux9020}{%
\section{19.1 有限折扣控制的完整构造}\label{ux6709ux9650ux6298ux6263ux63a7ux5236ux7684ux5b8cux6574ux6784ux9020}}

有限状态、每状态有限非空动作集，奖励 \(|r|\le R\)，折扣 \(0<\beta<1\)。定义 \[
(Tv)(x)=\max_a[r(x,a)+\beta P_{x,a}v].
\]

\textbf{P-CTL-01：Bellman 存在与最优性。} \(T\) 有唯一固定点 \(V^*\)，值迭代从任意有界初值收敛；存在平稳确定性策略在每个状态对所有因果策略最优。

\textbf{证明。} 对任意 \(v,w\)，最大值的差不超过各候选差的最大值，故 \[
\|Tv-Tw\|_\infty\le\beta\|v-w\|_\infty.
\] 有限维完备空间中 Banach 定理给唯一固定点及收敛。选择每点达到最大值的动作 \(\pi^*\)。对任意因果策略， \[
V^*(X_t)\ge r(X_t,A_t)+\beta E[V^*(X_{t+1})\mid H_t,A_t].
\] 逐步折扣求和和取期望，得其前 \(T\) 步收益不超过 \(V^*(x)-\beta^TEV^*(X_T)\)。有界性令末项消失。对 \(\pi^*\) 每步取等号，故它达到 \(V^*\)。\(\blacksquare\)

值迭代残差可停止： \[
\|v-V^*\|_\infty\le\frac{\|Tv-v\|_\infty}{1-\beta}.
\] 证明是在压缩性不等式中插入 \(Tv\) 并移项。

\hypertarget{ux6709ux9650ux65f6ux57dfux4e0eux751fux5b58ux8bc4ux4ef7}{%
\section{19.2 有限时域与生存评价}\label{ux6709ux9650ux65f6ux57dfux4e0eux751fux5b58ux8bc4ux4ef7}}

给终端奖励 \(G\)， \[
V_T=G,\qquad
V_t(x)=\max_a[r_t(x,a)+E V_{t+1}(X')].
\] 逐步后向归纳给最优值与时变策略。

若评价仅为在域 \(\mathcal V\) 内生存 \(H\) 步， \[
V_0(x)=\mathbf1_{\mathcal V}(x),\quad
V_{h+1}(x)=\mathbf1_{\mathcal V}(x)\max_a\int V_h(y)P(dy\mid x,a).
\] 同一归纳证明它等于最大生存概率。

生存最优动作一般依赖时域：某次进入路线 A 后生存率为 \(0.9^H\)，路线 B 的进入成功率 \(0.8\)、此后每步 \(0.99\)，其 \(H\) 步生存为 \(0.8(0.99)^{H-1}\)。短时 A 更好，长时 B 更好。路径模型必须将路线状态纳入，不能只比较无记忆标签。

\hypertarget{ux7d27-feller-ux63a7ux5236}{%
\section{19.3 紧 Feller 控制}\label{ux7d27-feller-ux63a7ux5236}}

\textbf{K-CTL-01 / P-CTL-02。} 状态、动作均为紧度量空间，动作集合固定非空，奖励连续，转移核对状态动作弱连续。则 \(T:C(\mathcal X)\to C(\mathcal X)\) 为 \(\beta\)-压缩，有唯一连续固定点，并有可测平稳最优选择。

\textbf{证明。} 对连续有界 \(v\)，积分项连续；紧动作上的 maximum theorem 给最大值连续且 argmax 非空紧。可测选择定理给选择器；压缩与最优性证明同上。\(\blacksquare\)

状态依赖动作集需要闭图、紧值和相应连续性。可测选择不必连续，因而最优闭环不自动 Feller；长期占据结论要另验。

\hypertarget{ux8fdeux7eedux65f6ux95f4-hjb-ux9a8cux8bc1}{%
\section{19.4 连续时间 HJB 验证}\label{ux8fdeux7eedux65f6ux95f4-hjb-ux9a8cux8bc1}}

设可实施扩散 \(dX_t=b(X_t,a_t)dt+\sigma(X_t,a_t)dW_t\)， 收益 \(E\int_0^\infty e^{-\rho t}r(X_t,a_t)dt\)，\(\rho>0\)。

\textbf{P-CTL-03。} 若有有界 \(C^2\) 函数 \(V\)，满足 \[
\rho V(x)=\sup_a\{r(x,a)+\mathcal L^aV(x)\},
\] 存在达到上确界且闭环 SDE 良定义的可测选择器；所有有关局部化、可积与随机积分条件满足，则 \(V\) 为最优值，该选择器最优。

\textbf{证明。} 对 \(e^{-\rho t}V(X_t)\) 用 Itô 公式。HJB 给任意控制下 \(r+\mathcal L^aV-\rho V\le0\)。积分取期望得到 \[
V(x)\ge E\left[\int_0^T e^{-\rho t}r\,dt+e^{-\rho T}V(X_T)\right].
\] 有界性使终端项消失。最优选择器使不等式处处取等号。\(\blacksquare\)

若值函数只有 viscosity 正则性，采用相应标准控制理论；不能把不存在的经典梯度直接代入。

\hypertarget{god-ux7684ux6b63ux5f0fux7c7bux578b}{%
\section{19.5 GOD 的正式类型}\label{god-ux7684ux6b63ux5f0fux7c7bux578b}}

定义 \[
Q^*(x,a)=r(x,a)+\beta P_{x,a}V^*,\qquad
\mathcal A^*(x)=\arg\max_aQ^*(x,a).
\] \textbf{D-GOD。} GOD 首先是最优动作对应 \(\mathcal A^*\)。有局部漂移时可映为 \[
\mathcal D^*(x)=\{b(x,a):a\in\mathcal A^*(x)\}.
\] 跳过程可用最优生成元，离散系统可用最优转移核，无需强造欧氏方向。

\(\epsilon\)-GOD 包含满足 \(Q^*(x,a)\ge V^*(x)-\epsilon\) 的动作。若每步使用这样的动作，折扣总后悔至多 \(\epsilon/(1-\beta)\)：把单步 Bellman 缺损沿实际轨迹折扣求和即可。此处单步 \(\epsilon\) 与整条策略的 \(\epsilon\)-最优不是同一数值。

\hypertarget{hjb-ux7684ux6b63ux5f0fux524dux63d0ux4e0eux751fux5b58ux7ea6ux675fux6700ux4f18ux63a7ux5236}{%
\section{19.6 HJB 的正式前提与生存约束最优控制}\label{hjb-ux7684ux6b63ux5f0fux524dux63d0ux4e0eux751fux5b58ux7ea6ux675fux6700ux4f18ux63a7ux5236}}

\nolinkurl{DiffusionHJBVerification.p_ctl_03} 的 Lean 实现消费过程特定的 Model.ItoRun：折扣占据测度、终端概率律、奖励与漂移可积性、有限时域 Itô 期望恒等式，以及折扣累计奖励的极限。它随后证明 HJB 残差非正、选择器残差为零、有界终端项消失和最优性。该证书没有存入最优值结论，但目前也没有从任意 SDE 自动构造完整受控随机分析底座。P-CTL-03 在本节这些过程条件下成立；不能删去 ItoRun 的来源义务。

有限模型可以得到另一条已完成的加强：先由生存固定点构造保留源核的动作子类型，再对同一奖励和转移的受限模型做 Bellman。\nolinkurl{SafeOptimalControlFromViability.viable_restricted_optimal_dominates_all_causal_policies} 比较所有遵守该约束的历史依赖随机策略；其最优性不是在未受限策略全集上宣称。第 35.3 节给完整证明。

\hypertarget{ux5bf9ux8c61ux5546ux5982ux4f55ux4fddux6301ux51b3ux7b56ux4e0eux4ef7ux503c}{%
\chapter{20. 对象商如何保持决策与价值}\label{ux5bf9ux8c61ux5546ux5982ux4f55ux4fddux6301ux51b3ux7b56ux4e0eux4ef7ux503c}}

本章采用第 19 章已经保证 Bellman 压缩、最优值和所用最优策略存在的有限模型，或满足该章相应紧性与可测选择条件的模型。奖励有界、折扣相同；一般非紧模型须另证这些前提，不能只凭以下商等式取得最优策略存在性。

\hypertarget{ux7cbeux786eux63a7ux5236ux5546}{%
\section{20.1 精确控制商}\label{ux7cbeux786eux63a7ux5236ux5546}}

\textbf{P-QUO-01。} 满射 \(f:\mathcal X\to\mathcal M\) 满足同纤维动作集相同、 \[
r(x,a)=\bar r(f(x),a),\qquad
f_\#P_{x,a}=\bar P_{f(x),a},
\] 则 \[
V^*=\bar V^*\circ f,
\] 宏观最优策略提升后微观最优。

\textbf{证明。} 对任意宏观函数 \(\bar v\)， \(T(\bar v\circ f)=(\bar T\bar v)\circ f\)。 取宏观固定点，得到微观 Bellman 固定点；由唯一性得到价值等式。动作值逐项相同，故 argmax 提升。\(\blacksquare\)

``预测自身准确''不保证奖励闭合，``某策略下 Markov''不保证所有动作闭合。这两项是对象表示进入控制的必要接口条件。

\hypertarget{ux8de8ux5ea6ux654fux611fux7684ux8fd1ux4f3cux5546ux754c}{%
\section{20.2 跨度敏感的近似商界}\label{ux8de8ux5ea6ux654fux611fux7684ux8fd1ux4f3cux5546ux754c}}

设奖励与转移误差分别不超过 \(\epsilon_r,\epsilon_p\)，且动作对应同前。令 \[
w=\bar V^*\circ f,\quad
\delta=\epsilon_r+\beta\epsilon_p\operatorname{sp}(\bar V^*),\quad
D=\delta/(1-\beta).
\]

\textbf{P-QUO-02。} \[
\|V^*-w\|_\infty\le D,
\] 宏观最优策略 \(\bar\pi^*\) 的提升 \(\widehat\pi\) 满足 \[
0\le V^*-V^{\widehat\pi}\le2D.
\]

\textbf{证明。} P-MET-02 给每个动作的 Bellman 表达误差至多 \(\delta\)，故 \(w\) 的最优 Bellman 残差至多 \(\delta\)。压缩残差界给第一式。固定提升策略下，因宏观最优动作达到宏观值，同一残差界也给 \(\|V^{\widehat\pi}-w\|\le D\)。两式相加。\(\blacksquare\)

若 \(|\bar r|\le R\)，可进一步粗化为 \[
D\le\frac{\epsilon_r}{1-\beta}
+\frac{2\beta R\epsilon_p}{(1-\beta)^2}.
\] 优先使用实际跨度。这项方法在数学 v2.1 已有，不在本版重新声称原创。

\hypertarget{ux52a8ux4f5cux95f4ux9694ux4e0eux8fd1ux6700ux4f18ux8befux5dee}{%
\section{20.3 动作间隔与近最优误差}\label{ux52a8ux4f5cux95f4ux9694ux4e0eux8fd1ux6700ux4f18ux8befux5dee}}

\textbf{P-QUO-03。} 若 \(\max_a|Q^*(x,a)-\widehat Q(x,a)|\le\eta\)，且估计最优动作与其余动作间隔大于 \(2\eta\)，则它也是真实唯一最优动作。

\textbf{证明。} 估计最佳动作的真实值最多下降 \(\eta\)，竞争动作最多上升 \(\eta\)，严格差仍为正。\(\blacksquare\)

近似商中可取 \(\eta=D\)：动作值误差由局部误差 \(\delta\) 加未来值误差 \(\beta D\)，等于 \(D\)。间隔不足只意味着该证书未决，不意味着动作一定改变。

若宏观策略只有逐点价值误差 \(\epsilon_{\rm opt}\)，应重新使用其固定策略价值跨度界，或在逐步动作缺损账本中加入算法误差；不能仍把它当作精确宏观最优策略。

\hypertarget{ux5360ux636eux52a0ux6743ux7684ux7cbeux786eux6052ux7b49ux5f0f}{%
\section{20.4 占据加权的精确恒等式}\label{ux5360ux636eux52a0ux6743ux7684ux7cbeux786eux6052ux7b49ux5f0f}}

固定策略 \(\pi\)，任意有界参考函数 \(w\)，定义 \[
b_\pi=r^\pi+\beta P^\pi w-w,\quad
d_{\mu}^\pi=(1-\beta)\sum_{t\ge0}\beta^t\mu(P^\pi)^t.
\]

\textbf{P-QUO-04。} \[
V^\pi-w=(I-\beta P^\pi)^{-1}b_\pi,
\qquad
\mu(V^\pi-w)=\frac{E_{d_\mu^\pi}b_\pi}{1-\beta}.
\]

\textbf{证明。} 从固定策略 Bellman 方程相减，得到线性方程；Neumann 级数收敛，因为 \(\|\beta P^\pi\|_\infty<1\)。左乘初始分布即得。\(\blacksquare\)

这说明实际决策误差与策略访问的状态有关，而不只由最坏误差支配。

\hypertarget{ux4eceux5c40ux90e8ux6a21ux578bux8befux5deeux5230ux7b56ux7565ux540eux6094}{%
\section{20.5 从局部模型误差到策略后悔}\label{ux4eceux5c40ux90e8ux6a21ux578bux8befux5deeux5230ux7b56ux7565ux540eux6094}}

用状态---动作折扣占据分布 \(d_\mu^\pi(x,a)\)。允许局部误差 \[
\delta(x,a)=\epsilon_r(x,a)+
\beta\epsilon_p(x,a)\operatorname{sp}(\bar V^*).
\]

\textbf{P-QUO-05：占据加权的商策略后悔。} \[
J(\pi^*;\mu)-J(\widehat\pi;\mu)
\le\frac{
E_{d_\mu^{\pi^*}}\delta+
E_{d_\mu^{\widehat\pi}}\delta}{1-\beta}.
\]

\textbf{证明。} 取 \(w=\bar V^*\circ f\)，单步参考残差 \[
b(x,a)=r(x,a)+\beta P_{x,a}w-w(x).
\] 任意动作的宏观 Bellman 优势不超过零，所以 \(b(x,a)\le\delta(x,a)\)。对提升的宏观最优动作，宏观残差为零，因此 \(b(x,\widehat\pi(x))\ge-\delta(x,\widehat\pi(x))\)。分别将 P-QUO-04 应用于真实最优策略和提升策略，相减得到结论。\(\blacksquare\)

若两占据分布分别被验证分布 \(\nu\) 以密度比 \(C_*,C_{\widehat\pi}\) 支配，则上界进一步为 \[
\frac{(C_*+C_{\widehat\pi})E_\nu\delta}{1-\beta}.
\] 这是第 16 章覆盖条件与控制商的具体连接。未知 \(C_*\) 时不能把该表达报告为已计算的实际保证。

\hypertarget{ux540cux6e90ux5546ux4e0eux53d7ux9650ux63a7ux5236ux7684ux7ec4ux5408ux987aux5e8f}{%
\section{20.6 同源商与受限控制的组合顺序}\label{ux540cux6e90ux5546ux4e0eux53d7ux9650ux63a7ux5236ux7684ux7ec4ux5408ux987aux5e8f}}

实际统一核先用同一个有限模型、注册标签和奖励构造稳定分割，再通过精确核因子化获得控制商。若进一步加生存约束，应证明允许动作子类型确实保持原始转移，而不是另造一个较容易安全的模型。

这也解释占据界为什么必须使用实际策略：估计策略和真实最优策略通常访问不同状态。P-QUO-05 的两个占据项各保留其来源；本轮有限风险中的 \(\nu_n\) 则由给定实际前缀和同一粗历史策略生成。不同公式的误差槽位必须按状态空间、时域和策略类型分别对齐。

\hypertarget{goaux957fux671fux7ed3ux6784ux7684ux5b58ux5728ux4e0eux7a33ux5b9aux6027}{%
\chapter{21. GOA：长期结构的存在与稳定性}\label{goaux957fux671fux7ed3ux6784ux7684ux5b58ux5728ux4e0eux7a33ux5b9aux6027}}

\hypertarget{ux7c7bux578bux5316ux7684-goa-ux5bb6ux65cf}{%
\section{21.1 类型化的 GOA 家族}\label{ux7c7bux578bux5316ux7684-goa-ux5bb6ux65cf}}

GOD 决定当下最优动作，最优策略生成路径，GOA 描述这些路径的长期结构。按模型分别定义：

\begin{itemize}
\tightlist
\item
  确定性：满足预紧条件的最优闭环轨迹的 \(\omega\)-极限集；
\item
  随机性：最优平稳策略的 Cesàro 时刻分布的弱极限；
\item
  多复归类：初态决定的类权重及类内不变分布；
\item
  亚稳性：QSD 与 Q-process 不变占据分别记录；
\item
  非平稳环境：增广环境状态，或使用另行定义的非自治长期结构。
\end{itemize}

最优策略族可按``所有初态同时最优''或``给定初始分布最优''定义，使用时固定一种。后者在从该初态永不可达的状态上可任意选择，不应要求全部状态都属于全局 argmax。

\hypertarget{ux6709ux9650ux94feux7684ux5360ux636eux6781ux9650}{%
\section{21.2 有限链的占据极限}\label{ux6709ux9650ux94feux7684ux5360ux636eux6781ux9650}}

\textbf{P-GOA-01。} 固定有限随机核 \(P\)，时间平均 \[
\bar\mu_N=N^{-1}\sum_{t=0}^{N-1}\mu_0P^t
\] 至少有收敛子列，每个极限 \(\nu\) 满足 \(\nu P=\nu\)。

\textbf{证明。} 概率单纯形紧，且 \[
\bar\mu_NP-\bar\mu_N=(\mu_0P^N-\mu_0)/N\to0.
\] 取极限。\(\blacksquare\)

有限链进一步有完整 Cesàro 极限，可由瞬态---复归分解和各不可约类的 Cesàro 定理得到。一般 Feller 过程采用 P-PER-02。极限集合不等于具有开放吸引域的吸引子；``Attractor''在 GOA 名称中不能替代吸引性证明。

\hypertarget{dobrushin-ux6df7ux5408ux88d5ux91cf}{%
\section{21.3 Dobrushin 混合裕量}\label{dobrushin-ux6df7ux5408ux88d5ux91cf}}

有限随机核定义 \[
\alpha(P)=\max_{x,x'}D_{\rm TV}(P_x,P_{x'}).
\] 标准有限核收缩系数满足 \(D_{\rm TV}(\mu P,\nu P)\le\alpha(P)D_{\rm TV}(\mu,\nu)\)； 可将输入有符号差分为等质量正负部分并对行混合使用凸性证明。

\textbf{P-GOA-02。} 若 \(\alpha(P)<1\)，不变概率唯一。若 \(\mu P=\mu,\widehat\mu\widehat P=\widehat\mu\)，且行误差至多 \(\epsilon\)，则 \[
D_{\rm TV}(\mu,\widehat\mu)\le\frac{\epsilon}{1-\alpha(P)}.
\]

\textbf{证明。} 唯一性由收缩立即得到。插入 \(\widehat\mu P\)，有 \(d\le\alpha(P)d+\epsilon\)，移项。\(\blacksquare\)

\(\alpha(\widehat P)\le\alpha(P)+2\epsilon\)，所以若右侧小于 1，估计核也具有唯一不变概率。没有混合裕量，接近的核可有完全不同长期结构：单位矩阵与 \(\bigl(\begin{smallmatrix}1-\epsilon&\epsilon\\0&1\end{smallmatrix}\bigr)\) 的复归类数量不同。

\hypertarget{ux4fddux6301ux590dux5f52ux5206ux89e3ux7684ux591aux7ed3ux6784ux7a33ux5b9aux6027}{%
\section{21.4 保持复归分解的多结构稳定性}\label{ux4fddux6301ux590dux5f52ux5206ux89e3ux7684ux591aux7ed3ux6784ux7a33ux5b9aux6027}}

假定扰动前后相同瞬态集和相同复归类。令瞬态转移块 \(Q\)，直接进入各类的概率矩阵 \(R\)， \[
N=(I-Q)^{-1},\quad H=NR.
\] 使用矩阵最大绝对行和范数，设 \(\|\widehat Q-Q\|_\infty\le\epsilon_Q\)， \(\|\widehat R-R\|_\infty\le\epsilon_R\)， \(\|N\|_\infty\epsilon_Q<1\)。

\textbf{P-GOA-03。} \[
\|\widehat H-H\|_\infty\le
B_H=
\frac{\|N\|_\infty^2\epsilon_Q+
\|N\|_\infty\epsilon_R}
{1-\|N\|_\infty\epsilon_Q}.
\] 若类内不变分布 TV 误差为 \(\epsilon_j^{\rm stat}\)，同初态诱导的 Cesàro 极限满足 \[
D_{\rm TV}(\nu,\widehat\nu)
\le\tfrac12 B_H+\sum_j\widehat w_j\epsilon_j^{\rm stat}.
\]

\textbf{证明。} resolvent 恒等式和 Neumann 级数给 \(\|\widehat N\|\le\|N\|/(1-\|N\|\epsilon_Q)\)、 \(\|\widehat N-N\|\le\|N\|^2\epsilon_Q/(1-\|N\|\epsilon_Q)\)。 由 \(\|R\|_\infty\le1\) 估计 \((\widehat N-N)R+\widehat N(\widehat R-R)\) 得第一式。 初态吸收权重的总 \(L^1\) 差不超过 \(B_H\)，混合分布三角不等式得第二式。\(\blacksquare\)

复归支撑改变时不适用该定理。

\hypertarget{ux5bf9ux79f0-killed-ux6838ux7684ux8c31ux7a33ux5b9aux6027}{%
\section{21.5 对称 killed 核的谱稳定性}\label{ux5bf9ux79f0-killed-ux6838ux7684ux8c31ux7a33ux5b9aux6027}}

设 \(Q\) 对称、非负、不可约、次随机，单位正主向量 \(u\)，谱间隔 \(g=\rho-\lambda_2>0\)。对称扰动 \(E\) 满足 \(\|E\|_2\le\eta<g/2\)，且 \(\widehat Q=Q+E\) 保持同类性质。

\textbf{P-GOA-04。} \[
|\widehat\rho-\rho|\le\eta,\qquad
D_{\rm TV}(u^2,\widehat u^2)\le\frac{2\sqrt2\,\eta}{g}.
\]

\textbf{证明。} Rayleigh 极值给特征值误差。将 \(\widehat u\) 投影到 \(u^\perp\)，特征方程得其正交分量范数至多 \(\eta/(g-\eta)\le2\eta/g\)。两正主向量夹角不超过直角，故 \(\|u-\widehat u\|_2\le\sqrt2\sin\theta\le2\sqrt2\eta/g\)。 最后 \[
\tfrac12\sum_i|u_i^2-\widehat u_i^2|
\le\tfrac12\|u-\widehat u\|_2\|u+\widehat u\|_2
\le\|u-\widehat u\|_2.
\] 对称 killed 核的 Q-process 占据为 \(u_i^2\)，得结论。\(\blacksquare\)

该界属于特殊谱稳定性，不是任意非自伴 QSD 的普遍界。

\hypertarget{ux6709ux9650ux98ceux9669ux9884ux7b97ux4e0dux80fdux65e0ux6761ux4ef6ux53d8ux6210ux957fux65f6ux5438ux5f15}{%
\section{21.6 有限风险预算不能无条件变成长时吸引}\label{ux6709ux9650ux98ceux9669ux9884ux7b97ux4e0dux80fdux65e0ux6761ux4ef6ux53d8ux6210ux957fux65f6ux5438ux5f15}}

即使每步误差很小，\(\sum_{t<T}\epsilon_t\) 也可能随 \(T\) 增长到平凡上界一。没有收缩、谱或常返条件时，不能对有限路径误差界直接令 \(T\to\infty\)，便声称 GOA 稳定。

M-OI-01 统一的是``有极限且可观测残差趋零→极限不变''的核心论证；M-PE-01 统一的是给定正特征数据的归一化后果。它们不自动给出 P-GOA-02 的混合裕量、P-GOA-03 的同复归分解或 P-GOA-04 的谱隙。第 32 章据此保留 GOA 家族中的领域适配。

\hypertarget{ux540cux4e00ux53d7ux63a7ux6e90ux7684ux7a33ux6001ux6761ux4ef6ux751fux5b58ux4e0eux4ef7ux503c}{%
\section{21.7 同一受控源的稳态、条件生存与价值}\label{ux540cux4e00ux53d7ux63a7ux6e90ux7684ux7a33ux6001ux6761ux4ef6ux751fux5b58ux4e0eux4ef7ux503c}}

本节把第 10、19、21 章的一般数学机制实例化到同一个实际有限受控源，并与第 35.7 节的对象划分、支持和恢复相接。源核、精确控制商与安全固定点已有 Lean 实现；以下参数化稳态、QSD 及完整受限最优闭环的显式实例给出新的书面证明，尚不声称这些新数值和实例已有专门 Lean 端点。第 35.4 节的 8/5 模型及第 38.6 节的 9/64 注册实验保持各自独立来源。

\hypertarget{ux552fux4e00ux7684ux53d7ux63a7ux6e90ux4e0eux5b8cux6574ux7b56ux7565}{%
\subsection{21.7.1 唯一的受控源与完整策略}\label{ux552fux4e00ux7684ux53d7ux63a7ux6e90ux4e0eux5b8cux6574ux7b56ux7565}}

取实际定义 \[
\begin{gathered}
X=\mathrm{Bool}\times\mathrm{Bool},\qquad A=\mathrm{Bool},\\
M=\mathrm{conflictRewardControl.micro},\\
P(x,a)=M.\mathrm{transitionPMF}(x,a).
\end{gathered}
\]

以下以 0、1 表示 false、true，写状态为 \(x=(b,u)\)。观察是 \(q(b,u)=b\)，安全域是 \(K=\{(1,0),(1,1)\}\)。初始律明确取 \(\nu_0=\delta_{(1,0)}\)。转移、奖励、折扣全部来自同一个 M： \[
P((b,u),a)\{(b',u')\}=\tfrac12{\bf1}_{b'=a},\qquad
r((b,u),a)=1-a,\qquad \gamma=\tfrac12.                 \tag{21.7-1}
\]

式（21.7-1）的 PMF 质量应由 M 的实数转移经已有 transitionPMF 转换定理取得，后续形式化不能另造同值表而省去相等证明。内部位按该转移重新取值，允许保持原值。

对每个 \(\theta\in[0,1]\)，定义同一个完整因果策略 \(\pi_\theta\)：对每一时刻及每个可见历史，其动作 PMF 恒为 \[
\pi_\theta(0)=\theta,\qquad \pi_\theta(1)=1-\theta.     \tag{21.7-2}
\]

这一定义也覆盖零概率历史。实际有限路径律按生产路径 API 的因果乘积／bind 递推构造。因而式（21.7-2）是在\textbf{每个历史条件下}的动作分布，不只是无条件动作频率；不需要从某个边际频率假定独立性。

令 \[
p_\theta(0)=\theta,\quad p_\theta(1)=1-\theta,\qquad
\mu_\theta(b,u)=p_\theta(b)/2.
\]

\hypertarget{ux65e0ux6761ux4ef6ux95edux73afux6240ux6709ux53c2ux6570ux5747ux4e00ux6b65ux5230ux8fbeux552fux4e00ux4e0dux53d8ux5f8b}{%
\subsection{21.7.2 无条件闭环：所有参数均一步到达唯一不变律}\label{ux65e0ux6761ux4ef6ux95edux73afux6240ux6709ux53c2ux6570ux5747ux4e00ux6b65ux5230ux8fbeux552fux4e00ux4e0dux53d8ux5f8b}}

\textbf{M-21.1。} 对所有 \(\theta\in[0,1]\)，从（21.7-1）及（21.7-2）诱导的真实闭环转移为 \[
L_\theta(x,y)
=\sum_{a\in A}\pi_\theta(a)P(x,a)\{y\}
=\mu_\theta(y).                                      \tag{21.7-3}
\] 每个概率行 \(\eta\) 满足 \(\eta L_\theta=\mu_\theta\)。因此 \(\mu_\theta\) 是唯一不变概率行，所有初始分布在一步后即与它相等；Dobrushin 系数为零。

\textbf{证明。} 固定 \(y=(b',u')\)，求和中只有 \(a=b'\) 一项可能非零，得到 \(p_\theta(b')/2\)。对 \(x\) 求和， \[
(\eta L_\theta)(y)=\sum_x\eta(x)\mu_\theta(y)=\mu_\theta(y).
\] 取 \(\eta=\mu_\theta\) 得不变性；若 \(\eta L_\theta=\eta\)，则同一等式给 \(\eta=\mu_\theta\)。各行相同，行间 TV 均为零。\(\blacksquare\)

\textbf{M-21.2。} 记实际第 t 步边际为 \(\nu_t=\nu_0L_\theta^t\)。对正整数 \(N\)， \[
\bar\nu_N=\frac1N\sum_{t=0}^{N-1}\nu_t
=\frac1N\nu_0+\frac{N-1}{N}\mu_\theta.                \tag{21.7-4}
\] 从而 \[
D_{\rm TV}(\bar\nu_N,\mu_\theta)
=\frac1N D_{\rm TV}(\nu_0,\mu_\theta)
=\frac{1+\theta}{2N}\longrightarrow0,                \tag{21.7-5}
\] \[
\bar\nu_NL_\theta-\bar\nu_N
=\frac{\mu_\theta-\nu_0}{N}.                         \tag{21.7-6}
\]

\textbf{证明。} M-21.1 给 \(t\ge1\) 时 \(\nu_t=\mu_\theta\)。代入平均得（21.7-4）。TV 对同一有符号差的正标量齐次，且 \(D_{\rm TV}(\delta_x,\mu)=1-\mu(x)\)，得到（21.7-5）；代入一次转移得（21.7-6）。\(\blacksquare\)

（21.7-4）是概率律的 Cesàro 平均，尚未宣称某条样本轨迹的经验频率定理。生产 cesaroRow 用参数 \(n\) 表示 \(N=n+1\)，适配时必须处理这个下标偏移。

在本例中，M-OI 的``残差趋零且存在极限''条件由显式计算给出，有限单纯形的子列抽取只需复用。通用 P-GOA-01 本身仍只声称子列存在及每个子列极限不变；本例更强的全序列收敛和唯一性来自（21.7-3），不能反推通用定理也有这些结论。

\hypertarget{ux6740ux6b7bux6838ux6765ux81eaux540cux4e00ux4e2aux5b9eux9645ux8fc7ux7a0b}{%
\subsection{21.7.3 杀死核来自同一个实际过程}\label{ux6740ux6b7bux6838ux6765ux81eaux540cux4e00ux4e2aux5b9eux9645ux8fc7ux7a0b}}

定义首次离开安全域的时刻 \[
\tau=\inf\{t\ge1:X_t\notin K\},
\] 并令 \(S_T=\{X_1,\ldots,X_T\in K\}\)，\(S_0\) 为整个样本空间。

把 K 通过 \(u\mapsto(1,u)\) 与 Bool 对应。取同一个闭环 \(L_\theta\) 在 K 上的限制，而不另给一个独立矩阵： \[
Q_\theta(u,u')=L_\theta((1,u),(1,u'))=\frac{1-\theta}{2}. \tag{21.7-7}
\]

它的行和为 \(1-\theta\)，缺失质量 \(\theta\) 正是原过程这一步离开 K 的概率。若只删除 K 外状态后重新归一化，得到的是另一核，不能称作（21.7-7）的原始 killed 核。

设 \[
H(u,u')=\tfrac12,\quad \rho=1-\theta,\quad
q_*(u)=\tfrac12,\quad h_*(u)=1.
\] 于是 \(Q_\theta=\rho H\)，且 \(H^2=H\)。

\textbf{M-21.3。} 当 \(0\le\theta<1\) 时，以上数据逐项满足现有 BiEigenData 与 KilledQSDGOA 的要求： \[
Q_\theta\ge0,\quad Q_\theta{\bf1}=\rho{\bf1},\quad
0<\rho\le1,\quad q_*Q_\theta=\rho q_*,
\quad \sum_uq_*(u)=\sum_uq_*(u)h_*(u)=1.              \tag{21.7-8}
\] 其 Doob 变换及不变权重为 \[
Q_\theta^{h_*}(u,u')
=\frac{Q_\theta(u,u')h_*(u')}{\rho h_*(u)}
=\tfrac12=H(u,u'),\qquad q_*h_*=q_*.                 \tag{21.7-9}
\]

\textbf{证明。} 所有非负性来自 \(\theta\in[0,1)\)。两项求和使 \(Q_\theta{\bf1}\) 等于 \(1-\theta\)。对每个 \(u'\)， \[
(q_*Q_\theta)(u')=2\cdot\tfrac12\cdot\tfrac{\rho}{2}
=\rho/2=\rho q_*(u').
\] 两项归一化直接求和。因为 \(\rho>0\)，（21.7-9）的除法合法。\(\blacksquare\)

不用 Perron--Frobenius 存在定理来假设一个未知正特征向量；这个有限实例把正特征数据明确算出，然后复用正特征结构生成器的后果。

\textbf{M-21.4。} 对任意 K 内初始概率行 \(\eta\)，以及任意 \(T\ge1\)， \[
\eta Q_\theta^T=\rho^Tq_*,\qquad
\sum_{u'}(\eta Q_\theta^T)(u')=\rho^T.                \tag{21.7-10}
\] 当 \(\theta<1\) 时，存活到 T 条件下的终端内部位分布恰为 \(q_*\)。

\textbf{证明。} \(H^2=H\) 给 \(Q_\theta^T=\rho^TH\)；任意概率行满足 \(\eta H=q_*\)。质量求和得（21.7-10），并在 \(\rho^T>0\) 时归一化。\(\blacksquare\)

这一步不能直接用 survival\_mass\_pow 的 QSD 初始律代替 \(\nu_0=\delta_{(1,0)}\)：该通用定理从 q\_* 开始。式（21.7-10）是所需的任意 K 内初始律适配。对于 \(T=0\)，实际条件终端内部位仍为 \(\delta_0\)，并非 q\_\emph{。从 q\_} 出发的 all-steps QSD 命题包括零步，两者量词不同。

\hypertarget{ux5b9eux9645ux6709ux9650ux98ceux9669ux4e0eux6761ux4ef6ux6837ux672cux7a7aux95f4}{%
\subsection{21.7.4 实际有限风险与条件样本空间}\label{ux5b9eux9645ux6709ux9650ux98ceux9669ux4e0eux6761ux4ef6ux6837ux672cux7a7aux95f4}}

\textbf{M-21.5。} 对 \(\theta\in[0,1]\) 和 \(T\ge0\)，同一个 \((P,\pi_\theta,\nu_0)\) 的有限路径满足 \[
\Pr_{\pi_\theta}(S_T)=(1-\theta)^T,\qquad
\Pr_{\pi_\theta}(\tau\le T)=1-(1-\theta)^T.           \tag{21.7-11}
\] 对 \(T\ge1\)，其终端失效概率却是 \[
\Pr_{\pi_\theta}(X_T\notin K)=\theta.                \tag{21.7-12}
\]

\textbf{证明。} 在任意已经存活的历史之后，下一步仍在 K 当且仅当选动作 1；其条件概率由（21.7-2）恒为 \(1-\theta\)。因此 \(s_{T+1}=(1-\theta)s_T\)，\(s_0=1\)，归纳得（21.7-11）。这也与实际 killed 核的质量（21.7-10）一致。式（21.7-12）由M-21.1 的真实边际 \(\nu_T=\mu_\theta\) 得出。\(\blacksquare\)

以 \(\theta=1/8,T=2\) 为例： \[
\Pr(\tau\le2)=15/64,\qquad \Pr(X_2\notin K)=1/8.
\] 两者之差 \(7/64\) 正是先离开 K、后又返回 K 的概率。在这两个运行步骤内的``返回''是原物理转移的恢复，不等于已经运行了独立程序恢复协议。

若认证行谓词取 \(C(y,a)\iff a=1\)，每个决策时刻的实际未认证行占据为 \(\theta\)。只有 \(\theta=0\) 时它为零；同一个 \(\pi_\theta\) 不能借用 \(\pi_0\) 的零占据。若选择全部行认证并使用精确参考核，占据虽然可为零，参考模型的真实物理失败风险仍为（21.7-11）。这是两个不同的证书设计，不能挑各自有利的项拼接。

\hypertarget{ux5168ux8defux5f84ux6761ux4ef6ux5316ux7684ux7cbeux786eux7b49ux5f0f}{%
\subsubsection{21.7.4.1 全路径条件化的精确等式}\label{ux5168ux8defux5f84ux6761ux4ef6ux5316ux7684ux7cbeux786eux7b49ux5f0f}}

对 \(\theta<1\)，在任意固定 T 下，以 \(S_T\) 为条件的完整有限状态---动作路径律等于 \(\pi_0\) 的同长路径律。

证明如下：每一条存活路径的动作全部为 1，故其 \(\pi_\theta\) 路径质量等于相同路径的 \(\pi_0\) 质量乘以 \((1-\theta)^T\)；非存活路径的条件质量为零；最后除以（21.7-11）。这说明条件化怎样产生（21.7-9）的公平内部核，同时也说明两种无条件部署律并不相等。这个简单常危险率实例不提供一般 Q-process 与有限时窗条件化恒等的定理。

\hypertarget{ux65e0ux9650ux8defux5f84ux7684ux6570ux5b66ux89e3ux91ca}{%
\subsubsection{21.7.4.2 无限路径的数学解释}\label{ux65e0ux9650ux8defux5f84ux7684ux6570ux5b66ux89e3ux91ca}}

有限式（21.7-11）已经给出 \[
0<\theta\le1\quad\Longrightarrow\quad
\lim_{T\to\infty}\Pr_{\pi_\theta}(\tau\le T)=1.         \tag{21.7-13}
\]

如需把（21.7-13）表述为同一无限实验中的``几乎必然最终失效''，可以明确构造：在动作序列与内部位序列的乘积概率空间上，让各动作坐标服从（21.7-2）、内部位坐标为公平 Bool，两个坐标族独立；置 \(X_0=(1,0)\)，\(X_{n+1}=(A_n,U_{n+1})\)。其每个有限圆柱分布逐项等于上述实际因果乘积，所以它实现了同一源与同一策略。由事件 \(S_T\) 的递减连续性， \[
\Pr(\tau=\infty)=\lim_T(1-\theta)^T=0
\] 对 \(0<\theta\le1\) 成立。对于 \(\theta=0\)，该概率为一。

这里给出的是本例明确乘积空间上的数学延伸；对应无限路径 Lean 对象及其与生产有限 PMF 的一致性尚未形式化。通用``每个有限 T 存在可行策略''的命题仍不能直接换量词成``存在一个策略无限存活''。

\hypertarget{ux6700ux4f18ux6027ux53efux884cux6027ux4e0eux5168ux7a7aux95f4ux4e0dux53d8ux5f8b}{%
\subsection{21.7.5 最优性、可行性与全空间不变律}\label{ux6700ux4f18ux6027ux53efux884cux6027ux4e0eux5168ux7a7aux95f4ux4e0dux53d8ux5f8b}}

\textbf{M-21.6。} 对所有 \(\theta\in[0,1]\)， \[
J_M(\pi_\theta;\nu_0)
=\mathbb E\!\left[\sum_{n=0}^\infty2^{-n}(1-A_n)\right]
=2\theta.                                          \tag{21.7-14}
\] 原模型 M 的无限折扣最优值在每个状态均为 2，唯一 Bellman greedy 动作为 0。其闭环是 \(L_1\)，唯一不变律为 \(\mu_1\)，完全支持在 K 外。

\textbf{证明。} 每步奖励非负有界，有限期望和为 \(\theta\sum_{n<T}2^{-n}\)，单调取极限得到（21.7-14）。任意因果策略的逐步奖励至多 1，所以总值至多 2；恒选 0 达到此值。常值函数 2 的两动作 Bellman 回报分别为 2 和 1，因此 greedy 唯一为 0。\(\blacksquare\)

可见，所有 \(L_\theta\) 都具有唯一且一步吸引的不变律，这一性质本身不能识别哪一策略优化 M 的给定奖励。最优策略也未因此保持预先登记的安全条件。

\hypertarget{ux5fc5ux987bux4fddux7559-k-ux5916ux7684ux5b9eux9645ux52a8ux4f5cux57df}{%
\subsubsection{21.7.5.1 必须保留 K 外的实际动作域}\label{ux5fc5ux987bux4fddux7559-k-ux5916ux7684ux5b9eux9645ux52a8ux4f5cux57df}}

当前生产 viableRestrictedModel 的限制为 \[
A_K(x)=\{a:x\in K\Rightarrow \operatorname{supp}P(x,a)\subseteq K\}.
\] 由（21.7-1）， \[
A_K(x)=\{1\}\ \ (x\in K),\qquad A_K(x)=\{0,1\}\ \ (x\notin K). \tag{21.7-15}
\]

令 R 为这个实际受限模型，它保留 M 的转移、奖励、折扣。\textbf{不能把（21.7-15）改写为对所有 x 都只许动作 1。}

\textbf{M-21.7。} R 的最优值及唯一 greedy 为 \[
V_R(b,u)=2(1-b),\qquad g_R(b,u)=b.                   \tag{21.7-16}
\] 其完整闭环矩阵为 \[
D((b,u),(b',u'))=\tfrac12{\bf1}_{b'=b},              \tag{21.7-17}
\] 全部不变概率恰好为 \[
\{\lambda\mu_0+(1-\lambda)\mu_1:0\le\lambda\le1\}.   \tag{21.7-18}
\] 从 \(\nu_0=\delta_{(1,0)}\) 出发，其一步后的边际与长期极限均为 \(\mu_0\)，整条实际路径与 \(\pi_0\) 相同。

\textbf{证明。} K 内唯一合法动作是 1，此后保持 K，所有奖励为零。K 外恒选 0 保持 K 外且每步奖励 1，达到全局上界 2。代入候选值可直接比较：K 外动作 0 的 Bellman 回报为 2，动作 1 的回报为 0；K 内只有动作 1。得到（21.7-16）及（21.7-17）。

对任意概率行 \(\eta\)，记其可见位质量为 \(\eta_b=\sum_u\eta(b,u)\)。一次作用 D 后， \[
(\eta D)(b,u)=\eta_b/2.
\] 所以不变性恰要求每个可见纤维内为公平内部位分布，而两个纤维的总权重任意。这给出（21.7-18）。从 K 内初始，所有实际后继仍在 K，greedy 恒选 1，因果递推逐步与 \(\pi_0\) 一致。\(\blacksquare\)

故 R 的完整 GreedyInvariantGOASet \textbf{不是单点}；其支持在 K 的成员只有 \(\mu_0\)，而从指定初态所得 CesàroLimitGOA 也只有 \(\mu_0\)。完整矩阵 D 的 Dobrushin 系数为 1，因为两个可见纤维的行支撑互不相交。不能调用要求全空间严格收缩的定理来证明 D 的全局唯一不变律。

若另行把状态类型限制为 K，得到的核是 H，确有唯一不变律 q\_*。这是带明确状态包含映射的限制操作，不能悄悄替换原 R 的状态类型。也不能把完整 D 与 \(L_0\) 视为相同：它们只在 K 内一致，K 外的动作与后继不同。

\hypertarget{ux7aefux70b9ux4e0eux4e0dux6210ux7acbux7684ux63a8ux8bba}{%
\subsection{21.7.6 端点与不成立的推论}\label{ux7aefux70b9ux4e0eux4e0dux6210ux7acbux7684ux63a8ux8bba}}

\begin{longtable}[]{@{}
  >{\raggedright\arraybackslash}p{(\columnwidth - 6\tabcolsep) * \real{0.2500}}
  >{\raggedright\arraybackslash}p{(\columnwidth - 6\tabcolsep) * \real{0.2500}}
  >{\raggedright\arraybackslash}p{(\columnwidth - 6\tabcolsep) * \real{0.2500}}
  >{\raggedright\arraybackslash}p{(\columnwidth - 6\tabcolsep) * \real{0.2500}}@{}}
\toprule\noalign{}
\begin{minipage}[b]{\linewidth}\raggedright
情形
\end{minipage} & \begin{minipage}[b]{\linewidth}\raggedright
\(\theta=0\)
\end{minipage} & \begin{minipage}[b]{\linewidth}\raggedright
\(0<\theta<1\)
\end{minipage} & \begin{minipage}[b]{\linewidth}\raggedright
\(\theta=1\)
\end{minipage} \\
\midrule\noalign{}
\endhead
\bottomrule\noalign{}
\endlastfoot
\(L_\theta\) 的唯一不变律 & \(\mu_0\)，在 K 内 & \(\mu_\theta\)，两纤维都有质量 & \(\mu_1\)，在 K 外 \\
正 T 的曾失败概率 & 0 & \(1-(1-\theta)^T\) & 1 \\
正 T 的终端失败概率 & 0 & \(\theta\) & 1 \\
killed 行和 \(\rho\) & 1 & 严格介于 0、1 & 0 \\
BiEigenData / 正生存 QSD & 可用；实际没有杀死质量 & 可用 & 不可用，scale\_pos 失败 \\
Doob 核 & H & H & 归一化分母为零，不能定义为本定理的 Doob 核 \\
\(J_M(\pi_\theta)\) & 0 & \(2\theta\) & 2 \\
\end{longtable}

T=0 时所有参数的存活概率为 1、曾失败概率为 0；使用约定 \(0^0=1\) 的自然数幂。实际初始内部位是 \(\delta_0\)，所以``条件终端内部位为 q\_*''严格保留 T≥1。

本节反驳以下无条件推论：唯一稳态推出零失效；QSD 支持在 K 推出未条件化部署安全；从安全初态的唯一极限推出全状态空间唯一 GOA；奖励最优推出可行；同一状态类型或同一个不变律推出同一实际政策／路径来源。它没有反驳附有真实最优性、保持域、初态、收缩或条件化前提的既有定理。

\hypertarget{lean-ux5bf9ux5e94ux53caux4e0dux80fdux7701ux7565ux7684ux6765ux6e90ux6865}{%
\subsection{21.7.7 Lean 对应及不能省略的来源桥}\label{lean-ux5bf9ux5e94ux53caux4e0dux80fdux7701ux7565ux7684ux6765ux6e90ux6865}}

InvariantLawGOA、CesaroLimitGOA 和 KilledQSDGOA 描述给定核的长期对象，本身没有最优性字段。GreedyInvariantGOASet 则明确使用所给控制模型的 Bellman-greedy 闭环。本节每个名称都按这个区别使用，不因一个稳定分布存在而给它附加奖励最优或规范价值。

现有 \href{../formalization/ueot-core/UEOT/V3/FiniteCesaroInvariant.lean}{FiniteCesaroInvariant} 与 \href{../formalization/ueot-core/UEOT/V3/Compression/OccupationLimitInvariance.lean}{\nolinkurl{OccupationLimitInvariance}} 可消费本节的随机矩阵、初始单纯形和 Cesàro 极限。现有 \href{../formalization/ueot-core/UEOT/V3/Compression/PositiveEigenstructure.lean}{PositiveEigenstructure} 的 BiEigenData、qsdDoob\_bundle 与 killed wrapper 可消费本节明确算出的正特征数据。一般接口已经证明；把 Lθ 与真实策略路径、Qθ 与实际 killed 源、R 的 greedy 与其真实状态相关动作域逐项对齐，是本实例仍需补的形式化。

\href{../formalization/ueot-core/UEOT/V3/Compression/TheoryCompletion/SameObjectViableGOA.lean}{SameObjectViableGOA} 已证明：真实同对象 greedy 保持非空持续域时，存在支持在该域内的不变律。这个接口还需要完整 Objecthood parent、control spec 及其来源等式，本节没有仅凭四状态类型构造这些字段。它的结论也不要求完整状态空间只有一个不变律。更强的全局唯一与混合端点另外要求真实收缩或孤立条件，不能用于本节拥有两个闭类的 D。

同源实例由此把 M-QD 的准确观察商、有限路径传播、M-OI 的无条件占据和 M-PE 的条件存活结构接在一起。它没有改变四生成机制及 106 个 P-ID 的分类，也没有由此补齐物理程序来源、父对象形成、谱系或经验双驱动的其他条件。

\hypertarget{ux8defux5f84ux4fe1ux606fux4ee3ux4ef7ux53efux8fbeux4f5cux7528ux4e0eux53d8ux5206ux5b9eux73b0}{%
\chapter{22. 路径信息代价、可达作用与变分实现}\label{ux8defux5f84ux4fe1ux606fux4ee3ux4ef7ux53efux8fbeux4f5cux7528ux4e0eux53d8ux5206ux5b9eux73b0}}

\hypertarget{ux751fux5b58ux4e8bux4ef6ux7684ux4fe1ux606fux4e0bux754c}{%
\section{22.1 生存事件的信息下界}\label{ux751fux5b58ux4e8bux4ef6ux7684ux4fe1ux606fux4e0bux754c}}

给基线路径律 \(P^0\)、可达路径律 \(Q\ll P^0\) 和事件 \(A\)，记 \(q=P^0(A),p=Q(A)\)。

\textbf{P-KL-01。} \[
D_{\rm KL}(Q\|P^0)\ge
d_{\rm Bern}(p\|q).
\] \textbf{证明。} 对路径作可测映射 \(\gamma\mapsto\mathbf1_A\)，用 KL 数据处理。\(\blacksquare\)

若 \(q=0\)，绝对连续迫使 \(p=0\)；要达到正概率必须改变可达支撑或使该比较的 KL 为无穷。若 \(q=1\)，绝对连续也迫使 \(p=1\)。

\hypertarget{ux7cbeux786eux4e8bux4ef6-i-projection}{%
\section{22.2 精确事件 I-projection}\label{ux7cbeux786eux4e8bux4ef6-i-projection}}

\textbf{P-KL-02。} 给 \(0<q<1\)、目标 \(0\le p\le1\)，在全部绝对连续路径律中， \[
\inf_{Q:Q(A)\ge p}D_{\rm KL}(Q\|P^0)
=
\begin{cases}
0,&p\le q,\\
d_{\rm Bern}(p\|q),&p>q.
\end{cases}
\] 当 \(p>q\)，达到者为 \[
\frac{dQ^*}{dP^0}
=\frac pq\mathbf1_A+
\frac{1-p}{1-q}\mathbf1_{A^c}.
\]

\textbf{证明。} \(p\le q\) 时基线本身可行。\(p>q\) 时，Bernoulli KL 对事件概率 \(r\ge p\) 严格递增，因为导数为 \(\log[r(1-q)/(q(1-r))]>0\)。 给定密度只改变两事件的质量，保留各自条件路径律，因此达到下界。边界 \(p=1\) 按 \(0\log0=0\) 解释。\(\blacksquare\)

这是在全部路径律中的最小信息改变量。受限反馈控制器未必能实现 \(Q^*\)，因此真实控制最小代价只会更高。指数倾斜的存在不单独证明自然系统实施了这个优化。

\hypertarget{ux79bbux6563-markov-ux8defux5f84ux7684-kl}{%
\section{22.3 离散 Markov 路径的 KL}\label{ux79bbux6563-markov-ux8defux5f84ux7684-kl}}

同一初态，两核 \(P_t,Q_t\) 满足适当逐行绝对连续，则 \[
D_{\rm KL}(P_{0:T}\|Q_{0:T})
=\sum_{t=0}^{T-1}E_P
D_{\rm KL}(P_t(\cdot\mid H_t)\|Q_t(\cdot\mid H_t)).
\]

\textbf{P-KL-03，证明。} 路径密度比为逐步条件密度比之积。取对数变成和，再逐步条件期望。不同初态另加初始 KL。\(\blacksquare\)

该链式法则也适用于历史依赖核，是一般过程的可执行误差账本。

\hypertarget{ctmc-ux7684ux663eux5f0fux516cux5f0f}{%
\section{22.4 CTMC 的显式公式}\label{ctmc-ux7684ux663eux5f0fux516cux5f0f}}

对有限非爆炸 CTMC、相同初态，受控速率 \(l^u_{ij}\) 的正支撑包含于基线 \(l^0_{ij}\) 的正支撑，记 \(\lambda_i^u=\sum_{j\ne i}l^u_{ij}\)。则

\textbf{P-KL-04。} \[
D_{\rm KL}(P^u_{[0,T]}\|P^0_{[0,T]})
=
E_u\int_0^T
\sum_{j\ne X_t}
\left[
l^u_{X_tj}\log\frac{l^u_{X_tj}}{l^0_{X_tj}}
-l^u_{X_tj}+l^0_{X_tj}
\right]dt.
\]

\textbf{证明。} 路径 log 似然比是每次跳变的 log 速率比之和，减去两退出速率差的时间积分。对受控律取期望，用跳计数补偿将跳变和换为相应速率积分。\(\blacksquare\)

零受控速率项用 \(0\log0=0\)。时间或历史依赖速率需要相应可预测性与补偿可积条件。

\hypertarget{ux6269ux6563ux566aux58f0ux7a7aux95f4ux4e0eux72b6ux6001ux8defux5f84ux987bux5206ux5f00}{%
\section{22.5 扩散：噪声空间与状态路径须分开}\label{ux6269ux6563ux566aux58f0ux7a7aux95f4ux4e0eux72b6ux6001ux8defux5f84ux987bux5206ux5f00}}

在参考概率空间上， \[
dX_t=b^0(t,X_t)dt+\sigma(t,X_t)dW_t^0.
\] 令 \(u_t\) 渐进可测，指数 \[
Z_T=\exp\left(\int_0^T u_t\,dW_t^0-
\tfrac12\int_0^T\|u_t\|^2dt\right)
\] 是真鞅且 \(EZ_T=1\)，例如满足适用的 Novikov 条件。用 \(d\mathbf Q=Z_Td\mathbf P^0\) 变换后，\(W^u=W^0-\int u\,dt\) 为 Brownian motion，\(X\) 的漂移变为 \(b^0+\sigma u\)。

\textbf{P-KL-05：可观测性修正的控制能量关系。} 若 \(E_{\mathbf Q}\int_0^T\|u_t\|^2dt<\infty\)，则完整概率空间上的 \[
D_{\rm KL}(\mathbf Q\|\mathbf P^0)
=\tfrac12E_{\mathbf Q}\int_0^T\|u_t\|^2dt.
\] 对仅观察 \(X\) 的路径律，一般只能保证 \[
D_{\rm KL}(Q_X\|P_X^0)
\le\tfrac12E_{\mathbf Q}\int_0^T\|u_t\|^2dt.
\]

\textbf{证明。} 将 \(dW^0=dW^u+u\,dt\) 代入 \(\log Z_T\)。在 \(\mathbf Q\) 下随机积分因平方可积均值为零，得第一式。第二式由投影到 \(X\) 路径的数据处理。\(\blacksquare\)

若 \(Z_T\) 对 \(X\) 路径可测，则信息未在投影中丢失，第二式取等号。一个常用充分设置是：扩散矩阵方阵可逆、驱动和控制均由可观测状态历史决定、相应 SDE 与测度变换良定义，使 Brownian 增量和密度比可从状态路径重构。

\textbf{N-KL-01。} 取二维噪声、\(\sigma=(1,0)\)、\(u=(0,1)\)。改变第二个噪声坐标不改变 \(X=W_1\) 的状态路径律，所以状态 KL 为零；控制能量却为 \(T/2\)。这排除了任意退化扩散中``状态路径 KL 恒等于全部控制能量''的表述。

Girsanov 与熵的精确条件可参见原始数学研究；本文使用上述有真鞅与可观测性限定的形式。\href{https://leonard.perso.math.cnrs.fr/papers/2012-Leonard-Girsanov\%20theory\%20under\%20a\%20finite\%20entropy\%20condition.pdf}{Léonard 原文}

与 P-KL-01 联合，即使只知上界仍有 \[
\tfrac12E_{\mathbf Q}\int\|u\|^2dt
\ge D_{\rm KL}(Q_X\|P_X^0)
\ge d_{\rm Bern}(Q_X(A)\|P_X^0(A)).
\] 因此持续事件的控制能量必要下界得到保留。它不是自动的热力学功界；能量单位、温度和耗散对应需独立物理实现。

\hypertarget{ux4fe1ux606fux4ee3ux4ef7ux7ed3ux8bbaux4e0eux8fdeux7eedux8fc7ux7a0bux9002ux914d}{%
\section{22.6 信息代价结论与连续过程适配}\label{ux4fe1ux606fux4ee3ux4ef7ux7ed3ux8bbaux4e0eux8fdeux7eedux8fc7ux7a0bux9002ux914d}}

有限路径和有限 CTMC 的概率构造、相对熵链与实际同源核可以在相应模块内部展开。连续扩散分支则显式保留 TerminalGirsanovData；恢复矩保留 \nolinkurl{DynkinExpectationCertificate}，HJB 保留 ItoRun。它们描述由合法过程定理应提供的测度和期望恒等式，不是自动完成了全部随机微积分。

StochasticSubstrate 将这些接口接入 P11，仍不把控制能量变成无条件热力学功。若密度不对观察路径可测，第 22.5 节只用数据处理上界；第 33 章的单位校准也必须另行给定。

第 39.9---39.15 节给出实际全时间 Gaussian／连续 Brownian 源、自然过滤条件矩及真正 predictable L² 终点积分。第 39.17、39.19 节进一步构造同一输入的实际连续积分过程，并从真实时间能量和条件二阶关系推导平方补偿鞅；Brownian 分割平方和也有原源网格误差界。自然过滤的指数密度已规范化并与实际 RN 导数相接。第 39.22 节现已从同一完整可预测标量输入推导确定性零网格分割的实际二次变差及实际加权一、二阶随机和。这些结果仍未合成已公开的 Taylor／Itô 来源，也未给出受控扩散换测度漂移或全部 HJB 前提；原过程、维度／过滤与模型类来源责任继续保留。

\hypertarget{ux3c0ux3c6ux7ea6ux675fux51e0ux4f55ux4e0eux53ccux9a71ux52a8ux7684ux4e25ux683cux7248ux672c}{%
\chapter{23. Π/Φ、约束几何与双驱动的严格版本}\label{ux3c0ux3c6ux7ea6ux675fux51e0ux4f55ux4e0eux53ccux9a71ux52a8ux7684ux4e25ux683cux7248ux672c}}

\hypertarget{ux8d44ux6e90ux4e0eux4fe1ux606fux7684ux5171ux540cux63a5ux53e3}{%
\section{23.1 资源与信息的共同接口}\label{ux8d44ux6e90ux4e0eux4fe1ux606fux7684ux5171ux540cux63a5ux53e3}}

原始 \(\Pi\) 表示资源、收益或功能效率，\(\Phi\) 表示预测不足、决策相关不确定性或信息成本。可采用 \[
J=\Pi-\lambda_\Phi\Phi-C.
\] 这是可用模板，具体量纲、符号、测量与比较方式由桥接决定。若把 \(\Phi\) 定为``信息收益''而非``信息缺口''，符号应相应改变。

资源约束与信息约束确实可以共同决定可行策略；这不自动证明所有系统的有效响应维数恰为二。

\hypertarget{ux89c4ux8303ux81eaux7531ux5ea6}{%
\section{23.2 规范自由度}\label{ux89c4ux8303ux81eaux7531ux5ea6}}

\textbf{P-DDH-01。} \[
\Pi'=\Pi+\lambda\chi,\qquad
\Phi'=\Phi+\chi
\] 保持 \(\Pi'-\lambda\Phi'=\Pi-\lambda\Phi\)。

\textbf{证明。} 直接代入相消。\(\blacksquare\)

因此行为最多识别被它约束的目标等价类。独立资源测量、预测实验或选择性干预用于锚定语义轴；任意旋转后的两个潜变量不能自动称为物理轴和信息轴。

\hypertarget{ux6709ux9650ux6307ux6570ux65cfux4e0eux7ea6ux675fux5bf9ux5076}{%
\section{23.3 有限指数族与约束对偶}\label{ux6709ux9650ux6307ux6570ux65cfux4e0eux7ea6ux675fux5bf9ux5076}}

给有限路径集合 \(\Gamma\)、严格正基线 \(p_0\)、预先指定特征 \(F:\Gamma\to\mathbb R^k\)，定义 \[
\psi(\lambda)=\log\sum_\gamma p_0(\gamma)e^{\lambda^\top F(\gamma)},
\quad
q_\lambda(\gamma)=p_0(\gamma)e^{\lambda^\top F(\gamma)-\psi(\lambda)}.
\]

\textbf{P-DDH-02。} \[
\nabla\psi=E_{q_\lambda}F,\qquad
\nabla^2\psi=\operatorname{Cov}_{q_\lambda}(F).
\]

\textbf{证明。} 有限和逐项求导；一阶导为归一化加权均值，二阶导为二阶矩减去均值外积。\(\blacksquare\)

\textbf{P-DDH-03：矩约束下的变分最小性。} 若目标矩 \(b=E_{q_\lambda}F\)，则 \(q_\lambda\) 在所有 \(q\) 满足 \(E_qF=b\) 中唯一最小化 \(D_{\rm KL}(q\|p_0)\)。

\textbf{证明。} \[
D(q\|p_0)=D(q\|q_\lambda)+\lambda^\top E_qF-\psi(\lambda).
\] 固定矩后，后两项相同，第一项非负且仅在 \(q=q_\lambda\) 为零。\(\blacksquare\)

边界矩不一定由有限 \(\lambda\) 达到，可需要参数极限或支撑缩减。只有独立规定特征与预算，才有可检验内容；事后选 \(F=\log(q/p_0)\) 会使任何观测律变成一次倾斜。

有限指数族的协方差与变分对偶是标准理论。\href{https://www.emerald.com/ftmal/article/1/1-2/1/1332382/Graphical-Models-Exponential-Families-and}{Wainwright--Jordan}

\hypertarget{ux5171ux540cux4e8cux7ef4ux74f6ux9888}{%
\section{23.4 共同二维瓶颈}\label{ux5171ux540cux4e8cux7ef4ux74f6ux9888}}

\textbf{H-DDH。} 在固定预算坐标 \(B\) 和所有待比较环境中，存在共同可微映射 \(z(B)\in\mathbb R^2\)，使 \[
m_e(B)=g_e(z(B)).
\]

\textbf{P-DDH-04。} 在同一预算点 \(B\)，各环境响应 Jacobian 纵向堆叠的秩不超过 2。

\textbf{证明。} 链式法则给 \(J_e=Dg_e\,Dz\)。堆叠后仍右乘同一个 \(2\times k\) 矩阵 \(Dz\)。\(\blacksquare\)

不同预算点的 \(Dz(B)\) 可旋转；跨点再堆叠需额外共同线性子空间条件。不同环境各自秩二，也不推出共同秩二： \(\operatorname{diag}(1,1,0)\) 与 \(\operatorname{diag}(0,1,1)\) 堆叠后秩三。

一个标量目标 \[
J(a,B)=-\tfrac12\|a-B\|^2
\] 给最优响应 \(a^*=B\)，维数可任意高。因此``目标数''``部门数''``潜变量维数''和``响应秩''须分开。

\hypertarget{ux951aux70b9ux4e0eux8c31ux8ba4ux8bc1}{%
\section{23.5 锚点与谱认证}\label{ux951aux70b9ux4e0eux8c31ux8ba4ux8bc1}}

若响应秩二分解为 \(J=AB^\top\)，任意 \(R\in GL(2)\) 给同一矩阵 \[
J=(AR)(BR^{-\top})^\top.
\] 给两条独立、预先定义且定向的干预轴后，允许基变换缩减为各轴正比例缩放。若还给量纲单位，则可进一步固定尺度。

\textbf{P-DDH-05：有效维数阈值。} 若 \(\|\widehat S-S\|_2\le\eta\)，则奇异值误差至多 \(\eta\)。若 \[
\sigma_3(S)+\eta<\tau<\sigma_2(S)-\eta,
\] 则估计中超过阈值 \(\tau\) 的奇异值恰有两个。

\textbf{证明。} 奇异值的变分表征给 Lipschitz 扰动界，代入阈值即可。\(\blacksquare\)

这认证的是阈值意义的有效维数。只有另有结构前提 \(\sigma_3=0\) 时，才可称为精确秩二模型的恢复。有限样本中第三奇异值未显著离零，不证明它真实等于零。

\hypertarget{ux3c0ux3c6-ux5982ux4f55ux8fdbux5165ux5e94ux7528}{%
\section{23.6 Π/Φ 如何进入应用}\label{ux3c0ux3c6-ux5982ux4f55ux8fdbux5165ux5e94ux7528}}

资源侧可使用任务收益、材料/能量预算、时间、控制幅度或计算成本；信息侧可使用预测 KL、所需记忆率、贝叶斯风险或策略相关不确定性。必须说明每个量怎样测量、是否可独立干预、怎样决定行为。

若信息价值已完整进入 belief-space 控制，不再机械追加一个信息增益奖励，除非它确实表示新的偏好或成本。第 27 章给出精确区分。

\hypertarget{ux5b8cux6574ux89c4ux8303ux523bux753bux4e0eux79e9ux68c0ux9a8cux7684ux72ecux7acbux6027}{%
\section{23.7 完整规范刻画与秩检验的独立性}\label{ux5b8cux6574ux89c4ux8303ux523bux753bux4e0eux79e9ux68c0ux9a8cux7684ux72ecux7acbux6027}}

对于固定 \(\lambda\)，本章 P-DDH-01 的保持方向已有逆向加强：若 \(\Pi'-\lambda\Phi'=\Pi-\lambda\Phi\)，取 \(\chi=\Phi'-\Phi\) 即得 \(\Pi'=\Pi+\lambda\chi\)。\nolinkurl{DualDriveGaugeCompleteness.same_objective_iff_gauge_related} 因此完整刻画同一数值目标的规范轨道，锚定识别的精确条件见第 33 章。

这个代数 iff 不证明自然系统存在共同二维响应瓶颈。C5 的共同秩必要性、奇异值误差和三值结论承担另一任务：兼容表示``该测试未排除''，不能当作潜变量唯一识别；两条独立物理轴还需独立测量和干预。

\hypertarget{ux591aux5c42ux7ea7ux4ef7ux503cux53efux79efux6027ux4e0eux534fux8c03ux673aux5236}{%
\chapter{24. 多层级价值、可积性与协调机制}\label{ux591aux5c42ux7ea7ux4ef7ux503cux53efux79efux6027ux4e0eux534fux8c03ux673aux5236}}

\hypertarget{ux591aux76eeux6807ux4e0eux5171ux540cux6700ux4f18}{%
\section{24.1 多目标与共同最优}\label{ux591aux76eeux6807ux4e0eux5171ux540cux6700ux4f18}}

给各评价层 \(J_1,\ldots,J_m\)，默认保留 Pareto 偏序。若 \[
\bigcap_{\ell=1}^m\mathfrak P_\ell^*\ne\varnothing,
\] 则存在共同最优策略；交为空说明不存在精确共同最优，需协商、约束或另加权重。

对两个策略的价值向量 \((1,0),(0,1)\)，若两层没有优先级且要求交换对称，任何对称标量化都不能唯一择一。额外权重是评价桥接，不由对象性推出。

\hypertarget{ux65b9ux5411ux573aux4ec0ux4e48ux65f6ux5019ux6709ux5168ux5c40ux52bf}{%
\section{24.2 方向场什么时候有全局势}\label{ux65b9ux5411ux573aux4ec0ux4e48ux65f6ux5019ux6709ux5168ux5c40ux52bf}}

在连通光滑流形上，给 \(C^1\) 一形式 \(\omega\)。

\textbf{P-ALI-01：全局可积判据。} \[
\omega=dV
\iff
\oint_C\omega=0
\quad\text{对每条分段光滑闭曲线 }C.
\]

\textbf{证明。} 正向由微积分基本定理。反向固定基点 \(x_0\)，定义 \(V(x)=\int_{\gamma_{x_0,x}}\omega\)。闭路零积分保证路径无关；在局部坐标中对端点求导给 \(dV=\omega\)。\(\blacksquare\)

局部闭合 \(d\omega=0\) 在有拓扑孔洞的空间不保证全局势。给定度量后漂移可转成一形式；该度量是模型的一部分，不具有任意坐标下无需运输的欧氏特权。

\textbf{K-GEO-01。} 紧、定向、无边界 Riemann 流形上的适当光滑一形式有 Hodge 分解 \(\omega=dV+\delta B+h\)。有边界时须声明边界条件。它将梯度、循环与调和部分分开；非梯度运动并不意味没有长期控制价值。

\hypertarget{ux5b50ux5c42ux66f4ux65b0ux4e0eux7236ux5c42ux4ef7ux503c}{%
\section{24.3 子层更新与父层价值}\label{ux5b50ux5c42ux66f4ux65b0ux4e0eux7236ux5c42ux4ef7ux503c}}

在实 Hilbert 直和策略空间中，父价值 \(J_P\)、子价值 \(J_i\) 可微。令 \[
g_i=\nabla_iJ_P,\quad
e_i=\nabla_iJ_i-g_i,\quad
\dot a_i=\alpha_i\nabla_iJ_i,\quad\alpha_i>0.
\]

\textbf{P-ALI-02。} \[
\frac{dJ_P}{dt}
=\sum_i\alpha_i(\|g_i\|^2+\langle g_i,e_i\rangle)
\ge\sum_i\alpha_i\|g_i\|(\|g_i\|-\|e_i\|).
\]

\textbf{证明。} 链式法则代入更新式，最后用 Cauchy--Schwarz。\(\blacksquare\)

若全部非零父梯度分量满足 \(\|e_i\|<\|g_i\|\)，且至少一项非零，父价值严格上升。这个瞬时充分条件不保证全局最优收敛。

\hypertarget{ux534fux8c03ux9608ux503c}{%
\section{24.4 协调阈值}\label{ux534fux8c03ux9608ux503c}}

将联合子场 \(G\) 改为 \(G_\eta=(1-\eta)G+\eta g\)，\(g=\nabla J_P\)，\(0\le\eta\le1\)。

\textbf{P-ALI-03。} 若 \(\langle g,G\rangle<0\)，父价值的瞬时导数转正恰需 \[
\eta>
\eta_*=\frac{-\langle g,G\rangle}
{\|g\|^2-\langle g,G\rangle}.
\]

\textbf{证明。} \(\dot J_P=(1-\eta)\langle g,G\rangle+\eta\|g\|^2\)，解一次不等式。\(\blacksquare\)

这给可解释的机制设计模板。税收、激励、权限、通信或反馈控制是否实现这种场混合，需要单独建模；公式本身不构成某种制度的规范正当性。

\hypertarget{ux8bc4ux4ef7ux5c42ux7ea7ux4e0eux5176ux5b83ux5c42ux7ea7ux7684ux4ea4ux53c9ux6865}{%
\section{24.5 评价层级与其它层级的交叉桥}\label{ux8bc4ux4ef7ux5c42ux7ea7ux4e0eux5176ux5b83ux5c42ux7ea7ux7684ux4ea4ux53c9ux6865}}

H0-VAL 与预测商、物理分辨率和选择层级不同。本章父价值的梯度、子更新场和权重必须在同一个策略空间中比较；仅有父子对象映射不能自动生成内积、价值或 \(\eta_*\)。

P-ALI-01 的全局周期障碍尤其不能被局部``驱动方向''术语掩盖。数学上已给的漂移可来自梯度、循环与调和部分；哲学上要解释哪个部分与目的有关，还须有明确的评价机制。第 34 章保留这种忠实性端口。

\hypertarget{ux7e41ux6b96ux9057ux4f20ux4e0eux6b63ux7b97ux5b50ux6f14ux5316}{%
\chapter{25. 繁殖、遗传与正算子演化}\label{ux7e41ux6b96ux9057ux4f20ux4e0eux6b63ux7b97ux5b50ux6f14ux5316}}

\hypertarget{ux4eceux5bf9ux8c61ux5230ux540eux4ee3ux673aux5236}{%
\section{25.1 从对象到后代机制}\label{ux4eceux5bf9ux8c61ux5230ux540eux4ee3ux673aux5236}}

繁殖不是对象定义中的必要条件。进入本章前，父代和候选后代应已有第 12 章对象证书，以及第 11 章的历史关系；材料脱落不自动构成繁殖。固定有限类型集 \(\{1,\ldots,d\}\)、一个世代长度和计数单位。类型 \(i\) 的个体产生随机后代计数向量 \(N^{(i)}\in\mathbb N^d\)，定义 \[
M_{ij}=E N_j^{(i)}<\infty,\qquad b_i=\sum_jM_{ij}.
\] \(M\) 是平均繁殖算子，\(b_i\) 是平均后代总数。若 \(b_i>0\)，令 \(K_{ij}=M_{ij}/b_i\)，于是 \[
M=D_bK.
\] 这唯一分离了繁殖强度和\textbf{按后代强度抽样}的类型传递。若 \(b_i=0\)，该行没有后代，\(K_{i\cdot}\) 无法由观测定义；任取概率行不改变 \(M\)，但不得声称识别了它。

注意 \[
K_{ij}=\frac{EN_j^{(i)}}{E\sum_kN_k^{(i)}}
\] 通常不等于先均匀抽一个非空家庭、再在其家庭内均匀抽孩子所得的 \(E[N_j/\sum_kN_k\mid\sum_kN_k>0]\)。必须固定抽样单位，才谈得上传递概率和遗传信息。

\hypertarget{ux5e73ux5747ux9891ux7387ux52a8ux529bux5b66ux4e0e-price-ux6052ux7b49ux5f0f}{%
\section{25.2 平均频率动力学与 Price 恒等式}\label{ux5e73ux5747ux9891ux7387ux52a8ux529bux5b66ux4e0e-price-ux6052ux7b49ux5f0f}}

确定性的平均强度模型令类型频率为行向量 \(p\)，且 \(\bar b=pb>0\)，则 \[
p'=\frac{pM}{\bar b}.
\] 这不是一般有限随机种群中 \(E[Z_{n+1}/|Z_{n+1}|]\) 的精确公式；随机比值与期望比值必须分开。

\textbf{P-EVO-01：选择与传递的精确分解。} 对父类型数值性状 \(z_i\) 和后代类型性状 \(z'_j\)，令 \((Kz')_i=\sum_jK_{ij}z'_j\)。则 \[
p'z'-pz=
\frac{\operatorname{Cov}_p(b,z)}{\bar b}
+\frac{\sum_i p_ib_i[(Kz')_i-z_i]}{\bar b}.
\]

\textbf{证明。} \(p'z'=\sum_i p_ib_i(Kz')_i/\bar b\)。加减 \(\sum_i p_ib_iz_i/\bar b\)，余下第一项为 \((E_p[bz]-E_p[b]E_p[z])/\bar b\)。零繁殖行乘以零，不依赖未识别的 \(K\)。\(\blacksquare\)

这是代数恒等式，不是一个独立的因果选择模型。因果解释还要求 \(b,K\) 在相应干预下具有机制意义。

\hypertarget{ux9057ux4f20ux4fe1ux606fux4e0eux53efux9057ux4f20ux9884ux6d4bux6838}{%
\section{25.3 遗传信息与可遗传预测核}\label{ux9057ux4f20ux4fe1ux606fux4e0eux53efux9057ux4f20ux9884ux6d4bux6838}}

单纯增长不能推出遗传：若所有 \(K_{i\cdot}\) 相同，后代类型与亲本类型独立，即使 \(b_i\) 很大也不改变这一事实。在指定亲子联合分布下，可报告 \(I(F;F')\)，但它依赖类型划分。

\textbf{P-EVO-02：无关标签会膨胀遗传信息。} 若有限标签 \(T\) 独立于 \((F,F')\)，并被无差错复制，则 \[
I((F,T);(F',T))=I(F;F')+H(T).
\]

\textbf{证明。} 用链式法则展开；独立性消去交叉项，共享标签贡献 \(I(T;T)=H(T)\)。\(\blacksquare\)

因此本核心优先使用与后代行为相联系的规范对象：在预先固定的谱系抽样机制、观察目标和可数干预族下，将亲本历史映射到全部后代未来的条件分布族。按 P-PRED-01，这个\textbf{谱系预测状态}是该预测问题的最小充分表示。它排除对指定未来无预测贡献的标签，而不是宣布标签在所有其他任务中都无意义。

谱系抽样可以是固定一条继承规则的随机谱系，也可以是完整后代树。两者不是同一未来空间，不能混用其熵或增长率。

\hypertarget{perron-ux589eux957fux4e0eux7ec4ux6210}{%
\section{25.4 Perron 增长与组成}\label{perron-ux589eux957fux4e0eux7ec4ux6210}}

\textbf{K-PF-01（复用并展开）：有限本原非负矩阵定理。} 若某个正整数幂 \(M^k\) 的所有元素严格正，则存在 \(R>0\)、正列向量 \(r\)、正行向量 \(l\)，满足 \[
Mr=Rr,\quad lM=Rl,\quad l\mathbf1=1,\quad lr=1.
\] \(R\) 是单重主特征值，其他特征值模严格小于 \(R\)，且 \(R^{-n}M^n\to rl\)。本章仅在此明确有限本原假设下调用 Perron--Frobenius 定理。

\textbf{P-EVO-03：平均增长和繁殖价值。} 对非零非负初始计数行向量 \(z_0\)， \[
R^{-n}z_0M^n\to(z_0r)l,\qquad
\frac{z_0M^n}{z_0M^n\mathbf1}\to l.
\] 若 \(w>0\) 满足对每个非负 \(z\)，\((zM)w=\rho zw\)，则 \(\rho=R\)，且 \(w\) 是 \(r\) 的正倍数。

\textbf{证明。} 第一部分直接左乘 K-PF-01 的矩阵极限；分母极限为 \(z_0r>0\)。第二部分取各坐标单位向量得 \(Mw=\rho w\)，再由正特征向量的 Perron 唯一性得结论。\(\blacksquare\)

右向量 \(r\) 表示线性计数评价中的相对繁殖价值，左向量 \(l\) 表示长期平均类型组成。这里``规范''是相对于指定正线性繁殖机制的唯一性，不是所有对象或所有价值的普适唯一性。

\hypertarget{ux968fux673aux79cdux7fa4ux4e0eux5e73ux5747ux6a21ux578bux7684ux8fb9ux754c}{%
\section{25.5 随机种群与平均模型的边界}\label{ux968fux673aux79cdux7fa4ux4e0eux5e73ux5747ux6a21ux578bux7684ux8fb9ux754c}}

设各个体的后代在给定类型后独立，繁殖律跨世代固定，\(Z_n\) 为真实种群计数。于是 \[
E[Z_{n+1}\mid\mathcal F_n]=Z_nM.
\]

\textbf{P-EVO-04：繁殖价值鞅。} 有限初始计数下， \[
W_n=R^{-n}Z_nr
\] 是非负可积鞅，且 \(EW_n=Z_0r\)。

\textbf{证明。} 条件期望为 \(R^{-(n+1)}Z_nMr=R^{-n}Z_nr\)。非负性和有限一阶矩给可积性。\(\blacksquare\)

\textbf{K-MAR-01。} 非负鞅几乎处处收敛到有限 \(W_\infty\)。但不由此得到 \(EW_\infty=EW_0\)、\(L^1\) 收敛，或``未灭绝时 \(W_\infty>0\)''。这些结论需要额外的一致可积性和繁殖矩条件。本版不把相应分支过程深层定理作为无条件核心。随机灭绝、有限种群漂变和确定性平均增长必须分别报告。

\hypertarget{ux4ee5ux540cux4e00ux7e41ux6b96ux7b97ux5b50ux6865ux63a5ux5b9eux9645ux8c31ux7cfb}{%
\section{25.6 以同一繁殖算子桥接实际谱系}\label{ux4ee5ux540cux4e00ux7e41ux6b96ux7b97ux5b50ux6865ux63a5ux5b9eux9645ux8c31ux7cfb}}

LineageEvolution 的 lineage\_price\_bridge 消费与 Price 代数相同的传递核；lineage\_perron\_bridge 消费与 P-EVO-03 相同的矩阵及 KPF01Certificate，并要求其每条正繁殖边具有实际子代语义。这样谱结论不会来自与谱系记录无关的另一矩阵。

修复同一个身份与产生一个不同身份的后代被形式化为互斥角色。谱的存在性、增长与随机种群结果仍保持本章条件；``任意客体都自发形成满足 Perron 假设的谱系''不在结论中。

\hypertarget{ux4e2aux4f53ux6027ux8f6cux53d8ux4e0eux5185ux751fux76eeux7684ux6865}{%
\chapter{26. 个体性转变与内生目的桥}\label{ux4e2aux4f53ux6027ux8f6cux53d8ux4e0eux5185ux751fux76eeux7684ux6865}}

\hypertarget{ux5bf9ux8c61ux5c3aux5ea6ux4e0dux7b49ux4e8eux9009ux62e9ux5c42ux7ea7}{%
\section{26.1 对象尺度不等于选择层级}\label{ux5bf9ux8c61ux5c3aux5ea6ux4e0dux7b49ux4e8eux9009ux62e9ux5c42ux7ea7}}

对象尺度 \(s\) 描述候选父子分解和观察接口；选择层级 \(\ell\) 描述什么单位被计数、继承和差异复制。两者可一致，也可不一致。例如一个协作群体可在有限时域构成父对象，却没有群体层面的复制；复制子可以借用另一个对象的维持机制。

本核心采用四项\textbf{操作性 ETI 证书}，并不宣称它是所有自然个体性转变的逻辑充要定义：

\begin{enumerate}
\def\labelenumi{\arabic{enumi}.}
\tightlist
\item
  父层具有预测、持续与切割不可约证书。
\item
  子对象到父层有有效的实现及分辨率接口。
\item
  父层单位存在可识别的复制与传递机制。
\item
  群体内冲突、子层保真度和父层复制优势落在预先给定的允许域内。
\end{enumerate}

若只满足前两项，称父层对象形成，不直接称为演化个体性转变。若只满足第三项，不能据此宣告群体是不可约父对象。

\hypertarget{ux673aux5236ux5185ux5316ux4e0eux811aux624bux67b6ux64a4ux9664}{%
\section{26.2 机制内化与脚手架撤除}\label{ux673aux5236ux5185ux5316ux4e0eux811aux624bux67b6ux64a4ux9664}}

将维持机制分成候选内部 \(I\) 与外部脚手架 \(S\)。在相同资源预算、相同可比环境协议下，实施明确的脚手架撤除或替代干预，并重新计算对象证书和持续性。

``移除 \(S\) 后失效''证明的是\textbf{该干预和替代集下}的依赖；它未证明不存在任何其他外部实现。``撤除后仍成立''说明该脚手架不是此证书的必要条件，也不自动证明全部因果闭合。应同时报告撤除操作改变了哪些因果边和资源边，避免把断电与组织解体混为一谈。

\hypertarget{ux9009ux62e9ux5982ux4f55ux8bf1ux5bfcux63a7ux5236ux76eeux6807}{%
\section{26.3 选择如何诱导控制目标}\label{ux9009ux62e9ux5982ux4f55ux8bf1ux5bfcux63a7ux5236ux76eeux6807}}

设可遗传控制器类型为 \(\theta\)，在环境 \(e\) 中实施策略 \(\pi_\theta\)，生成路径律 \(Q_\theta^e\)。给定繁殖机制 \[
R_e(\theta)=\mathcal R_e(Q_\theta^e)>0.
\] 这里从行为到复制的 \(\mathcal R_e\) 是\textbf{目的桥}。它可以依据生存、产出或其他明确机制定义；对象证书本身并不唯一指定它。

\textbf{P-BRG-01：无突变的固定适合度集中。} 有限类型、固定 \(R_i>0\) 下， \[
p_i(n+1)=\frac{p_i(n)R_i}{\sum_jp_j(n)R_j}
\quad\Longrightarrow\quad
p_i(n)=\frac{p_i(0)R_i^n}{\sum_jp_j(0)R_j^n}.
\] 令 \(R_*=\max_{i:p_i(0)>0}R_i\)。低于 \(R_*\) 的类型总质量指数趋零，最大者之间保持初始相对比例。

\textbf{证明。} 归纳得到显式公式。分子分母除以 \(R_*^n\)，次优项以 \((R_i/R_*)^n\) 衰减，最优项不变。\(\blacksquare\)

因此演化在这些条件下诱导最大化 \(R\) 的行为筛选；它不会创造初始支持外的控制器。若引入突变核 \(K\)，长期组成由 \(D_RK\) 的 Perron 结构决定，而非通常集中在单一最大 \(R\) 类型。

\hypertarget{ux4e0dux5b58ux5728ux65e0ux6761ux4ef6ux7684ux6f14ux5316ux5373ux6700ux4f18ux5316}{%
\section{26.4 不存在无条件的``演化即最优化''}\label{ux4e0dux5b58ux5728ux65e0ux6761ux4ef6ux7684ux6f14ux5316ux5373ux6700ux4f18ux5316}}

\textbf{N-BRG-01：频率依赖循环。} 在内部单纯形上取反对称石头---剪刀---布矩阵 \(A\)，行列和均为零， \[
\dot p_i=p_i(Ap)_i.
\] 因为 \(p^\top Ap=0\)，这是复制动力学；并且 \[
\frac{d}{dt}\log(p_1p_2p_3)=\mathbf1^\top Ap=0.
\] 非中心内部轨道不可能趋向边界或唯一内部平衡，因为其乘积为固定的正数且小于中心值；因此``所有选择轨道都收敛到一个最优点''不成立。反对称性也阻止将该动力学一般解释为单一严格上升势函数流。

若两个环境交换 \(R_1,R_2\) 的大小，就不存在由这两个固定适合度排序共同指定的唯一最优类型。跨环境泛化必须明确环境分布、最坏情况目标或适应机制。

\textbf{P-BRG-02：行为等价的选择不可分辨性。} 在同一环境、复制完全通过 \(Q_\theta^e\) 决定且无类型依赖突变或其他复制通道时，若 \(Q_\theta^e=Q_{\theta'}^e\)，则 \(R_e(\theta)=R_e(\theta')\)，二者不能仅由该复制桥被差异筛选。

\textbf{证明。} 将相同自变量代入同一个 \(\mathcal R_e\)。\(\blacksquare\)

这是一条接口一致性结果：给控制器附加不同语义标签，不会自动改变已声明的选择动力学。引入不同突变邻域或材料成本后，模型改变，应重新检验。

\hypertarget{ux5185ux751fux76eeux7684ux7684ux95edux73afux67b6ux6784}{%
\section{26.5 内生目的的闭环架构}\label{ux5185ux751fux76eeux7684ux7684ux95edux73afux67b6ux6784}}

完整闭环是时间递推： \[
\theta_n \to \pi_{\theta_n}\to Q_{\theta_n}^{e_n}
\to\text{复制/学习反馈}\to\theta_{n+1}.
\] 每一箭头需要独立机制。得到特定环境中生存或增长的解释，可以加强原论文关于功能和目的的论证；它仍不等于从``存在''演绎出唯一的道德价值。跨层价值冲突按第 24 章保留向量或明确协调机制。

\hypertarget{ux6761ux4ef6ux6784ux9020ux95edux73afux4e2dux7684ux6765ux6e90ux95eeux9898}{%
\section{26.6 条件构造闭环中的来源问题}\label{ux6761ux4ef6ux6784ux9020ux95edux73afux4e2dux7684ux6765ux6e90ux95eeux9898}}

P7/P12 从包含物理载体和程序来源的 OntogeneticSeed 出发，经过可信底座构造组织、进入恢复与维持过程；现有 no-go 说明相同物理部分而不同程序来源无法由物理部分单独统一恢复。这个边界与``后代按哪个评价被选择''又不同，后者需要复制---目的机制桥。

因此内生可以表示机制在明确时间过程中逐步被当前系统维持或传递；它不能表示证明中出现过的程序、解释器、目标和环境规律从未需要来源。第 34.5 节把已经连接的生命周期及保留前提完整列出。

\hypertarget{ux53cdux8eabux7ed3ux6784ux4fe1ux5ff5ux5b66ux4e60ux4e0eux80fdux52a8ux6027}{%
\chapter{27. 反身结构、信念学习与能动性}\label{ux53cdux8eabux7ed3ux6784ux4fe1ux5ff5ux5b66ux4e60ux4e0eux80fdux52a8ux6027}}

\hypertarget{ux628aux53efux4feeux6539ux7ed3ux6784ux7eb3ux5165ux72b6ux6001}{%
\section{27.1 把可修改结构纳入状态}\label{ux628aux53efux4feeux6539ux7ed3ux6784ux7eb3ux5165ux72b6ux6001}}

令物理状态为 \(X_t\)，控制器/模型结构为 \(S_t\)。若 \[
P(X_{t+1},S_{t+1}\mid H_t,a_t)
=K(X_t,S_t,a_t;\cdot),
\] 则 \(Z_t=(X_t,S_t)\) 是受控 Markov 状态。自修改并不在数学上排斥状态空间描述，但可能使状态空间很大或不可有限压缩。

\textbf{P-REF-01：反身状态增广。} 在上述条件和可测核条件下，给定任意因果策略，\(Z\) 的路径律由初始律及 \(K\) 唯一确定；确定性修改 \(S_{t+1}=F(S_t,X_t,a_t,\xi_t)\) 是其特例。

\textbf{证明。} 条件独立给出一步核分解，再用第 2 章核扩张。若压缩为 \(\phi(Z_t)\)，仍须重新满足第 7 章的受控闭合条件；反身性不提供自动闭合。\(\blacksquare\)

\hypertarget{ux6709ux9650ux9690ux72b6ux6001ux6a21ux578bux7684-bayes-ux63a5ux53e3}{%
\section{27.2 有限隐状态模型的 Bayes 接口}\label{ux6709ux9650ux9690ux72b6ux6001ux6a21ux578bux7684-bayes-ux63a5ux53e3}}

取静态未知参数 \(\theta\in\Theta\)、隐状态 \(x\in\mathcal X\)，两者有限。已知模型族为 \(P_\theta(x'|x,a)\) 和观测核 \(O_\theta(y|x',a)\)。当前信念 \(b_t(\theta,x)=P(\theta,X_t=x\mid\text{已见历史})\)。则 \[
b_{t+1}(\theta,x')=
\frac{O_\theta(y_{t+1}|x',a_t)
\sum_xP_\theta(x'|x,a_t)b_t(\theta,x)}
{\sum_{\vartheta,z}O_\vartheta(y_{t+1}|z,a_t)
\sum_xP_\vartheta(z|x,a_t)b_t(\vartheta,x)}.
\] 分母为零的观测在当前模型下不可能：接口应返回模型冲突，不能任意填一个后验并继续声称已认证。

\textbf{P-REF-02：信念充分性。} 上述模型下，下一观测的条件律和更新信念的条件律只依赖 \(b_t,a_t\)。有界折扣任务因此可转写为信念状态控制问题。

\textbf{证明。} 对隐状态和参数求全概率得到分母中的预测观测律；Bayes 公式给更新。下一信念是当前信念、动作和观测的确定函数，故其核仅依赖 \((b_t,a_t)\)。有界奖励的条件期望同理。\(\blacksquare\)

这只是模型内充分性。若真实机制不在模型族、参数实际漂移，或修改动作改变了 \(P_\theta,O_\theta\)，必须增广相应机制状态，或以第 16 章误差接口处理。

\hypertarget{ux4fe1ux606fux4ef7ux503cux4e0dux662fux4fe1ux606fux589eux76ca}{%
\section{27.3 信息价值不是信息增益}\label{ux4fe1ux606fux4ef7ux503cux4e0dux662fux4fe1ux606fux589eux76ca}}

\textbf{P-REF-03：免费信息的决策价值非负。} 在一次决策问题中，信号 \(Y\) 不改变世界分布或动作集、无成本，且允许忽略它。对有限动作集和可积收益， \[
E_Y\max_a E[u(a,\theta)\mid Y]
\ge \max_a E[u(a,\theta)].
\]

\textbf{证明。} 固定任意动作 \(a\)，逐个 \(Y\) 的最大值不小于该动作的条件收益。取期望，再对固定动作取最大。\(\blacksquare\)

若信号揭示一个与所有行动收益无关的独立硬币，互信息可以为 \(\log2\)，决策价值仍为零。若信息有成本、改变物理状态或妨碍行动，净价值也可能为负。因此不能把 \(\Phi\) 无条件等同于``最大化任意可获得信息''。

\hypertarget{ux76eeux6807ux4feeux8ba2ux4e0eux8befux6307ux5b9aux98ceux9669}{%
\section{27.4 目标修订与误指定风险}\label{ux76eeux6807ux4feeux8ba2ux4e0eux8befux6307ux5b9aux98ceux9669}}

\textbf{P-REF-04：统一目标误差的后悔界。} 若同一策略集上 \[
\sup_\pi|J_g(\pi)-J_{\hat g}(\pi)|
\le Ld(g,\hat g),
\] 且 \(\pi_g,\pi_{\hat g}\) 分别最优，则 \[
J_g(\pi_g)-J_g(\pi_{\hat g})\le2L d(g,\hat g).
\]

\textbf{证明。} 加减 \(J_{\hat g}(\pi_g),J_{\hat g}(\pi_{\hat g})\)，中间项由最优性非正，其余各不超过 \(Ld\)。\(\blacksquare\)

该结论需要对全部相关策略的统一界。它不把规范分歧化为已测出的距离，也不保证不可逆行为的代价很小。若奖励逐步误差为 \(\epsilon\)，折扣回报可取 \(Ld=\epsilon/(1-\beta)\)。

\hypertarget{ux80fdux52a8ux6027ux7684ux64cdux4f5cux6027ux5206ux91cf}{%
\section{27.5 能动性的操作性分量}\label{ux80fdux52a8ux6027ux7684ux64cdux4f5cux6027ux5206ux91cf}}

在固定任务和初始条件下，可分别测量：

\begin{itemize}
\tightlist
\item
  可实施策略集合及其对未来分布的影响；
\item
  识别环境并使用反馈的能力；
\item
  维持自身实现、模型与权限边界的能力；
\item
  扩展或修订控制结构的能力；
\item
  在代价和安全约束下恢复、回滚和停止的能力。
\end{itemize}

\textbf{P-REF-05：可行选择扩展的单调性。} 若同一评价 \(J\) 下 \(\Pi_1\subseteq\Pi_2\)，则 \(\sup_{\Pi_2}J\ge\sup_{\Pi_1}J\)。

\textbf{证明。} 在较大的集合上取上确界。\(\blacksquare\)

如果新增能力强制占用资源或移除旧策略，包含关系可能不成立，此定理不适用。权限扩张不是免费的自由增益。

对于维持能力，可比较匹配干预下 \[
\Delta_{\rm self}(T)
=P^{\rm active}(\tau>T)-P^{\rm disabled}(\tau>T).
\] 只有两种干预除所检验机制外具有可比资源和外部条件，才有相应因果解释；它与第 22 章维持信息代价是不同量。

\hypertarget{ux610fux4e49ux7684ux4e09ux5c42ux63a5ux53e3}{%
\section{27.6 意义的三层接口}\label{ux610fux4e49ux7684ux4e09ux5c42ux63a5ux53e3}}

功能意义：一个机制对对象证书、持续性或任务价值的反事实贡献。意向/表征意义：内部表示与环境变量及行动后果的协议化关系。规范意义：哪些结果值得追求，须由额外的价值桥和主体间约束说明。

本核心可以形式化前两层的部分可观测关系，并表示第三层的约束；不能由预测压缩、复制优势或一条梯度公式推出体验、意识或唯一的终极价值。将这些边界明确化，才能使统一架构保留可验证内容。

\hypertarget{ux4e00ux822cux4fe1ux5ff5ux9012ux63a8ux4e0eux53efux5b9eux65bdux7b56ux7565ux5b58ux5728}{%
\section{27.7 一般信念递推与可实施策略存在}\label{ux4e00ux822cux4fe1ux5ff5ux9012ux63a8ux4e0eux53efux5b9eux65bdux7b56ux7565ux5b58ux5728}}

C1ParametricBelief 在标准 Borel 潜在空间及明确的可数生成条件下构造联合可测的参数化后验，再将其与既有实际联合律对齐。固定参数下有条件分布，不足以保证同时依赖全部信念和动作的选择已经可测；C1 专门补的是这个接口。

有限支持控制的策略存在还需另外的可行性递推。先固定一个策略验证其实际支持，是证明该策略安全；把这个支持当作不变输入去另选策略可能改变未来可达集。第 37 章的后向递推加前向实际支撑证明避免这种循环，且没有依赖未知最优策略的支持假设。

\hypertarget{ux5b8cux6574ux63a5ux53e3ux7b97ux6cd5ux4e0eux7ec4ux5408ux5951ux7ea6}{%
\chapter{28. 完整接口、算法与组合契约}\label{ux5b8cux6574ux63a5ux53e3ux7b97ux6cd5ux4e0eux7ec4ux5408ux5951ux7ea6}}

\hypertarget{ux6a21ux5757ux5951ux7ea6ux603bux8868}{%
\section{28.1 模块契约总表}\label{ux6a21ux5757ux5951ux7ea6ux603bux8868}}

接口既是数学类型，也是实现时必须提供的信息。没有被提供的前提不得由下游默认为真。

\begin{longtable}[]{@{}
  >{\raggedright\arraybackslash}p{(\columnwidth - 6\tabcolsep) * \real{0.2500}}
  >{\raggedright\arraybackslash}p{(\columnwidth - 6\tabcolsep) * \real{0.2500}}
  >{\raggedright\arraybackslash}p{(\columnwidth - 6\tabcolsep) * \real{0.2500}}
  >{\raggedright\arraybackslash}p{(\columnwidth - 6\tabcolsep) * \real{0.2500}}@{}}
\toprule\noalign{}
\begin{minipage}[b]{\linewidth}\raggedright
模块
\end{minipage} & \begin{minipage}[b]{\linewidth}\raggedright
必需输入
\end{minipage} & \begin{minipage}[b]{\linewidth}\raggedright
输出及证书
\end{minipage} & \begin{minipage}[b]{\linewidth}\raggedright
主要失效状态
\end{minipage} \\
\midrule\noalign{}
\endhead
\bottomrule\noalign{}
\endlastfoot
Process & 状态/历史空间、初始律、时间、可测转移/观测核 & 协议下的路径律 & 无可测核、支持不一致 \\
Protocol & 干预时序、可用信息、动作约束、环境版本 & 合法因果协议族 & 使用未来信息、不可实施动作 \\
Factor & 候选读出、所在世界、预测目标、时域 & 因子及其相对预测状态 & 只比较标签、改变了未来目标 \\
Interface & 内外部变量、干预语义、条件核 & 联合残差、分离闭合证书 & 缺失协同信息、缺少支持覆盖 \\
Dynamics & 动作一致性、商映射、推送核 & 精确或近似闭合 & 历史依赖未进入状态 \\
Persistence & 评估域、恢复集、初始律、退出机制 & 不变/恢复/生存/QSD 剖面 & 不满足紧性、吸收或谱条件 \\
Continuation & 历史边、发展运输、容许误差 & 有向身份路径、保真剖面 & 把相似性当谱系、运输未校准 \\
Realization & 物理读出集、访问操作、毁损操作 & 载体族、访问 blocker、机制失效族 & 将读出删除混同物理删除 \\
Resolution & 两尺度域、映射、闭包、可见块 & 运输规则、IR--UV 区间 & 不可测商、非自然删失 \\
Composition & 候选父未来、切割族、条件环境 & 父层不可约性和子层保真度 & 共同驱动、切割改变目标 \\
Statistics & 抽样设计、置信度、覆盖、模型类别 & 误差事件或未决结构族 & 数据依赖遗漏、分布外部署 \\
Control & 可行策略、奖励/风险、时域、折扣 & 价值、GOD、策略误差 & 未声明目的桥、平均界误用 \\
GOA & 已指定最优策略族、闭环模型 & 长期集合/分布及稳定性 & 近分岔、无混合、策略非唯一 \\
Evolution & 个体单位、复制律、抽样、突变 & 平均算子、遗传核、增长结构 & 混用平均与随机比值 \\
Learning & 模型族、观测核、结构更新、成本 & 信念更新和可行策略扩展 & 零似然、误指定、越权修改 \\
\end{longtable}

各模块允许不同数学实现。有限表格、随机微分方程、路径核或可执行模拟器都是可能后端；``同一接口''不要求``同一物理本体''。

\hypertarget{ux5bf9ux8c61ux8bc1ux4e66ux7684ux6570ux636eux7ed3ux6784}{%
\section{28.2 对象证书的数据结构}\label{ux5bf9ux8c61ux8bc1ux4e66ux7684ux6570ux636eux7ed3ux6784}}

以下是规范性结构示例，不是声称已有通用自动证明软件。数值项必须携带单位、度量、量词与适用域。

\begin{verbatim}
{
  "schema": "UEOT-Core/ObjectCertificate/3.0",
  "object_id": "example-object",
  "world_version": "specified-model-v1",
  "factor": {
    "readout": "f",
    "target": "factor-relative-future",
    "horizon": 8,
    "protocol_family": "declared-family"
  },
  "claims": [
    {
      "id": "predictive-closure",
      "status": "certified-conditional",
      "metric": "TV",
      "quantifier": "uniform-on-declared-support",
      "bound": 0.01,
      "failure_probability": 0.05,
      "assumptions": ["finite-state", "independent-repeated-sampling"],
      "evidence": ["model-id", "data-id", "proof-id"]
    }
  ],
  "continuation": {"incoming_edges": [], "transport_model": "declared"},
  "realization": {"universe": ["r1", "r2", "r3"], "operation": "access"},
  "persistence": {"type": "finite-horizon-survival", "initial_law": "mu"},
  "purpose_bridge": null,
  "unresolved": ["physical-deletion-equivalence"]
}
\end{verbatim}

推荐状态集合为：defined、proved-conditional、certified-conditional、estimated、hypothesized、refuted、unresolved、out-of-scope。数值为零的估计不能自动升级为 proved-conditional。默认不把 null 的目的桥补成生存最大化。

每条证书应保留：模型版本、数据范围、协议支持、证明编号、误差来源、有效期限或漂移检查条件。实例标识是追踪工具，不是身份的数学证明。

\hypertarget{ux8defux5f84ux63a5ux53e3ux7684ux53efux7ec4ux5408ux6027}{%
\section{28.3 路径接口的可组合性}\label{ux8defux5f84ux63a5ux53e3ux7684ux53efux7ec4ux5408ux6027}}

一个从模型 \(A\) 到模型 \(B\) 的过程接口由以下部分组成：

\begin{itemize}
\tightlist
\item
  协议提升 \(J_{AB}:\mathcal U_B\to\mathcal U_A\)；
\item
  可测路径读出 \(C_{AB}:\Omega_A\to\Omega_B\)；
\item
  在指定初始律和全部声明协议下的自然性： \[
  (C_{AB})_\#P_A^{J_{AB}u}=P_B^u.
  \]
\end{itemize}

\textbf{P-API-01：精确与近似过程接口的复合。} 若 \(A\to B\) 和 \(B\to C\) 具有上述类型，则 \[
J_{AC}=J_{AB}\circ J_{BC},\qquad
C_{AC}=C_{BC}\circ C_{AB}.
\] 精确接口复合仍精确。若两接口在全部相应协议下的 TV 缺陷分别不超过 \(\epsilon_{AB},\epsilon_{BC}\)，则复合缺陷不超过 \(\min(1,\epsilon_{AB}+\epsilon_{BC})\)。

\textbf{证明。} 精确情形依次代入两次自然性。近似情形在复合推送与目标之间插入 \((C_{BC})_\#P_B^{J_{BC}u}\)，用三角不等式及 P-MET-01。\(\blacksquare\)

单位接口取恒等映射；复合结合律来自函数复合。时间粗化、读出分辨率和协议限制可以落在这一形式中，但必须各自显式给出路径与协议映射。任意``从一个领域到另一个领域的相似类比''不是过程接口。

\hypertarget{ux63a7ux5236ux63a5ux53e3ux9700ux8981ux989dux5916ux7ed3ux6784}{%
\section{28.4 控制接口需要额外结构}\label{ux63a7ux5236ux63a5ux53e3ux9700ux8981ux989dux5916ux7ed3ux6784}}

为了把过程接口提升为控制接口，还须提供动作对应、可实施策略提升、初始条件对应以及收益/约束运输。仅有观测路径自然性不能推出所有细层策略都被粗层表示，也不能推出最优值相等。

精确商的标准契约是第 20 章：相同纤维的可行动作、奖励和商转移一致。若采用动作映射，必须逐状态验证动作映射对可行动作集满射，并保证相应策略有可测提升。近似接口应输出 \[
(\epsilon_r,\epsilon_p,\beta,\operatorname{span}\bar V^*,D)
\] 及适用域。若只有 MMD，应传出核与函数类，不能把其数值填入 TV 槽位。

\hypertarget{ux5df2ux77e5ux6709ux9650ux6a21ux578bux7684ux6700ux7c97ux7a33ux5b9aux5206ux5272ux7b97ux6cd5}{%
\section{28.5 已知有限模型的最粗稳定分割算法}\label{ux5df2ux77e5ux6709ux9650ux6a21ux578bux7684ux6700ux7c97ux7a33ux5b9aux5206ux5272ux7b97ux6cd5}}

给定有限状态集 \(X\)、有限且各状态相同的动作集 \(A\)、已知转移矩阵 \(P_a\)、奖励 \(r(x,a)\)，以及需要保留的输出标签 \(o(x)\)。初始分割 \(\mathcal P_0\) 将标签与全部动作奖励完全相同的状态分在同一块。递推：在当前每块内，按签名 \[
\operatorname{sig}_n(x)
=\left(\text{当前块}(x),
\left(\sum_{y\in B}P_a(x,y)\right)_{a\in A,B\in\mathcal P_n}\right)
\] 重新分割，直到不再改变。

\textbf{P-ALG-01：有限稳定分割的正确性与最粗性。} 算法有限终止，终止分割是所有细化 \(\mathcal P_0\) 的受控可合并分割中最粗者。因此其商保留原输出、奖励和所有动作下的一步分布，并由归纳保留相应全部有限时域输出律。

\textbf{证明。} 每轮只细分、不合并；每次严格改变至少增加一个块，故严格改变次数最多 \(|X|-|\mathcal P_0|\)。终止时，同块状态到每个商块的转移概率在每个动作下相同，因此闭合成立。若 \(\mathcal Q\) 是任一稳定且细化 \(\mathcal P_0\) 的分割，归纳假设其细化 \(\mathcal P_n\)。每个 \(\mathcal P_n\) 块是若干 \(\mathcal Q\) 块的并；同一 \(\mathcal Q\) 块内到这些并集的概率相同，故下一轮也不会将该 \(\mathcal Q\) 块拆开。因此 \(\mathcal Q\) 细化终止分割，证明最粗性。有限时域结论由核递推。\(\blacksquare\)

\begin{verbatim}
partition := group_by(output, complete_reward_vector)
repeat:
    signature[x] := (block_id[x],
                     transition_mass_to_each_block_for_each_action)
    refined := split_each_block_by(signature)
    if refined == partition: return quotient(partition)
    partition := refined
\end{verbatim}

这里得到的是指定 Markov 状态域上的最粗稳定商，不应无条件等同于全部历史空间上的 FFIPS；从历史到状态还需充分性证明。浮点容差不具有精确等式的传递性：近似算法必须输出缺陷与阈值证书，不能仅将舍入相同的数视作严格相同。

\hypertarget{ux7ec4ux5408ux8ba4ux8bc1ux4e0eux5931ux8d25ux4f20ux64ad}{%
\section{28.6 组合认证与失败传播}\label{ux7ec4ux5408ux8ba4ux8bc1ux4e0eux5931ux8d25ux4f20ux64ad}}

证明依赖图中，一项前提失效使相应下游结论变为 unresolved 或 invalidated，而不是将其数值误差设为零。例如：

\begin{itemize}
\tightlist
\item
  部署支持超出采样支持：撤回统一误差证书，保留模型内代数证明。
\item
  发现隐藏状态使商不闭合：有限一步读出预测仍可成立，但控制和全未来身份须重审。
\item
  毁损实际改变资源预算：访问 blocker 仍是组合对象，但原物理必要性解释失效。
\item
  目标变更：对象身份未必改变，GOD/GOA 必须依新任务重算。
\item
  混合裕量消失：有限时域误差界可保留，长期平稳误差界不能沿用。
\end{itemize}

\hypertarget{ux9886ux57dfux6269ux5c55ux7684ux6700ux5c0fux63a5ux53e3}{%
\section{28.7 领域扩展的最小接口}\label{ux9886ux57dfux6269ux5c55ux7684ux6700ux5c0fux63a5ux53e3}}

领域模块必须提供：变量与单位、状态/路径机制、实验或干预操作、读出和噪声、可实施控制、对象候选、领域守恒/约束，以及误差模型。它可调用本核心定理；若反过来提出新通用定理，应去除领域专属前提后重新给出证明和反例检查。

一个成功的跨领域统一首先是结构和接口统一，不是将所有领域常数设成同一个数，也不是用统一术语替代领域方程。

\hypertarget{ux5df2ux5b9eux73b0ux7684ux6709ux9650ux652fux6301ux5224ux5b9aux63a5ux53e3}{%
\section{28.8 已实现的有限支持判定接口}\label{ux5df2ux5b9eux73b0ux7684ux6709ux9650ux652fux6301ux5224ux5b9aux63a5ux53e3}}

新的有限判定器输入：有限状态与动作、观察映射、初始 PMF、逐时源支持关系、认证谓词和时域。数学正确性直接关于实际 PMF 支持；可执行版本允许用有限关系表，但必须证明该表恰等于源支持。

输出 True 的精确含义是所有非空初始观察支持可行，进而存在实际零未认证占据策略；False 只否定同一模型、同一时域、同一认证条件下的该性质。浮点阈值或有限样本未出现的边不提供``恰等于支持''的证书。第 37 章给完整递推、正确性证明、算法代码和真实死路例。

\hypertarget{ux5e94ux7528ux65b9ux6cd5ux4eceux6a21ux578bux5230ux53efux68c0ux9a8cux8bc1ux4e66}{%
\chapter{29. 应用方法：从模型到可检验证书}\label{ux5e94ux7528ux65b9ux6cd5ux4eceux6a21ux578bux5230ux53efux68c0ux9a8cux8bc1ux4e66}}

\hypertarget{ux5197ux4f59ux8bbeux5907ux4e0eux5bb9ux9519ux7ef4ux62a4}{%
\section{29.1 冗余设备与容错维护}\label{ux5197ux4f59ux8bbeux5907ux4e0eux5bb9ux9519ux7ef4ux62a4}}

\textbf{输入。} 部件状态 \(X_t\)、负载接口 \(U_t\)、维修动作 \(a_t\)、故障/修复转移，以及需要保持的服务输出 \(f(X_t)\)。先声明评价的是服务预测身份、材料个体，还是物理维持机制。

\textbf{构造。} 对服务未来做预测分割；枚举可访问读出 \(R_S\)，计算第 13 章的缺陷 \(e(S)\)，得到最小服务读出载体。另以真实部件失效干预生成物理机制失效族。二者只有在物理实现条件满足时才可合并解释。

\textbf{决策。} 奖励可以是有效服务减维修成本，并加硬安全约束。求商 MDP 的 GOD，检查第 20 章值误差和动作间隔。若闭环有混合裕量，报告服务状态 GOA；若有永久失效，采用生存或 QSD，不把吸收态的普通平稳分布当``健康稳态''。

\textbf{可检验输出。} 最小读出族、物理割集、维修策略、服务中断概率、误差界。与全状态模型、单部件启发式、无维修基线在同一预算下比较。

\textbf{失败标准。} 载荷改变导致转移漂移，维修动作影响未观测隐藏模式，或删去传感器被错误当作删去供能部件时，分别返回相应接口失效。

\hypertarget{ux7a33ux6001ux7ef4ux6301ux968fux673aux52a8ux529bux5b66ux4e0eux53d1ux80b2}{%
\section{29.2 稳态维持、随机动力学与发育}\label{ux7a33ux6001ux7ef4ux6301ux968fux673aux52a8ux529bux5b66ux4e0eux53d1ux80b2}}

\textbf{输入。} 连续状态扩散 \[
dX_t=b(X_t,u_t)dt+\sigma(X_t,u_t)dW_t,
\] 非爆炸条件、噪声单位、候选持续域 \(\mathcal V\) 和恢复集 \(A\)。

\textbf{构造。} 可先寻找满足 \(\mathcal LV\le-aV+b\) 的非负势函数，得到矩恢复界；若关心返回 \(A\)，使用 \(\mathcal LV\le-c\) 的区域外漂移，或解适当边界条件下的 \(\mathcal LE=-1\)。关心退出则分析 killed 半群、条件存活分布与混合时间，不要求每条轨道返回同一点。

\textbf{发育接口。} 将结构阶段 \(s\) 纳入状态，定义阶段间运输 \(\Gamma_{s\to s'}\)。用后续环境协议检查运输后的预测/接口保真度。几何形状改变本身不证明身份中断；运输未经校准也不证明身份延续。

\textbf{可检验输出。} 恢复速度、噪声底、平均返回时间界、存活曲线、阶段运输误差和身份分支。恢复集或退出阈值必须在测试之前固定。

\textbf{失败标准。} 势函数只在有限训练轨道下降、非爆炸未保证、把噪声热平衡宽度称为精确不变点，均不足以建立持续性定理。

\hypertarget{ux591aux4e3bux4f53ux534fux4f5cux4e0eux7236ux5bf9ux8c61ux5f62ux6210}{%
\section{29.3 多主体协作与父对象形成}\label{ux591aux4e3bux4f53ux534fux4f5cux4e0eux7236ux5bf9ux8c61ux5f62ux6210}}

\textbf{输入。} 子对象状态与策略、公共环境接口、通信与资源耦合，以及父层候选未来。先固定条件环境，再设计切割协议，避免共同外部驱动造成假不可约性。

\textbf{构造。} 第 15 章给出父层完整路径律与切割律之间的距离及不可约性裕量；第 24 章评估子更新场与父价值梯度的夹角和误差。扫描耦合强度时同时测量父层完整性和子层保真度，得到可能的耦合窗口。

\textbf{决策。} 协调机制需要改变真实更新场或激励，而不只是更改分析者写下的目标函数。若任务是多目标，保留 Pareto 集或预先给出权重选择规则。

\textbf{可检验输出。} 父对象证书、有效切割族、子层损失剖面、协调阈值、跨环境稳定区间。单纯通信频率或相关系数不是父对象强度的充分指标。

\hypertarget{ux5b66ux4e60ux7cfbux7edfux7684ux8bb0ux5fc6ux538bux7f29ux4e0eux95edux73afux51b3ux7b56}{%
\section{29.4 学习系统的记忆压缩与闭环决策}\label{ux5b66ux4e60ux7cfbux7edfux7684ux8bb0ux5fc6ux538bux7f29ux4e0eux95edux73afux51b3ux7b56}}

\textbf{输入。} 观察历史、动作日志、预测目标、部署环境族、允许的记忆状态和更新规则。

\textbf{构造。} 用预测充分性定义所需记忆，而非要求重建所有历史。信息界给出精确或近似记忆成本下限；有限上下文下用第 16 章检验预测类。对压缩状态还须验证动作闭合，否则``预测良好''不保证``控制正确''。

\textbf{决策。} 以信念更新或明确可学习状态构造控制器；误差经过 TV、跨度或占据加权接口传到策略后悔。信息获取只在真实影响任务决策时计算决策价值，不能把不断产生新奇标签当作改进。

\textbf{可检验输出。} 压缩率---预测损失---控制后悔曲线、部署覆盖证书、失效检测和回滚条件。保留未压缩模型及简单历史窗口作基线；基线和新方法使用相同数据与计算预算。

\textbf{安全边界。} 工具权限、资源预算、不可逆动作和停机机制属于可行策略约束，不能在``扩大选择空间''名义下被模型自行删除。数学上的能力扩展不构成外部授权。

\hypertarget{ux6f14ux5316ux4e2aux4f53ux6027ux4e0eux9057ux4f20ux673aux5236ux7814ux7a76}{%
\section{29.5 演化个体性与遗传机制研究}\label{ux6f14ux5316ux4e2aux4f53ux6027ux4e0eux9057ux4f20ux673aux5236ux7814ux7a76}}

\textbf{输入。} 亲子对应、单位计数、世代时长、环境干预和子/父层读出。必须记录零后代家庭，才能分别估计繁殖强度和条件传递。

\textbf{构造。} 先认证对象与谱系，再估计 \(M\)、\(b\)、\(K\)，使用与可遗传未来相关的状态，剔除对当前预测目标无关的标签。计算平均增长的左右 Perron 结构；真实有限种群另报告灭绝和方差。

\textbf{转变检验。} 对外部脚手架、子层冲突调节和父层复制机制实施匹配干预，按第 26 章四分量证书判断。不能用一个增长率跳变替代全部 ETI 证据。

\textbf{可检验输出。} 传递可识别性、父层不可约性、繁殖价值、机制内化程度，以及哪些判断仍依赖采样或环境假设。

\hypertarget{ux673aux5236ux8bc6ux522bux4e0eux79d1ux5b66ux6a21ux578bux6bd4ux8f83}{%
\section{29.6 机制识别与科学模型比较}\label{ux673aux5236ux8bc6ux522bux4e0eux79d1ux5b66ux6a21ux578bux6bd4ux8f83}}

\textbf{输入。} 候选模型族、可实施实验族、可观测变量和预算。先按观测等价关系取商，再问能否通过新实验分离剩余类。

\textbf{构造。} 在规则参数模型中检查不同协议 Fisher 核的交，在线性问题中检查累积 Gram 矩阵。不能把局部满秩直接当作全局可识别，也不能把奇异 Fisher 一律当作不存在高阶识别。

\textbf{实验选择。} 对有限候选类，可选择覆盖未分离模型对的协议；对连续参数，需使用额外的连续性、紧性或局部设计理论。协议必须可实施，且区分强度须进入样本复杂度。

\textbf{可检验输出。} 可识别商、无法分离的机制类、下一实验的目标模型对、有限样本区间和外推限制。统一理论的经验内容体现在这些可失败的接口，而非事后总能重新命名对象。

\hypertarget{ux4eceux5df2ux8bc1ux660eux6a21ux578bux5230ux5b9eux9a8cux7ed3ux8bbaux7684ux4ea4ux4ed8}{%
\section{29.7 从已证明模型到实验结论的交付}\label{ux4eceux5df2ux8bc1ux660eux6a21ux578bux5230ux5b9eux9a8cux7ed3ux8bbaux7684ux4ea4ux4ed8}}

应用报告应同时给出固定源或模型族、实际抽样单位、可用信息、选择规则、部署生成次序、公共失败事件和误差预算。只给一个``通过''标签，无法判断它指的是预测、访问、真实物理维持还是任务风险。

C2/C5 的三值方法允许认证、拒绝或未决；C7 将协议、模型、解释和核心数学失败分开。当前科学闭合状态仍保留现实支持未验证。下一步真实系统研究应以可失败的模型合同和竞争解释开展，不能把合成模型的正确性直接报告为自然界的普遍结论。

\hypertarget{ux5b8cux6574ux7b97ux4f8bux4e0eux72ecux7acbux53efux6267ux884cux68c0ux9a8c}{%
\chapter{30. 完整算例与独立可执行检验}\label{ux5b8cux6574ux7b97ux4f8bux4e0eux72ecux7acbux53efux6267ux884cux68c0ux9a8c}}

本章全部为明确给定的数学模型或合成实验，不是自然系统实证。附录 D 内嵌 v3.0 历史运行的完整代码和输出。数值检验是实现一致性和反例检查，不代替前文一般证明。

\hypertarget{ux4e00ux4e2aux8d2fux7a7fux9884ux6d4bux8f7dux4f53ux63a7ux5236ux4e0e-goa-ux7684ux6709ux9650ux6a21ux578b}{%
\section{30.1 一个贯穿预测、载体、控制与 GOA 的有限模型}\label{ux4e00ux4e2aux8d2fux7a7fux9884ux6d4bux8f7dux4f53ux63a7ux5236ux4e0e-goa-ux7684ux6709ux9650ux6a21ux578b}}

微观状态 \(x=(s,z)\in\{0,1\}^2\)，动作 \(a\in\{0,1\}\)。下一步 \(z'\) 是独立公平比特， \[
P(s'=1\mid s,a)=
\begin{pmatrix}
0.15&0.85\\
0.90&0.85
\end{pmatrix}_{s,a},
\qquad r(s,z,a)=s-0.12a,\qquad\beta=0.9.
\] 故微观核为 \[
P((s',z')\mid(s,z),a)=\tfrac12
p_{sa}^{s'}(1-p_{sa})^{1-s'}.
\]

预测目标是 \(s\) 的未来，协议允许两种动作及其因果组合。\(z\) 对该未来无贡献，\(s=0,1\) 在动作零下的响应差为 \(0.75\)，故两类不能合并。有限稳定分割算法得到 \(\{\{(0,0),(0,1)\},\{(1,0),(1,1)\}\}\)；\(f(s,z)=s\) 是精确控制商。

取三个读出 \(r_1=s,r_2=s,r_3=z\)，则 \[
\mathcal M_0=\{\{r_1\},\{r_2\}\},\qquad
b(\mathcal M_0)=\{\{r_1,r_2\}\}.
\] 含任一个 \(s\) 读出的集合缺损为零；其余集合缺损为 \(0.75\)。这说明删去两个冗余读出才会破坏访问充分性；它并未断言删去两个物理部件恰是毁损设备的唯一方式。

\hypertarget{ux72ecux7acbux91c7ux6837ux548cux7cbeux786eux65cfux6062ux590d}{%
\section{30.2 独立采样和精确族恢复}\label{ux72ecux7acbux91c7ux6837ux548cux7cbeux786eux65cfux6062ux590d}}

在八个 \((s,z,a)\) 单元各独立生成 \(N=20000\) 个 Bernoulli 响应，总计 160000 个合成样本，随机种子 30905，\(\alpha=0.05\)。由 P-STAT-01 的 Bernoulli 版本， \[
\eta=\sqrt{\frac{\log(2\cdot8/0.05)}{40000}}
=0.0120086646.
\] v3.0 记录的最大经验误差为 \(0.00545<\eta\)。取缺损阈值 \(\tau=0.2\)，则 \[
2\eta<0.2<0.75-2\eta,
\] 满足精确结构间隔条件，恢复正确载体族与 blocker。

这里零缺损间隔由所给模型证明，不是单凭数据估计出真实等式。若换成未知连续模型，仅有该置信半径不足以宣布精确载体恢复。

\hypertarget{ux6700ux4f18ux4ef7ux503cux4e0eux53efux8ba4ux8bc1ux52a8ux4f5c}{%
\section{30.3 最优价值与可认证动作}\label{ux6700ux4f18ux4ef7ux503cux4e0eux53efux8ba4ux8bc1ux52a8ux4f5c}}

真实最优策略和价值为 \[
\pi^*(0)=1,\quad\pi^*(1)=0,\qquad
V^*=(7.7717277487,\;8.9445026178).
\] 用各 \(s\) 的 \(z=0\) 代表单元构造估计商，得到 \[
\widehat V=(7.7698153658,\;8.9411000664),
\qquad \widehat\pi=(1,0).
\]

\begin{longtable}[]{@{}
  >{\raggedright\arraybackslash}p{(\columnwidth - 4\tabcolsep) * \real{0.3000}}
  >{\raggedleft\arraybackslash}p{(\columnwidth - 4\tabcolsep) * \real{0.4000}}
  >{\raggedright\arraybackslash}p{(\columnwidth - 4\tabcolsep) * \real{0.3000}}@{}}
\toprule\noalign{}
\begin{minipage}[b]{\linewidth}\raggedright
量
\end{minipage} & \begin{minipage}[b]{\linewidth}\raggedleft
v3.0 历史结果
\end{minipage} & \begin{minipage}[b]{\linewidth}\raggedright
解释
\end{minipage} \\
\midrule\noalign{}
\endhead
\bottomrule\noalign{}
\endlastfoot
实际最大价值误差 & 0.00340255 & 已知真模型才可直接计算 \\
用实际模型缺陷的跨度界 \(D\) & 0.05745151 & 确定性条件界 \\
用同时置信半径的跨度界 \(D_{\rm conf}\) & 0.12659009 & 至少 95\% 响应事件上的界 \\
提升策略真实后悔 & 0 & 本例真实策略恰好相同 \\
状态零的估计动作间隔 & 0.62344368 & 大于 \(2D_{\rm conf}\)，可认证 \\
状态一的估计动作间隔 & 0.17797859 & 小于 \(2D_{\rm conf}\)，此证书未决 \\
\end{longtable}

所有``95\%''均相对于声明的重复抽样实验。跨度使用数据决定的 \(\widehat V\) 不需要另一次并合：响应误差事件是统一核界，在该事件上跨度不等式对任意相应有界函数成立。

\hypertarget{ux957fux671fux7ed3ux6784ux4e0eux6709ux9650ux8defux5f84}{%
\section{30.4 长期结构与有限路径}\label{ux957fux671fux7ed3ux6784ux4e0eux6709ux9650ux8defux5f84}}

提升策略的真实闭环核为 \[
P^\pi=\begin{pmatrix}.15&.85\\.10&.90\end{pmatrix}.
\] 其 Dobrushin 系数为 \(0.05\)，不变分布为 \[
\mu=(0.1052631579,\;0.8947368421).
\] 估计闭环不变分布为 \((0.1055867977,0.8944132023)\)，TV 误差约 \(0.00032364\)，低于置信运输界 \[
\eta/(1-0.05)=0.01264070.
\] 这一长期结论使用真实闭环的混合条件；一般未知模型应另认证相应裕量。

从 \(s_0=0\) 出发，枚举八步路径得到真实与估计路径律 TV 差 \(0.00206291\)，低于基于统一实际模型缺陷 \(0.00545\) 的 \(1-(1-0.00545)^8=0.04277733\)。使用全部动作的统一缺陷较保守，但类型一致且有效。

\hypertarget{ux8fdeux7eedux6062ux590dornsteinuhlenbeck-ux6a21ux578b}{%
\section{30.5 连续恢复：Ornstein--Uhlenbeck 模型}\label{ux8fdeux7eedux6062ux590dornsteinuhlenbeck-ux6a21ux578b}}

设 \[
dX_t=-\kappa X_tdt+\sigma dW_t,\quad\kappa>0.
\] 积分因子给 \[
X_t=e^{-\kappa t}X_0+
\sigma\int_0^te^{-\kappa(t-u)}dW_u.
\] 若 \(X_0=x_0\)，由 Itô 等距， \[
EX_t^2=e^{-2\kappa t}x_0^2+
\frac{\sigma^2}{2\kappa}(1-e^{-2\kappa t}).
\] 这与 \(\mathcal L(x^2)=-2\kappa x^2+\sigma^2\) 的恢复界一致。令初始律为 \(N(0,\sigma^2/(2\kappa))\)，上述高斯卷积保持同一分布，故得到不变概率。共享噪声的两解差为 \(e^{-\kappa t}(x_0-y_0)\)，给出耦合收缩。

这里有严格恢复和噪声底，却没有状态始终固定于零。若另定义区间退出，需要新的 killed 分析；普通不变分布并不表示永不退出任何有界区间。

\hypertarget{ux6761ux4ef6ux751fux5b58ux4e0dux7b49ux4e8eux6761ux4ef6ux6c38ux7eed}{%
\section{30.6 条件生存不等于条件永续}\label{ux6761ux4ef6ux751fux5b58ux4e0dux7b49ux4e8eux6761ux4ef6ux6c38ux7eed}}

取 killed CTMC 生成元 \[
K=\begin{pmatrix}-2&1\\2&-4\end{pmatrix},
\] 两状态额外吸收率分别为 1 和 2。特征值为 \(-3\pm\sqrt3\)，故 \[
\lambda=3-\sqrt3=1.2679491924,
\quad q=(0.7320508076,0.2679491924).
\] 计算右正特征向量 \(h\)，按 \(qh=1\) 归一化，得到 \[
\varpi=q\odot h=(0.7886751346,0.2113248654).
\] \(q\) 是此刻仍存活的 QSD，\(\varpi\) 是 Doob 条件永续过程的不变律。代码分别验证 \(qK=-\lambda q\)、变换生成元行和零及 \(\varpi K^\uparrow=0\)。

对区间 \((0,L)\) 中生成元 \(\frac12\partial_{xx}\) 的吸收 Brownian 运动，直接解 Dirichlet 特征方程得 \[
\phi_n(x)=\sin(n\pi x/L),\quad
\lambda_n=\frac{n^2\pi^2}{2L^2},
\] \[
q(x)=\frac{\pi}{2L}\sin(\pi x/L),\quad
\varpi(x)=\frac2L\sin^2(\pi x/L).
\] 谱混合间隔 \(\lambda_2-\lambda_1=3\lambda_1\)。这个固定比例并不独自保证任意给定误差与存活要求有很宽的共同时间窗；初态系数和阈值仍要进入 P-QSD-03。

\hypertarget{ux4e09ux4e2aux9632ux6b62ux865aux5047ux7edfux4e00ux7684ux53cdux4f8b}{%
\section{30.7 三个防止虚假统一的反例}\label{ux4e09ux4e2aux9632ux6b62ux865aux5047ux7edfux4e00ux7684ux53cdux4f8b}}

\textbf{协同遗漏。} 独立公平 \(X,E\)，\(Y=X\oplus E\)。有 \(I(Y;X)=I(Y;E)=0\)，但 \(I(Y;X,E)=\log2\)。分别检查内部、环境边际残差不能代替联合闭合。

\textbf{不可观测控制代价。} 二维 Brownian 噪声，观察 \(X=W^1\)，控制仅向第二噪声坐标加入单位漂移。时长一的完整噪声路径 KL 为 \(1/2\)，观察到的状态路径 KL 为零。因此``状态路径 KL 恰等于全部控制能量''需要可观测性；数据处理给出的不等式仍成立。

\textbf{逐环境二维不等于共同二维。} 共同预算坐标 \(B\in\mathbb R^3\)，两环境响应 Jacobian 分别为 \(\operatorname{diag}(1,1,0)\) 和 \(\operatorname{diag}(0,1,1)\)。各自秩二，但纵向堆叠秩三，不存在在该点以同一二维光滑中间变量同时分解两者的表示。Π/Φ 的共同双驱动必须检验这种跨环境一致性。

\hypertarget{ux7e41ux6b96ux4e0eux534fux8c03ux7684ux6570ux503cux4f8b}{%
\section{30.8 繁殖与协调的数值例}\label{ux7e41ux6b96ux4e0eux534fux8c03ux7684ux6570ux503cux4f8b}}

令 \[
M=\begin{pmatrix}1.8&0.2\\0.4&0.8\end{pmatrix},\quad
p=(0.4,0.6),\quad z=(0,1).
\] 则 \(b=(2,1.2)\)，\(\bar b=1.52\)， \[
p'=(0.6315789474,0.3684210526).
\] 平均性状变化 \(-0.2315789474\) 分解为选择项 \(-0.1263157895\) 与传递项 \(-0.1052631579\)。主增长因子约为 \(1.8744562647\)，不等于简单取最大的行和 2。

层级协调中取父梯度 \(g=(1,0)\)、子联合场 \(G=(-1,1)\)，则 \(\eta_*=1/2\)。大于此阈值的混合使父价值瞬时上升，但不证明远期最优或所有子价值改善。

\hypertarget{v3.0-ux6570ux503cux68c0ux9a8cux7684ux80fdux529bux8fb9ux754c}{%
\section{30.9 v3.0 数值检验的能力边界}\label{v3.0-ux6570ux503cux68c0ux9a8cux7684ux80fdux529bux8fb9ux754c}}

v3.0 原文记录代码通过 2371 项断言，包括三顶点全部 20 个反链的 blocker 对合，以及三原子到二原子的全部 6 个满射下的 blocker 自然性、左右伴随、精确实现区间与活跃纤维唯一性。另检验载体恢复、精确商、近似价值界、动作认证、稳定分割、路径 TV、平稳扰动、QSD/Doob 区分、Price 恒等式、奖励塑形及反例。其中 2316 项是有限分辨率结构的枚举实例，不是 2316 个不同的一般定理。

这些检查没有证明任意大系统中所有结论，也没有进行 Lean/Coq 等形式化证明核验。一般结论依赖正文证明和明确标准数学前提；合成数据未验证 UEOT 对自然界的普适性。

\hypertarget{ux7ef4ux62a4ux7248ux65b0ux589eux7684ux6b63ux5f0fux6709ux9650ux89c1ux8bc1}{%
\section{30.10 维护版新增的正式有限见证}\label{ux7ef4ux62a4ux7248ux65b0ux589eux7684ux6b63ux5f0fux6709ux9650ux89c1ux8bc1}}

本章 v3.0 的数值记录完整保留。维护版另有真实 Lean 见证：第 37 章的后验支持严格弱化模型、存在安全策略却有已覆盖死路的 Bool 模型；第 38 章的两条真实注册源行各 256 样本、以首样本选择两条不同部署律的实验。

后一个模型的选择概率为 \(7/8,1/8\)，实际联合失败率精确为 \(9/64\)；Good-only 校准给通用保证 \(17/32<1\)。其罕见 all-true 记录具有正概率却不满足小预算，说明所有记录都必须通过数值检查的前提过强。精确值、通用界和 Lean 原件分别列在第 38 章，不能与本章旧浮点输出合并为同一次实验。

\hypertarget{ux7edfux4e00ux88c5ux914dux6838ux5fc3ux80fdux5171ux540cux63a8ux51faux4ec0ux4e48}{%
\chapter{31. 统一装配：核心能共同推出什么}\label{ux7edfux4e00ux88c5ux914dux6838ux5fc3ux80fdux5171ux540cux63a8ux51faux4ec0ux4e48}}

\hypertarget{ux6709ux9650ux64cdux4f5cux6027ux6838ux5fc3ux7684ux88c5ux914dux5b9aux7406}{%
\section{31.1 有限操作性核心的装配定理}\label{ux6709ux9650ux64cdux4f5cux6027ux6838ux5fc3ux7684ux88c5ux914dux5b9aux7406}}

将多个模块组织成一个可检验结论，比单独罗列相似数学工具更接近本项目所需的``统一''。

\textbf{P-CORE-01：有限对象---实现---控制认证装配。} 给定：

\begin{enumerate}
\def\labelenumi{\arabic{enumi}.}
\tightlist
\item
  有限受控过程，其转移、奖励、动作语义及观察目标已正确指定；候选 \(f\) 在真实模型中满足受控闭合和奖励闭合，或具有声明的统一缺陷。
\item
  有限历史---协议响应表，未来目标固定；每单元独立重复采样，得到同时响应误差事件 \(E_{\rm resp}\)，概率至少 \(1-\alpha_{\rm resp}\)。
\item
  真预测类间隔 \(\gamma>4\eta\)，访问载体非零缺损间隔 \(\Delta>4\eta\)；阈值分别满足第 16 章条件。
\item
  读出纤维、有限载体域和访问操作固定。若解释成物理毁损，另具备第 13 章声明的实现条件。
\item
  折扣有界控制问题及策略提升有效，奖励误差和转移误差在事件 \(E_{\rm ctl}\) 上分别不超过 \(\epsilon_r,\epsilon_p\)，其失效概率不超过 \(\alpha_{\rm ctl}\)。
\item
  若要求持续历史身份，则另有有效历史边、运输和对象证书；若要求长期稳态误差，则有相应混合或复归结构前提；这些附加认证事件失效概率总计不超过 \(\alpha_{\rm aux}\)。
\end{enumerate}

则以至少 \[
1-\alpha_{\rm resp}-\alpha_{\rm ctl}-\alpha_{\rm aux}
\] 的概率（同一事件复用时不重复计数），同时有：

\begin{itemize}
\tightlist
\item
  被测试响应问题的预测类、精确极小访问载体族及其 blocker 正确恢复；
\item
  近似商价值误差不超过 \[
  D=\frac{\epsilon_r+\beta\epsilon_p\operatorname{sp}(\bar V^*)}{1-\beta},
  \] 提升策略后悔不超过 \(2D\)，动作间隔大于 \(2D\) 的估计最佳动作可认证；
\item
  相同策略下 \(T\) 步路径误差不超过第 7 章逐步耦合误差合成界；
\item
  在额外长期假设下得到对应 GOA 稳定性界；在已认证运输链上，历史保真误差按第 11 章累积；
\item
  若固定失效集合的估计距离裕量大于对象表示误差，则该表示扰动未越过该失效集合。
\end{itemize}

\textbf{证明。} 并合界给共同事件概率，不要求不同模块事件独立。在共同事件上，P-STAT-05 恢复预测类，P-STAT-04 恢复载体族，有限 blocker 是确定性函数，故一并恢复。P-QUO-02 和 P-QUO-03 给价值、后悔和动作认证。P-DYN-03 给相同实施策略的有限路径界。长期项依其所声明的 P-GOA 条件调用，不从有限时域界直接取极限。历史运输由 P-ID 的相应逐边界递推。最后由固定非空失效集合的距离函数 1-Lipschitz 性得到稳健性。\(\blacksquare\)

本定理不声称从有限响应表识别了所有未测协议，也不从访问读出恢复了全部物理机制。它把已有证书在共同前提下装配，而不是增添一个未经证明的``对象存在万能判据''。

\hypertarget{ux6838ux5fc3ux95edux5408ux4e0eux660eux786eux5f00ux653eux7aefux53e3}{%
\section{31.2 核心闭合与明确开放端口}\label{ux6838ux5fc3ux95edux5408ux4e0eux660eux786eux5f00ux653eux7aefux53e3}}

到这里，``对象---持续---组合---目的---学习---演化''获得同一套过程、协议、因子、实现和证书语言。核心的统一性体现在：

\begin{itemize}
\tightlist
\item
  对象不再由名称或材料集合先验指定，而由候选未来、内部/环境接口和历史机制共同认证。
\item
  宏观压缩通过明确商条件连接微观动力学、有限统计和控制价值。
\item
  载体与毁损的组合对偶通过操作语义连接物理实现，而不混淆访问与破坏。
\item
  生存、恢复、条件稳定和最优长期行为拥有不同但兼容的数学类型。
\item
  遗传与学习可改变未来控制结构，而目的桥和规范桥仍保持显式。
\end{itemize}

以下仍是研究问题，而非隐含定理：

\begin{enumerate}
\def\labelenumi{\arabic{enumi}.}
\tightlist
\item
  无限、连续、非紧环境中预测商及最小实现的可测存在条件。
\item
  在受限实验、强相关采样和结构漂移下同时发现对象与协议的有效复杂度。
\item
  不枚举全部子集而恢复大规模载体超图与物理机制割集的可证明算法。
\item
  非平稳、多尺度与分岔附近可迁移的身份运输和 GOA 稳定性。
\item
  共同 Π/Φ 低维机制的独立实验支持及何时必须增加驱动维数。
\item
  从具体复制、学习和协调机制识别目的桥，而非用行为事后拟合任意目标。
\item
  为尚未关闭的一般过程适配继续补齐机器证明；对已完成的有限同源连接，开展与其假设吻合的真实系统预注册检验。
\end{enumerate}

统一``真理''理论的可持续路线，是让共同架构承担越来越多明确、可证明和可失败的连接；不是把尚未得到的结论写成已经成立的公理。

\hypertarget{ux5b9eux9645ux4f7fux7528ux672cux6587ux4ef6}{%
\section{31.3 实际使用本文件}\label{ux5b9eux9645ux4f7fux7528ux672cux6587ux4ef6}}

新任务先定位第 28 章相应输入契约，再读取该任务依赖的定义、定理和反例。推导新结论时保留量词、支撑、协议和误差类型；增添领域模块时不要反向污染通用核心的前提。发现真实反例时更新对应证书和证明依赖，而不是改换对象名称规避失败。

本文件是通用数学主版本；后续修订应在此连续维护，并保留本次版本与验证记录。它既不要求读者另找旧版补齐关键证明，也不把其他领域分支作为使用前提。

\hypertarget{ux7531ux6761ux4ef6ux88c5ux914dux63a8ux8fdbux5230ux771fux5b9eux540cux6e90ux95edux73af}{%
\section{31.4 由条件装配推进到真实同源闭环}\label{ux7531ux6761ux4ef6ux88c5ux914dux63a8ux8fdbux5230ux771fux5b9eux540cux6e90ux95edux73af}}

第 35--38 章将本章装配原则具体落实到同一个有限源。稳定细化保持注册接口；源固定点产生安全动作；受限 Bellman 比较所有允许因果策略；恢复程序返回相同安全源核。对不完整观察，真实历史后验支持等于递推支持，后向可行性恰好刻画零未认证占据策略存在；有限关系判定器通过证明接回真实支持。统计 Good 上的经验模型和所选策略再通过实际训练---部署联合律得到总体风险。

这些箭头均有具体假设和定理来源。它们没有消去奖励来源、程序来源、物理实现框架或外部部署条件律的校准。第 34 章的 P0--P12/C1--C7 边界因此仍与这条更强正向链共存。统一性的进展是减少不必要的前提、把可导出的条件真正导出，并保留不可消去的条件，而非把全部条件改名为一个``客体公理''。

第 21.7、35.7 节又把同一源的准确观察商、实际策略、占据极限与 killed 条件结构接到一起。结果既提供正向装配，也辨明其不同问题：同一策略可以具有唯一稳定分布而不断积累曾失败风险；安全受限最优可以从指定初态持续，却在整个状态空间拥有多个不变律。第 36.8、38.10 节分别检验有限表示是否保留全部逃逸质量、统计选择是否仍消费同一真实采样与部署。这些联系使 MINCORE 的定量责任、H0 的层级区分和具体定理在真实来源上彼此约束，并没有取消目的、形成与经验验证仍需提供的条件。

\hypertarget{ux6709ux7c7bux578bux7684ux7edfux4e00ux516bux7c7bux5c42ux7ea7ux4e0eux56dbux4e2aux751fux6210ux673aux5236}{%
\chapter{32. 有类型的统一：八类层级与四个生成机制}\label{ux6709ux7c7bux578bux7684ux7edfux4e00ux516bux7c7bux5c42ux7ea7ux4e0eux56dbux4e2aux751fux6210ux673aux5236}}

\hypertarget{ux7edfux4e00ux95eeux9898ux7684ux6570ux5b66ux5f62ux5f0f}{%
\section{32.1 统一问题的数学形式}\label{ux7edfux4e00ux95eeux9898ux7684ux6570ux5b66ux5f62ux5f0f}}

前 31 章的结果已经说明，多种问题可以共享概率核、因子和误差传输。但统一不是将不同类型的数据替换成同一个符号。例如，预测商 \(q:X\to Y\) 是状态映射；载体粗化 \(\pi:V\to W\) 作用于物理或访问原子；谱系关系作用于带身份的对象实例；价值排序作用于策略。三个映射都可能被称为``升到高层''，它们仍不能在一个公式里直接复合。

本章用有类型的关系组织它们。一个接口至少给出定义域、值域、保持的结构和比较准则。只有中间对象相同，或有已证明的适配，复合才有意义。这也是形式化审查中检查类型之外还要检查``是否同一源''的原因：两个核即便类型相同，仍可能描述不同过程。

\hypertarget{ux516bux7c7bux5c42ux7ea7ux5173ux7cfb}{%
\section{32.2 八类层级关系}\label{ux516bux7c7bux5c42ux7ea7ux5173ux7cfb}}

H0 的八类关系如下。表中的区分是数学内容，不是八个互不相干的学科。

\begin{longtable}[]{@{}
  >{\raggedright\arraybackslash}p{(\columnwidth - 4\tabcolsep) * \real{0.3333}}
  >{\raggedright\arraybackslash}p{(\columnwidth - 4\tabcolsep) * \real{0.3333}}
  >{\raggedright\arraybackslash}p{(\columnwidth - 4\tabcolsep) * \real{0.3333}}@{}}
\toprule\noalign{}
\begin{minipage}[b]{\linewidth}\raggedright
类型
\end{minipage} & \begin{minipage}[b]{\linewidth}\raggedright
核心数学对象
\end{minipage} & \begin{minipage}[b]{\linewidth}\raggedright
允许的典型结论
\end{minipage} \\
\midrule\noalign{}
\endhead
\bottomrule\noalign{}
\endlastfoot
PQ 预测与商 & 响应族、纤维相容、可测统计量、受控核 & 最小充分表示、核降落、控制商 \\
DYN 动力学跨尺度 & 同一源核、交换图、半群、逐步误差 & 跨尺度一致、有限路径误差 \\
RES 分辨率与实现 & 闭包系统、原子映射、上闭族、伴随 & 极小实现粗化、实现不确定区间 \\
COMP 组合与父对象 & 子路径联合律、切割、父预测核、覆盖 & 条件不可约性、父层新增信息、组合窗口 \\
HIST 历史延续 & 有向实例、核运输、发展模型 & 冻结核缺陷和发展误差的累计 \\
VAL 价值与评价 & 偏好、微分形式、策略梯度、权重 & 可积性、价值保持与协调阈值 \\
SEL 选择与繁殖 & 亲子关系、平均繁殖算子、传递核 & Price 分解、Perron 组成、选择不可分辨性 \\
GOA 长期结构 & 不变律、QSD、复归类、谱隙、占据极限 & 有条件的存在、唯一性、混合和稳定性 \\
\end{longtable}

例如：同一个精确受控商可连接 PQ 与 DYN；如果奖励和动作也保持，可以连接 VAL；如果还给出真实载体读出，可连接 RES。没有这些附加数据时，从商定义本身不能推出父对象形成，也不能推出长期不变律存在。

Lean 对应的代表源分别为 QuotientDescent、DynamicsKernel、PathError、ClosureResolution、CompositionPathSource、TransportDefect、AlignmentGlobalExactness、LineageEvolution、FiniteCesaroInvariant 等。H0\_HIERARCHY\_INVENTORY 是这些关系的目录；它没有定义一个替代全部八类关系的统一算子。

\hypertarget{m-qdux76f8ux5bb9ux6027ux4e0eux552fux4e00ux964dux843d}{%
\section{32.3 M-QD：相容性与唯一降落}\label{m-qdux76f8ux5bb9ux6027ux4e0eux552fux4e00ux964dux843d}}

设 \(q:X\twoheadrightarrow Y\) 满射，\(g:X\to Z\)。定义纤维相容：

\[q(x)=q(x')\Longrightarrow g(x)=g(x').\]

\textbf{M-32.1：集合层的唯一降落。} 纤维相容当且仅当存在唯一 \(\bar g:Y\to Z\) 使 \(g=\bar g\circ q\)。

\textbf{证明。} 相容时，对每个 \(y\) 选 \(x_y\) 满足 \(q(x_y)=y\)，令 \(\bar g(y)=g(x_y)\)。相容性使选取不影响值，并给出 \(\bar g(q(x))=g(x)\)。若 \(h\circ q=g\)，则对 \(x_y\) 代入得 \(h(y)=g(x_y)=\bar g(y)\)。反向由同一个 \(\bar g\) 直接得到纤维相容。\(\blacksquare\)

\textbf{Lean。} \nolinkurl{UEOT.V3.Compression.QuotientDescent.fiberCompatible_iff_existsUnique_descend}。

这条定理本身是集合论的。若 \(X,Y,Z\) 带可测结构，还必须证明 \(\bar g\) 可测。第 4 章采用概率测度值坐标及可数协议族，给出在共同概率一集合上的可测因子化；\nolinkurl{countableAEFamily_descend} 复用这一共同机制。对一般不可数版本，第 4.3 节的单点零集反例继续有效。

双侧分离情形也可用同一机制。若 \(K:X\times E\to Z\) 对内部等价关系与环境等价关系分别不变，先替换 \(x\)、再替换 \(e\)，得到良定义的 \(\bar K\)；任意其他精确分离表示都必须细于响应等价。Lean 的 \nolinkurl{existsUnique_twoSidedDescend} 和 \nolinkurl{exactSeparatedFactorization_refines_responseEq} 对应这一证明。

计数压缩中，M-QD-01 生成 P-PRED-01 和 P-INT-02。可测 Markov 核、控制值保持、算法终止和条件解体的其余结构仍分别存在于保留适配中。

\hypertarget{m-tcux8befux5deeux5728ux540cux4e00ux4f20ux8f93ux94feux4e0aux7684ux590dux7528}{%
\section{32.4 M-TC：误差在同一传输链上的复用}\label{m-tcux8befux5deeux5728ux540cux4e00ux4f20ux8f93ux94feux4e0aux7684ux590dux7528}}

设各阶段比较量 \(e_n\ge0\)，传输常数 \(L_n\ge0\)，局部缺陷 \(\epsilon_n\ge0\)，并有

\[e_{n+1}\le L_ne_n+\epsilon_n.\]

\textbf{M-32.2：加权误差链。} 对任意 \(N\)，

\[e_N\le \left(\prod_{j<N}L_j\right)e_0+
\sum_{k<N}\epsilon_k\prod_{k<j<N}L_j.\]

空积取一。

\textbf{证明。} \(N=0\) 立即成立。将 \(N\) 阶界乘以非负 \(L_N\)，再加 \(\epsilon_N\)，把末项写成带空积的 \(k=N\) 项，即得 \(N+1\) 阶界。\(\blacksquare\)

\textbf{Lean。} \nolinkurl{UEOT.V3.Compression.TransportCertificate.weighted_chain_bound}；发展源包装为 \nolinkurl{development_pipeline_via_weighted_chain}。

两阶段的三角关系、精确交换图复合和 \(L_n=1\) 的链是同一机制的特例。概率路径还可使用乘法版本：若每阶段``未分歧的共同质量''至少保留比例 \(1-\epsilon_n\)，则全程共同质量至少为 \(\prod_n(1-\epsilon_n)\)，从而得到 P-DYN-03 的路径 TV 界。PathError 中真正的 PMF 延长与共同质量计算证明了这些比例来自实际核，而不只是一个抽象递推符号。

M-TC-01 的五个计数端点为 P-API-01、P-ID-01、P-DYN-04、P-DYN-03、P-ID-02。它统一计算机制，不消去各源的概率、可测性、奖励或局部缺陷条件。

\hypertarget{m-oiux6781ux9650ux4e2dux6d88ux5931ux7684ux6b8bux5dee}{%
\section{32.5 M-OI：极限中消失的残差}\label{m-oiux6781ux9650ux4e2dux6d88ux5931ux7684ux6b8bux5dee}}

设 \(\mu_k\) 弱收敛到 \(\mu\)，\(P\) 将所用有界连续测试函数保持为同类函数，并且对每个这样的 \(f\)，

\[\mu_k(Pf)-\mu_k(f)\longrightarrow0.\]

\textbf{M-32.3：可观测残差推出不变性。} 若该测试类决定概率测度，则 \(\mu P=\mu\)。

\textbf{证明。} 弱收敛分别作用于 \(f\) 与 \(Pf\)，给出 \(\mu(Pf)-\mu(f)=0\)。测试类决定测度，所以两测度相同。\(\blacksquare\)

\textbf{Lean。} \nolinkurl{UEOT.V3.Compression.OccupationLimitInvariance.invariant_of_continuous_observable_residual}。

对有限 Cesàro 占据，残差是 \((\mu_0P^N-\mu_0)/N\)；对连续 Feller 占据，残差由两个长度为 \(s\) 的时间边界项控制。它们分别产生 P-GOA-01 和 P-PER-02。紧性、弱极限子列及闭支集的传承仍需各自证明。这个论证不允许在没有极限或紧性时宣布稳态存在。

\hypertarget{m-peux5f52ux4e00ux5316ux6b63ux7279ux5f81ux7ed3ux6784}{%
\section{32.6 M-PE：归一化正特征结构}\label{m-peux5f52ux4e00ux5316ux6b63ux7279ux5f81ux7ed3ux6784}}

给有限非负矩阵 \(M\)、\(\rho>0\) 和逐项正的行向量 \(q\)、列向量 \(h\)，满足

\[qM=\rho q,\qquad Mh=\rho h,\qquad q\mathbf1=1,\quad qh=1.\]

\textbf{M-32.4：QSD/Doob 与随机读出。} 令

\[M^h_{ij}=\frac{M_{ij}h_j}{\rho h_i},\qquad \varpi_i=q_ih_i.\]

则 \(M^h\) 为随机矩阵，\(\varpi\) 为其不变概率。若 \(M\) 是次随机核，\(qM^n=\rho^nq\) 给出相应 QSD 关系。若逐项非负的随机行向量过程 \(Z_n\) 适应且可积，并满足

\[E[Z_{n+1}\mid\mathcal F_n]=Z_nM,\]

则 \(W_n=\rho^{-n}Z_nh\) 是非负可积鞅。

\textbf{证明。} 行和为 \((Mh)_i/(\rho h_i)=1\)；平稳性为

\[\sum_iq_ih_iM^h_{ij}=\frac{h_j}{\rho}(qM)_j=q_jh_j.\]

总质量为 \(qh=1\)。QSD 关系由反复应用左特征方程。鞅关系由条件期望的线性性和右特征方程得到。\(\blacksquare\)

\textbf{Lean。} \nolinkurl{UEOT.V3.Compression.PositiveEigenstructure.qsdDoob_bundle}、\nolinkurl{stochasticRightEigen_bundle}；计数端点为 P-QSD-02 与 P-EVO-04。

本命题消费特征数据；没有证明任意 \(M\) 满足这些数据。Perron 存在、唯一性、严格谱支配与幂极限仍由适用的 K-PF-01 或 KPF01Certificate 负责。正特征结构与随机分支过程的深层矩条件也不能混为一谈。

\hypertarget{ux4e2aux6e90ux7ed3ux679cux7684ux771fux5b9eux538bux7f29ux5206ux89e3}{%
\section{32.7 106 个源结果的真实压缩分解}\label{ux4e2aux6e90ux7ed3ux679cux7684ux771fux5b9eux538bux7f29ux5206ux89e3}}

冻结压缩账本给出

\[106=11\text{ 个生成端点}+89\text{ 个领域适配}+6\text{ 个保留边界}.\]

11 个生成端点为 P-PRED-01、P-INT-02；P-API-01、P-ID-01、P-DYN-04、P-DYN-03、P-ID-02；P-PER-02、P-GOA-01；P-QSD-02、P-EVO-04。

6 个保留边界为 P-INFO-05、P-INV-01、P-INV-02、P-DDH-01、P-EVO-02、P-BRG-02。它们限制了压缩、识别、规范或选择的无条件推广。保留边界属于统一理论的一部分，不是待抹掉的``未统一余项''。

最小核心的最小性相对于登记的定理表面与允许复用路线：删除一个生成元时，指定见证不再可经剩余登记路线导出。消融禁止在过程中临时从更底层库合成替代生成元。因此它是架构非冗余，不是绝对逻辑独立性，也不是四条公理唯一推出全部 106 个命题。

早期 633 模块与后来的 650 模块均为源快照计数；模块数量不会提高定理的逻辑强度。面向读者的规范入口可以简化，但不能通过删除保留适配使其前提消失。

\hypertarget{ux5173ux7cfbux7f3aux9677ux5b9aux91cfux6821ux51c6ux4e0eux89c4ux8303ux534fux53d8}{%
\chapter{33. 关系缺陷、定量校准与规范协变}\label{ux5173ux7cfbux7f3aux9677ux5b9aux91cfux6821ux51c6ux4e0eux89c4ux8303ux534fux53d8}}

\hypertarget{ux96f6ux5173ux7cfbux4e0eux5b9aux91cfux7f3aux9677}{%
\section{33.1 零关系与定量缺陷}\label{ux96f6ux5173ux7cfbux4e0eux5b9aux91cfux7f3aux9677}}

给集合 \(X\) 上的关系 \(R\)，一个缺陷表示为非负函数 \(d:X\times X\to\mathbb R_{\ge0}\)，满足

\[d(x,y)=0\iff R(x,y).\]

这里只要求零集对应；对称性、三角不等式与动力学收缩必须另外证明。\nolinkurl{RelationDefectPresentation} 正是这种最小接口。

\textbf{M-33.1：关系不能唯一恢复定量缺陷。} 若 \(d(x_0,y_0)>0\)，则对每个 \(c>0\)，\(d_c=cd\) 有相同零关系，而 \(c\ne1\) 时 \(d_c(x_0,y_0)\ne d(x_0,y_0)\)。因此不存在只读取关系、却能对所有这样的表示恢复其原始缺陷数值的规则。

\textbf{证明。} 正数乘积为零恰好等价于原数为零。若规则同时恢复 \(d\) 与 \(2d\)，相同关系会给出同一输出，在 \((x_0,y_0)\) 却应给两个不同数值，矛盾。\(\blacksquare\)

\textbf{Lean。} \nolinkurl{MinCore.RelationDefectPresentation.rescale_preserves_exact_relation}、\nolinkurl{rescale_changes_nonzero_defect}；\nolinkurl{MinCore.no_relation_only_quantitative_recovery}。

另取三点 \(0,1,2\)，令对角缺陷为零，\(d(0,1)=d(1,2)=1\)、\(d(0,2)=3\)，并对称。其零关系仍是相等，三角不等式却失败。对应 \nolinkurl{relationDefectPresentation_does_not_force_triangle}。这个反例说明：给某个定量指标取名为``结构距离''不能替代证明。

\hypertarget{ux53ccux8fb9ux6821ux51c6}{%
\section{33.2 双边校准}\label{ux53ccux8fb9ux6821ux51c6}}

设 \(d_P,d_Q\) 为同一对输入上的缺陷。给 \(a>0,b\ge0\)，校准条件是

\[a d_P(x,y)\le d_Q(x,y)\le b d_P(x,y)\quad\forall x,y.\]

\textbf{M-33.2：零关系、预算和趋零的运输。} 在校准条件下：

\begin{enumerate}
\def\labelenumi{\arabic{enumi}.}
\tightlist
\item
  \(d_P(x,y)=0\iff d_Q(x,y)=0\)；
\item
  \(d_P\le\epsilon\) 推出 \(d_Q\le b\epsilon\)；
\item
  \(d_Q\le\delta\) 推出 \(d_P\le\delta/a\)；
\item
  对任意序列或相应滤子，\(d_P\to0\iff d_Q\to0\)。
\end{enumerate}

\textbf{证明。} 第一项用非负性夹逼及 \(a>0\)。第二、三项分别使用上界和正数除法。第四项使用 \(0\le d_Q\le bd_P\)，以及 \(0\le d_P\le d_Q/a\) 的夹逼。\(\blacksquare\)

\textbf{Lean。} \nolinkurl{MinCore.DefectCalibration.zero_defect_iff}、rel\_iff、forward\_budget、reverse\_budget、tendsto\_zero\_iff。

不需要另外把``同一零关系''存入校准字段，因为它已由不等式推出。也不能仅凭零关系推出校准常数。例如 \(d\) 与 \(cd\) 可以使常数 \(c\) 任意大。

\hypertarget{ux51c6ux4e09ux89d2ux4e0eux8de8ux9636ux6bb5ux6536ux7f29}{%
\section{33.3 准三角与跨阶段收缩}\label{ux51c6ux4e09ux89d2ux4e0eux8de8ux9636ux6bb5ux6536ux7f29}}

\textbf{M-33.3：校准后的三角常数。} 若 \(d_P\) 满足三角不等式，则

\[d_Q(x,z)\le\frac ba\bigl(d_Q(x,y)+d_Q(y,z)\bigr).\]

\textbf{证明。} 依次应用目标上校准、源三角和两次源下校准：

\[d_Q(x,z)\le b d_P(x,z)
\le b[d_P(x,y)+d_P(y,z)]
\le\frac ba[d_Q(x,y)+d_Q(y,z)].\quad\square\]

这不是宣称 \(b/a\) 总为最优常数，也不是宣称目标仍具有原来常数一的三角性。

对于映射 \(F:X\to Y\)，若源度量满足

\[d_{P,Y}(Fx,Fx')\le Ld_{P,X}(x,x'),\qquad L\ge0,\]

并且输入校准下常数为 \(a_X>0\)、输出校准上常数为 \(b_Y\ge0\)，相同计算给

\[d_{Q,Y}(Fx,Fx')\le\frac{b_YL}{a_X}d_{Q,X}(x,x').\]

\textbf{Lean。} \nolinkurl{MinCore.DefectCalibration.triangle_quasi} 与 \nolinkurl{MinCore.calibrated_cross_stage_contraction}。

\hypertarget{ux7cbeux786eux6807ux5ea6ux53d8ux6362ux7684ux94feux534fux53d8}{%
\section{33.4 精确标度变换的链协变}\label{ux7cbeux786eux6807ux5ea6ux53d8ux6362ux7684ux94feux534fux53d8}}

一般双边校准只给界；若每阶段确有精确标度 \(d'_n=c_nd_n\)，\(c_n>0\)，则系数可精确运输。

\textbf{M-33.4：阶段标度与误差链。} 从 \(e_{n+1}\le L_ne_n+\epsilon_n\) 及 \(e'_n=c_ne_n\) 得

\[e'_{n+1}\le L'_ne'_n+\epsilon'_n,
\quad L'_n=\frac{c_{n+1}}{c_n}L_n,
\quad\epsilon'_n=c_{n+1}\epsilon_n.\]

M-32.2 的展开随之变成原展开整体乘 \(c_N\)。

\textbf{证明。} 将原一步界乘 \(c_{n+1}\)，代入 \(e_n=e'_n/c_n\)。展开中每串比例 \(c_{j+1}/c_j\) 望远镜相消，初始项与每个误差项都留下正确的末阶段标度。\(\blacksquare\)

\textbf{Lean。} \nolinkurl{MinCore.rescale_contraction_transformed}、\nolinkurl{rescale_local_budget_transformed}、\nolinkurl{rescaled_weighted_chain_bound}。精确同阶段正比例也保持三角常数，见 \nolinkurl{rescale_triangle_preserved}。

若更换单位却继续使用旧误差阈值、旧 Lipschitz 常数，就在比较不同量。这个协变定理修复的是该数学问题；是否某个 \(c_n\) 来自真实物理重整化，需要另行的模型相容关系。

\hypertarget{ux540cux4e00ux53ccux9a71ux52a8ux76eeux6807ux7684ux5b8cux6574ux89c4ux8303}{%
\section{33.5 同一双驱动目标的完整规范}\label{ux540cux4e00ux53ccux9a71ux52a8ux76eeux6807ux7684ux5b8cux6574ux89c4ux8303}}

在同一策略集 \(\mathfrak P\) 上给实函数 \(\Pi,\Phi\)，固定 \(\lambda\in\mathbb R\)，令

\[J(\pi)=\Pi(\pi)-\lambda\Phi(\pi).\]

\textbf{M-33.5：完整规范刻画。} 对另一对 \(\Pi',\Phi'\)，

\[J'=J\iff\exists\chi:\mathfrak P\to\mathbb R,
\quad\Pi'=\Pi+\lambda\chi,\quad\Phi'=\Phi+\chi.\]

\textbf{证明。} 右向左代入相消。左向右取 \(\chi=\Phi'-\Phi\)，由 \(J'=J\) 移项得到 \(\Pi'-\Pi=\lambda\chi\)。\(\blacksquare\)

\textbf{Lean。} \nolinkurl{UEOT.V3.Compression.TheoryCompletion.UnifiedClosure.same_objective_iff_gauge_related}。

若 \(\Phi'=\Phi\)，同目标立即给 \(\Pi'=\Pi\)。若 \(\Pi'=\Pi\) 且 \(\lambda\ne0\)，则 \(\Phi'=\Phi\)。\(\lambda=0\) 时 \(\Phi\) 不进入目标，不能识别它。对应 \nolinkurl{pi_identified_of_phi_anchor}、\nolinkurl{phi_identified_of_pi_anchor}、\nolinkurl{unanchored_drive_is_not_uniquely_identified}。

同一数值目标强于同一策略排序。三种策略上的目标值 \((0,1,2)\) 与 \((0,1,4)\) 排序相同，却不数值相等，因此不能用``保持 \(J\) 完全相等''的代数规范把二者联系起来。TeleologicalGauge 的目的等价与代数轨道必须保留这一差别。

\hypertarget{ux4eceux5f62ux5f0fux6821ux51c6ux5230ux53efux8bc6ux522bux673aux5236}{%
\section{33.6 从形式校准到可识别机制}\label{ux4eceux5f62ux5f0fux6821ux51c6ux5230ux53efux8bc6ux522bux673aux5236}}

完整规范定理既不决定 \(\Pi,\Phi\) 的物理量纲，也不证明跨环境响应只需要二维。资源测量、信息任务、干预轴、共同预算和误差模型仍是机制输入。第 23 章的秩必要条件是独立的可反驳命题：若在同一预算点的共同响应秩明显大于二，可排除所声明的共同二维因子化；秩兼容并不唯一确定两个潜变量。

所以本章带来两个互补结论：没有校准时哪些量不能识别；有明确校准时哪些预算、极限和规范可以合法运输。二者共同约束统一理论的定量表达。

\hypertarget{ux64cdux4f5cux5bf9ux8c61ux76eeux7684ux4e0eux6761ux4ef6ux751fux547dux5468ux671fux7684ux5b8cux6574ux63a5ux53e3}{%
\chapter{34. 操作对象、目的与条件生命周期的完整接口}\label{ux64cdux4f5cux5bf9ux8c61ux76eeux7684ux4e0eux6761ux4ef6ux751fux547dux5468ux671fux7684ux5b8cux6574ux63a5ux53e3}}

\hypertarget{ux516dux79cdux5bf9ux8c61ux89d2ux8272}{%
\section{34.1 六种对象角色}\label{ux516dux79cdux5bf9ux8c61ux89d2ux8272}}

给定候选域、实验协议及过程，至少要分清六种判断。

\begin{longtable}[]{@{}
  >{\raggedright\arraybackslash}p{(\columnwidth - 4\tabcolsep) * \real{0.3333}}
  >{\raggedright\arraybackslash}p{(\columnwidth - 4\tabcolsep) * \real{0.3333}}
  >{\raggedright\arraybackslash}p{(\columnwidth - 4\tabcolsep) * \real{0.3333}}@{}}
\toprule\noalign{}
\begin{minipage}[b]{\linewidth}\raggedright
角色
\end{minipage} & \begin{minipage}[b]{\linewidth}\raggedright
数学问题
\end{minipage} & \begin{minipage}[b]{\linewidth}\raggedright
证据责任
\end{minipage} \\
\midrule\noalign{}
\endhead
\bottomrule\noalign{}
\endlastfoot
预测身份 & 同一表示是否保留指定未来响应 & 纤维相容、可测因子化 \\
形成出处 & 候选是否由给定形成过程产生 & 属于真实形成关系或生成像 \\
相互作用区分 & 与其它候选的差别是否在注册探针中显现 & 逐对可区分或相应界 \\
组成持续 & 动力学是否保留组成条件 & 真实源与持续证书 \\
可恢复实现 & 受扰后是否存在有效修复路径或程序 & 恢复策略、载体及可信实现 \\
语义身份 & 哪些实例被解释为同一对象 & 指定身份关系及保持桥 \\
\end{longtable}

OperationalObject 的 PredictiveIdentity、FormedProvenance、InteractionDelimitedAt、ConstitutivelyPersists、PhysicallyRecoverableOn、SemanticFiberIdentity 对应这些角色。它们不共享一个自动推出其余五项的万能字段。

统计识别与物理形成尤其不同：算法可以以高概率识别一个已有父对象的类，而没有在世界中构造它。反过来，一个生成机制可以稳定地产生某结构，却未必使当前观察协议能够唯一识别它。

\hypertarget{ux76eeux7684ux5408ux540cux7684ux5fc5ux8981ux6761ux4ef6ux4e0eux673aux5236ux6865}{%
\section{34.2 目的合同的必要条件与机制桥}\label{ux76eeux7684ux5408ux540cux7684ux5fc5ux8981ux6761ux4ef6ux4e0eux673aux5236ux6865}}

给策略集 \(\mathfrak P\) 上的比较关系 \(\preceq\)。若要求实值 \(J\) 完全表示它，即

\[\pi\preceq\pi'\iff J(\pi)\le J(\pi'),\]

实数序的总性和传递性就成为 \(\preceq\) 必须满足的条件。

\textbf{M-34.1：不可比较合同的数值表示障碍。} 两个不同策略只各自与自身可比较时，不存在这样的实值完全表示。

\textbf{证明。} 两个实数必有一个不大于另一个，完全表示会强制至少一个交叉比较成立，与合同矛盾。\(\blacksquare\)

\textbf{Lean。} \nolinkurl{TheoryCompletion.incomparableBoolContract_has_no_numerical_representation}，以及不可有有效双驱动代表的对应定理。

现在给机制集合 \(M\)、观测 \(o:M\to O\)、机制目标 \(J_m\)。机制---目的忠实性要求

\[o(m)=o(m')\Longrightarrow J_m,J_{m'}\text{ 在同一策略域诱导相同排序}.\]

在此前提下，同观测机制的目的等价类相同。证明是直接应用忠实性再转入等价类；准确端点是 \nolinkurl{ScientificClosure.mechanismObservation_identifies_teleologicalClass}。

必要性反例取恒定观测 \(o(m)=*\)，两个机制分别给 Bool 策略相反目标。此时观察相同，目的相反。\nolinkurl{equalMechanismObservation_does_not_determine_purpose} 因而排除无忠实性桥的识别。这个反例不否认具体机制能支持目的解释，而是要求解释具有独立内容。

\hypertarget{p0p12-ux7684ux9010ux5c42ux804cux8d23}{%
\section{34.3 P0--P12 的逐层职责}\label{p0p12-ux7684ux9010ux5c42ux804cux8d23}}

\begin{longtable}[]{@{}
  >{\raggedright\arraybackslash}p{(\columnwidth - 4\tabcolsep) * \real{0.3333}}
  >{\raggedright\arraybackslash}p{(\columnwidth - 4\tabcolsep) * \real{0.3333}}
  >{\raggedright\arraybackslash}p{(\columnwidth - 4\tabcolsep) * \real{0.3333}}@{}}
\toprule\noalign{}
\begin{minipage}[b]{\linewidth}\raggedright
包
\end{minipage} & \begin{minipage}[b]{\linewidth}\raggedright
已组织的主要问题
\end{minipage} & \begin{minipage}[b]{\linewidth}\raggedright
必须保留的边界
\end{minipage} \\
\midrule\noalign{}
\endhead
\bottomrule\noalign{}
\endlastfoot
P0 & 操作客体、目的、GOD、GOA、尺度的语义区分 & 定义和同义词不是外部事实 \\
P1 & 目的合同、数值代表、规范与排序 & 表示存在、忠实性和规范识别分开 \\
P2 & 同一对象上的 agency、控制、GOD、GOA & 真实受控源、目的实现与生存前提 \\
P3/P4 & 有限数据与变化对象类的识别 & 共同 Good、间隔、协议与样本来源 \\
P5 & 组织故障下的常返与平均稳态 & 修复漂移、故障包络和同一混合核 \\
P6 & 资源累计与维持预算 & 初始储备、供给、成本和所用平均语义 \\
P7 & 从种子组装修复组织 & 程序来源、可信解释器、编解码和物理源 \\
P8 & 谱系、繁殖、Price/Perron 桥 & 同一算子、真实子代边及谱条件 \\
P9 & 尺度复合和性质运输 & 六类性质分别保持；未定义物理 RG 相容律 \\
P10 & 一般观察、核与信念接口 & 相应可测与拓扑正则性 \\
P11 & 所需随机分析接口 & ItoRun、Dynkin 与 Girsanov 的过程义务 \\
P12 & 将已给桥连接成条件生命周期 & PARTIAL / EXPLICIT BOUNDARY，不合成全部来源 \\
\end{longtable}

CLOSED 表示该包的约定交付已完成，可以包括条件命题和 no-go。它不把此表右列变成已自动解决的问题。

\hypertarget{ux5f62ux6210bindingux4f20ux8f93ux7684ux5b9aux91cfux7ec4ux5408}{%
\section{34.4 形成---binding---传输的定量组合}\label{ux5f62ux6210bindingux4f20ux8f93ux7684ux5b9aux91cfux7ec4ux5408}}

令 \(X_0,X_1\) 为底层输入域，\(F_0,F_1\) 为形成输出域，\(P_0,P_1\) 为父实现域。给

\[f_i:X_i\to F_i,\quad B_i:F_i\to P_i,\]

以及运输 \(t_X:X_0\to X_1\)、\(t_F:F_0\to F_1\)、\(t_P:P_0\to P_1\)。在 \(F_1,P_1\) 上给相应距离，假设

\[d_F(f_1(t_Xx),t_F(f_0x))\le\epsilon_F,\]

\[d_P(B_1(t_Fu),t_P(B_0u))\le\epsilon_B,\]

\[d_P(B_1u,B_1v)\le Ld_F(u,v),\qquad L\ge0.\]

\textbf{M-34.2：实现结构的 FBT 界。}

\[d_P(B_1(f_1(t_Xx)),t_P(B_0(f_0x)))\le L\epsilon_F+\epsilon_B.\]

若两端输出还分别满足实际形成关系，且右侧不超过预先指定的结构容限，则得到相应已形成结构的延续证书。

\textbf{证明。} 在两端之间插入 \(B_1(t_F(f_0x))\)。第一段用 \(B_1\) 的 Lipschitz 与形成误差；第二段用 binding 传输误差。形成关系和容限分别保留，不由三角不等式创造。\(\blacksquare\)

\textbf{Lean。} \nolinkurl{ScientificClosure.approximate_fbt_realized_bound}、\nolinkurl{certified_structural_continuation}、\nolinkurl{certified_formed_structural_continuation}。

当缺少 binding 传输条件时，可以令形成映射完全可交换而任意改变父输出，结论失败；\nolinkurl{formationTransport_does_not_imply_bindingTransport} 给出确切 Bool 反例。结构距离界也不等于父身份 token 相同。

\hypertarget{p12-ux7684ux5b9eux9645ux7ec8ux7aefux4e0eux7a0bux5e8fux6765ux6e90ux53cdux4f8b}{%
\section{34.5 P12 的实际终端与程序来源反例}\label{p12-ux7684ux5b9eux9645ux7ec8ux7aefux4e0eux7a0bux5e8fux6765ux6e90ux53cdux4f8b}}

P12 较强终端 \nolinkurl{p12_terminal_conditional_finite_lifecycle} 消费的是一组已经相互对齐的数据：有限物理 PMF 源 \(P\)、生存固定点 \(K\)、物理修复证书、可信修复底座、有限程序表示编解码、具有物理部分与程序来源的种子，以及实际程序实施修复后保持策略的证明。故障核必须落在声明的联合故障包络中，故障比例 \(\epsilon<1\)。

令 \(S\) 为由这些数据构造的组织维持系统，\(\mu_0\) 为实际组装组织的 Dirac 初始律。记 \(D_n\) 为其第 \(n\) 个真实边缘落在非法组织集合的概率；设有效修复漂移为 \(\kappa>0\)，平均故障注入负担为 \(\lambda\ge0\)。在源码中它们由同一故障比例与同一系统给出，而不是自由选择的两个参数。

平均稳态的结论为：对每个 \(\eta>0\)，最终所有充分大 \(N\) 满足

\[\frac1N\sum_{n<N}D_n\le\frac\lambda\kappa+\eta.\]

这不表示每条路径从某时起永不出错。若初始势为 \(V_0\)、可变成本为 \(c_v\)、固定维持成本为 \(c_m\)，资源前提具有以下形式：

\[c_v\frac{V_0}{\kappa}\le R_0,\qquad
c_m+c_v\frac\lambda\kappa\le s.\]

它们推出由真实组织边缘计算的期望成本满足所有累计资源预算

\[\sum_{n<N}E[c(Z_n)]\le R_0+Ns.\]

这里保证的是定义中的期望成本账本；没有从它推出每条随机路径都不破产。

同一个生命周期终端还包括：变化类发现器对同一父对象的错误概率趋零；组装初始律支撑在组织 carrier；每个支持边缘经显式 SameObject 桥保持同一 parent；真实 parent 与给定 child 身份不同；同一繁殖核的 Price 分解；以及给定尺度保持桥下的目标性质。证明分别调用真实组装入口、平均稳态、期望资源、身份、谱系与尺度端点，再组合其结论。它没有把这些结论作为一个任意 record 的字段直接返回。

这个终端的``有限''首先是有限状态、动作和程序表示的合同。结论中含所有自然数边缘和渐近平均；它\textbf{并不证明该组织必在有限时间死亡}，也不证明任意条件下可以永续。有限寿命、吸收、资源耗竭等更强主张要另加条件。

\textbf{M-34.3：程序来源不能从不可区分物理部分恢复。} 若两个合法种子具有同一物理部分却要求不同程序输出，则不存在只读取物理部分而同时正确输出二者的确定性合成器。

\textbf{证明。} 同一输入只能有同一函数值，与两个不同正确输出矛盾。\(\blacksquare\)

\textbf{Lean。} \nolinkurl{TheoryCompletion.no_physicalOnly_synthesizer_for_samePhysical_distinctSeeds}；\nolinkurl{UnifiedClosure.same_physical_seed_program_reconstruction_no_go}。

对已给编码 \(e:\mathrm{Program}\to\mathrm{Information}\)，若程序域有 fallback 值，则``存在完整左逆解码器''恰好等价于 \(e\) 单射：左逆推出单射；单射时在像内取唯一前像，像外取 fallback。\nolinkurl{perfect_program_decoder_iff_injective_information} 证明这一正反连接。编码的物理存在仍是单独的问题。

\hypertarget{c1c7-ux7684ux79d1ux5b66ux63a5ux53e3}{%
\section{34.6 C1--C7 的科学接口}\label{c1c7-ux7684ux79d1ux5b66ux63a5ux53e3}}

\textbf{C1：参数化信念。} 在标准 Borel 潜在空间及相应可数生成条件下，从给定信念、动作与观察联合核构造联合可测的后验核。条件概率因子分解保证给定参数的后验截面与既有后验几乎处处一致，再由观察分布推送得到信念转移。\nolinkurl{c1_parameterized_belief_recursion}、\nolinkurl{c1_countablyGenerated_belief_recursion} 将这些构造接回同一实际联合律；零质量观察仍有版本自由，不能无条件推广到任意可测空间。

\textbf{C2：有漂移的三值决策。} 若真实量 \(d\) 与估计 \(\widehat d\) 的误差半径和漂移总和为 \(r\)，则真实量在 \([\widehat d-r,\widehat d+r]\) 内。上端不超过阈值时可认证，下端超过阈值时可拒绝，其余未决。证明只用区间包含。\nolinkurl{IntervalCertificate.decideAt_certified_sound}、decideAt\_rejected\_sound、decideAt\_ambiguous\_iff 保留三个输出。

\textbf{C3：共同通道校准。} 对同一个 Markov 通道，数据处理证明输出比较不放大所用 TV，给特定 parent-binding 的常数一校准。\nolinkurl{commonMarkov_calibrates_parentBinding} 不能用于每个样本都重新改变通道却不计误差的场景。

\textbf{C4：条件 FBT。} 已由 M-34.2 给出；它不自动生成父形成关系和 binding 相容性。

\textbf{C5：共同二驱动的必要秩条件。} 若每个环境 \(m_e(B)=g_e(z(B))\) 共享同一二维 \(z\)，则同预算点堆叠 Jacobian 右乘同一个 \(Dz\)，故秩至多二。\nolinkurl{commonTwoDrive_rankNecessary} 与奇异值扰动结果给可拒绝、兼容、未决的判断；兼容不是唯一识别。

\textbf{C6：目的忠实性。} 已由 M-34.1 后的机制桥给出。识别排序类以后，控制实现和生存条件仍须独立验证。

\textbf{C7：证据与失败类型。} 同一个观测失败可能否定采样协议、源模型、物理解释或形式命题的应用前提；不能不经分型就声称数学内核已被证伪，反过来也不能用内核编译通过保护错误的世界模型。旧证据存在无法验证全部中途终止尝试是否保留的历史边界；新的持久记录机制不追溯修复过去。当前科学接口记录的现实支持仍为 UNVERIFIED。

\hypertarget{ux5c3aux5ea6ux6027ux8d28ux4e0eux6ca1ux6709ux5b8cux6210ux7684-rg-ux7aefux53e3}{%
\section{34.7 尺度性质与没有完成的 RG 端口}\label{ux5c3aux5ea6ux6027ux8d28ux4e0eux6ca1ux6709ux5b8cux6210ux7684-rg-ux7aefux53e3}}

ObjectScaleMap 可复合；若某谓词在第一个映射和第二个映射下都保持，依次应用两项证明即得复合保持。\nolinkurl{p9_terminal_objectScale_semantic_transport} 分别消费身份、目的、GOD、GOA、修复和稳态六项证书。

当前 \nolinkurl{wilsonianCouplingFlow_nontrivial_of_twoCouplings} 仅说明：耦合域有两个不同点，就能构造两个不同常值流记录。该记录没有与实际商映射、动力学或缺陷校准的相容方程。因此本书不将它升级为``给定尺度映射存在两个合法 Wilson 流''。新研究若要提出这一命题，必须先给出相容定义，再证明相容流存在或非唯一。

\hypertarget{ux540cux4e00ux6709ux9650ux6e90ux4e0aux7684ux9884ux6d4bux8f7dux4f53ux751fux5b58ux4feeux590dux4e0eux6700ux4f18ux63a7ux5236}{%
\chapter{35. 同一有限源上的预测、载体、生存、修复与最优控制}\label{ux540cux4e00ux6709ux9650ux6e90ux4e0aux7684ux9884ux6d4bux8f7dux4f53ux751fux5b58ux4feeux590dux4e0eux6700ux4f18ux63a7ux5236}}

\hypertarget{ux5171ux540cux6570ux5b66ux5e95ux5ea7}{%
\section{35.1 共同数学底座}\label{ux5171ux540cux6570ux5b66ux5e95ux5ea7}}

本章固定有限状态 \(X\)、有限非空动作 \(A\)、同一源核 \(P(x,a)\in\mathrm{PMF}(X)\)、奖励 \(r(x,a)\)、折扣 \(0<\beta<1\) 与注册输出 \(o(x)\)。需要保持的组织、程序、资源与安全标签可以作为 \(o\) 的分量。每一后续构造必须保留它所声称使用的这同一个 \(P\)；类型相同的另一核不能代替。

按第 28 章稳定细化得到最粗稳定分割 \(\sim\)，商映射为 \(q:X\to X/{\sim}\)。初始分割保留 \(o\) 和完整动作奖励向量；稳定性保证 \(q_\#P(x,a)\) 只依赖 \(q(x),a\)。于是得精确商核与奖励，P-QUO-01 给相同价值及宏观最优策略的提升。

\textbf{Lean。} \nolinkurl{OrganizationalCoreAssembly} 中 \nolinkurl{exactControlFromStablePartition}、\nolinkurl{organizationalTerminalControl}；\nolinkurl{UnifiedCoreClosureCertificate} 中 \nolinkurl{registered_organizational_attributes_preserved}、\nolinkurl{unified_finite_core_control_and_viability}。

\hypertarget{ux4eceux540cux4e00ux6838ux4ea7ux751fux5b89ux5168ux533aux57dfux4e0eux5b9eux9645ux8defux5f84}{%
\section{35.2 从同一核产生安全区域与实际路径}\label{ux4eceux540cux4e00ux6838ux4ea7ux751fux5b89ux5168ux533aux57dfux4e0eux5b9eux9645ux8defux5f84}}

给安全集合 \(V\subseteq X\)，使用第 8 章删除递推得到最大固定点 \(K\)。其动作条件为

\[x\in K\Longrightarrow\exists a:\operatorname{supp}P(x,a)\subseteq K.\]

\textbf{M-35.1：同源所有时刻安全。} 对任何 \(\mu_0\) 支撑在 \(K\)，存在确定性平稳策略 \(\pi\)，其由原始 \(P\) 和 \(\mu_0\) 构造的无限轨迹满足

\[\Pr_{P,\pi,\mu_0}(\forall n, X_n\in K)=1.\]

\textbf{证明。} 逐点选满足固定点的动作。若第 \(n\) 个边缘支撑在 \(K\)，原始 \(P\) 在这些动作下的支持也在 \(K\)，所以第 \(n+1\) 个边缘同样成立。真实路径扩张的边缘等于此递推；对所有 \(n\) 的概率一事件取可数交得到结论。\(\blacksquare\)

\textbf{Lean。} ViabilitySource.p\_per\_03；同源统一端点为 \nolinkurl{UnifiedClosure.single_source_terminal_object_viability_all_times}。

\(K\) 可以为空。这时不存在支撑在 \(K\) 的概率初始律，命题没有保证某现实初态可行。第 30.1 节的旧模型若要求服务位始终为一，每个动作仍有正概率离开，故相应概率一生存核为空，尽管该模型有高平均服务质量和普通不变律。

\hypertarget{ux7531ux5b89ux5168ux56faux5b9aux70b9ux4ea7ux751fux53d7ux9650ux6700ux4f18ux63a7ux5236}{%
\section{35.3 由安全固定点产生受限最优控制}\label{ux7531ux5b89ux5168ux56faux5b9aux70b9ux4ea7ux751fux53d7ux9650ux6700ux4f18ux63a7ux5236}}

定义每个状态的允许动作

\[A_K(x)=\{a\in A: x\in K\Rightarrow\operatorname{supp}P(x,a)\subseteq K\}.\]

在 \(K\) 内非空性来自固定点；在 \(K\) 外任意动作满足蕴涵。受限模型保留原奖励与原核，只缩小动作域。

\textbf{M-35.2：安全约束内对所有因果策略的最优性。} 受限 Bellman 算子

\[(T_Kv)(x)=\max_{a\in A_K(x)}[r(x,a)+\beta P(x,a)v]\]

有唯一固定点 \(V_K^*\)，存在确定性平稳达到策略。对任何每一步只选择 \(A_K(X_n)\) 内动作的历史依赖随机策略，其折扣回报不超过 \(V_K^*\)。若初始律支撑在 \(K\)，它们均满足同源概率一安全。

\textbf{证明。} 动作集有限非空，所以第 19 章 Bellman 压缩给固定点和达到者。对任意受限策略，Bellman 上界对其每个支持动作成立；按其实际历史条件期望、折扣求和、有界终端项取极限，得到全因果策略支配。安全由 M-35.1 的支持归纳，并不要求这些竞争策略平稳。\(\blacksquare\)

\textbf{Lean。} \nolinkurl{UnifiedClosure.exists_viable_source_action}、\nolinkurl{viable_restricted_greedy_preserves_source_kernel}、\nolinkurl{viable_restricted_optimal_dominates_all_causal_policies}。

这里最优策略与安全约束使用相同核。它没有把未受限策略的最优性偷偷写进安全结论。

\hypertarget{ux4e00ux4e2aux533aux5206ux5b89ux5168ux4e0eux672aux53d7ux9650ux6700ux4f18ux7684ux6709ux9650ux4f8b}{%
\section{35.4 一个区分安全与未受限最优的有限例}\label{ux4e00ux4e2aux533aux5206ux5b89ux5168ux4e0eux672aux53d7ux9650ux6700ux4f18ux7684ux6709ux9650ux4f8b}}

取状态 \(x=(s,z)\in\{0,1\}^2\)，\(z'\) 每步独立公平；动作为等待 \(a=0\) 和修复 \(a=1\)。修复总令 \(s'=1\)；等待在 \(s=0\) 时保持零，在 \(s=1\) 时以 \(3/4\) 留在一、以 \(1/4\) 变零。令

\[r(s,z,a)=s-a,\qquad\beta=1/2.\]

映射 \(q(s,z)=s\) 保持所有动作核和奖励。对安全集合 \(V=\{s=1\}\)，固定点为 \(K=V\)，而 \(K\) 内唯一允许动作是修复。受限安全值为零，因为每步奖励 \(1-1=0\)。

未受限模型的值为 \(V(0)=0,V(1)=8/5\)：在零状态，修复动作值为 \(-1+(1/2)(8/5)=-1/5\)，等待为零；在一状态，等待值为 \(1+(1/2)(3/4)(8/5)=8/5\)，修复为 \(4/5\)。故未受限最优策略等待，它在健康状态有正失败风险。

这是由前述有限 Bellman 方程直接检验的数学例，不另声称它有专门的 Lean 见证端点。它说明``未受限最优''与``安全最优''可以不同；注册输出和精确商都不能抹去这种价值约束差别。

\hypertarget{ux6062ux590dux7a0bux5e8fux8fd4ux56deux540cux4e00ux5b89ux5168ux6838}{%
\section{35.5 恢复程序返回同一安全核}\label{ux6062ux590dux7a0bux5e8fux8fd4ux56deux540cux4e00ux5b89ux5168ux6838}}

设 \(T\) 为可信程序执行底座，程序 \(p\) 对上节某个安全动作选择有效。修复过程联合更新物理状态与程序表示；在已经证明的编解码及恢复条件下，投影到下一物理状态的核恰好是 \(P(x,a_p(x))\)。

\textbf{M-35.3：程序恢复的源核一致性桥。} 若恢复程序的执行核与原安全源动作核相同，则从 \(K\) 内状态执行恢复程序后仍支撑在 \(K\)。

\textbf{证明。} 恢复定理先给核的等式，将已证明的 \(\operatorname{supp}P(x,a_p(x))\subseteq K\) 沿该等式改写。程序恢复与物理安全分别由它们自己的前提取得，然后经同一核连接。\(\blacksquare\)

\textbf{Lean。} \nolinkurl{UnifiedClosure.reconstructed_program_source_kernel_agrees_with_safe_action}、\nolinkurl{reconstructed_program_preserves_source_viability}、\nolinkurl{reconstructed_program_joint_and_source_safety}。

这条桥没有从未知物理材料合成 \(p\)，也没有断言恢复的程序达到 M-35.2 的受限最优值。若需要安全最优程序，要再证明该程序实施受限 Bellman 选择器；不能只因它安全便赋予最优标签。

\hypertarget{ux672cux7ae0ux5f97ux5230ux7684ux7edfux4e00ux4e0eux4fddux7559ux7684ux8f93ux5165}{%
\section{35.6 本章得到的统一与保留的输入}\label{ux672cux7ae0ux5f97ux5230ux7684ux7edfux4e00ux4e0eux4fddux7559ux7684ux8f93ux5165}}

同一有限源能够贯通：注册输出→稳定预测/控制商→访问目标与载体→生存固定点→实际无限安全路径→受限最优控制→恢复程序的源核一致性。每一步都有独立证明，不要求读者把几个不同的模型视为同一物理世界。

保留的输入包括奖励与安全任务、注册组织属性的物理意义、真实源模型、初始可行性和程序底座。它们与第 34 章的来源 no-go 相容。后三章进一步研究：在不完整观察下能否实际构造控制策略，以及在有限训练后使用所选模型时怎样形成真实风险保证。

\hypertarget{ux540cux4e00ux5b9eux4f8bux4e2dux7684ux5bf9ux8c61ux5212ux5206ux884cux52a8ux652fux6301ux4e0eux6062ux590d}{%
\section{35.7 同一实例中的对象划分、行动、支持与恢复}\label{ux540cux4e00ux5b9eux4f8bux4e2dux7684ux5bf9ux8c61ux5212ux5206ux884cux52a8ux652fux6301ux4e0eux6062ux590d}}

第 21.7 节固定的四状态源来自现有 \nolinkurl{conflictRewardControl.micro}：状态为 \((b,u)\in\{0,1\}^2\)，动作 \(a\) 令下一可见位等于 \(a\)，内部位公平重新取值；奖励为 \(1-a\)，折扣为 \(1/2\)。本节始终使用这个完整 M，初始律为 \(\delta_{(1,0)}\)，不使用第 35.4 节的另一转移或奖励。

对注册读出 \(q(b,u)=b\)，逐个动作有 \(q_\#P(x,a)=\delta_a\)，奖励也不区分同一可见纤维中的状态。已有 \nolinkurl{survivalRead_strongLumpable} 与 conflictRewardControl 因此提供非零奖励的精确控制商。宏观类型实际是 q 的 ReachableState 子类型；它与 Bool 的对应要有显式映射。若把内部位也登记为必须保留的程序属性，原两类分割就不够；既有 \nolinkurl{hiddenProgramProtectedModel} 保留了这种区别。

P-ALG 的 registeredSurvivalModel 具体见证使用同一个物理核，但奖励字段为零。它与本节完整 M 不是同一条模型记录。本节奖励与当前状态无关，所以书面上不会继续细化按 b 的初始分割；把这个非零奖励注册模型写入并证明其终端分割仍为 q 的核，是另一个具体适配。不能由已存在的精确控制商反向宣称这个 P-ALG 实例也已形式化。

真实支持为

\[
\operatorname{supp}P(x,1)=K,\qquad
\operatorname{supp}P(x,0)=K^c,\qquad K=\{b=1\}.
\]

在 K 内，保持域的动作恰为 1。按本章受限动作定义，K 外两动作均合法。M-21.7 已逐状态给出受限最优值：K 内为零、K 外为二；完整 greedy 保持当前可见位，拥有两个闭类及其所有混合不变律。从指定安全初态出发，它才与恒选 1 的策略拥有相同的实际路径。这样得到的持续与最优有真实共同来源，也有不可省略的初态和比较范围。

本源的后验支持提供一个清楚的适用边界：在初始时刻，正概率观察只有 1，其真实支持是单点 \((1,0)\)；在任意正时间、正概率实际历史上，最后一步的内部位公平重新取值，所以终端支持恰为该历史的当前可见纤维。这个实例没有展示``历史严格缩小当前纤维''的记忆增益。第 37 章的独立严格后验与死路见证仍有不可替代的作用。

对在每个历史下以 \(\theta\) 选择动作 0 的同一策略，若认证条件是 \(a=1\)，各步实际未认证占据均为 \(\theta\)；其 T 步曾失败风险为 \(1-(1-\theta)^T\)。若把全部真实行认证，并选真实精确参考核，则未认证占据可为零，参考风险仍为这个正数。另一证书可以只认证动作 1、使用零失败参考轨迹，但要支付本策略自己的占据预算。这些参数属于各自完整证书，不能混搭零参考风险与另一证书的零占据。

恢复程序还能接回同一个 P。一个明确保存可信程序来源的恢复步骤可以先重建程序、暂时保持物理状态，下一组织步骤执行恢复出的动作 1，再由 P 进入 K。这里的程序表示是联合恢复状态中的附加内容，不等于本例每步重新取值的内部位 u；两个组织步骤也不等于两个未经改变的物理 P 步。已有 JointRecovery 与本章 M-35.3 提供相应核一致性接口，具体两时钟实例仍需自己的适配。一次恢复重新开放了未来安全运行，却没有把先前中断从实际历史中删除。

这个共同实例连接了对象描述、可行行动、受限价值、支持、概率风险和恢复的来源条件。它不替代载体毁损的物理框架、父对象形成、谱系和外部实验。各条数学联系成立的范围，恰是其能为哲学的对象、持续和意义讨论提供清楚内容的原因。

\textbf{来源。} \href{../formalization/ueot-core/docs/compression/theory_completion/unified_closure/interface_research_v1/UNIFIED_FINITE_SOURCE_EXAMPLE.md}{统一有限源研究及逐箭头表}、\href{../formalization/ueot-core/UEOT/V3/Compression/TheoryCompletion/UnifiedClosure/SafeRewardConflictWitness.lean}{\nolinkurl{SafeRewardConflictWitness}}、\href{../formalization/ueot-core/UEOT/V3/Compression/TheoryCompletion/UnifiedClosure/OrganizationalSufficiencyWitness.lean}{\nolinkurl{OrganizationalSufficiencyWitness}}。新数值、参数化策略、完整 greedy 不变律与额外实例的形式化范围按第 21.7 节和附录 F 阅读。

\hypertarget{ux5b9eux9645ux6709ux9650ux8defux5f84ux540eux9a8cux5206ux89e3ux4e0eux90e8ux7f72ux8986ux76d6}{%
\chapter{36. 实际有限路径、后验分解与部署覆盖}\label{ux5b9eux9645ux6709ux9650ux8defux5f84ux540eux9a8cux5206ux89e3ux4e0eux90e8ux7f72ux8986ux76d6}}

\hypertarget{ux5148ux786eux5b9aux5b9eux9a8cux7684ux751fux6210ux5f8b}{%
\section{36.1 先确定实验的生成律}\label{ux5148ux786eux5b9aux5b9eux9a8cux7684ux751fux6210ux5f8b}}

本章固定有限非空状态空间 \(X\)、有限非空动作空间 \(A\)、观察映射 \(q:X\to Y\)、初始概率 \(p_0\) 和时变 Markov 源

\[P_n(x,a)\in\mathrm{PMF}(X).\]

观察历史不是单个当前观察。令

\[H_0^Y=Y,\qquad H_{n+1}^Y=H_n^Y\times(A\times Y),\]

微观历史 \(H_n^X\) 同样递归定义。\(Q_n:H_n^X\to H_n^Y\) 逐个保留动作、把每个状态变为其 \(q\) 观察。这里用大写 \(Q_n\) 表示历史投影；下文参考动力学核另记 \(\widehat P_n\)。

策略 \(\pi_n(g)\in\mathrm{PMF}(A)\) 读取完整粗历史 \(g\)。给微观前缀 \(h\)，先按 \(\pi_n(Q_nh)\) 生成动作 \(a\)，再按 \(P_n(\operatorname{last}h,a)\) 生成下一状态。由 \(p_0\) 和此递推产生实际微观路径律 \(p_n\)，实际粗路径律为 \(\bar p_n=(Q_n)_\#p_n\)。

整个后续证明都使用这个生成关系。即使训练数据、参考模型或另一个 PMF 具有相同的状态类型，也不能把它们改称 \(p_n\)。概率计算使用有限支撑 PMF；涉及测度时给有限空间以离散可测结构，在一般观察空间上保留可测单点条件。有限源映射与实际有限支撑路径据此可测，而不是声称任意无限历史策略自动可测。PMF 扩张是这里的实质构造，不只是某个抽象联合律的存在断言。

\textbf{Lean。} PathError 的 causalLaw、pmfExtend；UnifiedClosure 的 riskMicroControlKernel、riskHistoryMap；DeploymentCoverage 中 coverageDeploymentRowLaw。

\hypertarget{ux540cux4e00ux5386ux53f2ux7ea4ux7ef4ux4e0aux7684ux771fux5b9eux540eux9a8c}{%
\section{36.2 同一历史纤维上的真实后验}\label{ux540cux4e00ux5386ux53f2ux7ea4ux7ef4ux4e0aux7684ux771fux5b9eux540eux9a8c}}

当 \(\bar p_n(g)>0\) 时，定义末状态后验

\[
\rho_n^g(x)=
\frac{\displaystyle\sum_{\substack{h\in H_n^X\\Q_nh=g\\\operatorname{last}h=x}}p_n(h)}
{\bar p_n(g)}.
\]

它支撑在当前观察纤维 \(q^{-1}(\operatorname{last}g)\) 中，但可能远小于整个纤维，因为过去的观察和动作排除了部分状态。给粗历史和动作，实际下一观察核为

\[
K_n(g,a)(y)=\sum_x\rho_n^g(x)\,[q_\#P_n(x,a)](y).
\]

动作与后验的结合有一个明确前提：同一粗历史纤维中的全部微观前缀都使用 \(\pi_n(g)\)。于是对正质量 \(g\)，

\[
\Pr(X_n=x,A_n=a\mid G_n=g)=\rho_n^g(x)\pi_n(g)(a).
\]

这不是任意训练---部署耦合下的独立性声明。它由本节已指定的动作生成顺序推出。策略如果还读取未包含在 \(g\) 中的隐藏微观信息，此因子分解一般不成立。

\textbf{M-36.1：有限纤维分解与零质量分支。} 对正质量粗历史，\(\rho_n^g\) 是概率分布，并满足

\[
\bar p_n(g)\rho_n^g(x)
=\sum_{\substack{Q_nh=g\\\operatorname{last}h=x}}p_n(h).
\]

给零质量粗历史任取一个概率补全，以上加权等式仍成立。用这些后验构造的 \(K_n\) 与同一个 \(p_n\) 的下一粗历史扩张一致。

\textbf{证明。} 正质量时分母可消去，按末状态求和后分子就是 \(\bar p_n(g)\)。零质量时，同一纤维中各项非负且总和为零，因此每项为零；任意补全乘以零仍为零。最后按 \(g,a,x,x'\) 重排有限和，将同一纤维的共同因子 \(\pi_n(g)(a)\) 提出，即得到投影后下一步路径律。\(\blacksquare\)

\textbf{Lean。} FinitePMFProjection 中 riskFiberPosterior、\nolinkurl{risk_map_mul_fiberPosterior}、risk\_conditionedRow\_mass、risk\_map\_extend\_apply、risk\_map\_extend\_eq。末状态后验是既有历史纤维后验沿 last 的推送；新支持桥复用该分解，没有再假设一个与实际律匹配的后验。

零质量历史上的后验只是版本补全。不能因为选定补全含某个状态，反向要求该状态在实际实验中接受认证。

\hypertarget{ux4fddux7559ux524dux7f00ux65f6ux7684ux7cbeux786e-tv-ux5206ux89e3}{%
\section{36.3 保留前缀时的精确 TV 分解}\label{ux4fddux7559ux524dux7f00ux65f6ux7684ux7cbeux786e-tv-ux5206ux89e3}}

对有限集合上的概率分布，统一使用

\[
\operatorname{TV}(p,r)=\frac12\sum_z|p(z)-r(z)|
=\sup_E|p(E)-r(E)|.
\]

给共同前缀律 \(p\) 以及两组后继核 \(K,L\)，记扩张联合律为

\[
(p\ltimes K)(h,z)=p(h)K(h)(z).
\]

\textbf{M-36.2：共同前缀的精确行误差公式。}

\[
\operatorname{TV}(p\ltimes K,p\ltimes L)
=\sum_h p(h)\operatorname{TV}(K(h),L(h)).
\]

若使用共同后继核，则

\[
\operatorname{TV}(p\ltimes K,r\ltimes K)\le\operatorname{TV}(p,r).
\]

\textbf{证明。} 第一式逐项提出非负因子 \(p(h)\)，再交换有限和。第二式可用 Markov 数据处理；保留整个前缀时事实上可得到等式，上界形式已足够递推。若把前缀丢弃，只比较混合后的末状态，第一式通常只剩``\(\le\)''，因为不同前缀的有符号误差可以抵消。\(\blacksquare\)

\textbf{Lean。} \nolinkurl{coverage_tv_extend_same_initial}、\nolinkurl{coverage_tv_extend_same_kernel}；相应混合界复用 \nolinkurl{risk_tv_bind_same_weights}。第一条的新价值是把后续误差的权重固定为真实前缀，不能把任意训练行分布代入等号右侧。

\hypertarget{ux4eceux9010ux884cux8befux5deeux5230ux5b9eux9645ux8defux5f84ux8befux5dee}{%
\section{36.4 从逐行误差到实际路径误差}\label{ux4eceux9010ux884cux8befux5deeux5230ux5b9eux9645ux8defux5f84ux8befux5dee}}

给参考核 \(\widehat P_n(y,a)\)，其初始观察律为 \(q_\#p_0\)，动作仍由同一个 \(\pi_n(g)\) 产生。记参考路径律为 \(\widehat p_n\)，定义

\[
d_n(x,a)=\operatorname{TV}\bigl(q_\#P_n(x,a),\widehat P_n(qx,a)\bigr).
\]

实际部署的状态---动作占据律为

\[
\nu_n(x,a)=
\sum_{\operatorname{last}h=x}p_n(h)\pi_n(Q_nh)(a).
\]

它是实际前缀和动作联合 PMF 沿 \((h,a)\mapsto(\operatorname{last}h,a)\) 的推送，因而总质量为一。

\textbf{M-36.3：同策略、实际占据加权的路径界。}

\[
\operatorname{TV}(\bar p_T,\widehat p_T)
\le
\min\left\{1,\ \sum_{n<T}\sum_{x,a}\nu_n(x,a)d_n(x,a)\right\}.
\]

\textbf{证明。} 设 \(\delta_n=\operatorname{TV}(\bar p_n,\widehat p_n)\)。插入中间律：使用实际粗前缀 \(\bar p_n\)，但使用参考后继核。三角不等式将下一步误差分成``共同前缀、改变核''和``共同核、改变前缀''两项。第二项不超过 \(\delta_n\)；第一项由 M-36.2、同历史动作因子和 TV 的凸性，不超过

\[
\sum_g\bar p_n(g)\sum_a\pi_n(g)(a)\sum_x\rho_n^g(x)d_n(x,a).
\]

用 M-36.1 改写为 \(\sum_{x,a}\nu_n(x,a)d_n(x,a)\)。初始误差为零，归纳求和，再用概率 TV 不超过一。\(\blacksquare\)

\textbf{Lean。} \nolinkurl{coverage_projected_extension_error}、\nolinkurl{coverage_projected_causalLaw_error}、\nolinkurl{coverage_rowLaw_expectation}、\nolinkurl{coverage_same_policy_row_tv}、\nolinkurl{coverage_same_policy_path_error}。它们使用原始 causalLaw 的 actual prefix；没有额外输入``所用占据等于实际占据''的结论字段。

\hypertarget{ux4ec5ux8ba4ux8bc1ux4e00ux90e8ux5206ux90e8ux7f72ux884c}{%
\section{36.5 仅认证一部分部署行}\label{ux4ec5ux8ba4ux8bc1ux4e00ux90e8ux5206ux90e8ux7f72ux884c}}

令 \(C_n(x,a)\) 为被认证的源行。设 \(0\le\varepsilon_n\le1\)，且在这些行上 \(d_n(x,a)\le\varepsilon_n\)；对其余行只使用概率距离的一般上界 \(d_n\le1\)。

\textbf{M-36.4：未认证占据的精确线性收费。}

\[
\sum_{x,a}\nu_n(x,a)d_n(x,a)
\le
\varepsilon_n+(1-\varepsilon_n)\nu_n(C_n^c).
\]

若实际未认证质量还满足 \(\nu_n(C_n^c)\le\kappa_n\)，则可进一步用

\[
\varepsilon_n+(1-\varepsilon_n)\kappa_n
\]

代替。对微观失败事件 \(E\subseteq H_T^X\) 和粗事件 \(F\subseteq H_T^Y\)，若 \(E\subseteq Q_T^{-1}F\)，便得到

\[
\Pr_{p_T}(E)\le
\min\left\{1,\ \Pr_{\widehat p_T}(F)
+\sum_{n<T}\bigl[\varepsilon_n+(1-\varepsilon_n)\kappa_n\bigr]\right\}.
\]

\textbf{证明。} 将有限和分为 \(C_n\) 与 \(C_n^c\)。第一部分至多 \(\varepsilon_n\nu_n(C_n)\)，第二部分至多 \(\nu_n(C_n^c)\)，相加得第一式。替换上界需要 \(1-\varepsilon_n\ge0\)。最后由事件包含、TV 的事件界以及 M-36.3 得风险界。\(\blacksquare\)

\textbf{Lean。} \nolinkurl{coverage_uncovered_indicator_expectation}、\nolinkurl{coverage_expected_error_le}、\nolinkurl{coverage_same_policy_certified_path_error}、\nolinkurl{coverage_same_policy_microscopic_failure_risk}。本轮 \nolinkurl{ObservableActionCoverageContinuationCheck.deployment_coverage_budget_le_of_actual_occupancy_cap} 的预算单调性接口把实际 \(\nu\) 上界替换为给定 \(\kappa\)；这是已证明的单调代换，不会产生缺失的覆盖证据。

如果训练只观测到 \(C_n\) 中的源行，却没有控制部署进入 \(C_n^c\) 的概率，这个风险预算仍含未知的实际占据。样本数量再大也不能单独消除此缺项。

\hypertarget{ux52a8ux4f5cux8ba4ux8bc1ux600eux6837ux7ed9ux5b9eux9645ux8986ux76d6ux4e0aux754c}{%
\section{36.6 动作认证怎样给实际覆盖上界}\label{ux52a8ux4f5cux8ba4ux8bc1ux600eux6837ux7ed9ux5b9eux9645ux8986ux76d6ux4e0aux754c}}

给每个实际正质量粗历史一个允许动作集合 \(\Gamma_n(g)\)。假定

\[
a\in\Gamma_n(g),\quad \rho_n^g(x)>0
\quad\Longrightarrow\quad C_n(x,a),
\]

且策略在这些历史上满足

\[
\pi_n(g)(\Gamma_n(g)^c)\le\kappa_n.
\]

\textbf{M-36.5：支持认证的动作预算推出实际覆盖预算。}

\[
\nu_n(C_n^c)\le\kappa_n.
\]

\textbf{证明。} 按 M-36.1 将 \(\nu_n(C_n^c)\) 写成粗历史质量、动作质量和后验状态质量的有限三重和。正历史上，属于 \(\Gamma_n(g)\) 的动作对所有正后验状态都已认证，故只可能由允许集外动作贡献。每个正历史的条件贡献不超过其动作漏出质量 \(\kappa_n\)。零质量粗历史不贡献，零质量动作或后验状态也不贡献。最后对概率权重 \(\bar p_n(g)\) 求和。\(\blacksquare\)

\textbf{Lean。} \nolinkurl{PosteriorSupportCoverageCheck} 的 \nolinkurl{supported_history_actions_bound_actual_coverage}、\nolinkurl{posterior_supported_actions_bound_actual_coverage}；已通过的最小检查直接消费既有 riskFiberPosterior，而没有重建条件概率库。

一种更强的充分条件是用整个当前观察纤维：

\[
\Gamma_n(y)\subseteq\bigcap_{q(x)=y}\{a:C_n(x,a)\}.
\]

它不依赖历史后验，容易作为固定动作掩码使用。但它可能排除实际上可行的历史策略；第 37.6 节给出严格见证。

\hypertarget{ux8badux7ec3ux5206ux5e03ux4e0dux80fdux66ffux4ee3ux90e8ux7f72ux540eux9a8c}{%
\section{36.7 训练分布不能替代部署后验}\label{ux8badux7ec3ux5206ux5e03ux4e0dux80fdux66ffux4ee3ux90e8ux7f72ux540eux9a8c}}

本章 \(\nu_n\) 位于 \(X\times A\)，训练实验律 \(\mu\) 通常位于样本记录空间 \(\Omega\)，二者甚至不具有相同的类型。即使另给一个定义在 \(X\times A\) 上的训练行分布 \(\lambda_n\)，也需要额外关系才能替换期望。

例如，对 \(0\le d\le1\)，若 \(\operatorname{TV}(\nu_n,\lambda_n)\le\delta_n\)，则

\[
\mathbb E_{\nu_n}d\le\mathbb E_{\lambda_n}d+\delta_n.
\]

若已知 \(\nu_n\ll\lambda_n\) 且密度不超过 \(M_n\)，则

\[
\mathbb E_{\nu_n}d\le M_n\mathbb E_{\lambda_n}d.
\]

这些是第 16 章已有分布偏移界的同型应用。响应空间上``一个源行接近它的经验行''的 TV 校准，不蕴含``部署选择哪些源行''的分布接近训练设计；两个距离的空间和数学作用不同。

\hypertarget{ux6709ux9650ux5b9eux9645ux89c2ux5bdfux50cfux4e0eux5fc5ux987bux4fddux7559ux7684ux9003ux9038ux8d28ux91cf}{%
\section{36.8 有限实际观察像与必须保留的逃逸质量}\label{ux6709ux9650ux5b9eux9645ux89c2ux5bdfux50cfux4e0eux5fc5ux987bux4fddux7559ux7684ux9003ux9038ux8d28ux91cf}}

有限 X 的读出 \(q:X\to Y\) 具有有限完整像 \(S=q(X)\)，即使 Y 本身无限。参考核可能把质量送到 S 外，因此不能直接删除这些点并重新归一化。本节给出保留原参考过程、原实际策略及占据的书面扩展；完整路径适配的 Lean 状态与较基础的有限编码引理分开记录。

\textbf{M-36.6：一侧有限支撑的精确 TV 编码。} 设 Y 的单点可测，S 有限，概率测度 \(\mu,\nu\) 满足 \(\mu(S^c)=0\)。把每个 S 内点独立保留，并把全部外部点送到一个额外符号 \(\bot\)，得到可测映射 \(c:Y\to S\sqcup\{\bot\}\)。则

\[
\operatorname{TV}(\mu,\nu)=\operatorname{TV}(c_\#\mu,c_\#\nu).
\]

\textbf{证明。} 可测性来自有限 S 的可测单点纤维和可测补集。数据处理给右侧不大于左侧。反向取任意可测 A，令 B 仅包含 \(A\cap S\) 对应的保留点，则

\[
\mu(A)-\nu(A)\le\mu(A\cap S)-\nu(A\cap S)
=(c_\#\mu)(B)-(c_\#\nu)(B).
\]

以编码 TV 控制右侧，再对 \(A^c\) 使用同一不等式。两侧总质量同为一，故也控制 \(\nu(A)-\mu(A)\)。对 A 取确界。\(\blacksquare\)

若 \(\mu=\delta_0,\nu=(\delta_0+\delta_1)/2,S=\{0\}\)，逃逸质量为 \(1/2\)，编码后 TV 仍为 \(1/2\)；删除 1 后重新归一化却会错误给零。若两侧都可能在 S 外，例如 \(\mu=\delta_1,\nu=\delta_2\)，合并外部点会把 TV 从一降至零。一侧集中条件不能删除。

在无限环境 \(Y=\mathbb N\) 中还可同时区分源误差与两个参考的差别。取 \(S=\{0,2\}\)，

\[
\mu=\tfrac13\delta_0+\tfrac23\delta_2,\qquad
\nu_e=\tfrac12\delta_0+\tfrac14\delta_2+\tfrac14\delta_e,
\quad e\in\{5,7\}.
\]

两参考的逃逸质量均为 \(1/4\)，但 \(\operatorname{TV}(\mu,\nu_e)=5/12\)：有限原子上的半 \(L^1\) 差为 \(\tfrac12(1/6+5/12+1/4)\)。同时 \(\operatorname{TV}(\nu_5,\nu_7)=1/4\)，编码后两参考相等，TV 为零。源与参考的距离因而不只是逃逸概率；相对于固定有限源的精确比较，也不把编码变成所有参考之间的全局等距。这个联合边界已有 \nolinkurl{FiniteObservationEscapeBoundary} 的形式证明。

\hypertarget{ux4eceux6709ux9650ux7f16ux7801ux5230ux540cux4e00ux5b9eux9645ux8defux5f84}{%
\subsection{36.8.1 从有限编码到同一实际路径}\label{ux4eceux6709ux9650ux7f16ux7801ux5230ux540cux4e00ux5b9eux9645ux8defux5f84}}

取 \(Z=S\sqcup\{\bot\}\)。用初始 PMF 的非空支持选取 \(x_\star\)，令 \(y_\star=q(x_\star)\)；解码 d 将保留点还原，\(\bot\) 送到 \(y_\star\)。对完整 Z 历史逐观察解码、保留动作后，应用原完整策略 π，得到 \(\bar\pi\)。它也在不可能历史与逃逸历史上定义，不需凭空构造动作分布。

在保留点 \(s\in S\) 上令 \(\bar Q_n(s,a)=c_\#Q_n(s,a)\)，在 \(\bot\) 上令 \(\bar Q_n(\bot,a)=\delta_\bot\)。这些是归一化的真实 PMF 行。令 \(\bar q=c\circ q\)；任何实际微观历史经 \(\bar q\) 后再解码都等于它原来的 q 历史，因此 \(\bar\pi\) 在该历史上给出的动作 PMF 与 π 完全相同。原实际微观路径律和状态---动作占据均严格保持。

原参考过程可能离开 S 后又返回，有限参考则在 \(\bot\) 吸收。两者并非逐坐标编码后的同一个过程；需要比较完整历史。

设 \(L_n\) 为所有观察均在 S 内的完整 n 步观察---动作历史，它有限。向原 Y 历史及 Z 历史分别作保留动作的自然注入 \(i_n,j_n\)。由初始律相同及路径质量的实际乘法递推，对每个 \(\ell\in L_n\)，

\[
\bar\nu_n(j_n\ell)=\nu_n(i_n\ell).
\]

归纳步分别使用前缀质量相等、保留历史上的策略相等、参考行在保留后继点上的质量相等。它不是一个另行假设的整路径匹配条件。

再把两个完整历史空间都编码到同一个 \(L_n\sqcup\{\bot\}\)：保留点各自对应同一 \(\ell\)，非保留历史统一送到 \(\bot\)。两参考在所有保留点质量相等，归一化给剩余质量也相等。实际观察路径全部落在保留历史上，分别应用 M-36.6，得到原路径 TV 与有限参考比较 TV 精确相等。这里把``曾逃逸后返回''仍视为非保留完整历史，不能只看末观察。

\textbf{M-36.7：保持原来源的有限观察像路径界。} 在上述条件下，沿用原源 P、完整策略 π、初始律和原参考 Q，记实际状态---动作占据为 \(w_n\)，则

\[
\operatorname{TV}(\mu_T,\nu_T)
\le \sum_{n<T}\sum_{x,a}w_n(x,a)
\operatorname{TV}(q_\#P_n(x,a),Q_n(q(x),a)).
\]

\textbf{证明。} 每条实际源观察行集中于 S，M-36.6 因而逐行保持与相应参考行的 TV。将有限 Z 上的实际路径界应用于 \(\bar q,\bar Q,\bar\pi\)，然后用刚证明的路径 TV、逐行 TV 及实际占据等式改写。\(\blacksquare\)

在认证行上有误差 \(\varepsilon_n\) 时，M-36.4 的同一个实际占据预算继续有效。若再以 \(\kappa_n\) 代替实际未认证占据，仍保留 \(\varepsilon_n\le1\)。给原 Y 历史上的\textbf{可测}事件 F 及同一微观 E 的包络，事件转移给

\[
\Pr_{\rm actual}(E)\le
\min\{1,\nu_T(F)+B_T\}.
\]

其中 \(\nu_T(F)\) 是原 Q、原 π 的参考事件概率；F 的可测性替代了旧有限 Y 下自动得到的可测性。

M-36.6 的基本编码结论适用于一般概率测度，可以包括非原子参考；本章路径层的源与参考仍是 PMF。允许 Y 无限也允许参考 PMF 有可数无限支撑，但不等于已经构造任意非原子连续核的路径接口。

第一版取完整 \(q(X)\)，使所有源状态和全局认证行都在保留范围内。仅取某个策略的较小可达像时，还需证明实际支持归纳、改写认证域，并证明新旧占据相等；不能用``这些状态不可达''省略来源证明。

\textbf{Lean 范围。} 本版生产接口包括：\nolinkurl{FiniteObservationEscapeCheck} 的精确编码，\nolinkurl{PathObservationEscapeCheck.escape_path_tv} 的同策略完整路径等式，以及 \nolinkurl{ArbitraryObservationRiskCheck} 的原空间路径、认证占据与微观事件风险端点。\nolinkurl{FiniteObservationEscapeBoundary} 同时证明上述正逃逸数值边界。原有限 Y 接口仍保留，新端点沿本节的来源等式扩展其范围。编码使用非计算定义，不新增一个已经执行的通用算法。书面证明来源见 \href{../formalization/ueot-core/docs/compression/theory_completion/unified_closure/interface_research_v1/FINITE_OBSERVATION_ESCAPE.md}{\nolinkurl{FINITE_OBSERVATION_ESCAPE}}，当前声明和验收范围见 \href{../formalization/ueot-core/docs/compression/theory_completion/unified_closure/REGISTERED_INTERFACES_CLOSURE_V1.md}{接口闭环} 及 \href{LEAN_API_SCOPE.md}{LEAN\_API\_SCOPE}。

\hypertarget{ux90e8ux5206ux89c2ux5bdfux4e0bux7684ux6709ux9650ux652fux6301ux751fux5b58ux53ccux5411ux5b9aux7406ux4e0eux53efux6267ux884cux68c0ux9a8c}{%
\chapter{37. 部分观察下的有限支持生存：双向定理与可执行检验}\label{ux90e8ux5206ux89c2ux5bdfux4e0bux7684ux6709ux9650ux652fux6301ux751fux5b58ux53ccux5411ux5b9aux7406ux4e0eux53efux6267ux884cux68c0ux9a8c}}

\hypertarget{ux652fux6301ux9012ux63a8ux7684ux5bf9ux8c61ux4e0eux5fc5ux8981ux8fb9ux754c}{%
\section{37.1 支持递推的对象与必要边界}\label{ux652fux6301ux9012ux63a8ux7684ux5bf9ux8c61ux4e0eux5fc5ux8981ux8fb9ux754c}}

继续固定第 36 章的 \(p_0,q,P\)。本章认证的是有限时域 \(0\le n<T\) 中每一步使用的状态---动作行 \(C_n(x,a)\)。终点默认不附加新的安全状态条件；如果任务要求终端集合，就应修改递推的终值，不能由当前定理凭空获得。

给可能末状态集合 \(S\subseteq X\)、已记录动作 \(a\) 与下一观察 \(y\)，定义

\[
\operatorname{Succ}_n(S,a,y)=
\{x':q(x')=y\ \land\
\exists x\in S,\ P_n(x,a)(x')>0\}.
\]

初始观察支持为

\[
S_0(y)=\{x:p_0(x)>0,\ q(x)=y\}.
\]

对完整粗历史递归定义

\[
S_{n+1}(g,(a,y))=\operatorname{Succ}_n(S_n(g),a,y).
\]

这个 trace support 由 \(p_0,q,P\) 和记录中的动作确定，不把策略的概率值放进定义。它包含完整观察约束，不能只取某个任意代表或整个末观察纤维。

\textbf{Lean。} SupportViabilityClosure 中 initialSupport、successorSupport、traceSupport。数学支持桥不要求把 \(Y\) 枚举为有限集；由于 \(X\) 有限，有实际质量的观察总在有限像 \(q(X)\) 中。

\hypertarget{ux652fux6301ux9012ux63a8ux4e0eux5b9eux9645ux6982ux7387ux7684ux7b49ux4ef7}{%
\section{37.2 支持递推与实际概率的等价}\label{ux652fux6301ux9012ux63a8ux4e0eux5b9eux9645ux6982ux7387ux7684ux7b49ux4ef7}}

定义与第 36 章相同的 microLaw、coarseLaw。以下所有结论针对任意一个给定完整策略 \(\pi\)。

\textbf{M-37.1：正历史下的实际支持桥。}

\begin{enumerate}
\def\labelenumi{\arabic{enumi}.}
\item
  初始粗观察有正质量，当且仅当 \(S_0(y)\) 非空。
\item
  若粗历史 \(g\) 有正质量，则

  \[
  S_n(g)=\operatorname{supp}\rho_n^g.
  \]
\item
  对下一粗历史 \((g,(a,y))\)，

  \[
  \bar p_{n+1}(g,(a,y))>0
  \iff
  \bar p_n(g)>0\ \land\ \pi_n(g)(a)>0\
  \land\operatorname{Succ}_n(S_n(g),a,y)\ne\varnothing.
  \]
\item
  在上述正后继历史上，实际末状态后验支持正好等于 \(\operatorname{Succ}_n(S_n(g),a,y)\)。
\end{enumerate}

\textbf{证明。} 有限非负和为正等价于至少一项为正，两个有限非负概率因子的积为正等价于两个因子都为正。微观下一步支持据此分解为旧微观支持、策略动作支持和源行支持。对粗历史归纳：一个 \(x\in S_n(g)\) 恰好能由实际正质量、投影为 \(g\)、末状态为 \(x\) 的微观历史实现。

归纳的逆向构造尤其需要同一粗历史策略。换用任意另一个同纤维微观前缀时，原动作仍由同一个 \(\pi_n(g)\) 给出正质量，故可与支持中的源后继重新拼接。最后用已有纤维后验的支持定理把这些实现集合改写成后验支持。\(\blacksquare\)

\textbf{Lean。} \nolinkurl{mem_microLaw_succ_support_iff}、\nolinkurl{mem_coarseLaw_zero_support_iff}、\nolinkurl{mem_traceSupport_iff_realized}、\nolinkurl{mem_coarseLaw_succ_support_iff}、\nolinkurl{traceSupport_eq_fiberTerminalPosterior_support}、\nolinkurl{successorSupport_eq_fiberTerminalPosterior_support}。

定理只识别支持，不声称不同策略在同一可达粗历史上的完整后验概率相同。对零质量历史，trace support 仍有组合定义，但实际后验补全可以不同；当前第二、第四项没有向零质量历史延伸。

\textbf{M-37.2：零未认证占据的逐支持刻画。}

\[
\nu_n(C_n^c)=0
\iff
\forall g\in\operatorname{supp}\bar p_n\
\forall a\in\operatorname{supp}\pi_n(g)\
\forall x\in S_n(g),\ C_n(x,a).
\]

\textbf{证明。} 首先由实际前缀---动作 PMF 的推送定义证明

\[
(x,a)\in\operatorname{supp}\nu_n
\iff
\exists g\in\operatorname{supp}\bar p_n,\
a\in\operatorname{supp}\pi_n(g),\ x\in S_n(g).
\]

有限非负和 \(\nu_n(C_n^c)\) 为零，当且仅当每个正质量状态---动作对都已认证。代入支持刻画即可。\(\blacksquare\)

\textbf{Lean。} \nolinkurl{mem_deploymentRowLaw_support_iff}、\nolinkurl{zero_uncertified_mass_iff_support}、\nolinkurl{zero_uncovered_iff_support_certified}。这里的实际支持正确性是内部证明，不是最终定理另行要求用户提供的匹配假设。

\hypertarget{ux540eux5411ux652fux6301ux751fux5b58ux9012ux63a8}{%
\section{37.3 后向支持生存递推}\label{ux540eux5411ux652fux6301ux751fux5b58ux9012ux63a8}}

令 \(V(n,k,S)\) 表示从时刻 \(n\) 起、还有 \(k\) 个动作步骤、当前可能末状态集合为 \(S\) 时的支持可行性：

\[
V(n,0,S):=\mathrm{True},
\]

\[
\begin{aligned}
V(n,k+1,S):\iff \exists a\in A:\;&
\bigl[\forall x\in S,\ C_n(x,a)\bigr]\\
&\land
\bigl[\forall y,\ \operatorname{Succ}_n(S,a,y)\ne\varnothing
\\
&\qquad \Rightarrow V(n+1,k,\operatorname{Succ}_n(S,a,y))\bigr].
\end{aligned}
\]

动作必须是对全部当前可能状态共同可用的\textbf{同一个动作}。后继条件必须覆盖\textbf{每一个有正概率可能性的观察分支}。把第一项改成 \(\forall x\exists a\)，会准许读取未观测隐藏状态；把第二项改成存在某个良好后继，则会放过其他有正质量的失败分支。

\textbf{Lean。} viableWith、Viable。递推用剩余步数 \(k\) 做结构递归，并允许源和认证随绝对时间 \(n\) 改变。

为得到处处定义的完整策略，显式给一个基准策略 \(\pi_{\rm base}\)。定义 \(\pi^V\)：若 \(n<T\) 且 \(V(n,T-n,S_n(g))\)，选择上述存在量词中的某个共同动作并给它 Dirac 分布；否则使用 \(\pi_{\rm base}(n,g)\)。

这是一项数学选择构造。它没有声称所有零质量历史都走 fallback；零质量但满足 \(V\) 的历史也可以选递推动作。所有实际正历史上的性质不受这些分支决定影响。非空动作或一个实际 PMF fallback 是完整策略存在所需的输入，不应把空动作空间的问题藏在``任取一个动作''里。

\hypertarget{ux5fc5ux8981ux4e14ux5145ux5206ux7684ux7b56ux7565ux5b58ux5728ux6027}{%
\section{37.4 必要且充分的策略存在性}\label{ux5fc5ux8981ux4e14ux5145ux5206ux7684ux7b56ux7565ux5b58ux5728ux6027}}

\textbf{M-37.3：有限支持生存与真实零占据策略等价。} 给定完整 \(\pi_{\rm base}\)，以下两项等价：

\[
\forall y,\ S_0(y)\ne\varnothing\Rightarrow V(0,T,S_0(y));
\]

\[
\exists\pi,\ \forall n<T,\ \nu_n^\pi(C_n^c)=0.
\]

右侧 \(\nu_n^\pi\) 是同一 \(p_0,q,P\) 与这个所量化策略 \(\pi\) 实际生成的部署占据律；不是先指定一条参考律再任意改变控制器。

\textbf{充分性证明。} 使用上一节定义的 \(\pi^V\)。初始正粗历史满足假设。若时刻 \(n\) 的实际正历史满足相应 \(V\)，策略支持中的唯一动作就是递推选出的共同认证动作。任何实际正后继，按 M-37.1 都对应一个非空的 Succ，因此递推的全称后继条件保证该后继继续满足 \(V\)。前向归纳覆盖所有实际正历史；在每一步用 M-37.2 得 \(\nu_n(C_n^c)=0\)。

\textbf{必要性证明。} 给右侧策略。对任意实际正历史，它的动作 PMF 至少有一个正质量动作。由 M-37.2，这个动作对所有 \(S_n(g)\) 内状态共同认证。对每个非空 Succ，M-37.1 保证对应粗后继有实际正质量，因此可对剩余步数使用归纳假设。得到 \(V(n,T-n,S_n(g))\)，再应用于初始支持。

两方向都使用同一实际 law 的支持桥，因此没有循环地假设未来所构造策略拥有参考策略的后验。\(\blacksquare\)

\textbf{Lean。} \nolinkurl{viablePolicy_support_step}、\nolinkurl{viablePolicy_supported_histories_viable}、\nolinkurl{viablePolicy_zero_uncovered}、\nolinkurl{supported_viable_of_zero_uncovered}、\nolinkurl{initial_viability_iff_exists_zero_uncovered_policy}。

当 \(T=0\) 时，左侧终值为真，右侧任取显式基准策略且没有需认证的时刻，两侧相符。该定理也说明：对这一``概率一逐步认证''目标，若某个随机历史策略成功，则存在由支持递推选出的确定性历史策略成功。它不推出无限时域生存、奖励最优性、学习最优性或对近似支持判断的鲁棒性。

\hypertarget{ux6709ux9650ux5173ux7cfbux8868ux4e0aux7684ux53efux6267ux884cux5224ux5b9a}{%
\section{37.5 有限关系表上的可执行判定}\label{ux6709ux9650ux5173ux7cfbux8868ux4e0aux7684ux53efux6267ux884cux5224ux5b9a}}

原始源 \(P_n\) 作为 PMF 数学对象可以用非可计算的表示定义，不能只因 \(X\) 有限，就宣称其支持判断能直接执行。执行层因此显式提供布尔关系表 \(R_n(x,a,x')\)，并证明

\[
R_n(x,a,x')\iff P_n(x,a)(x')>0.
\]

这是一项确切表示定理。有限样本中``没有看见 \(x'\)''不满足这个等价式。

对有限集合 \(S\)，先按 \(R\) 枚举所有后继，按观察相等合并；再对动作枚举做``存在一个共同认证动作且所有非空后继都通过''的递归检验。即使 \(Y\) 未作为有限类型枚举，也只需遍历 \(q(X)\)。

\textbf{M-37.4：可执行检验的正确性。} 若关系表精确表示实际 \(P\) 支持，则有限关系检验返回 true 当且仅当 \(V(n,k,S)\) 成立。因此对所有非空初始支持返回 true，当且仅当 M-37.3 的真实零占据策略存在。

\textbf{证明。} 对 \(k\) 归纳。关系表精确性给每个构造的有限后继集合与 Succ 相同；有限动作枚举对应存在量词，有限观察像枚举对应所有非空后继。将布尔与逻辑连接逐层改写即得检验与 \(V\) 等价，再用 M-37.3。\(\blacksquare\)

\textbf{Lean。} finiteViabilityCheck、\nolinkurl{finiteViabilityCheck_correct}；实际可执行层为 finiteRelationSuccessor、\nolinkurl{finiteRelationViabilityCheck}、\nolinkurl{finiteRelationViabilityCheck_correct}。后者先证明其输出与前者相同，再复用前者正确性；没有重复建设策略存在证明。

下列伪代码给出完整控制结构。它要求 exactSupport 为已证明精确的关系表，observe 的像可枚举，certified 可判定。

\begin{verbatim}
ViabilityCheck(n, k, S):
    if k == 0:
        return true
    for a in Actions:
        if not all(certified(n, x, a) for x in S):
            continue
        passes = true
        for y in image(observe, States):
            Next = {z in States |
                    observe(z) == y and
                    exists x in S: exactSupport(n, x, a, z)}
            if Next is nonempty and not ViabilityCheck(n+1, k-1, Next):
                passes = false
                break
        if passes:
            return true
    return false
\end{verbatim}

可按 \((n,k,S)\) 缓存结果；状态支持至多有 \(2^{|X|}\) 种。指数级支持空间是部分观察问题的实质成本，当前 Lean 检验正确性没有声称一个新的多项式复杂度结果，也没有以符号算法代替大实例性能实测。

\hypertarget{ux4e25ux683cux653eux677eux6574ux4e2aux5f53ux524dux89c2ux5bdfux7ea4ux7ef4ux7684ux89c1ux8bc1}{%
\section{37.6 严格放松整个当前观察纤维的见证}\label{ux4e25ux683cux653eux677eux6574ux4e2aux5f53ux524dux89c2ux5bdfux7ea4ux7ef4ux7684ux89c1ux8bc1}}

取四个微观状态 \((b,s)\in\{0,1\}^2\)：\(b=0\) 为起始阶段，\(b=1\) 为后续阶段；\(s\) 是隐藏侧别。初始分布在 \((0,0),(0,1)\) 上各为 \(1/2\)。源的两种动作都把起始阶段送到同侧的后续阶段；后续阶段执行匹配动作 \(a=s\) 时保持状态，执行不匹配动作时回到同侧起始阶段。观察取三个值：

\[
q(0,0)=\mathrm{some}(0),\qquad
q(0,1)=\mathrm{some}(1),\qquad
q(1,0)=q(1,1)=\mathrm{none}.
\]

第一动作固定为零；第二动作记住初始观察并选侧别 \(s\)。起始阶段所有动作都认证，后续阶段仅 \(a=s\) 认证。

\textbf{M-37.5：真实后验支持严格优于整个当前纤维。} 两个实际正质量粗历史

\[
g_L=(\mathrm{some}(0),(0,\mathrm{none})),\qquad
g_R=(\mathrm{some}(1),(0,\mathrm{none}))
\]

各有质量 \(1/2\)。其实际末状态后验支持分别为 \(\{(1,0)\}\) 与 \(\{(1,1)\}\)，所用记忆策略满足 \(\nu_1(C_1^c)=0\)。但对整个当前观察纤维 \(q^{-1}(\mathrm{none})\)，不存在同时认证的共同动作。

\textbf{证明。} 初始两支各 \(1/2\)，首动作与转移确定，所以恰好得到上述两条正历史。第一条只能来自左状态，第二条只能来自右状态；已有纤维后验支持定理分别给单点支持。策略在各自支持上选择对应侧别，M-36.5 的零漏出情形给实际覆盖为零。整个当前纤维同时含两个状态；同一动作不可能既为零又为一。\(\blacksquare\)

此外，\(g_\varnothing=(\mathrm{none},(0,\mathrm{none}))\) 初始观察不可能发生，故质量为零。即使在这个分支设置空允许动作集，使条件动作漏出为一，它也不影响实际 \(\nu\)。这同时检查了零质量边界。

\textbf{Lean。} \nolinkurl{PosteriorSupportStrictWitness} 中 \nolinkurl{two_positive_coarse_histories}、\nolinkurl{posterior_supports_are_singletons}、\nolinkurl{both_terminal_states_reachable}、\nolinkurl{actual_uncovered_probability_zero}、\nolinkurl{whole_current_fiber_has_no_common_action}、null\_history\_boundary。当前公开端点形式化的是后验支持为单点；本书不把它记成另有一个``后验 PMF 等于 Dirac''的独立已审计声明。也不将手算的无记忆策略风险推论冒称本轮已编译端点。

\hypertarget{ux53efux884cux95eeux9898ux4e2dux4ecdux53efux5b58ux5728ux5c40ux90e8ux5b89ux5168ux6b7bux8def}{%
\section{37.7 可行问题中仍可存在局部安全死路}\label{ux53efux884cux95eeux9898ux4e2dux4ecdux53efux5b58ux5728ux5c40ux90e8ux5b89ux5168ux6b7bux8def}}

取 \(X=A=\{0,1\}\)，观察恒为单点。初始状态公平。第一步实际转移 \(X_1=A_0\)，此后保持状态。认证谓词是

\[
C_n(x,a)\iff(n=0)\ \lor\ (x=0).
\]

时域 \(T=2\)。第一步所有动作都认证；第二步仅状态零上的动作认证。

\textbf{M-37.6：存在完整可行策略，不蕴含每个已认证前缀都可延长。} 总选动作零的完整策略使两步未认证占据都为零。但是第一步选择动作一也是一个已认证前缀；它实际到达支持 \(\{1\}\)，第二步没有任何认证动作。

\textbf{证明。} 第一动作零导致 \(X_1=0\)，故第二步所有动作认证。第一动作一则导致 \(X_1=1\)，且 \(C_1(1,a)\) 对全部动作为假。第一步本身的认证与动作无关，故坏前缀确为``已认证但不能安全延长''。\(\blacksquare\)

\textbf{Lean。} SafePrefixDeadend 中 \nolinkurl{a_complete_zero_uncovered_policy_exists}、\nolinkurl{deadend_policy_second_step_not_covered}、\nolinkurl{viable_source_with_a_covered_deadend}。已通过的关系检验实际执行：

\[
\begin{array}{c|c}
\text{输入}&\text{输出}\\ \hline
(n=0,k=2,S=\{0,1\})&\mathrm{true}\\
(n=1,k=1,S=\{1\})&\mathrm{false}\\
(n=1,k=1,S=\{0\})&\mathrm{true}
\end{array}
\]

这个反例区分两种量词：要求``每个已经覆盖的部分策略都存在下一步延长''是较强的充分合同；后向 \(V\) 则在当前就排除会进入死路的动作，并给必要充分条件。它不能被改述为``一次只检查当前动作安全就足够''。

\hypertarget{ux4eceux6ce8ux518cux91c7ux6837ux5230ux6570ux636eux9009ux62e9ux540eux7684ux771fux5b9eux8054ux5408ux98ceux9669}{%
\chapter{38. 从注册采样到数据选择后的真实联合风险}\label{ux4eceux6ce8ux518cux91c7ux6837ux5230ux6570ux636eux9009ux62e9ux540eux7684ux771fux5b9eux8054ux5408ux98ceux9669}}

\hypertarget{ux8badux7ec3ux8bb0ux5f55ux4e0eux672aux6765ux8defux5f84ux7684ux8054ux5408ux5b9eux9a8c}{%
\section{38.1 训练记录与未来路径的联合实验}\label{ux8badux7ec3ux8bb0ux5f55ux4e0eux672aux6765ux8defux5f84ux7684ux8054ux5408ux5b9eux9a8c}}

设训练空间 \(\Omega\) 和本章所用未来路径空间 \(H\) 有限，训练概率为 \(\mu\in\mathrm{PMF}(\Omega)\)。模型候选以 \(c\in\mathcal C\) 索引，每个候选规定其实际未来部署 PMF \(K_c\in\mathrm{PMF}(H)\)。选择器 \(s:\Omega\to\mathcal C\) 可使用整个训练记录。

直接定义训练---未来联合律

\[
J(\omega,h)=\mu(\omega)K_{s(\omega)}(h).
\]

它表示先实际生成训练记录，再按所选条件部署核生成未来路径。训练和所选模型当然可以相关；不要求选择器独立于训练。这里有限 \(\Omega\) 保证所选有限未来律可测，候选索引本身不必强加有限性。

\textbf{M-38.1：真实联合律的边缘与塔式公式。} \(J\) 是概率分布，训练边缘为 \(\mu\)，未来边缘为 \(\sum_\omega\mu(\omega)K_{s(\omega)}\)。对共同未来事件 \(E\subseteq H\)，

\[
J(\Omega\times E)
=\sum_\omega\mu(\omega)r(\omega),
\qquad r(\omega)=K_{s(\omega)}(E).
\]

\textbf{证明。} 对 \(h\) 求和使用每个 \(K_c\) 总质量一；对 \(\omega\) 求和得到未来边缘；对事件指标做两次有限求和得到塔式公式。若训练原先以有限空间上的概率测度给出，Measure.toPMF 及其 toMeasure 逆等式保证使用的是同一个训练律。\(\blacksquare\)

\textbf{Lean。} \nolinkurl{selectedTrainingDeployment_mass}、\nolinkurl{selectedTrainingDeployment_training_marginal}、\nolinkurl{selectedTrainingDeployment_path_marginal}、\nolinkurl{selectedTrainingDeployment_failure_tower}、\nolinkurl{selectedTrainingDeployment_failure_integral}；\nolinkurl{finite_training_selection_measurable}、\nolinkurl{selectedDeployment_measure_measurable} 给所用可测接口。

共同事件 \(E\) 的固定很重要。可以让不同宏观表示使用不同 \(F_c\)，但必须各自覆盖同一个微观失败 \(E\)；选择模型不能同时重新定义何谓失败。

\hypertarget{ux53eaux5728ux6821ux51c6ux597dux4e8bux4ef6ux4e0aux68c0ux67e5ux9884ux7b97}{%
\section{38.2 只在校准好事件上检查预算}\label{ux53eaux5728ux6821ux51c6ux597dux4e8bux4ef6ux4e0aux68c0ux67e5ux9884ux7b97}}

设 \(G\subseteq\Omega\) 是好事件。在 \(G\) 上已有真实条件风险上界 \(r(\omega)\le b(\omega)\)。

\textbf{M-38.2：数据依赖预算的联合界。}

\[
J(\Omega\times E)
\le
\mu(G^c)+\sum_{\omega\in G}\mu(\omega)b(\omega).
\]

这里不要求好事件外的 \(b\) 满足某个数值预算，也不需要额外假设 \(b\) 在整个 \(\Omega\) 上非负。证明只在 \(G\) 使用它。

\textbf{证明。} 将 M-38.1 的和分成 \(G\) 与 \(G^c\)。好事件上代入给定条件界；坏事件上用风险不超过一。\(\blacksquare\)

若在 \(G\) 上统一有 \(r\le\beta\)，\(0\le\beta\le1\)，并有 \(\mu(G^c)\le\alpha\)，则得到

\textbf{M-38.3：固定好事件预算的精确混合形式。}

\[
J(\Omega\times E)
\le\min\{1,\alpha+(1-\alpha)\beta\}.
\]

\textbf{证明。} 记实际坏质量为 \(q=\mu(G^c)\)。M-38.2 给 \(q+(1-q)\beta=\beta+(1-\beta)q\)。因 \(1-\beta\ge0\)，可用 \(q\le\alpha\)。再与概率上界一取最小值。若 \(\alpha>1\)，右侧截断为一；命题仍合法，但没有提供非平凡信息。\(\blacksquare\)

\textbf{Lean。} \nolinkurl{selectedTrainingDeployment_data_dependent_bound}、\nolinkurl{selectedTrainingDeployment_good_bad_bound}、\nolinkurl{selectedTrainingDeployment_failure_le}。这个概率代数本身是标准混合界；本项目新增连接的实质是让旧统计证书给出的 \(G\) 与所构造实际 \(J\) 共用源、训练律、选定策略和失败事件。

\hypertarget{ux56faux5b9aux6ce8ux518cux884cux4e0eux89c4ux8303ux6821ux51c6ux4e8bux4ef6}{%
\section{38.3 固定注册行与规范校准事件}\label{ux56faux5b9aux6ce8ux518cux884cux4e0eux89c4ux8303ux6821ux51c6ux4e8bux4ef6}}

训练开始前固定注册源行集合 \(\mathcal R\)。每条注册行 \(j\) 有真实有限响应律 \(p_j\)，从该行取得 \(N_j\ge1\) 个独立同分布响应，经验律为 \(\widehat p_j(\omega)\)。不同源行间的独立性不是有限并合界本身的必要前提；具体产品采样模型可以额外提供它。

给半径 \(\rho_j\ge0\)，规范好事件为

\[
G=\{\omega:\forall j\in\mathcal R,\
\operatorname{TV}(p_j,\widehat p_j(\omega))\le\rho_j\}.
\]

对响应空间大小 \(m\)，有限事件并合与 Hoeffding 给

\[
\mu(G^c)\le
\alpha:=\sum_{j\in\mathcal R}
2\cdot2^m\exp(-2N_j\rho_j^2).
\]

上界可以大于一；实际使用时需检查数值非平凡性。这里枚举所有响应子集是与当前形式证书匹配的直接界，并不声称它是最优集中常数。

若旧统计证书先给某个满足这些逐行性质的好事件 \(G_0\)，则 \(G_0\subseteq G\)，从而 \(\mu(G^c)\le\mu(G_0^c)\)，不需要重新收取一次失败概率。

\textbf{Lean。} 原统一样本数版本为 \nolinkurl{measure_finite_row_calibration_bad_le}、\nolinkurl{coverage_registered_source_calibration} 及 RegisteredGoodJointCheck 的规范 Good。第 38.10 节将逐行预算与 \nolinkurl{VariableRegisteredSamples}、\nolinkurl{VariableRegisteredJointCheck} 的依赖行索引样本类型对应；任意 Y 扩展的精确接口范围见附录 F.3.1。

规范 \(G\) 的定义含未知真实源 \(p_j\)。因此它是分析中的真实校准事件，不能因为写成有限合取就称为``训练后可直接观察到是否成立''。可执行选择规则应对全部记录定义；定理只在它恰好属于 \(G\) 时使用校准保证。

\hypertarget{ux4eceux9010ux884cux62dfux5408ux5230ux9009ux5b9aux6761ux4ef6ux98ceux9669}{%
\section{38.4 从逐行拟合到选定条件风险}\label{ux4eceux9010ux884cux62dfux5408ux5230ux9009ux5b9aux6761ux4ef6ux98ceux9669}}

对每个训练记录，可以选择商映射 \(q_\omega\)、初始律、参考核 \(\widehat P^\omega\) 和因果策略 \(\pi^\omega\)。微观源 \(P\) 保持同一身份，注册行也仍是训练前固定的行。

假定仅在 \(G\) 上：

\begin{enumerate}
\def\labelenumi{\arabic{enumi}.}
\tightlist
\item
  在注册源行上，经验微观响应经 \(q_\omega\) 推送后，与所选参考行的 TV 拟合误差不超过 \(\eta_n(\omega)\)。
\item
  所选实际策略的未认证占据满足 \(\nu_n^\omega(C_n^c)\le\kappa_n(\omega)\)。
\item
  宏观事件 \(F_\omega\) 覆盖共同微观失败 \(E\)，并且参考风险加路径预算不超过某个 \(\beta\)。
\end{enumerate}

单个真实行到所选参考行的比较只需要

\[
\varepsilon_n(\omega)=\min\{1,\rho_n+\eta_n(\omega)\}.
\]

第一段是真实微观行与经验行的校准，经 Markov 推送不放大 TV；第二段是经验推送与参考行的拟合。因此这条路线使用 \(\rho+\eta\)，不应无理由再加第二个 \(\rho\)。若实际采用不同的代表行比较链，必须按其真实链条重新计费。

由 M-36.4，在 \(G\) 上实际条件风险满足

\[
r(\omega)\le b(\omega):=
\min\left\{1,\
\Pr_{\widehat p_T^\omega}(F_\omega)+
\sum_{n<T}\bigl[
\varepsilon_n(\omega)+
(1-\varepsilon_n(\omega))\kappa_n(\omega)
\bigr]\right\}.
\]

若 \(b(\omega)\le\beta\) 对所有 \(\omega\in G\) 成立，M-38.3 给总风险界。选择器可以在训练后改变模型，但同一个 \(G\) 已对全部注册行同时成立，所以不因``选择了一个候选''再做一次候选并合。

\textbf{Lean。} \nolinkurl{coverage_empirical_registered_microscopic_failure_risk}；\nolinkurl{registered_selected_joint_failure_of_cap_on_good}。此端点的 hfit、hcap、hchecked 只在 Good 上要求。数据依赖版本则直接消费 M-38.2，不需要把所有训练记录压成同一个 \(\beta\)。

参考事件概率与初始 PMF 仍是所选模型的输入。除非另给计算或区间认证程序，不能把此数学表达式整体称为``完全由可观察数据算出的量''。

\hypertarget{ux7528ux652fux6301ux751fux5b58ux5b9aux7406ux6d88ux9664ux8986ux76d6ux8f93ux5165}{%
\section{38.5 用支持生存定理消除覆盖输入}\label{ux7528ux652fux6301ux751fux5b58ux5b9aux7406ux6d88ux9664ux8986ux76d6ux8f93ux5165}}

\textbf{M-38.4：支持可行性接入统计联合风险。} 假定在同一个 \(G\) 上，对每个选定候选的实际 \(p_0^\omega,q_\omega,P,C\)，其全部非空初始观察支持满足 M-37.3 的后向可行性；使用对应 viablePolicy，且保留完整 fallback 策略作为输入。则不再需要另行输入 \(\kappa\) 或实际 \(\nu\) 上界：定理内部得到

\[
\forall n<T,\quad \nu_n^\omega(C_n^c)=0.
\]

若同一所选策略的经验拟合与参考预算在 \(G\) 上满足

\[
\min\{1,\Pr_{\widehat p_T^\omega}(F_\omega)
+\sum_{n<T}\varepsilon_n(\omega)\}\le\beta,
\]

则真实联合风险仍满足 \(\min\{1,\alpha+(1-\alpha)\beta\}\)。

\textbf{证明。} 逐记录应用 \nolinkurl{viablePolicy_zero_uncovered}，得到与其实际源核和策略一致的零覆盖缺口；以 \(\kappa=0\) 代入上一节 Good-only 端点。好事件始终是同一个 \(G\)，所以只计一次 \(\alpha\)。\(\blacksquare\)

\textbf{Lean。} 原 \nolinkurl{ViableRegisteredJointCheck} 中 viableFamilyPolicy、\nolinkurl{viable_registered_selected_joint_failure_on_good} 已生产成立。第 38.10 节的 \nolinkurl{ViableVariableRegisteredJointCheck} 复用同一个 viableFamilyPolicy，将逐行样本与选中 Good 包络相接，并在同一 Good 上内部推出自己的零占据；其具体版本范围见该节。第 38.6 节的具体采样例则用直接的实际零覆盖证明，不冒充该选择策略的专门实例。

\hypertarget{ux4e00ux4e2aux5b8cux6574ux7684ux6570ux636eux9009ux62e9ux771fux5b9eux90e8ux7f72ux4e0eux8054ux5408ux98ceux9669ux89c1ux8bc1}{%
\section{38.6 一个完整的数据选择、真实部署与联合风险见证}\label{ux4e00ux4e2aux5b8cux6574ux7684ux6570ux636eux9009ux62e9ux771fux5b9eux90e8ux7f72ux4e0eux8054ux5408ux98ceux9669ux89c1ux8bc1}}

取 \(X=A=\{0,1\}\)、\(q=\mathrm{id}\)、时域 \(T=1\)，初始状态固定为零。实际下一状态为一是失败，两个动作的源行分别满足

\[
p_0(1)=\frac18,\qquad p_1(1)=\frac14.
\]

这里 \(p_a\) 的下标表示动作，不是初始律。源行与当前状态无关；注册认证只覆盖时刻零、当前状态零上的两个动作行。

训练对每条注册行取 \(N=256\) 个独立响应，并用两个行产品构造完整训练 PMF。因此样本记录有 \(512\) 个布尔坐标。定义选择规则

\[
A(\omega)=\text{动作零源行的第一个响应}.
\]

选择器使用真实训练数据，并不声称它最大化某个学习目标。参考模型采用每个动作行的全部 \(256\) 个响应得到的经验律，而部署动作只依赖上述第一个坐标。

\hypertarget{ux9009ux62e9ux786eux5b9eux6539ux53d8ux5b9eux9645ux90e8ux7f72ux5f8b}{%
\subsection{38.6.1 选择确实改变实际部署律}\label{ux9009ux62e9ux786eux5b9eux6539ux53d8ux5b9eux9645ux90e8ux7f72ux5f8b}}

第一个坐标的真实边缘是 \(p_0\)，故

\[
\Pr(A=0)=\frac78,\qquad \Pr(A=1)=\frac18.
\]

两个选择概率都严格正，所选实际未来失败概率分别为 \(1/8\) 和 \(1/4\)，确为不同部署 PMF。候选索引可以包含整个记录，因为参考模型读取全部经验数据；这不改变实际部署由 \(A(\omega)\) 决定的事实。

\hypertarget{ux6ce8ux518cux8986ux76d6ux6765ux81eaux540cux4e00ux5b9eux9645ux5b9eux9a8c}{%
\subsection{38.6.2 注册覆盖来自同一实际实验}\label{ux6ce8ux518cux8986ux76d6ux6765ux81eaux540cux4e00ux5b9eux9645ux5b9eux9a8c}}

时域只有一次动作，初始状态必为零，两动作源行都已注册，故对每个训练记录

\[
\nu_0^\omega(C_0^c)=0.
\]

这不是关于训练行分布的推测，而是直接由实际初始 PMF 和策略扩张所得。经验参考正是对应经验行，故拟合余项 \(\eta=0\)。

\hypertarget{ux597dux4e8bux4ef6ux4e0eux9884ux7b97}{%
\subsection{38.6.3 好事件与预算}\label{ux597dux4e8bux4ef6ux4e0eux9884ux7b97}}

取 \(\rho=1/8\)。在规范好事件上，两条经验行的失败概率分别至多

\[
\widehat p_0(1)\le\frac18+\frac18=\frac14,\qquad
\widehat p_1(1)\le\frac14+\frac18=\frac38.
\]

因此无论数据选择哪个动作，参考失败不超过 \(3/8\)，加路径误差 \(1/8\) 后不超过

\[
\beta=\frac12.
\]

两个注册行、布尔响应的四个子集以及 Hoeffding 的两侧系数给

\[
\alpha
=2\cdot4\cdot2\exp\!\left(-2\cdot256\cdot(1/8)^2\right)
=16e^{-8}\le\frac1{16}.
\]

已有指数数值界给最后一步；本轮没有凭浮点近似替代形式不等式。因此

\[
\Pr_J(E)\le\frac12+8e^{-8}\le\frac{17}{32}<1.
\]

\hypertarget{ux7cbeux786eux5b9eux9645ux98ceux9669}{%
\subsection{38.6.4 精确实际风险}\label{ux7cbeux786eux5b9eux9645ux98ceux9669}}

由 M-38.1 和真实选择边缘，

\[
\Pr_J(E)=
\frac78\cdot\frac18+\frac18\cdot\frac14
=\frac9{64}.
\]

这证明形成的确是真实训练---所选未来联合律，而不是把一个恒定策略风险重新命名为``适应选择''。\(17/32\) 是通用校准方法给出的保守证书，\(9/64\) 是这个完全指定有限模型的精确值；两者不应混称最优或紧界。对已完全知道的这两个源行，直接使用条件风险不超过 \(1/4\) 也能给更好的通用上界，说明统计证书的目的在于验证其数据校准链条，并非击败已知模型的精确计算。

\hypertarget{ux4e3aux4ec0ux4e48ux4e0dux80fdux8981ux6c42ux5168ux90e8ux8bb0ux5f55ux7684ux9884ux7b97ux90fdux5c0f}{%
\subsection{38.6.5 为什么不能要求全部记录的预算都小}\label{ux4e3aux4ec0ux4e48ux4e0dux80fdux8981ux6c42ux5168ux90e8ux8bb0ux5f55ux7684ux9884ux7b97ux90fdux5c0f}}

考虑所有 \(512\) 个训练坐标均为一的记录 \(\omega_{\rm all}\)。其实际概率严格正；在产品模型中可直接算为

\[
\left(\frac18\right)^{256}\cdot
\left(\frac14\right)^{256}>0.
\]

在该记录上，选择动作一，两条经验行都成为纯失败分布。因此截断后的参考风险加校准预算为一，不满足 \(\beta=1/2\)。

这说明旧的``对所有训练记录都通过小预算检查''的假设对这个真实采样实验不成立。新 Good-only 端点只在校准事件使用预算，坏记录由实际坏质量收费，恰好覆盖此情形。全一记录也确实不属于规范 \(G\)。

\textbf{Lean。} \nolinkurl{RegisteredSelectionWitness} 中 selectedAction\_law、selectedAction\_pmf、selection\_probability、\nolinkurl{actual_deployments_distinct}、\nolinkurl{two_positive_deployment_choices}、checkedBudget\_on\_good、confidenceBudget\_eq、confidenceBudget\_small、\nolinkurl{registered_selected_joint_bound}、\nolinkurl{actual_joint_failure_exact}、allTrue\_checkedBudget、all\_record\_budget\_false、\nolinkurl{allTrue_positive_probability}、\nolinkurl{registered_selection_joint_witness}。产品采样及各坐标的实际边缘、独立性复用既有 StatisticalRiskWitness；不重新假设样本来自目标源。

\hypertarget{ux6df7ux5408ux754cux7684ux7d27ux6027ux4e0eux5fc5ux8981ux6761ux4ef6}{%
\section{38.7 混合界的紧性与必要条件}\label{ux6df7ux5408ux754cux7684ux7d27ux6027ux4e0eux5fc5ux8981ux6761ux4ef6}}

\textbf{一个紧的两分支例。} 训练坏事件质量 \(1/4\)，好事件质量 \(3/4\)；好分支未来失败概率 \(1/2\)，坏分支为一。选择器在两个正质量分支选到不同实际部署律，则

\[
\Pr_J(E)=\frac14+\frac34\cdot\frac12=\frac58.
\]

这达到 M-38.3 的 \(\alpha+(1-\alpha)\beta\)。Lean 的 \nolinkurl{sharp_joint_failure_probability}、\nolinkurl{sharp_mixture_is_attained} 和 \nolinkurl{selected_deployments_are_nonconstant} 给该例。因此仅凭一般的 \(\alpha,\beta\) 两项信息，不能删去坏事件收费。

\textbf{固定策略的更强结论。} 若实际未来律完全不依赖训练记录，而且 \(G\) 具有正质量，那么 \(G\) 非空，取一个好记录便得到该固定律本身的风险上界。不必再为同一个恒定风险付一次 \(\alpha\)。旧 StatisticalRiskWitness 的固定部署风险为 \(143/256\) 的证书，经新联合律仍可保持 \(143/256\)；不把较弱的 \(685/1024\) 当作主要进步。

\textbf{条件部署律必须真实。} 如果证书用的未来核为纯安全律，却把外部实际部署改成纯失败律，前者当然不能约束后者。\nolinkurl{wrong_conditional_law_counterexample} 用真实构造展示此必要边界。即使分别知道训练边缘和某些候选的无条件边缘，也不能无条件识别数据选择后的条件核。当前正向定理解决的是已明确构造的实验；把外部系统识别为这个生成实验仍需模型证据。

\hypertarget{ux5b9eux9645ux5f62ux5f0fux68c0ux67e5ux4e0eux672cux4e66ux7684ux8bc1ux660eux5730ux4f4d}{%
\section{38.8 实际形式检查与本书的证明地位}\label{ux5b9eux9645ux5f62ux5f0fux68c0ux67e5ux4e0eux672cux4e66ux7684ux8bc1ux660eux5730ux4f4d}}

新增数学的检查来自固定源码的真实 Lean 编译与打印公理结果。下表是本轮数学草稿检查的历史事实，不代替后续正式模块拆分的构建，也不提前宣称它们已进入主干。

\begin{longtable}[]{@{}
  >{\raggedright\arraybackslash}p{(\columnwidth - 4\tabcolsep) * \real{0.3333}}
  >{\raggedright\arraybackslash}p{(\columnwidth - 4\tabcolsep) * \real{0.3333}}
  >{\raggedright\arraybackslash}p{(\columnwidth - 4\tabcolsep) * \real{0.3333}}@{}}
\toprule\noalign{}
\begin{minipage}[b]{\linewidth}\raggedright
检查
\end{minipage} & \begin{minipage}[b]{\linewidth}\raggedright
实际结果
\end{minipage} & \begin{minipage}[b]{\linewidth}\raggedright
可据此使用的结论
\end{minipage} \\
\midrule\noalign{}
\endhead
\bottomrule\noalign{}
\endlastfoot
check-03，支持桥与策略双向定理 & 完整退出 0；12 项打印仅标准三公理 & M-37.1--M-37.3 的实际 PMF 接通与 iff \\
check-04，注册 Good 端点 & 该单项退出 0；同批有限检验第一稿失败 & Good-only 端点已通过；整批保持历史失败状态 \\
check-06，关系检验及薄 A→B 桥 & 两项均退出 0；检验实际输出 true、false、true & M-37.4、死路见证与 M-38.4 \\
check-08，完整注册选择见证第三稿 & 完整退出 0；12 项标准公理，无警告 & 第 38.6 节的实际采样、选择、联合风险和全一记录边界 \\
前一轮严格后验支持见证第二稿 & 该单项退出 0；当时另一独立扩展第一稿失败 & M-37.5；不重写混合批次的失败历史 \\
\end{longtable}

标准公理集合是 propext、Classical.choice、Quot.sound 的子集；这些打印结果没有 sorryAx。它们说明相应声明在所列 Lean 环境和假设下经过内核检查，不说明真实物理模型已经被实验证实，也不说明每项声明都有构造性执行程序。确切源路径、主定理全限定名、编译源码哈希与历史输入关系见附录 F 和配套来源索引。

\hypertarget{ux6570ux5b66ux95edux73afux7684ux5b8cux6210ux8303ux56f4}{%
\section{38.9 数学闭环的完成范围}\label{ux6570ux5b66ux95edux73afux7684ux5b8cux6210ux8303ux56f4}}

本书现已把下列关系逐项接通：给定实际源和注册信息，定义真实有限路径；从该路径产生条件后验与状态---动作占据；以支持递推判定部分观察下的有限可行性并构造策略；将所得实际覆盖与注册经验校准组合；对训练后选择的条件部署核构造真正联合律，最后得到共同微观失败事件的概率界。

这个闭环仍以明确数学输入为起点：实际源、观察协议、精确支持或其有效表示、训练生成协议、失败事件和任务认证条件。将它用于开放世界时，必须检验这些输入。有限支持可行性不是未知转移概率的自动学习器，有限时域零占据不是无限寿命定理，通用风险证书不是目的来源的唯一识别。正因为这些边界已经进入定理而不是隐藏在解释中，后续经验和哲学研究才有明确可检验的接口。

\hypertarget{ux9010ux884cux771fux5b9eux91c7ux6837ux4e0eux540cux4e00ux826fux4e8bux4ef6ux4e0aux7684ux9009ux4e2dux5305ux7edc}{%
\section{38.10 逐行真实采样与同一良事件上的选中包络}\label{ux9010ux884cux771fux5b9eux91c7ux6837ux4e0eux540cux4e00ux826fux4e8bux4ef6ux4e0aux7684ux9009ux4e2dux5305ux7edc}}

本节将前三节的数学范围与新的实际类型相接。\nolinkurl{VariableRegisteredSamples}、VariableRegisteredJoint、\nolinkurl{ViableVariableRegisteredJoint} 已在本版生产模块中形式化；联合端点允许任意 Y，并明示同一 Good 上选中事件的可测性。它们与第 36.8 节的四个模块构成本版所述的概率接口扩展。原统一 Fin N 接口、已有见证与首版有限 Y 的历史编译记录保留。

\hypertarget{ux7528ux771fux5b9eux5750ux6807ux5b9eux9a8cux5b9eux73b0ux4e0dux7b49ux6837ux672cux91cf}{%
\subsection{38.10.1 用真实坐标实验实现不等样本量}\label{ux7528ux771fux5b9eux5750ux6807ux5b9eux9a8cux5b9eux73b0ux4e0dux7b49ux6837ux672cux91cf}}

固定有限注册行集 J、有限响应空间 X、真实行 \(p_j\in\mathrm{PMF}(X)\)、样本分配 \(N:J\to\mathbb N\) 和半径 \(\rho:J\to\mathbb R\)，其中每条实际行满足 \(N_j>0,\rho_j\ge0\)。这些输入在观察数据前确定。定义坐标类型及产品实验

\[
I_N=\{(j,k):j\in J,\ 0\le k<N_j\},\qquad
\Omega_N=X^{I_N},\qquad
\mu_N=\bigotimes_{(j,k)\in I_N}p_j.
\]

样本 \(S_{j,k}(\omega)=\omega(j,k)\) 的真实边缘恰为 \(p_j\)，每行样本独立。于是实际经验律由属于该行的 \(N_j\) 个坐标组成，不是给统一长度数组补上没有采样含义的重复项。

\textbf{M-38.5：逐行真实注册校准。} 上述实验满足

\[
\mu_N(G^c)\le
\alpha_N:=\sum_{j\in J}2\cdot2^{|X|}
\exp(-2N_j\rho_j^2),
\quad
G=\{\omega:\forall j,\operatorname{TV}(p_j,\widehat p_j(\omega))\le\rho_j\}.
\]

同一结论适用于任何真实训练测度及样本族，只要逐行样本可测、独立且边缘等于相应真实行；不同注册行之间不要求独立。

\textbf{证明。} 对每个 j 使用其实际 \(N_j\) 的单行校准，再取好事件交集并对补集作有限并合。该交集包含于规范 G，所以规范 G 的坏质量不超过同一个预算。产品实验的坐标投影定理分别给可测性、正确边缘及所需独立性。\(\blacksquare\)

\textbf{Lean。} \nolinkurl{VariableRegisteredSamples.productRowLaw}、productRowSample\_law、\nolinkurl{productRowSample_independent}、\nolinkurl{registered_row_calibration}、\nolinkurl{product_experiment_calibration}。一般单行组合不要求训练空间有限；后续联合部署端点仍采用有限训练空间。若 J 为空，坐标实验仍良定义、预算为零、所有逐行条件为空合取；这不替部署自动提供非空可行性或零物理失败风险。

\hypertarget{ux53eaux5728-good-ux4e2dux5b9eux9645ux9009ux4e2dux7684ux5019ux9009ux4e0aux4f7fux7528ux5305ux7edc}{%
\subsection{38.10.2 只在 Good 中实际选中的候选上使用包络}\label{ux53eaux5728-good-ux4e2dux5b9eux9645ux9009ux4e2dux7684ux5019ux9009ux4e0aux4f7fux7528ux5305ux7edc}}

现固定生产注册类型 \(J=\mathrm{CoverageRegisteredRow}(C,T)\)，其中每个 j 指明实际时刻 \(n_j<T\)、状态、动作与认证证明。真实源 P、注册条件 C、时域 T、N、误差半径及微观失败事件 E 都在数据前固定。候选 c 可带观察映射、初始律、参考核、因果策略及粗事件 F；选取 \(c=s(\omega)\) 可以依赖全部训练记录。

联合端点目前使用对每个 \(n<T\) 非负的时间半径 \(\rho_n\)，因此其预算是 M-38.5 中取 \(\rho_j=\rho_{n_j}\) 的特例。这里不能把 Samples 层更一般的逐行半径也直接写成当前 Joint 的完整签名。

在同一个规范 G 上要求所选候选的经验拟合、实际占据上界及数值预算成立。采用第 36.8 节的一般 Y 路线时，还要求该记录选中的粗事件 F 可测；事件包络只要求

\[
\forall\omega\in G,\qquad
E\subseteq H_{q_{s(\omega)}}^{-1}(F_{s(\omega)}).
\]

这不是对所有候选 c 的包络要求，也不要求坏训练记录上的所选候选满足包络。失败事件 E 仍是同一个固定微观事件；不能在看到记录以后重定义何谓失败。

\textbf{M-38.6：不等样本的选中 Good 联合风险。} 若上述共同 G 上所选策略的参考事件风险加实际路径预算经截断后至多 \(\beta\in[0,1]\)，则真实训练---选择---部署联合律满足

\[
J(E)\le \min\{1,\alpha_N+(1-\alpha_N)\beta\}.
\]

\textbf{证明。} M-38.5 给 \(\mu(G^c)\le\alpha_N\)。对每个 \(\omega\in G\)，从该记录的真实行校准与经验拟合局部导出源行误差，再用该记录选中策略自己的实际占据和包络，得到其真实条件部署失败概率至多 β。最后分解同一个联合律在 G 与补集上的质量，应用 M-38.1 的混合界。不存在第二个需要重新收费的``包络好事件''。\(\blacksquare\)

\textbf{Lean。} \nolinkurl{VariableRegisteredJointCheck.RegisteredCalibrationGood}、\nolinkurl{registeredConfidenceBudget}、\nolinkurl{registered_selected_joint_failure_of_cap_on_good}。样本类型为 \(\forall j,\mathrm{Fin}(N_j)\to\Omega\to X\)；hfit、hcap、hF、henvelope、hchecked 均限定于同一个 \nolinkurl{RegisteredCalibrationGood}。该接口保留 Ω、X、A 有限，Y 具有可测空间及可测单点，hF 明示 Good 上所选粗事件的可测性。候选类型不要求有限。

置信预算的来源因而十分具体：它由固定注册行、实际 \(N_j\)、源响应空间 X 与半径决定。原始行先在 X 上校准，训练后选出的 q 再通过 TV 数据处理使用同一校准；不会仅因候选类型很大，或沿第 36.8 节数学路线允许整体 Y 无限，而自动产生另一次选择失败预算。这个结论不表示任意选后模型都免费正确：经验拟合、所选事件包络、原事件的可测性以及数值风险检查仍各自承担其明确条件，并且要在同一个 G 上相接。

\hypertarget{ux540cux4e00ux652fux6301ux751fux5b58ux7b56ux7565ux5185ux90e8ux6d88ux53bbux5360ux636e}{%
\subsection{38.10.3 同一支持生存策略内部消去占据}\label{ux540cux4e00ux652fux6301ux751fux5b58ux7b56ux7565ux5185ux90e8ux6d88ux53bbux5360ux636e}}

\textbf{M-38.7：可行策略接入逐行 Good 联合风险。} 若 M-38.6 所选候选的全部非空初始观察支持在同一个 G 上满足第 37 章的 Viable 条件，并以显式 fallback 构造 viableFamilyPolicy，则该策略自己的未认证占据在 G 上每步为零。其联合风险只需使用参考事件风险加

\[
\sum_{n<T}\min\{1,\rho_n+\eta_n(\omega)\},
\]

并满足同一个混合界。

\textbf{证明。} 固定好记录，应用实际支持生存定理得到 viableFamilyPolicy 在各 \(n<T\) 的零占据；以 \(\kappa=0\) 代入 M-38.6。参考路径、真实部署核和内部覆盖证明全部使用该同一策略。\(\blacksquare\)

\textbf{Lean。} \nolinkurl{ViableVariableRegisteredJointCheck.viable_registered_selected_joint_failure_on_good}。hinit、hfit、hF、henvelope、hchecked 共享 \nolinkurl{VariableRegisteredJointCheck.RegisteredCalibrationGood}；调用者不再提供 hcap。该接口与上一节具有相同的观察及训练空间范围。这个消去不包含零物理失败的结论，也没有为空动作类型创造完整策略；fallback 仍是明确输入。

这些新连接不把相关单轨迹采样、训练后自行改变 N、额外未注册行或任意非原子参考核纳入结论。它们把已经明确的实际采样与部署条件进一步接紧，使第 31.1 节的共同证书装配少一组重复前提，而不是通过重命名输入消去真实研究问题。

\hypertarget{ux6267ux884cux6545ux969cux590dux5236ux4e0eux76eeux7684ux7684ux540cux6e90ux6784ux9020}{%
\chapter{39. 执行、故障、复制与目的的同源构造}\label{ux6267ux884cux6545ux969cux590dux5236ux4e0eux76eeux7684ux7684ux540cux6e90ux6784ux9020}}

\hypertarget{ux54eaux4e9bux6761ux4ef6ux4ecdux7531ux6a21ux578bux627fux62c5}{%
\section{39.1 哪些条件仍由模型承担}\label{ux54eaux4e9bux6761ux4ef6ux4ecdux7531ux6a21ux578bux627fux62c5}}

前面的接口把程序执行、修复、父对象、资源和评价分开陈述。本章给出一个共同的非退化源族，实际推导其中若干原先独立的桥。它不是对任意世界的自创生证明：执行器及解释器可靠，保留程序受保护，复制按声明规则继承程序；环境模式与出生标签按给定的独立概率律产生。动作权限、价格、折扣、初始历史与资源供给也明确登记。

取有限寄存器集 \(I\)、指定执行位 \(i\in I\)，物理状态为 \(x\in\{0,1\}^I\)。保留程序位为 \(b\)。允许动作只有 idle、write-false、write-true；\(U_i(x,a)\) 按动作保持或写入第 \(i\) 位。实际解释器先读取 \(x_i\)：若 \(x_i=b\) 就 idle，否则 write-\(b\)。有限指令的解释证明其动作就是此策略 \(\pi_b\)，不是另用选择公理寻找一个实现相同目标的程序。

每步模式 \(R_t\) 的修复概率为 \(r=1-\epsilon\)，故障概率为 \(\epsilon\in[0,1]\)。物理转移是 \[
X_{t+1}=\begin{cases}
U_i(X_t,A_t),&R_t=1,\\
X_t[i\leftarrow\neg b],&R_t=0.
\end{cases}
\] 故障把指定寄存器置成程序目标的反值，保留程序和控制器。它不是对任意受损位再翻转一次；后一规则可能偶然修好损伤，属于不同源。源模式的 PMF 推送等于 P5 的物理故障/修复混合行，也等于本章控制模型的完整物理行。核相等由执行定义证明，没有作为调用者输入。

\hypertarget{ux4feeux590dux80fdux529bux5b9eux9645ux51faux751fux4e0eux5386ux53f2}{%
\section{39.2 修复能力、实际出生与历史}\label{ux4feeux590dux80fdux529bux5b9eux9645ux51faux751fux4e0eux5386ux53f2}}

联合组织状态中，有限修复载体恰由保留程序正确刻画。实际势能与故障包络由上述转移构造，故 P5 所需的执行和故障条件得到来源。这里的载体进入表示仍有恢复能力；寄存器在故障时可以不满足物理目标，不能从载体进入推出它以后每刻健康。

在修复模式中，动作后的目标位正确才产生一个实际子代；失败记为 None。子代字面复制程序及动作后物理状态，其历史为亲本历史加上本次标签。由实际出生可推出子代满足目标、继承程序，并把同一个子代作为种子接入同一 \(\epsilon\) 的 P5。没有重试，也没有按成功删去失败。

对 repair、damage、birth 指令组成的实际日志，执行后的历史等于原历史连接日志中按次序出现的出生标签。因此父身份桥来自真实执行历史。修复不增加出生标签，出生则严格增加代数。两个未标记物理/程序记录可以相同而真实历史不同，所以仅凭前者不能恢复后者。标签历史描述这一谱系，不保证跨时刻、并发人口中每个子代都有全局唯一标识。

这一点也修订了一个容易混淆的有限模型。若有限非空的永久个体表上每个正繁殖边都严格增加代数，就不可能每一行的繁殖质量都为一：取代数最大的个体，其全部出边只能非正，与行和一矛盾。正确的有限 Price 消费者以亲本类型和子代事件标签作矩形索引，从实际出生标签律给出归一化传递核；亲本频率、繁殖强度及性状赋值仍是输入。类型和事件标签不等于永久个体。严格个体谱系不能静默替换成类型链。 第 39.20---39.21 节在另登记的串行人口模型中进一步从实际新生槽位推导个体父身份、完整旧亲本队列及真实子代频率；它没有沿用这个不可能的有限归一严格谱系组合。

\hypertarget{ux540cux4e00ux8bb0ux5f55ux7684ux6536ux76caux4e0eux8d44ux6e90}{%
\section{39.3 同一记录的收益与资源}\label{ux540cux4e00ux8bb0ux5f55ux7684ux6536ux76caux4e0eux8d44ux6e90}}

设每次受控写入费用为 \(c\ge0\)，idle 费用为零，登记价格为 \(\lambda\ge0\)。实际记录同时保留动作、模式、标签、亲本、子代及费用。记 \[
B_t=\mathbf1_{\{R_t=1\}}\mathbf1_{\{U_i(X_t,A_t)_i=b\}},\qquad
C_t=c\,\mathbf1_{\{R_t=1\}}\mathbf1_{\{A_t\ne\mathrm{idle}\}}.
\] \(B_t\) 是本次实际出生计数。故障时没有受控写费用，也没有出生。物理损伤、代码拷贝与重新绑定在这个费用模型中不另收费，不能把 \(C_t\) 称为全部能耗。给定 \(0<d<1\)，有限评价为 \[
J_T(\pi)=E\!\left[\sum_{t<T}d^t(B_t-\lambda C_t)\right].
\] 初始策略记忆中的当前状态须与亲本实际物理状态一致，亲本保留程序须为登记目标；编译策略的指定种子自动满足这些条件。完整有限记录的 PMF 由同一策略、模式和标签逐步生成。直接证明其记录积分可积，并等于同一策略在上述物理控制模型中的截断值。策略可有原有的物理/动作因果记忆；模式和标签虽然被记录，不因此升级为策略观察。

实际付费记录还接入原 P6 资源账。对任意实际动作与模式序列，每步 \(0\le C_t\le c\)，所以初始储备非负、每步供给 \(s\ge c\) 时，所有累计余额非负。这是逐路径的充分供给条件，供给仍是设计输入，并未从平均收益或繁殖成功中创造出来。原有限记录之外补零只定义有限时域的费用消费者。新增无限解释器按实际事件递次延长同一历史，得到概率一的前缀相容性；其完整无限记录律的每个有限前缀恰是上述 recordLaw。对该实际无限费用流，P6 的同一供给条件也逐路径成立，无需补零。

\hypertarget{ux89e3ux91caux51faux7684ux540cux4e00ux7b56ux7565ux786eux5b9eux6700ux4f18}{%
\section{39.4 解释出的同一策略确实最优}\label{ux89e3ux91caux51faux7684ux540cux4e00ux7b56ux7565ux786eux5b9eux6700ux4f18}}

在充分条件 \(\lambda c\le1\) 下，\(\pi_b\) 在原全部因果策略类中最优，且对所有 \(\epsilon\in[0,1]\) 使用同一个程序。其健康与受损状态值分别为 \[
G=\frac{r(1-d\epsilon\lambda c)}{1-d},\qquad H=G-r\lambda c.
\] 证明先按同一模式律求实际单步出生减费用，再核对 Bellman 不等式以及 \(\pi_b\) 对该不等式取等，最后调用原因果最优性端点。条件只声称充分性。这里的新故障目标重新从实际记录推导，没有把旧无故障模型的路径积分搬过来。

取单寄存器、目标 \(b=\mathrm{true}\)、公平出生标签，\(\epsilon=1/4,c=1,\lambda=d=1/2\) 时，\(G=45/32,H=33/32\)。从受损初态出发，指定首步 write-true、次步 idle、先修复后故障的模式及标签 true、false 的完整两步 tape 事件质量为 \(3/64\)，折扣净收益为 \(1/2\)（实际写费用为 \(1\)，价格折算为 \(1/2\)）；第二步的失败仍留在记录中。同一程序的有限记录期望趋于最优控制模型的值。进一步，对上述实际无限解释器，单步收益绝对值不超过 \(1+\lambda c\)；几何折扣给绝对可和及可积性。支配收敛与已证明的前缀律据此推出实际无限记录积分恰为同一原因果策略的无限值，编译程序的实际积分即为上述最优值。原策略的任意 Type0 内部记忆由可数事件字重放，不增加有限记忆限制；模式和标签始终不进入原策略的信息权限。

\hypertarget{ux672aux77e5ux6545ux969cux7387ux4e0bux7684ux76eeux7684ux6765ux6e90}{%
\section{39.5 未知故障率下的目的来源}\label{ux672aux77e5ux6545ux969cux7387ux4e0bux7684ux76eeux7684ux6765ux6e90}}

真实 write-false、write-true 的出生读出及受控写费用读出给出三个精确均值 \[
\bigl(r\mathbf1_{\{b=\mathrm{false}\}},\ r\mathbf1_{\{b=\mathrm{true}\}},\ rc\bigr).
\] 当 \(r>0\) 时，前两项之和识别 \(r\)，非零项识别 \(b\)，第三项除以 \(r\) 识别 \(c\)。这些均值从动作后的实际出生/付款计算，且与亲本当时物理状态和历史无关；无需额外授予重置权限。由此先识别源参数，再推出完整记录评价的策略排序相同，供应原 C6 的忠实性消费者。价格、折扣和评价合同仍为登记输入，机制没有独自决定伦理评价。

当 \(\epsilon=1\) 时，这些均值全零，源参数不能这样识别；但登记控制值也全零。这并不证明全无差别的目的排序类不可识别。精确实验均值访问也不等于有限采样估计、已执行反演算法或选后风险保证；将本章机制接到统计校准仍需自己的共同消费者。保留程序的控制器本身不读取真实故障率。

\hypertarget{ux5df2ux6838ux9a8cux6765ux6e90ux4e0eux4ecdux987bux5b8cux6210ux7684ux8fdeux63a5}{%
\section{39.6 已核验来源与仍须完成的连接}\label{ux5df2ux6838ux9a8cux6765ux6e90ux4e0eux4ecdux987bux5b8cux6210ux7684ux8fdeux63a5}}

上述来源分别见 RegisterRepairSource、RegisterRepairHistory、RegisterFaultSource、RegisterResourceSource、\nolinkurl{RegisterFaultControlSource}、\nolinkurl{RegisterFaultControlModel}、RegisterFaultJointRecord、RegisterBirthMaintenance、\nolinkurl{RegisterFaultPurposeInformation} 和 \nolinkurl{RegisterFaultInfiniteRecord}。它们调用原 P5/P6/P12/C6 与因果控制端点。直接源码对应及核验回执见\href{../formalization/ueot-core/docs/compression/theory_completion/unified_closure/register_fault_resources_v1/RESULTS_AND_OPEN_OBLIGATIONS.md}{故障来源包}。

一般预测方面，\nolinkurl{MeasurablePredictiveSource/Sensor} 已从共享噪声的可测更新及不读取隐藏坐标的观察器推导实际后验的几乎处处投影和同策略递归；允许连续噪声及部分、有噪观察。隐藏坐标控制下一可见状态、破坏可见自治时，已有严格泄漏反例。在此基础上，\nolinkurl{MeasurablePredictiveJointPath} 保持同一时变可见 belief 策略的完整有限动作---观察---后验记录。\nolinkurl{MeasurablePredictiveLatentPath} 从初始潜状态 \(X_0\sim b_0\) 出发，按原 P/O 保留并转移实际潜状态，同时由实际观察更新同一 posterior。一步解体恒等式给随机 lift 的相容性；前缀归纳证明完整观察记录与末态潜变量的联合律，经最后实际 posterior 分解。由此取得实际潜状态运行的有限观察记录律，并以可数前缀唯一性推广到完整可数路径。它没有逐点假设任意 belief 与潜状态配对都满足 lumpability，也没有每步从 belief 重抽物理状态。这里的可见 belief 由声明源给出，不把它说成任意过程的最小可测预测像。

这些结果使对象、程序、动作、实际历史、出生、资源和注册目的之间多了一条非空的共同来源链。下面继续给出实际有限协议的最小预测来源、完整历史策略、时戳身份和连续 Gaussian 来源。它们各自只覆盖声明的模型与策略类；整个理论仍未闭合，未覆盖协议与动作类、一般连续随机分析、相关/适应校准与发现、未覆盖模型中的形成／谱系和尺度运输、完整长期资源与全局组合仍有必需义务。编译、数学成立及现实经验支持分别报告；本章不宣称已确认自然界的普适机制。

\hypertarget{ux4eceux4e00ux822cux5b9eux9645ux6e90ux5bfcux51faux6709ux9650ux534fux8baeux9884ux6d4bux9012ux5f52}{%
\section{39.7 从一般实际源导出有限协议预测递归}\label{ux4eceux4e00ux822cux5b9eux9645ux6e90ux5bfcux51faux6709ux9650ux534fux8baeux9884ux6d4bux9012ux5f52}}

现在不要求可见动力学自治。给定任意非空标准 Borel 潜状态 \(Z\)、原转移核 \(P\) 及观察核 \(O\)，令 \(K_u\) 为允许动作词 \(u\) 的完整有序观察律：先按 \(P\) 转移，再在同一步下一潜状态上按 \(O\) 观察，随后执行余下动作。重复动作仍占不同位置，空词输出唯一空元组。对先验 \(b\) 定义 \[
M_u(b)=K_u\circ b,\qquad C(b)=(M_u(b))_{u\in A^*}.
\] 这保存每个有限动作词下的整个观察概率律，未把它缩成几个经验评分。由同一真实 P/O 联合律的解体，直接推出 \[
M_{a::u}(b)=\operatorname{Obs}_{b,a}\otimes
  \bigl(K_u\circ\operatorname{Posterior}_{b,a}\bigr).
\] 分子及延续因子来自实际长词执行，不再由调用者另给 Bayes 因子化。若 \(C(b)=C(c)\)，则同一首动作的实际观察律相等，且每个延续词的下一后验响应几乎处处相等。动作可数时，对全部有限动作词取可数交，得到一个共同概率一集合。此集合针对固定的 \(b,c,a\)，没有偷换成对不可数先验同时成立的版本。

源响应的上述递归只需观察空间可数生成。若观察空间进一步为标准 Borel，可对响应坐标 \(M_{a::u}(b)\) 本身采用统一的可测解体，构造总的响应数组更新；它不读取潜状态、源参数或一个到潜信念的 section。动作可数且单点可测时，更新对状态、动作和观察联合可测，并得到一个实际 Markov 核：同一 C1 后验转移经 \(C\) 推送，恰等于由响应数组计算的下一响应核。任意不在源像内的数组只获数学延拓，不被说成物理可实现。零观察概率分支也不因总核存在而得到经验认证。

已有连续 Gaussian 源取明确允许的整数漂移动作，给出非原子的共同消费者，原点首步观察律为 \(N(0,1)\)。这项可数动作结论没有静默覆盖全部实数动作，也没有独自证明任意自适应协议或一般最小预测像。其源码为 \nolinkurl{MeasurablePredictiveResponseRecursion}，真实消费者为 \nolinkurl{actual_C1_response_kernel_factor}。

\nolinkurl{MeasurablePredictiveCanonicalSource} 又从已经生成的实际过去和保留的末态潜变量出发，执行同一 P/O 的未来动作词，推导过去---未来联合律等于实际过去律乘以最后后验给出的未来核。这是明确的未来干预，未来动作词可以不同于生成过去所用策略的继续动作。由这条来源再调用原 PredictionDependent 的预测充分性、事件条件期望以及模零最小性；竞争统计量能够预测这些真实未来律，是该统计量充分性的定义，不是另给一个实现证书。最小性针对实际过去律并取模零意义，动作可数时适用。概率响应数组到原 Measure 数组的嵌入只涉及全部已归一化数组，不要求物理可实现像为 Borel，也不构造到潜信念的截面。

\hypertarget{ux540cux4e00ux4eb2ux672cux5728ux4e0dux540cux65f6ux523bux4ea7ux751fux7684ux540eux4ee3}{%
\section{39.8 同一亲本在不同时刻产生的后代}\label{ux540cux4e00ux4eb2ux672cux5728ux4e0dux540cux65f6ux523bux4ea7ux751fux7684ux540eux4ee3}}

原 Bool 标签可以重复：同一亲本两次成功出生后，子代的旧标签历史可以完全相同。新增来源在原事务解释外保存一个实际计数器；每次事务，包括故障和失败，都使亲本计数加一。实际子代身份为亲本身份连接本次计数和标签，子代自己的局部计数从零开始。于是可从已执行前缀推出，同一保留亲本不同事务时刻的成功出生具有不同身份，并从真实子代身份的前缀解释其亲本。

整个带时戳记录律擦除新增元数据后，严格等于原来同一策略、初态、故障率和标签源的完整记录律；动作、物理状态、程序、付费和失败都保持。指定受损单寄存器、true 程序、故障率 \(1/4\)、公平标签，三步 \nolinkurl{write-true/idle/write-true}，模式修复/故障/修复且三个已采样标签均为 true 的完整 tape 质量为 \(9/512\)。中间失败仍在，两个实际子代身份分别为 \([(0,\mathrm{true})]\) 和 \([(2,\mathrm{true})]\)，其旧 Bool 历史却同为 \([\mathrm{true}]\)。

\nolinkurl{RegisterTimestampedBirthSource} 实际消费原 P8 的修复/后代区分、原 P5 的同一子代维持，以及原 P12 的父身份发现。后者在成功出生时解释 Option(parent ID)，失败保留 None，不识别失败请求的亲本。可靠且可访问的串行元数据、初始 founder 身份与计数仍是声明条件；元数据在原受控写费用模型中不另收费。它既不是从无标签物理状态中自动生成的历史知识，也不解决任意并发人口的 founder 分配、全局身份或繁殖强度来源。

\hypertarget{ux4eceux6709ux9650-gaussian-ux5f8bux6784ux9020ux5b9eux9645ux5168ux65f6ux95f4ux6765ux6e90}{%
\section{39.9 从有限 Gaussian 律构造实际全时间来源}\label{ux4eceux6709ux9650-gaussian-ux5f8bux6784ux9020ux5b9eux9645ux5168ux65f6ux95f4ux6765ux6e90}}

恢复的连续时间接口还需要真实随机来源。对满足 Brownian 有限维律的源，可以从时间变换 \[
\tau(t)=\frac{e^{2kt}-1}{2k},\qquad
X_t=e^{-kt}\bigl(x_0+sB_{\tau(t)}\bigr),\quad k>0,
\] 推导实际推送律、确定初态、二次可积性以及 \[
E[X_t^2]=e^{-2kt}x_0^2+\frac{s^2}{2k}(1-e^{-2kt}).
\] 此期望绝对连续，几乎处处导数等于 \(E[-2kX_t^2+s^2]\)；原 p\_rec\_02 据此消费实际构造的 Dynkin 证书。这个二次观测的右时间变化率也从源导出，没有直接推出一般生成元定义域、条件转移半群或 SDE 恒等式。

给定源的前提本身又有了一个实际构造。BrownianCountableSource 从已有有限 Gaussian 律出发，解体相邻有限前缀，再迭代现有轨迹核，构造任意可数时间列的一个概率空间。每个有限坐标集合具有原 Gaussian 律，重复时间共享坐标；同一个稠密时间源有几乎处处零初态和单位时间单位方差。

BrownianProjectiveSource 进一步处理全非负实时间。每个可数柱集族及其柱形并集只用到可数多个时间：在同一个上述可数源上实现它们，用真实概率测度的并集界证明柱集 content 的 \(\sigma\) 次可加性，再调用现有 Carathéodory 扩张。所得一个全时间概率律具有所有原有限维 Brownian Gaussian 律，实际给出 IsPreBrownianReal；未把 projective limit、次可加性或 Brownian 来源当成输入。这个同一实际源接入上述 OU 恢复，并从任意确定初态给出正噪声、正时刻的严格正方差。

有限维律、样本路径连续性和联合时间可测性分别承担证明责任。原始乘积坐标过程没有被默认几乎处处连续；下一节给出实际连续修改。一般 Itô、受控 SDE/HJB 和 Girsanov 仍是全理论必需义务。数学来源构造不等于对现实噪声模型的经验确认。

\hypertarget{ux540cux4e00ux9ad8ux65afux6982ux7387ux6e90ux7684ux5b9eux9645ux8fdeux7eedux4feeux6539}{%
\section{39.10 同一高斯概率源的实际连续修改}\label{ux540cux4e00ux9ad8ux65afux6982ux7387ux6e90ux7684ux5b9eux9645ux8fdeux7eedux4feeux6539}}

BrownianContinuousSource 从上述源的 Gaussian 增量推导 \[
E[|B_t-B_s|^4]=3|t-s|^2.
\] 在每个整数紧区间取步长 \(32^{-n}\) 的网格，以相邻增量阈值 \(2^{-n}\) 和 Markov 界得到可求和的坏事件概率。对全部整数区间与全部层数构造一个可测概率一事件，格点近似在此事件上一致收敛，细分与振荡估计推出连续极限。事件外统一取零路径，因而所构造过程的每条路径都连续；每个固定时刻的极限可测，再由连续路径与逐时可测性推出联合时间可测。

还须说明它与原源的关系。Gaussian 二阶矩保证固定时刻的格点近似依概率趋于原坐标；概率一事件上的极限及依概率极限唯一性，使每个固定时刻的新过程与原坐标几乎处处相等。这里没有对不可数时刻取零集交。有限维分布只需有限交，故完整转移所有原 Gaussian 有限维律，并得到一个不再以 Brownian 证书为输入的 IsBrownianReal。其零初态和单位时间单位方差保留，非零噪声没有因构造而消失。

同一概率律和连续过程实际接回上一节的 OU：每条 OU 路径连续，任意确定初态在正噪声和正时刻下有严格正方差；每个固定时刻的 OU 推送律与原坐标构造完全相同，故原恢复界仍由同一个源消费。这没有证明一般随机积分、Itô 公式、受控 SDE 或换测度 Girsanov；自然过滤下的条件增量与这些一般消费者仍各有来源责任。

\hypertarget{ux539fux7b56ux7565ux8bfbux53d6ux5b8cux6574ux5b9eux9645ux5386ux53f2}{%
\section{39.11 原策略读取完整实际历史}\label{ux539fux7b56ux7565ux8bfbux53d6ux5b8cux6574ux5b9eux9645ux5386ux53f2}}

令控制器为完整有限历史 \(h\in(A\times W)^*\) 上的可测 Markov 核，时间可由历史长度取得。历史可测空间由长度和每个有缺位标记的有序事件坐标生成，允许非原子观察；没有给所有历史强加离散可测空间。控制器读取同一真实观察历史，不读取潜状态、后验、响应数组或模型 oracle。

\nolinkurl{MeasurablePredictiveHistoryControl} 先在登记初始信念下采样一次潜状态，随后每步只按原 P/O 保留并推进它。实际后验解体给一步一致性，完整前缀归纳给末态潜变量的同一后验分布。再把当前后验变成上述完整响应数组、同时原样保留历史，得到同一控制器的有限与可数完整动作---观察历史路径律。潜状态标准 Borel 且非空，观察可数生成时已能证明源与 filter 的路径相容；响应过程再要求观察标准 Borel、动作可数且单点可测。

完整历史在响应旁边保留很重要。一个实际 Bool 源的前两步动作固定，后续动作读取第一次观察。两条时间二历史各有实际概率 \(1/4\)，长度和最后动作---观察事件相同，下一动作律却不同。这个反例没有挑选成功运行，也没有证明响应本身足以实现所有有记忆控制器。共同 Gaussian 源则给每个上述整数动作、完整实观察历史控制器的非原子消费者。进一步比较不同初始信念时，仍须固定同一个控制器及其同一初始记忆；不同控制器历史不能静默合并。

\hypertarget{ux4eceux771fux5b9eux5386ux53f2ux7a0bux5e8fux5bfcux51faux81eaux9002ux5e94ux9884ux6d4bux7684ux6700ux5c0fux6027}{%
\section{39.12 从真实历史程序导出自适应预测的最小性}\label{ux4eceux771fux5b9eux5386ux53f2ux7a0bux5e8fux5bfcux51faux81eaux9002ux5e94ux9884ux6d4bux7684ux6700ux5c0fux6027}}

前面的有限词响应还需要回答两个实际问题：未来动作会读取观察后再决定；已经发生的过去，也可能由完整历史程序产生。现在固定同一 P/O 源，过去程序只读取实际动作---观察历史，潜状态从登记初始信念采样一次后持续保留。前缀归纳推出：整个实际过去与末态潜变量的联合律，由实际过去律及同一个末态后验分解。未来程序从这个保留的真实潜变量执行；没有重新创造一个与过去独立的物理状态。

本节实际过去／保留潜态联合律只需观察空间可数生成。后文由响应数组构造实际可测的 fresh／记忆预测核，还要求观察标准 Borel、动作可数且单点可测。若把每个新未来程序的控制记忆明确初始化为空，其完整可数动作---观察历史律有一个这样的预测核，以当前完整响应数组为输入。于是同一响应、同一未来程序和同一初始控制记忆给出相同完整未来律。若未来程序继承的记忆是过去的可测函数 \(m(h)\)，相应预测器保留 \((C(h),m(h))\)；它没有把不同记忆中的行动规则强行等同。

最小性也需要实际实验的包含关系。对每个有限动作词 \(u\)，构造真实历史程序：按当前历史长度执行 \(u\) 的相应允许动作，词结束后使用登记的允许默认动作。在长度 \(|u|\) 的实际历史上，明确的有序观察投影恰好给原 \(M_u\)。重复动作的观察位置仍然不同，缺位默认读出不会改变任何实际词观察。由此，能够预测所有这些真实 fresh adaptive 程序的统计量，必然能预测每个原有限词；原 canonical 最小性才适用于它，而不是另外假定一个词预测证书。

更准确地，设 \(\mu_n\) 为任意完整历史过去程序的实际前缀律。若统计量 \(S\) 对每个 fresh adaptive 程序都具有可测 Markov 预测核，分别在 \(\mu_n\) 下几乎处处给出其真实未来律，则 \(C\) 在 \(\mu_n\) 下几乎处处可测因子化经过 \(S\)，并有模零信息包含 \(\sigma(C)\preceq\sigma(S)\)。共同零集只对原可数动作词族取交；没有对不可数的全部程序族取交。这里保留原有限／可数规范协议约定及观察空间条件，不扩大为没有共同可测结构的不可数版本。

\nolinkurl{MeasurablePredictiveAdaptiveCanonical}、\nolinkurl{MeasurablePredictiveWordController} 与 \nolinkurl{MeasurablePredictiveHistoryCanonical} 分别承担实际可测未来核、物理词投影和任意真实过去的来源。最后一个模块使过去也不再受可见 belief 策略限制；控制程序始终只读取其允许的观察历史。真实过去---未来联合律和这些充分性／最小性消费者使用同一源。

\hypertarget{ux6574ux4e2aux81eaux7136ux8fc7ux53bbux4e0aux7684ux6761ux4ef6ux566aux58f0}{%
\section{39.13 整个自然过去上的条件噪声}\label{ux6574ux4e2aux81eaux7136ux8fc7ux53bbux4e0aux7684ux6761ux4ef6ux566aux58f0}}

令 \(B\) 为第 39.10 节实际构造的连续过程，\(\mathcal F_s=\sigma(B_r:r\le s)\) 为其自然过滤。BrownianMartingaleSource 从同一源的独立增量证明，未来增量与整个自然过去独立，而不只是与单个当前坐标独立。对 \(0\le s\le t\)，实际得到 \[
E[B_t-B_s\mid\mathcal F_s]=0,\qquad
E[(B_t-B_s)^2\mid\mathcal F_s]=t-s.
\] 相应可积性由 Gaussian 矩分别推出。于是 \(B_t\) 和 \(B_t^2-t\) 是该自然过滤下的真实 martingale；没有把它们作为来源字段输入。

对每个固定实数 \(u\)，同源 Gaussian MGF 又给 \[
Z_t^u=\exp(uB_t-u^2t/2),\qquad E[Z_t^u]=1,
\qquad E[Z_t^u\mid\mathcal F_s]=Z_s^u.
\] 有限时域测度 \(Q_T^u=Z_T^uP\) 因而是概率律，对每个固定 \(u,T\)，其实际 Radon--Nikodym 导数在 \(P\) 下几乎处处为同一个显式 \(Z_T^u\)；对 \(A\in\mathcal F_s\)、\(s\le t\) 有 \(Q_t^u(A)=Q_s^u(A)\)。这是有限时域的规范化和过去事件一致性，不等于证明换测度后的漂移平移过程是 Brownian，也不等于一般 Girsanov 或无限时域的绝对连续性。

\hypertarget{ux540cux4e00ux4e2aux65f6ux95f4ux8f93ux5165ux8fd4ux56deux540cux4e00ux4e2aux968fux673aux79efux5206}{%
\section{39.14 同一个时间输入返回同一个随机积分}\label{ux540cux4e00ux4e2aux65f6ux95f4ux8f93ux5165ux8fd4ux56deux540cux4e00ux4e2aux968fux673aux79efux5206}}

取有限非降确定分割 \(0=t_0\le\cdots\le t_N=T\)。每个实系数 \(H_i\) 只对区间起点的 \(\mathcal F_{t_i}\) 强可测，并在同一 \(P\) 下本质有界。实际输出和时间输入分别定义为 \[
I(H)=\sum_{i<N}H_i(B_{t_{i+1}}-B_{t_i}),\qquad
h(t,\omega)=\sum_{i<N}H_i(\omega)\mathbf1_{(t_i,t_{i+1}]}(t).
\] 原自然过去条件矩推出一步零均值和二阶矩；一个先前输出在其终点可测，才可以进入后一区间起点的条件期望。这样推导交叉项正交，并取得 \[
E[I(H)^2]=\sum_{i<N}(t_{i+1}-t_i)E[H_i^2]
=\|h\|_{L^2([0,T]\times\Omega,dt\,dP)}^2.
\] 重复时刻的增量和输入区间严格为零，空分割与 \(T=0\) 也保留。系数看未来时，不能继续使用这个条件均值论证。

还不能因为每张分割表有相同形式的能量公式，就认为它们已经定义同一个积分。对两张表，BrownianCommonRefinement 实际构造区间交集的共同细分，逐点证明原时间输入与原执行输出都不变。空交集统一用零系数，因为另一表的系数可能要到更晚才可见；正长度交集只沿过滤的前向包含传递权限。同源共同差的能量等式据此推出：原生时间输入 L² 类相同，原 literal 输出便几乎处处相同，其原生输出类也相同；不同输入间的全部 L² 距离同样保持。

BrownianIntegralClosure 从已证距离和范数推导 Gram 恒等式，再用原生有限线性组合构造 elementary 输入张成空间上的连续线性等距映射，并连续延拓到这个空间在原生 L² 中的拓扑闭包。每个原有限输入仍返回原 \(I(H)\) 的同一输出类；原有限输入在时间---源 L² 中的收敛，推动它们原来的输出类收敛到该实际积分。单位系数在 \((0,T]\) 上的积分能量为 \(T\)，非零执行噪声没有在延拓中被删除。

这一步的积分域明确只是 elementary 张成空间的闭包。下一节补足它与完整 predictable L² 类的关系；该关系需要独立证明，不能用``闭包''这个定义替代。随时间变化的积分过程及其 martingale／二次变差、原一般 Itô／Girsanov／受控 SDE 消费者仍各有实际义务。

\hypertarget{ux79efux5206ux57dfux786eux5b9eux662fux539fux751f-predictable-luxb2}{%
\section{39.15 积分域确实是原生 predictable L²}\label{ux79efux5206ux57dfux786eux5b9eux662fux539fux751f-predictable-luxb2}}

上述闭包与可预测输入的关系现在有了独立证明。先采用原生自然过滤的 predictable σ 代数，由零时截面与 \((s,\infty)\times A\)、\(A\in\mathcal F_s\) 的矩形生成，再把它拉回本书使用的实时间坐标。完整 predictable L² 空间由这一 σ 代数下具有几乎处处强可测版本的原生 L² 类定义，未使用 elementary 闭包作为定义。

每个实际区间输入都可预测，故 elementary 闭包包含于该完整空间。反方向也从真实输入推导：在同一 \([0,T]\) 测度下，生成矩形由已有有限区间记录实现；超出时域的矩形及零时截面是零类。生成族的 π 系统性质、σ 代数归纳、同一乘积测度下指标函数的 L² 收敛及原生 L² 强可测归纳，据此把所有完整 predictable L² 类送入原闭包。两个空间因此相等，\(T=0\) 也包含在证明中。

\nolinkurl{BrownianPredictableInputs} 与 \nolinkurl{BrownianPredictableDensity} 将这条双向包含接到原积分，给完整 predictable L² 上的实际连续线性等距积分。每个原有限输入仍返回同一 literal 输出，单位系数能量仍为 \(T\)。这是终点积分的来源；随时间的自然过滤 martingale、连续路径版本、二次变差及原一般 Itô／Girsanov 消费者仍须分别推导。

\hypertarget{ux6709ux9650ux5b9eux9645ux51faux751fux4e0eux8d39ux7528ux6570ux636eux7684ux76eeux7684ux8bc1ux4e66}{%
\section{39.16 有限实际出生与费用数据的目的证书}\label{ux6709ux9650ux5b9eux9645ux51faux751fux4e0eux8d39ux7528ux6570ux636eux7684ux76eeux7684ux8bc1ux4e66}}

第 39.5 节的精确均值访问并非有限数据。现在固定完整预算：先执行 \(N\) 次 write-false，再执行 \(N\) 次 write-true；费用读出按预先登记的有理宽度 \(\tau>0\) 变成整数箱，价格 \(\lambda\ge0\) 也登记为有理数。允许动作、原未知实故障率、程序目标及实费用保持，训练观察只含出生位与费用箱。完整事务序列保留同一亲本，失败、费用和全部 \(2N\) 坐标都留下。源模式与标签的同一产品律导出实际时序样本的独立性及每行观测律，没有重置到成功状态，也没有重新抽样。

若某次请求位 \(b\) 产生实际子代，则 \(b\) 就是源目标，并且本次实际费用为 \(c\)。整数箱 \(k\) 满足 \[
k\tau\le c<(k+1)\tau.
\] 可执行选择器只读取完整 Bool／Nat 批次及登记的有理参数：若 \(\lambda(k+1)\tau<1\)，返回对应修复控制器；若 \(1<\lambda k\tau\)，在明确 \(\epsilon<1\) 及实际坏态的条件下给出一步 idle 优于修复的负判定；其余返回 UNKNOWN。源参数、真实均值及校准好事件不进入选择器。区间与目标证书在这个精确读出模型中没有错误；这不表示物理执行没有故障。

在正输出下，学得位等于原受保护目标，费用证书供应原 \(\lambda c\le1\) 的充分条件。新消费者从完成全部探针后的实际末父对象及其真实物理状态出发，执行同一返回控制器。一个真实无限记录律同时保留原全部有限前缀、失败及费用；其可积无限收益等于该返回策略自身的值，并等于原全部因果策略的最优值。没有另选一个存在的策略，也没有回到首次成功时的中间状态。训练已发生的费用仍在原完整记录中；最优结论评价登记的未来部署，不声称训练加未来的总净收益最优或训练免费。

若有信息的箱仍不能裁决，则 \(|\lambda c-1|\le\lambda\tau\)。在真实箱已与阈值分离的条件下，对任何仅用于分析、满足 \(0\le u<1-\epsilon\) 的半径，原同源标量 Hoeffding 校准给 \[
P(\mathrm{UNKNOWN})\le\min\{1,4e^{-2Nu^2}\}.
\] 这不是把未知修复概率的下界交给选择器。固定 \(N=64,\lambda=1/2,\tau=1/4,c=1,\epsilon=1/4\) 的源条件下，取分析半径 \(u=1/4\)，界为 \(4e^{-8}<1\)。这是固定预算模型的形式结论，执行的选择器 fixtures 也是确定性核验；本包没有生成新的现实采样批次。零预算总返回 UNKNOWN；总故障 \(\epsilon=1\) 时几乎处处返回 UNKNOWN，未从零质量的假想成功元组推出正修复概率。

必要负界也保留。故障率 \(1/2\) 下的共同全故障样本，具有正质量 \(2^{-2N}\)；价格一、目标 true、费用零与二的源在该样本上无法区分，但同一坏态的实际一步排序相反。因此确定性 data-only 且零错误的程序不能对这两个源都总是作出严格裁决；这里未声称随机判定器版本。费用宽度二还会使费用 \(1/2\) 和 \(3/2\) 的全部粗读出相同，这是该固定粗读出的相对障碍。

负输出的时域不能扩大。取 \(\epsilon=1/4,\lambda=d=1/2,c=3\)，坏态修复的一步差值为 \(-3/8\)，但同一编译策略自身的无限值为 \(3/32\)，idle 自身的无限值为零。这里使用策略评价，未调用不适用的 \(\lambda c\le1\) 最优性条件。有限资料的部分目的证书、完整源参数识别和任意时域的全部策略排序是不同责任。源码为 \nolinkurl{RegisterFaultFinitePurposeProbe} 与 \nolinkurl{RegisterFaultFinitePurposeInfinite}；新的资源维持通道在第 39.18 节另行登记；第 39.24 节在另登记的完整 idle／write 协议中接入实际训练账本；原旧协议与成本更全面的评价仍保留其各自范围。

\hypertarget{ux540cux4e00ux79efux5206ux7684ux771fux5b9eux8fdeux7eedux8fc7ux7a0b}{%
\section{39.17 同一积分的真实连续过程}\label{ux540cux4e00ux79efux5206ux7684ux771fux5b9eux8fdeux7eedux8fc7ux7a0b}}

固定 \(T\ge0\)，始终使用第 39.15 节同一个原生 predictable \(L^2\) 域与同一 \(dt|_{[0,T]}\otimes P\)。对输入 \(f\)，定义实际截断输入和输出类 \[
f_s=\mathbf1_{(0,\min(s,T)]}f,\qquad I_s(f)=I_T(f_s).
\] 这不是随 \(s\) 更换时间测度。原输入上的指示函数乘法保持真正 predictable 域，且是收缩。实际区间输入 \(H\mathbf1_{(a,b]}\) 在时刻 \(u\) 的输出为 \(H(B_{\max(a,\min(b,u))}-B_a)\)。当 \(u<a\) 时，条件期望通过 \(\mathcal F_a\) 的塔式法则归零，没有把未来的 \(H\) 当成 \(\mathcal F_u\) 已知量。有限分解和已证明的原生稠密性将条件投影恒等式推广到每个 \(f\)；随后才取得与这个实际截断输出在各固定时刻几乎处处相同的自然过滤鞅。

连续性另行推导。有限字面路径连续，平方条件 Jensen 与原生最大不等式给一个网格最大偏差界；排序与从下连续性将它推广为单个可数最大事件的界。原有限输入的几何精度逼近和 Borel--Cantelli 给出一个共同的全时间一致收敛事件。实际连续版本 \(M_s(f)\) 始终按时刻 \(s\) 的有限适应值取极限，不根据这个未来全局事件选择输出。因此它保持原自然过滤，具有联合与渐进可测性，且路径连续几乎处处成立。每个固定时刻，\(M_s(f)\) 与原 \(I_s(f)\) 几乎处处相同；每个原有限输入的完整字面路径则在一个共同满测事件上对所有时间返回。没有对不可数时间逐项取几乎处处交集，也没有声明一个事件同时适用于所有输入。实际单位输入的末端二阶能量仍为 \(T\)，零时域也包含在同一结论中。

直接来源为 \nolinkurl{BrownianIntegralTimeTruncation}、\nolinkurl{BrownianIntegralMartingale}、\nolinkurl{BrownianIntegralFinitePaths}、\nolinkurl{BrownianIntegralMaximalControl} 和 \nolinkurl{BrownianIntegralContinuousVersion}。一般 Itô、换测度和任意过滤上的运输仍有独立义务。

\hypertarget{ux53cdux9988ux8d44ux6e90ux6765ux6e90ux4e0eux5c01ux95edux50a8ux5907ux969cux788d}{%
\section{39.18 反馈资源来源与封闭储备障碍}\label{ux53cdux9988ux8d44ux6e90ux6765ux6e90ux4e0eux5c01ux95edux50a8ux5907ux969cux788d}}

第 39.3 节的每步供给不小于写费用，是合法但较强的条件。现在另登记一个开放环境端口：每次请求 idle 时供给 \(g\) 单位资源，即使该次物理执行被故障抢先；请求写入时供给零。这是新外部通道，没有从原封闭源内部创造资源，也没有更改允许动作、旧物理事务、出生或费用投影。若部署时储备 \(R\ge c\) 且 \(g\ge c\)，原编译修复程序的实际状态转移保持一个由健康／受损状态决定的储备垫。真正付费修复的同一步供给可以为零；修复执行后，同一个程序下一步请求 idle。因此原每个正质量有限记录有非负实际余额，原完整无限记录的前缀一致性进一步给 P6 的逐路径资源可行性几乎处处成立。

与此相对，在原单寄存器封闭源中，固定出生标签 false，令 \(c>0\)、\(0<\epsilon<1\)。从受损态开始，实际执行修复、故障 idle 的交替字有质量 \[
\bigl((1-\epsilon)\epsilon\bigr)^k>0,
\] 并实际支付 \(kc\)。每个有限储备都会被某个正质量前缀超支。因此这一个原编译策略的完整无限律不是几乎处处资源可行。它不是几乎必然破产定理，也不覆盖零费用、退化故障率或任意策略的普遍不可能性。

有限资料的正输出沿用原返回程序及完成全部原探针后的实际父对象：同一个真实无限记录既达到登记的原最优未来收益，也在新端口条件下满足 P6。这里的储备起算点是部署，过去探针的全部费用仍保留；训练加部署的整体资金来源尚需另证。来源为 \nolinkurl{RegisterFaultFeedbackResource}。

\hypertarget{ux5b9eux9645ux5e73ux65b9ux589eux91cfux4e0eux540cux8f93ux5165ux65f6ux95f4ux80fdux91cf}{%
\section{39.19 实际平方增量与同输入时间能量}\label{ux5b9eux9645ux5e73ux65b9ux589eux91cfux4e0eux540cux8f93ux5165ux65f6ux95f4ux80fdux91cf}}

对原 Brownian 源的非递减确定分割 \(0=t_0\le\cdots\le t_n=T\)，允许重复时间，记 \(Q=\sum_j(B_{t_{j+1}}-B_{t_j})^2\)。原高斯四阶矩和整个自然过去的条件独立性导出中心平方增量的正交性，因而 \[
E[(Q-T)^2]=2\sum_j(t_{j+1}-t_j)^2
\le2\,\operatorname{mesh}(t)\,T.
\] 网格从调用者的实际有限时钟计算。每个实际网格趋零的分割族均得到均方及原生 \(L^2\) 收敛；等分时钟族自身也证明了网格趋零。原过去可测且一致有界的权重可随分割变化，相同证明把加权平方和接到同一个加权时间和；此处没有直接将后者称为尚未证明的时间积分。来源为 \nolinkurl{BrownianQuadraticVariationSteps} 和 \nolinkurl{BrownianQuadraticVariation}。

对任意原 full predictable \(L^2\) 输入 \(f\)，再定义真正的路径时间能量 \[
A_s(f)=\int_0^{\min(s,T)}|f(t,\omega)|^2\,dt,
\] 其中使用同一个输入类的原生 predictable 代表。原生可预测权限与 Fubini 导出 \(A_s\) 对原 \(\mathcal F_s\) 强适应且可积；一个共同的可积截面事件给其所有时间上的单调与连续性。平方输出的自然尺度是 \(L^1\)： \[
\|u^2-v^2\|_1\le\|u-v\|_2(\|u\|_2+\|v\|_2).
\] 不为一般 \(L^2\) 输入偷偷增加四阶矩条件。

关键的条件二阶能量也由实际输入推导。对 \(A\in\mathcal F_s\)，只在未来窗口 \((s,\min(t,T)]\) 将 \(f\) 乘 \(\mathbf1_A\)；原区间输出和原生稠密性证明其实际积分恰为 \(\mathbf1_A(M_t-M_s)\)。等距性与同一个时间积分的 Fubini 恒等式遂给出所有此类事件上的二阶能量关系。实际鞅的交叉项条件均值为零，原生条件期望唯一性最终证明 \[
M_s(f)^2-A_s(f)
\] 是原自然过滤上的真实连续鞅，零时刻为零几乎处处。\(M,A\) 和补偿鞅均不是调用者字段。对应 \nolinkurl{BrownianIntegralTimeEnergy}、BrownianIntegralSquareL1、\nolinkurl{BrownianIntegralPastLocalization}、\nolinkurl{BrownianIntegralCompensator}。此结论本身仍不等于已经证明任意 \(f\) 的分割二次变差、一般 C² Itô、Girsanov 或受控 SDE。

\hypertarget{ux4eceux771fux5b9eux4ebaux53e3ux8c03ux5ea6ux5bfcux51faux7236ux8eabux4efd}{%
\section{39.20 从真实人口调度导出父身份}\label{ux4eceux771fux5b9eux4ebaux53e3ux8c03ux5ea6ux5bfcux51faux7236ux8eabux4efd}}

时戳身份现在扩展到另一个明确登记的串行人口模型。有限 founder 表的不同槽位代表不同所有者，即使其程序、物理值或出生标签相同。每次请求先按地址读取实际亲本，调用原事务并在原槽位更新它，再追加一个新的可选槽位；实际出生占据新槽，失败或无效地址留下空槽并消耗一个时钟。没有借用另一个亲本、重置或按成功删记录。

新生 ID 等于该次执行前的表长，因此跨亲本、跨代且标签相同的成功事务仍给不同 ID。实际出生还确定当时真正存在的唯一亲本，供应原 P12 父桥及 P8 的修复／后代区别。子代复制的实际程序也有来源；在所选亲本程序等于登记机制目标的条件下，该子代接原 P5。混合程序人口不因此自动获得同一个 P5 保障。调度地址与可靠日志是新增且明确的权限，不是从未标记物理状态反推出祖先。对应 \nolinkurl{RegisterPopulationHistorySource}；统计发现的好事件算法仍需另行来源。

\hypertarget{ux539f-price-ux5b9aux7406ux6d88ux8d39ux5b9eux9645ux5b50ux4ee3}{%
\section{39.21 原 Price 定理消费实际子代}\label{ux539f-price-ux5b9aux7406ux6d88ux8d39ux5b9eux9645ux5b50ux4ee3}}

冻结执行前的地址表，每个旧地址只尝试一次，新生个体不加入本代调度。实际执行证明所读亲本仍是原对象，末表的新槽正是这次产生的子代。亲本频率只计非空旧槽，繁殖率由实际是否出生计算。仅在实际出生数 \(B>0\) 时，构造实际子代上的归一传递行。无子代的行可补到一个已经存在的子代，但其繁殖率为零；只有 \(bK>0\) 的项代表实际谱系边。不能沿用第 39.2 节已被反例否定的``有限永久个体每行都归一且每条正边严格增代''组合。

源计算的 \(f bK\) 更新归一化后，恰等于实际子代的均匀经验频率。因此原 p\_evo\_01 的频率、繁殖率和传递矩阵均来自同一次实际出生。对子代继承的程序性状，传递项为零，得到实际子代均值变化等于选择协方差除以平均繁殖率。两种程序并存而标签相同的实例中，亲本 true 型比例为 \(1/2\)，实际子代中为零，协方差除以平均繁殖率的项为 \(-1/2\)；同一调度的全故障实现则是零子代，不归一化。

这是逐实现的经验 Price 恒等式，没有把随机比的期望换成期望之比。旧亲本仍保留在末表中，所更新的频率明确只属于子代普查，不是整个末表人口比例。它也不证明选择必然趋于最优。来源为 \nolinkurl{RegisterPopulationTypePrice}；形成、依赖统计发现、尺度、长期资金闭环及全局组合仍继续承担各自义务。

\hypertarget{ux540cux4e00ux5b8cux6574ux79efux5206ux7684ux5206ux5272ux5e73ux65b9ux548cux4e0eux5b9eux9645ux52a0ux6743ux8868}{%
\section{39.22 同一完整积分的分割平方和与实际加权表}\label{ux540cux4e00ux5b8cux6574ux79efux5206ux7684ux5206ux5272ux5e73ux65b9ux548cux4e0eux5b9eux9645ux52a0ux6743ux8868}}

第 39.19 节的平方补偿鞅并未单独证明分割二次变差。新的来源链固定原实时间乘原概率律上的完整可预测 \(L^2\) 输入 \(f\)，使用同一连续积分过程 \(M_f\) 和同一实际时间能量 \(A_f(t)=\int_0^{t\wedge T}f_s^2\,ds\)。任意声明的确定性有限非降分割 \(\pi_m\) 从 \(0\) 到 \(T\)，只要计算出的最大时间间隔趋零，就有

\[
\sum_{[u,v]\in\pi_m}(M_f(v)-M_f(u))^2\longrightarrow A_f(T)
\quad\text{于原概率源的 }L^1\text{，从而依概率。}
\]

这里没有换成另一个存在的输入。证明从原有限张成空间的实际左采样出发，经真实时间边缘积分和对所有分割一致的平方和稳定界扩展到完整输入；只要求原 \(L^2\)，没有另加 \(L^4\)。重复端点、空时间域均保留。结论不含随机停时分割，也不声称一条共同事件同时保证所有分割族的整列几乎处处收敛。

对同一 Brownian 源上的两个输入 \(f,g\)，实际交叉增量和的极限是真实的 \(\int_0^T f_sg_s\,ds\)；相应的 \(M_fM_g-\int fg\,ds\) 是原自然过滤上的可积鞅。这里的两个输入不是两个独立噪声。对一致有界、联合可测、适应且几乎处处连续的实际系数过程 \(H\)，原 \(F_u\) 简单函数逼近、未来区间定位和原输入的支配收敛进一步给出

\[
\begin{aligned}
\sum_{[u,v]\in\pi_m}H_u\,(M_f(v)-M_f(u)) &\longrightarrow I(Hf)_T,\\
\sum_{[u,v]\in\pi_m}H_u\,(M_f(v)-M_f(u))^2 &\longrightarrow \int_0^T H_sf_s^2\,ds.
\end{aligned}
\]

于原源的概率。每张表逐项使用同一实际 \(M_f\)；目标由原时间积分计算，未作为待证极限输入。这提供 Itô 所需的一、二阶随机和来源，尚须 Taylor 余项与实际望远镜和的合成。任意有界 \(C^2\) 函数并不自动有有界一、二阶导数；后续消费者须逐项提供自身所需的正则性和可积性。当前源是标量规范源，不免除一般维度、过滤、受控 SDE 与换测度的原义务。

\hypertarget{ux8bfbux53d6ux6574ux4e2aux8054ux5408ux5386ux53f2ux7684ux6709ux9650ux7ef4ux566aux58f0ux4e0eux539fux751fux8f93ux5165ux8eabux4efd}{%
\section{39.23 读取整个联合历史的有限维噪声与原生输入身份}\label{ux8bfbux53d6ux6574ux4e2aux8054ux5408ux5386ux53f2ux7684ux6709ux9650ux7ef4ux566aux58f0ux4e0eux539fux751fux8f93ux5165ux8eabux4efd}}

另登记有限个原标量 Brownian 源的原生乘积，向量使用真正的欧氏范数。未来向量增量与整个联合自然历史的独立性由坐标成对律和乘积重排推导；不能仅凭各坐标与自身过去独立就授予联合历史权限。因而一个实际系数可以读取其他坐标的过去，例如读取第二坐标后构造余弦／正弦权重。在每个确定 START 点，各坐标系数对整个联合自然历史强可测并在同一源上本质有界时，实际点积增量和具有由同源条件协方差推导的零均值与欧氏等距。

实际输入位于未归一的实时间乘同一向量概率源的欧氏值 \(L^2\)。不同确定性分割的公共细化保持原始 START 权限；塌缩公共时间单元的系数取零，不能把未来可测函数挪到更早的起点。由此推导原生输入、输出之间全部距离相等：同一输入类返回同一张有限程序的实际输出类和几乎处处相同的原输出。原生有限支撑线性组合及连续线性延拓再构造该张成空间闭包上的积分，并证明它返回每个原有限输入的原输出。

这个定义域明确是实际基本输入张成空间的闭包。它与全部联合可预测欧氏值 \(L^2\) 的相等仍需生成集／密度证明，不能作为字段补进来。零维、零时域和重复时间都包含在准确范围内；一般向量 Itô、控制及 Girsanov 不由有限和等距自动成立。

\hypertarget{ux5b8cux6574ux8badux7ec3ux8d39ux7528ux4e0eux5b9eux9645ux8fd4ux56deux7b56ux7565ux7684ux4e00ux6761ux8d26ux672c}{%
\section{39.24 完整训练费用与实际返回策略的一条账本}\label{ux5b8cux6574ux8badux7ec3ux8d39ux7528ux4e0eux5b9eux9645ux8fd4ux56deux7b56ux7565ux7684ux4e00ux6761ux8d26ux672c}}

第 39.18 节的外部保护 idle 端口现接入一个另登记的完整训练协议：每次原允许写动作前先请求一次 idle，最后再请求一次 idle。原 \(2N\) 次写动作因此对应完整 \(4N+1\) 次 primitive 抽取。全部故障、出生、费用和 idle 记录保留；没有重置、程序安装、筛选成功或重抽。所选写坐标的边缘律等于原训练源，完整轨迹和实际末亲本则确实不同，不能沿用旧轨迹的同一性。

在声明的端口输入 \(g\ge c\)、实际初始储备 \(R_0\ge0\) 下，每个实际训练前缀账面非负，末储备至少为 \(c\)。同一计数选择器只读实际允许观测；在原程序匹配、登记价格、正读出宽度与 \(0<d<1\) 折扣条件下，其正修复裁决返回的原策略从新的实际末亲本状态出发，接入原未来最优性和资源 P6。将完整训练记录字面前接到该实际未来记录律，得到同一初始储备的一条完整账本几乎处处可行。最优性仍从实际部署起点计，不冒称``训练加未来总净收益''最优；端口是另登记的外部供给，不是内部产生能源。

\hypertarget{ux6df7ux5408ux7a0bux5e8fux7684ux81eaux8eabux7ef4ux6301ux5b9eux9645ux5b50ux4ee3ux4e0eux5b8cux6574ux4eb2ux672cux53d1ux73b0}{%
\section{39.25 混合程序的自身维持、实际子代与完整亲本发现}\label{ux6df7ux5408ux7a0bux5e8fux7684ux81eaux8eabux7ef4ux6301ux5b9eux9645ux5b50ux4ee3ux4e0eux5b8cux6574ux4eb2ux672cux53d1ux73b0}}

原固定全局目标机制下，混合程序群体不能自动共享同一个 P5。新的模型在每次真实亲本查找后，以该亲本保护程序计算本次目标；故障改为写入其自身程序的反值。这是明确不同的环境语义，旧源和旧限制保留，未增加生物动作。实际出生的物理种子与继承程序相合，因而各类型子对象分别进入其自身的原 P5，无须再外给``亲本程序等于全局目标''。

冻结旧地址队列、每个原亲本只访问一次时，新模型的完整实际子对象选项与原子代普查投影相等，故能复用原 Price 的实际频率、繁殖数和传递项。故障后保留的亲本和整个末表可以不同；子代频率不是整个末表频率。实际新生槽位、唯一出生 ID 和亲本由同一分配器推导。一个三亲本源产生两类实际子代、一个保留故障及随后真实付费修复；其四步 primitive 字的原概率为 \(27/4096\)。新生子对象零局部 cushion 与声明的 idle 端口还给出其实际第一步资金桥，未给出全群体银行或全局调度最优性。

亲本发现另登记 marked journal 的完整噪声读权限。计数选择器只使用实际 ID 读出、有限 registry 和固定 fallback，隐蔽目标、噪声参数、正确标记及 Good 不进入选择器。在 \(\eta<1/2\) 时，真实独立读源以正确概率至少 \(1-\eta\) 给出严格多数事件，从源均值和 Hoeffding 推出其补集概率至多 \(\exp[-2N(1/2-\eta)^2]\)。实际 ID 类到亲本的桥把这一结论送入原 P12。在 \(\eta=1/2\)、未知任意 confuser 类中，不同实际亲本可有相等的完整读出律；任一确定性估计器在两源的错误风险之和至少为一。

更强的共同消费者覆盖全部四步群体抽取与全部 \(N\) 次读抽取的乘积律。第四动作对每个群体原始字都从真实前三步状态计算；实际 key 和亲本可随原始字变化。在群体与读噪声均为 \(1/4\) 的登记源上，联合实际亲本风险至多 \(\exp(-N/8)\) 并趋零。没有对正例、成功出生或接受结果条件化；质量 \(1/4096\) 的全故障无出生字仍在完整源中。只有在源推导的 Good 且目标确实出生时，返回的同一个实际子对象才进入自身 P5。None 表示无子代亲本，不是失败交易尝试者的身份。读取成本、全群体资金与全局目的最优仍承担各自来源义务，不能从元数据权限中免费取得。

\hypertarget{ux9644ux5f55-a.-ux7b26ux53f7ux4e0eux7c7bux578bux7d22ux5f15}{%
\chapter{附录 A. 符号与类型索引}\label{ux9644ux5f55-a.-ux7b26ux53f7ux4e0eux7c7bux578bux7d22ux5f15}}

同一字符在不同章节可能是局部变量；下表给主类型。实际引用必须以所在定义为准。

\begin{longtable}[]{@{}
  >{\raggedright\arraybackslash}p{(\columnwidth - 4\tabcolsep) * \real{0.3333}}
  >{\raggedright\arraybackslash}p{(\columnwidth - 4\tabcolsep) * \real{0.3333}}
  >{\raggedright\arraybackslash}p{(\columnwidth - 4\tabcolsep) * \real{0.3333}}@{}}
\toprule\noalign{}
\begin{minipage}[b]{\linewidth}\raggedright
符号/名称
\end{minipage} & \begin{minipage}[b]{\linewidth}\raggedright
数学类型与含义
\end{minipage} & \begin{minipage}[b]{\linewidth}\raggedright
入口
\end{minipage} \\
\midrule\noalign{}
\endhead
\bottomrule\noalign{}
\endlastfoot
\(X_t,H_t\) & 状态、截至 \(t\) 的历史 & 第 2 章 \\
\(P^u,K_i\) & 协议路径律、条件响应核 & 第 2、4 章 \\
\(f,\sigma(f)\) & 可测读出及其嵌入因子 & 第 3 章 \\
\(Y_f^+\) & 候选因子相对未来 & 第 4 章 \\
\(C^f,C\) & 概率核族值的规范预测状态 & 第 4 章 \\
FFIPS & 固定未来与协议的完整预测充分状态；使用时明确该状态域 & 第 4、16 章 \\
\(M,U\) & 内部表示、环境接口；不是默认物质/能量 & 第 6 章 \\
\(I(\cdot;\cdot\mid\cdot)\) & 条件互信息，单位 nat & 第 1、5 章 \\
\(\epsilon,\eta,\alpha\) & 模型缺陷、估计半径、置信失效概率；分别记录 & 第 16 章 \\
\(\mathcal V,\tau_{\mathcal V}\) & 持续域、退出时间 & 第 8 章 \\
\(V_A,\mathcal L\) & 平均返回时间势、生成元 & 第 9 章 \\
\(q,h,\varpi\) & QSD、正右特征函数、Doob 不变律 & 第 10 章 \\
\(\lambda,\gamma\) & 吸收指数率、条件混合率 & 第 10 章 \\
\(\Gamma,\Theta,\Psi\) & 历史运输、结构描述子、发展机制 & 第 11 章 \\
TCIC / DTRC & 冻结核交换 / 发展运输一致性 & 第 11 章 \\
\(\Omega\) & 已指定证书的因果完整性结构；不是预设几何环 & 第 12 章 \\
\(V,\mathcal M_0,b_V(\mathcal M_0)\) & 载体全集、极小充分族、极小击中族 & 第 13 章 \\
\(e(S),r(S)\) & 响应直径缺损、最优解码半径 & 第 13 章 \\
\(\mathcal L_s,\operatorname{cl}_s\) & 分辨率闭包系统、闭包 & 第 14 章 \\
\(\pi_*,J^-,J^+\) & 极小族粗像、实现区间下/上端 & 第 14 章 \\
\(\Delta_P\) & 父层超过子核心的条件预测熵 & 第 15 章 \\
\(G_N,\mathcal I\) & 设计 Gram、Fisher 信息矩阵 & 第 17 章 \\
\(J,\Pi,\beta\) & 任务评价、策略集合、折扣 & 第 18 章 \\
\(V^*,Q^*,\mathcal A^*\) & 最优值、动作值、最优动作对应 & 第 19 章 \\
GOD / GOA & 最优动作/生成机制、其长期结构家族 & 第 19、21 章 \\
TVRC & 相对于指定任务最优路径律族的目的实现证书 & 第 18 章 \\
\(\operatorname{sp}(v)\) & \(\sup v-\inf v\)，函数跨度 & 第 1、20 章 \\
\(d_\mu^\pi\) & 归一化折扣状态---动作占据分布 & 第 20 章 \\
\(\Pi,\Phi\) 双驱动 & 独立锚定的资源/信息坐标，勿与策略集混用 & 第 23 章 \\
\(B,z(B),S\) & 预算、共同低维响应变量、堆叠 Jacobian & 第 23 章 \\
\(M=D_bK\) & 平均繁殖算子、繁殖强度、传递核 & 第 25 章 \\
\(R,l,r\) & 主增长因子、长期组成、繁殖价值 & 第 25 章 \\
ETI & 本文四分量的演化个体性转变证书 & 第 26 章 \\
\(b_t(\theta,x)\) & 参数---状态后验信念 & 第 27 章 \\
\end{longtable}

\hypertarget{ux9644ux5f55-b.-ux8bc1ux660eux7ed3ux679cux6ce8ux518cux8868}{%
\chapter{附录 B. 证明结果注册表}\label{ux9644ux5f55-b.-ux8bc1ux660eux7ed3ux679cux6ce8ux518cux8868}}

第 0--31 章保留原 v3.0 的 106 个 P 编号结果，均有原有证明或明确组合。编号本身不表示原创或无条件性。现有 Core 对这套冻结源端点有独立形式化追踪；具体 Lean 假设与本文的标准数学条件仍须按逐章对应阅读。新增 M 编号结果另见附录 F，不改变下面这张 106 项稳定索引。点击编号可回到陈述。

\begin{longtable}[]{@{}lrl@{}}
\toprule\noalign{}
编号 & 章 & 结果/所在节 \\
\midrule\noalign{}
\endhead
\bottomrule\noalign{}
\endlastfoot
\protect\hyperlink{p-met-01}{P-MET-01} & 1 & 核与读出的数据处理 \\
\protect\hyperlink{p-met-02}{P-MET-02} & 1 & 跨度估计 \\
\protect\hyperlink{p-proc-01}{P-PROC-01} & 2 & 历史状态化 \\
\protect\hyperlink{p-fac-01}{P-FAC-01} & 3 & 3.2 表示协变 \\
\protect\hyperlink{p-pred-01}{P-PRED-01} & 4 & 共同协议的规范最小性 \\
\protect\hyperlink{p-pred-02}{P-PRED-02} & 4 & 目标推送与协议细化 \\
\protect\hyperlink{p-pred-03}{P-PRED-03} & 4 & 有限字母过程的递归更新 \\
\protect\hyperlink{p-info-01}{P-INFO-01} & 5 & 5.1 信息保留恒等式 \\
\protect\hyperlink{p-info-02}{P-INFO-02} & 5 & 5.2 从信息残差到预测误差 \\
\protect\hyperlink{p-info-03}{P-INFO-03} & 5 & 零失真复杂度 \\
\protect\hyperlink{p-info-04}{P-INFO-04} & 5 & Fano 型身份下界 \\
\protect\hyperlink{p-info-05}{P-INFO-05} & 5 & 5.5 遗漏与 packing 下界 \\
\protect\hyperlink{p-int-01}{P-INT-01} & 6 & 结构化充分性的充要条件 \\
\protect\hyperlink{p-int-02}{P-INT-02} & 6 & 6.2 有限结构核的最小内部---环境分解 \\
\protect\hyperlink{p-int-03}{P-INT-03} & 6 & 6.5 最小 Markov 边界的唯一性条件 \\
\protect\hyperlink{p-dyn-01}{P-DYN-01} & 7 & strong lumpability \\
\protect\hyperlink{p-dyn-02}{P-DYN-02} & 7 & 7.2 连续时间生成元判据 \\
\protect\hyperlink{p-dyn-03}{P-DYN-03} & 7 & 同策略路径误差 \\
\protect\hyperlink{p-dyn-04}{P-DYN-04} & 7 & 7.4 同一底过程的跨尺度商 \\
\protect\hyperlink{p-per-01}{P-PER-01} & 8 & 8.2 确定性持续的必要不变核 \\
\protect\hyperlink{p-per-02}{P-PER-02} & 8 & 8.4 随机持续与占据核 \\
\protect\hyperlink{p-per-03}{P-PER-03} & 8 & 8.5 有限状态的可行核 \\
\protect\hyperlink{p-per-04}{P-PER-04} & 8 & 切向必要性 \\
\protect\hyperlink{p-rec-01}{P-REC-01} & 9 & 9.1 离散恢复管道 \\
\protect\hyperlink{p-rec-02}{P-REC-02} & 9 & 9.2 连续随机恢复 \\
\protect\hyperlink{p-rec-03}{P-REC-03} & 9 & 9.3 恢复时间生成规范势 \\
\protect\hyperlink{p-rec-04}{P-REC-04} & 9 & 漂移推出恢复时间 \\
\protect\hyperlink{p-qsd-01}{P-QSD-01} & 10 & 条件稳定化的逆向结论 \\
\protect\hyperlink{p-qsd-02}{P-QSD-02} & 10 & 10.2 有限 Perron 表示 \\
\protect\hyperlink{p-qsd-03}{P-QSD-03} & 10 & 10.3 同初态的持续时间窗 \\
\protect\hyperlink{p-qsd-04}{P-QSD-04} & 10 & 10.4 可逆扩散的谱版本 \\
\protect\hyperlink{p-id-01}{P-ID-01} & 11 & 运输缺损累计 \\
\protect\hyperlink{p-id-02}{P-ID-02} & 11 & 发展管道 \\
\protect\hyperlink{p-omg-01}{P-OMG-01} & 12 & 极小毁损重建 \\
\protect\hyperlink{p-omg-02}{P-OMG-02} & 12 & 12.3 完整性裕量 \\
\protect\hyperlink{p-car-01}{P-CAR-01} & 13 & 13.1 单调访问系统 \\
\protect\hyperlink{p-car-02}{P-CAR-02} & 13 & 13.2 访问擦除与 blocker \\
\protect\hyperlink{p-car-03}{P-CAR-03} & 13 & 双 blocker \\
\protect\hyperlink{p-car-04}{P-CAR-04} & 13 & 13.4 缺损直径与解码半径 \\
\protect\hyperlink{p-res-01}{P-RES-01} & 14 & 极小性质的粗粒化 \\
\protect\hyperlink{p-res-02}{P-RES-02} & 14 & 复合律 \\
\protect\hyperlink{p-res-03}{P-RES-03} & 14 & blocker 的尺度自然性 \\
\protect\hyperlink{p-res-04}{P-RES-04} & 14 & 伴随三元组 \\
\protect\hyperlink{p-res-05}{P-RES-05} & 14 & 14.4 精确微观不确定区间 \\
\protect\hyperlink{p-res-06}{P-RES-06} & 14 & 对象相关的唯一性 \\
\protect\hyperlink{p-comp-01}{P-COMP-01} & 15 & 15.1 三种共同性 \\
\protect\hyperlink{p-comp-02}{P-COMP-02} & 15 & 15.2 干预不可约性 \\
\protect\hyperlink{p-comp-03}{P-COMP-03} & 15 & 15.3 组合裕量 \\
\protect\hyperlink{p-comp-04}{P-COMP-04} & 15 & 15.4 父层新增预测信息 \\
\protect\hyperlink{p-comp-05}{P-COMP-05} & 15 & 15.6 重叠几何的规范 Boolean 化 \\
\protect\hyperlink{p-comp-06}{P-COMP-06} & 15 & 物理载体到联盟的提升 \\
\protect\hyperlink{p-comp-07}{P-COMP-07} & 15 & 15.8 保留个体性的组合窗口 \\
\protect\hyperlink{p-stat-01}{P-STAT-01} & 16 & 有限字母 TV 同时界 \\
\protect\hyperlink{p-stat-02}{P-STAT-02} & 16 & 16.2 缺损误差不额外依赖载体候选数量 \\
\protect\hyperlink{p-stat-03}{P-STAT-03} & 16 & 固定容差族恢复 \\
\protect\hyperlink{p-stat-04}{P-STAT-04} & 16 & 精确族恢复 \\
\protect\hyperlink{p-stat-05}{P-STAT-05} & 16 & 16.5 规范预测类恢复 \\
\protect\hyperlink{p-stat-06}{P-STAT-06} & 16 & 嵌入同时误差 \\
\protect\hyperlink{p-stat-07}{P-STAT-07} & 16 & 核运输协变 \\
\protect\hyperlink{p-stat-08}{P-STAT-08} & 16 & 有限候选发现 \\
\protect\hyperlink{p-stat-09}{P-STAT-09} & 16 & 平均误差运输 \\
\protect\hyperlink{p-inv-01}{P-INV-01} & 17 & 二元检验下界 \\
\protect\hyperlink{p-inv-02}{P-INV-02} & 17 & 17.2 Fisher 的规范零方向 \\
\protect\hyperlink{p-inv-03}{P-INV-03} & 17 & 17.2 Fisher 的规范零方向 \\
\protect\hyperlink{p-inv-04}{P-INV-04} & 17 & 17.3 发展定律的设计识别 \\
\protect\hyperlink{p-inv-05}{P-INV-05} & 17 & 17.4 可预测设计的 OLS 界 \\
\protect\hyperlink{p-tel-01}{P-TEL-01} & 18 & 折扣势函数等价 \\
\protect\hyperlink{p-ctl-01}{P-CTL-01} & 19 & Bellman 存在与最优性 \\
\protect\hyperlink{p-ctl-02}{P-CTL-02} & 19 & 19.3 紧 Feller 控制 \\
\protect\hyperlink{p-ctl-03}{P-CTL-03} & 19 & 19.4 连续时间 HJB 验证 \\
\protect\hyperlink{p-quo-01}{P-QUO-01} & 20 & 20.1 精确控制商 \\
\protect\hyperlink{p-quo-02}{P-QUO-02} & 20 & 20.2 跨度敏感的近似商界 \\
\protect\hyperlink{p-quo-03}{P-QUO-03} & 20 & 20.3 动作间隔与近最优误差 \\
\protect\hyperlink{p-quo-04}{P-QUO-04} & 20 & 20.4 占据加权的精确恒等式 \\
\protect\hyperlink{p-quo-05}{P-QUO-05} & 20 & 占据加权的商策略后悔 \\
\protect\hyperlink{p-goa-01}{P-GOA-01} & 21 & 21.2 有限链的占据极限 \\
\protect\hyperlink{p-goa-02}{P-GOA-02} & 21 & 21.3 Dobrushin 混合裕量 \\
\protect\hyperlink{p-goa-03}{P-GOA-03} & 21 & 21.4 保持复归分解的多结构稳定性 \\
\protect\hyperlink{p-goa-04}{P-GOA-04} & 21 & 21.5 对称 killed 核的谱稳定性 \\
\protect\hyperlink{p-kl-01}{P-KL-01} & 22 & 22.1 生存事件的信息下界 \\
\protect\hyperlink{p-kl-02}{P-KL-02} & 22 & 22.2 精确事件 I-projection \\
\protect\hyperlink{p-kl-03}{P-KL-03} & 22 & 22.3 离散 Markov 路径的 KL \\
\protect\hyperlink{p-kl-04}{P-KL-04} & 22 & 22.4 CTMC 的显式公式 \\
\protect\hyperlink{p-kl-05}{P-KL-05} & 22 & 可观测性修正的控制能量关系 \\
\protect\hyperlink{p-ddh-01}{P-DDH-01} & 23 & 23.2 规范自由度 \\
\protect\hyperlink{p-ddh-02}{P-DDH-02} & 23 & 23.3 有限指数族与约束对偶 \\
\protect\hyperlink{p-ddh-03}{P-DDH-03} & 23 & 矩约束下的变分最小性 \\
\protect\hyperlink{p-ddh-04}{P-DDH-04} & 23 & 23.4 共同二维瓶颈 \\
\protect\hyperlink{p-ddh-05}{P-DDH-05} & 23 & 有效维数阈值 \\
\protect\hyperlink{p-ali-01}{P-ALI-01} & 24 & 全局可积判据 \\
\protect\hyperlink{p-ali-02}{P-ALI-02} & 24 & 24.3 子层更新与父层价值 \\
\protect\hyperlink{p-ali-03}{P-ALI-03} & 24 & 24.4 协调阈值 \\
\protect\hyperlink{p-evo-01}{P-EVO-01} & 25 & 选择与传递的精确分解 \\
\protect\hyperlink{p-evo-02}{P-EVO-02} & 25 & 无关标签会膨胀遗传信息 \\
\protect\hyperlink{p-evo-03}{P-EVO-03} & 25 & 平均增长和繁殖价值 \\
\protect\hyperlink{p-evo-04}{P-EVO-04} & 25 & 繁殖价值鞅 \\
\protect\hyperlink{p-brg-01}{P-BRG-01} & 26 & 无突变的固定适合度集中 \\
\protect\hyperlink{p-brg-02}{P-BRG-02} & 26 & 行为等价的选择不可分辨性 \\
\protect\hyperlink{p-ref-01}{P-REF-01} & 27 & 反身状态增广 \\
\protect\hyperlink{p-ref-02}{P-REF-02} & 27 & 信念充分性 \\
\protect\hyperlink{p-ref-03}{P-REF-03} & 27 & 免费信息的决策价值非负 \\
\protect\hyperlink{p-ref-04}{P-REF-04} & 27 & 统一目标误差的后悔界 \\
\protect\hyperlink{p-ref-05}{P-REF-05} & 27 & 可行选择扩展的单调性 \\
\protect\hyperlink{p-api-01}{P-API-01} & 28 & 精确与近似过程接口的复合 \\
\protect\hyperlink{p-alg-01}{P-ALG-01} & 28 & 有限稳定分割的正确性与最粗性 \\
\protect\hyperlink{p-core-01}{P-CORE-01} & 31 & 有限对象---实现---控制认证装配 \\
\end{longtable}

\hypertarget{ux9644ux5f55-c.-ux6765ux6e90ux56deux5f52ux6838ux5fc3ux7684ux6539ux8fdbux4e0eux6807ux51c6ux4f9dux8d56}{%
\chapter{附录 C. 来源、回归核心的改进与标准依赖}\label{ux9644ux5f55-c.-ux6765ux6e90ux56deux5f52ux6838ux5fc3ux7684ux6539ux8fdbux4e0eux6807ux51c6ux4f9dux8d56}}

\hypertarget{c.1-ux6765ux6e90ux7684ux89d2ux8272}{%
\section{C.1 来源的角色}\label{c.1-ux6765ux6e90ux7684ux89d2ux8272}}

本版的理论起点是用户提供的原始论文与早期数学文件，不把后来专门分支替换为新的哲学起点。后期材料用于寻找可回归通用核心的证明、反例和接口修正。历史对话中的研究意图用于解释任务范围，不代替数学证据；无法访问或尚未核验的历史内容不作为已证前提。

以下为直接影响本版结构的来源谱系。表中``归入''表示提取相应数学并在正文重新陈述，不表示旧版全部主张均被接受，也不声称所有历史文件或全部对话逐字完成了独立审计。

\begin{longtable}[]{@{}
  >{\raggedright\arraybackslash}p{(\columnwidth - 4\tabcolsep) * \real{0.3333}}
  >{\raggedright\arraybackslash}p{(\columnwidth - 4\tabcolsep) * \real{0.3333}}
  >{\raggedright\arraybackslash}p{(\columnwidth - 4\tabcolsep) * \real{0.3333}}@{}}
\toprule\noalign{}
\begin{minipage}[b]{\linewidth}\raggedright
来源
\end{minipage} & \begin{minipage}[b]{\linewidth}\raggedright
对本版的作用
\end{minipage} & \begin{minipage}[b]{\linewidth}\raggedright
归入章节
\end{minipage} \\
\midrule\noalign{}
\endhead
\bottomrule\noalign{}
\endlastfoot
UEOT\_v3\_top.pdf & 原始 Being/Becoming、持续身份、层级、Π/Φ、GOD/GOA 问题架构 & 0、8、12、18--27 \\
UEOT\_v1.0\_核心理论正式版.md & Ω、结构化开放、对象层级、适应与意义的展开 & 0、6、11--12、24--27 \\
UEOT\_数学基础\_统一正式版\_v1.1(1).md & 早期统一数学中的预测、信息、持续、组合、控制、反身结构 & 1--12、15、18--27 \\
UEOT\_数学基础\_完整系统正式版\_v2.1.md & I--IV 主链汇合、对象相对未来、严格化和跨度控制界 & 3--27 \\
UEOT\_数学扩展\_v2.1\_五大开放主线系统推进.md & 发展识别、连续响应、GOA 与多层目标扩展 & 10、16--17、21、24 \\
UEOT\_数学基础IV\_Dynamic-Composition与Multi-GOA\_v0.4.md & 动态组合与多长期结构区分 & 11、15、21 \\
UEOT\_数学基础IV\_Metastable-Object-GOA与Teleology-Bridge\_v0.5.md & 条件稳定、QSD 与目的桥的分离 & 10、18、21 \\
UEOT\_数学基础\nolinkurl{IV_Persistence-IProjection}与\nolinkurl{Anchored-Dual-Drives_v0.6.md} & 生存信息投影与独立锚定双驱动 & 22--23 \\
UEOT\_数学基础IV\_Lineage-IProjection与\nolinkurl{Positive-Operator-Evolution_v0.7.md} & 谱系评价、正算子与增长桥 & 25--26 \\
UEOT\_数学基础\nolinkurl{IV_Realized-Lineage_Heredity_Object-Selection_v0.8.md} & 实现谱系、遗传抽样与对象选择层级 & 25--26 \\
UEOT\_数学基础\nolinkurl{IV_Endogenous-Teleological-Bridge_v0.9.md} & 由明确选择机制诱导目的的条件 & 26 \\
UEOT\_数学基础\nolinkurl{IV_Causal-Omega_Teleonomy_Variational-Realization_v0.10.md} & 因果 Ω 与变分实现 & 12、22 \\
UEOT\_数学基础\nolinkurl{IV_Protocol-Saturated-Objecthood_v0.11.md} & 协议相对完整性与饱和边界 & 4、7、16 \\
UEOT\_数学基础IV\_Factor-Ontology与\nolinkurl{Directed-Identity_v0.12.md} & 嵌入因子与有向历史身份 & 3、11 \\
UEOT-QOD\_Interface\_v1.0\_通用对象实现-分辨率-组合接口\_严格版.tex & 通用载体、blocker、分辨率伴随与重叠组合证明 & 13--15、28 \\
UEOT\_Core\_v2.2\_QOD-G融合\_双主线收敛核心\_严格版.tex & 对象实现与控制主线接口 & 13--21、28、31 \\
前次 v2.3 核心与独立核对工作 & 条件修正、有限算例及严格表述参考；本版不以其短文替代完整主文 & 全文结构与检验 \\
\end{longtable}

来源条目按概念依赖排列，不凭文件名中的版本号推断绝对写作日期。本文 v3.0 表示一次完整架构重编与通用证明整合，不是宣布所有旧分支均已收敛。

\hypertarget{c.2-ux672cux7248ux660eux786eux5438ux6536ux548cux4feeux6b63ux7684ux8fdeux63a5}{%
\section{C.2 本版明确吸收和修正的连接}\label{c.2-ux672cux7248ux660eux786eux5438ux6536ux548cux4feeux6b63ux7684ux8fdeux63a5}}

\begin{longtable}[]{@{}
  >{\raggedright\arraybackslash}p{(\columnwidth - 4\tabcolsep) * \real{0.3333}}
  >{\raggedright\arraybackslash}p{(\columnwidth - 4\tabcolsep) * \real{0.3333}}
  >{\raggedright\arraybackslash}p{(\columnwidth - 4\tabcolsep) * \real{0.3333}}@{}}
\toprule\noalign{}
\begin{minipage}[b]{\linewidth}\raggedright
项目
\end{minipage} & \begin{minipage}[b]{\linewidth}\raggedright
本版处理
\end{minipage} & \begin{minipage}[b]{\linewidth}\raggedright
性质
\end{minipage} \\
\midrule\noalign{}
\endhead
\bottomrule\noalign{}
\endlastfoot
预测身份 & 固定对象相对未来、协议、共同条件版本 & 保留并统一前提 \\
内外部闭合 & 联合残差替代两个未联合残差；给 XOR 反例 & 逻辑修正 \\
历史身份 & 有向实例延续与无向谱系、类型相似分离 & 后期证明回归 \\
恢复与持续 & 确定不变核、随机占据、恢复势、QSD 分型 & 恢复原数学广度 \\
载体缺损 & 区分响应直径和最优解码半径，证明 \(e/2\le r\le e\) & 接口优化 \\
实现与毁损 & 访问 blocker 与物理机制失效分离 & 操作语义修正 \\
分辨率 & 闭包、Boolean 自然性、伴随三元组、精确微观区间 & 通用证明回归 \\
统计恢复 & 同时响应事件覆盖全部载体，精确结构需间隔 & 误差链统一 \\
核嵌入 & characteristic 不等于 TV 误差可控 & 度量修正 \\
商控制 & 保留既有跨度界，增加占据加权后悔接口 & 保留与加强 \\
目的实现 & 对完整最优路径律族取距离 & 非唯一最优修正 \\
Girsanov & 完整噪声 KL 与可观测状态路径 KL 分离 & 等式前提修正 \\
双驱动 & 独立锚定、共同预算、堆叠秩和有效维数 & 可检验化 \\
演化 & 繁殖强度、传递核、随机种群与平均模型分离 & 类型修正 \\
有限算法 & 最粗受控稳定分割的终止、正确性与最粗性证明 & 可执行接口补全 \\
\end{longtable}

这些修订的价值在于提高一般连接的有效性，不在于把标准数学重新标为原创。P 编号仅表示本文给出证明。

\hypertarget{c.3-ux76f4ux63a5ux8c03ux7528ux7684ux57faux7840ux6570ux5b66}{%
\section{C.3 直接调用的基础数学}\label{c.3-ux76f4ux63a5ux8c03ux7528ux7684ux57faux7840ux6570ux5b66}}

本书假定通常的测度概率、线性代数、微积分和条件期望基础。以下较重结果不在本文重新证明；其作用和适用条件不得被省略。

\begin{longtable}[]{@{}
  >{\raggedright\arraybackslash}p{(\columnwidth - 4\tabcolsep) * \real{0.3333}}
  >{\raggedright\arraybackslash}p{(\columnwidth - 4\tabcolsep) * \real{0.3333}}
  >{\raggedright\arraybackslash}p{(\columnwidth - 4\tabcolsep) * \real{0.3333}}@{}}
\toprule\noalign{}
\begin{minipage}[b]{\linewidth}\raggedright
标准结果
\end{minipage} & \begin{minipage}[b]{\linewidth}\raggedright
本文需要的条件/内容
\end{minipage} & \begin{minipage}[b]{\linewidth}\raggedright
用途
\end{minipage} \\
\midrule\noalign{}
\endhead
\bottomrule\noalign{}
\endlastfoot
正则条件概率、Doob--Dynkin & 标准 Borel 目标及相应可测性；模零结论使用共同完成空间 & 条件核与最小充分表示 \\
Ionescu--Tulcea & 可数时序的可测概率核与初始律 & 路径律一致存在 \\
KL 链式法则、Pinsker & 使用扩展非负 KL；相减表达另需有限性 & 信息到预测及路径代价 \\
Prokhorov、Portmanteau & Polish 空间上的紧概率族和弱收敛 & 持续占据的不变概率 \\
Marchaud viability / Nagumo & 闭域、指定包含或局部 Lipschitz 非爆炸 ODE & 切向条件与可行性 \\
自伴紧 resolvent 谱分解 & 第 10.4 节列明条件，正首函数与严格首间隔 & 扩散 QSD 的谱证明 \\
Perron--Frobenius & 有限本原非负矩阵；对称不可约处使用相应简单主根版本 & QSD、Doob、繁殖增长 \\
可测最大值/选择 & 紧度量动作、连续收益、弱连续核等第 19.3 节条件 & 连续状态最优策略存在 \\
Itô/Dynkin 与 Girsanov & 定义域、可积性、非爆炸、真实指数鞅、可观测性分别注明 & 连续恢复、HJB、能量界 \\
Hoeffding、McDiarmid & 独立样本、相应有界差条件 & 有限数据认证 \\
遍历定理、绝对正则耦合 & 平稳遍历用于一致性；混合率额外用于速率 & 相关数据 \\
Hodge 分解 & 紧定向无边界流形、Riemann 度量及相应光滑性 & 局部目标场的非梯度分量 \\
非负鞅收敛 & 适应、可积、非负及鞅条件 & 随机繁殖价值极限 \\
\end{longtable}

基础数学的使用不把其本身变成 UEOT 的新物理定律。完整应用前还需验证具体模型满足条件。

\hypertarget{c.4-v3.0-ux4fddux7559ux7684ux539fux59cbux7814ux7a76ux5165ux53e3}{%
\section{C.4 v3.0 保留的原始研究入口}\label{c.4-v3.0-ux4fddux7559ux7684ux539fux59cbux7814ux7a76ux5165ux53e3}}

以下条目继承自 v3.0 的标准模块来源；本轮没有重新运行其文献核验。它们是原始研究或作者文本，不是本版全部内容的外部来源，也不替代正文条件。

\begin{itemize}
\tightlist
\item
  预测等价与计算力学：Shalizi--Crutchfield 的原始论文讨论最优预测状态与相关结构，可与本版的协议相对推广对照。\href{https://arxiv.org/abs/cond-mat/9907176}{原文}
\item
  可行性：Aubin 的早期研究是闭域与微分包含可行性模块的来源之一。\href{https://pure.iiasa.ac.at/1685/1/WP-81-076.pdf}{原始工作}
\item
  QSD：Champagnat--Villemonais 研究统一指数收敛到准平稳分布及 Q-process；本版仅在明确谱或有限条件下给出所需证明。\href{https://arxiv.org/abs/1404.1349}{原文}
\item
  核两样本与均值嵌入：Gretton 等的原始研究提供 MMD 统计背景；本版不将其无条件转换成 TV。\href{https://www.jmlr.org/papers/volume13/gretton12a/gretton12a.pdf}{原文}
\item
  自适应线性估计：Abbasi-Yadkori 等给出自归一化过程工具；正文使用另行证明的保守固定时点界。\href{https://arxiv.org/abs/1102.2670}{原文}
\item
  MDP 近似同态：Ravindran--Barto 的研究是动作与状态抽象的既有理论背景。\href{https://all.cs.umass.edu/pubs/2004/ravindran_b_KBCS04.pdf}{原文}
\item
  有限熵 Girsanov：Léonard 的研究提供相对熵条件下变换的标准背景；观察映射仍受数据处理约束。\href{https://leonard.perso.math.cnrs.fr/papers/2012-Leonard-Girsanov\%20theory\%20under\%20a\%20finite\%20entropy\%20condition.pdf}{作者原文}
\end{itemize}

\hypertarget{ux9644ux5f55-d.-v3.0-ux5b8cux6574ux53efux6267ux884cux9a8cux8bc1ux4e0eux5386ux53f2ux8f93ux51fa}{%
\chapter{附录 D. v3.0 完整可执行验证与历史输出}\label{ux9644ux5f55-d.-v3.0-ux5b8cux6574ux53efux6267ux884cux9a8cux8bc1ux4e0eux5386ux53f2ux8f93ux51fa}}

\hypertarget{d.1-ux8fd0ux884cux65b9ux6cd5ux4e0eux4ee3ux7801ux8eabux4efd}{%
\section{D.1 运行方法与代码身份}\label{d.1-ux8fd0ux884cux65b9ux6cd5ux4e0eux4ee3ux7801ux8eabux4efd}}

将下方 Python 代码保存为 verify\_ueot\_core\_v3.py，在已安装 NumPy 的 Python 3 环境运行：

\begin{verbatim}
python3 verify_ueot_core_v3.py
\end{verbatim}

无需访问原始数据或其他文档。合成实验种子固定为 30905。不同平台的末位浮点数字可能变化，断言包含相应数值容差。算法的精确等式仅用于本例明确相同的有限参数，不构成任意浮点模型的证明规则。

代码 SHA-256：\nolinkurl{4d4072377cbae485f595eddc239ea6962a7bdc5cf7ad400e8d75086ce485c5f7}

\hypertarget{d.2-ux5b8cux6574ux4ee3ux7801}{%
\section{D.2 完整代码}\label{d.2-ux5b8cux6574ux4ee3ux7801}}

\begin{verbatim}
"""UEOT Core v3.0: finite-model checks, not a proof assistant.
Run: python3 verify_ueot_core_v3.py
Dependency: NumPy. No network, external datasets, or quantum modules.
"""
import itertools
import json
import math
import numpy as np

results = {}
checks = 0

def check(condition, message):
    global checks
    checks += 1
    if not bool(condition):
        raise AssertionError(message)

def tv(p, q):
    return float(np.abs(np.asarray(p) - np.asarray(q)).sum() / 2)

def minimal(family):
    return frozenset(s for s in family if not any(t < s for t in family))

def powerset(n):
    return [frozenset(i for i in range(n) if mask & (1 << i))
            for mask in range(1 << n)]

def blocker(family, n):
    return minimal([s for s in powerset(n)
                    if all(bool(s & edge) for edge in family)])

# All hypergraphs on three vertices, reduced to antichains.
universe = powerset(3)
antichains = set()
for mask in range(1 << len(universe)):
    antichains.add(minimal([universe[i] for i in range(len(universe))
                           if mask & (1 << i)]))
for family in antichains:
    check(blocker(blocker(family, 3), 3) == family, "blocker involution")
results["blocker_antichains_checked"] = len(antichains)

# Exhaustive resolution tests: all 3-to-2 surjections and both antichain spaces.
coarse_universe = powerset(2)
coarse_antichains = set()
for mask in range(1 << len(coarse_universe)):
    coarse_antichains.add(minimal([coarse_universe[i]
        for i in range(len(coarse_universe)) if mask & (1 << i)]))
surjections = [p for p in itertools.product(range(2), repeat=3) if set(p)=={0,1}]
resolution_checks_before = checks
for mapping in surjections:
    def image(s):
        return frozenset(mapping[i] for i in s)
    def preimage(s):
        return frozenset(i for i in range(3) if mapping[i] in s)
    def push(family):
        return minimal([image(s) for s in family])
    for family in antichains:
        check(blocker(push(family), 2) == push(blocker(family, 3)),
              "blocker quotient naturality")
        up_fine = frozenset(s for s in universe if any(e <= s for e in family))
        restriction = frozenset(s for s in coarse_universe
                                if preimage(s) in up_fine)
        for coarse in coarse_antichains:
            up_coarse = frozenset(s for s in coarse_universe
                                 if any(e <= s for e in coarse))
            ext_minus = frozenset(s for s in universe
                if any(preimage(b) <= s for b in up_coarse))
            ext_plus = frozenset(s for s in universe if image(s) in up_coarse)
            check((ext_minus <= up_fine) == (up_coarse <= restriction),
                  "left resolution adjunction")
            check((restriction <= up_coarse) == (up_fine <= ext_plus),
                  "right resolution adjunction")
            check((push(family) == coarse) ==
                  (ext_minus <= up_fine and up_fine <= ext_plus),
                  "exact realization interval")
    for coarse in coarse_antichains:
        up_coarse = frozenset(s for s in coarse_universe
                             if any(e <= s for e in coarse))
        ext_minus = frozenset(s for s in universe
                             if any(preimage(b) <= s for b in up_coarse))
        ext_plus = frozenset(s for s in universe if image(s) in up_coarse)
        active = frozenset().union(*coarse)
        singleton_active = all(len(preimage({w})) == 1 for w in active)
        check((ext_minus == ext_plus) == singleton_active,
              "active-fiber uniqueness")
results["resolution"] = dict(surjections=len(surjections),
    fine_antichains=len(antichains),coarse_antichains=len(coarse_antichains),
    assertions=checks-resolution_checks_before)

# Redundant representation and its response deficiency.
states = list(itertools.product(range(2), repeat=2))
p_one = np.array([[0.15, 0.85], [0.90, 0.85]])
readout = [(s, s, z) for s, z in states]
response = np.array([p_one[s] for s, z in states])
def deficiency(support, table):
    pairs = [(i, j) for i in range(4) for j in range(4)
             if all(readout[i][k] == readout[j][k] for k in support)]
    return max(float(np.max(np.abs(table[i] - table[j]))) for i, j in pairs)
exact_family = minimal([s for s in universe if deficiency(s, response) == 0])
check(exact_family == frozenset([frozenset([0]), frozenset([1])]),
      "minimal redundant carriers")
check(blocker(exact_family, 3) == frozenset([frozenset([0, 1])]),
      "redundancy blocker")
# Two Bernoulli rows: minimax decoding radius is half the diameter.
a, b = 0.2, 0.8
check(abs(max(abs(a - 0.5), abs(b - 0.5)) - abs(a-b)/2) < 1e-14,
      "radius versus diameter")

# Fresh independent repeated-sampling experiment.
rng = np.random.default_rng(30905)
N, L, alpha = 20000, 8, 0.05
emp = rng.binomial(N, response) / N
eta = math.sqrt(math.log(2 * L / alpha) / (2 * N))
observed = float(np.max(np.abs(emp-response)))
check(observed <= eta, "this run's simultaneous event")
for s in universe:
    check(abs(deficiency(s, emp)-deficiency(s, response)) <= 2*eta + 1e-14,
          "all carrier deficiencies")
tau = 0.2
emp_family = minimal([s for s in universe if deficiency(s, emp) <= tau])
check(emp_family == exact_family, "gap-certified carrier recovery")
results["sampling"] = dict(seed=30905, samples=N*L, eta=eta,
                           observed_max_error=observed,
                           empirical_response=emp.tolist(),
                           threshold=tau, nonzero_gap=0.75)

# Bellman solvers and exact lifting.
beta = 0.9
macro_P = np.stack([1-p_one, p_one], axis=-1)  # state, action, next
reward = np.array([[s-0.12*a for a in range(2)] for s in range(2)])
def solve(P, r):
    v = np.zeros(len(r))
    for _ in range(10000):
        q = r + beta*np.einsum("san,n->sa", P, v)
        new = q.max(axis=1)
        if np.max(np.abs(new-v)) < 1e-13:
            v = new
            break
        v = new
    else:
        raise AssertionError("value iteration did not converge")
    q = r + beta*np.einsum("san,n->sa", P, v)
    policy = q.argmax(axis=1)
    value = np.linalg.solve(np.eye(len(r))-beta*P[np.arange(len(r)), policy],
                            r[np.arange(len(r)), policy])
    return value, policy

true_v, true_policy = solve(macro_P, reward)
micro_P = np.array([[[((p_one[s,a] if sn else 1-p_one[s,a])/2)
                      for sn, zn in states] for a in range(2)]
                    for s, z in states])
micro_r = np.array([reward[s] for s, z in states])
micro_v, micro_policy = solve(micro_P, micro_r)
check(np.max(np.abs(micro_v-np.repeat(true_v, 2))) < 1e-11,
      "exact value lifting")
check(true_policy.tolist() == [1, 0], "true optimal policy")
# Learn macro rows at representative z=0, retaining uniform micro comparison.
hat_one = emp[[0, 2]]
hat_P = np.stack([1-hat_one, hat_one], axis=-1)
hat_v, hat_policy = solve(hat_P, reward)
epsilon_p = float(np.max(np.abs(response-hat_one[[0,0,1,1]])))
bound_D = beta*epsilon_p*float(np.ptp(hat_v))/(1-beta)
confidence_D = beta*eta*float(np.ptp(hat_v))/(1-beta)
value_error = float(np.max(np.abs(true_v-hat_v)))
check(value_error <= bound_D+1e-12, "span value bound")
closed_true = macro_P[np.arange(2), hat_policy]
lift_value = np.linalg.solve(np.eye(2)-beta*closed_true,
                             reward[np.arange(2), hat_policy])
regret = float(np.max(true_v-lift_value))
check(regret <= 2*bound_D+1e-12, "lifted policy regret bound")
def stationary(P):
    A = P.T-np.eye(len(P))
    A[-1] = 1
    bvec = np.zeros(len(P)); bvec[-1] = 1
    return np.linalg.solve(A, bvec)
closed_hat = hat_P[np.arange(2), hat_policy]
stationary_true, stationary_hat = stationary(closed_true), stationary(closed_hat)
mix = tv(closed_true[0], closed_true[1])
stat_error = tv(stationary_true, stationary_hat)
check(stat_error <= epsilon_p/(1-mix)+1e-12, "stationary perturbation bound")
hat_q = reward+beta*np.einsum("san,n->sa",hat_P,hat_v)
action_gaps = np.abs(hat_q[:,0]-hat_q[:,1])
results["control"] = dict(true_value=true_v.tolist(), estimated_value=hat_v.tolist(),
    true_policy=true_policy.tolist(), estimated_policy=hat_policy.tolist(),
    epsilon_p_actual=epsilon_p, value_error=value_error, span_bound_actual=bound_D,
    span_bound_confidence=confidence_D, policy_regret=regret,
    estimated_action_gaps=action_gaps.tolist(),
    certified_actions_at_confidence_radius=(action_gaps>2*confidence_D).tolist(),
    stationary_true=stationary_true.tolist(),stationary_estimated=stationary_hat.tolist(),
    stationary_error=stat_error,mixing_coefficient=mix,
    stationary_confidence_bound=eta/(1-mix))

# Coarsest finite controlled partition.
partition = [[0,1,2,3]]
def refine(blocks):
    refined = []
    for block in blocks:
        groups = {}
        for x in block:
            sig = tuple(micro_r[x]) + tuple(
                float(micro_P[x,a,B].sum()) for a in range(2) for B in blocks)
            groups.setdefault(sig, []).append(x)
        refined.extend(groups.values())
    return refined
for _ in range(5):
    refined = refine(partition)
    if refined == partition:
        break
    partition = refined
check(partition == [[0,1],[2,3]], "coarsest stable partition")
results["stable_partition"] = partition

# Path TV coupling bound, horizon eight, one fixed initial macro state.
T = 8
path_true, path_hat = [], []
for future in itertools.product(range(2), repeat=T):
    x = 0
    pt = ph = 1.0
    for xn in future:
        pt *= closed_true[x,xn]
        ph *= closed_hat[x,xn]
        x = xn
    path_true.append(pt); path_hat.append(ph)
path_error = tv(path_true, path_hat)
check(path_error <= 1-(1-epsilon_p)**T+1e-12, "path coupling bound")
results["path"] = dict(horizon=T, observed_tv=path_error,
                       actual_defect_bound=1-(1-epsilon_p)**T)

# Joint residual XOR: individual mutual information zero, joint ln 2.
def entropy(probabilities):
    return -sum(p*math.log(p) for p in probabilities if p>0)
hy = entropy([.5,.5])
conditional_single = .5*entropy([.5,.5])+.5*entropy([.5,.5])
check(abs(hy-conditional_single) < 1e-14, "XOR marginal information")
check(abs(hy-math.log(2)) < 1e-14, "XOR joint information")
results["xor"] = dict(I_Y_X=0.0,I_Y_E=0.0,I_Y_XE=hy)

# Continuous-time QSD, h-transform, and q versus q*h.
K = np.array([[-2.,1.],[2.,-4.]])
eigs, right = np.linalg.eig(K)
idx = int(np.argmax(eigs.real))
lam = float(-eigs[idx].real)
h = right[:,idx].real
if h.sum()<0: h=-h
leigs,left = np.linalg.eig(K.T)
q = left[:,int(np.argmax(leigs.real))].real
q = q/q.sum()
h = h/float(q@h)
varpi = q*h
K_up = np.diag(1/h)@K@np.diag(h)+lam*np.eye(2)
check(np.max(np.abs(q@K+lam*q)) < 1e-12, "QSD eigenlaw")
check(np.max(np.abs(K_up.sum(axis=1))) < 1e-12, "Doob generator")
check(np.max(np.abs(varpi@K_up)) < 1e-12, "Doob stationary law")
check(tv(q,varpi)>0.01, "QSD is not Doob equilibrium")
results["qsd"] = dict(exit_rate=lam,q=q.tolist(),doob_stationary=varpi.tolist())

# Price decomposition and Perron values.
M = np.array([[1.8,.2],[.4,.8]])
p = np.array([.4,.6]); z = np.array([0.,1.])
b = M.sum(axis=1); transmission = M/b[:,None]; mean_b=float(p@b)
pnew = p@M/mean_b
selection = float(p@(b*z)-mean_b*(p@z))/mean_b
change = float((p*b)@(transmission@z-z))/mean_b
check(abs(float(pnew@z-p@z)-selection-change)<1e-14, "Price identity")
ev,rv=np.linalg.eig(M); k=int(np.argmax(ev.real)); growth=float(ev[k].real)
results["evolution"] = dict(next_frequency=pnew.tolist(),selection=selection,
                            transmission=change,growth_rate=growth)

# Potential shaping, alignment, and a strict recovery counterexample.
psi=np.array([.3,-.2])
shaped_r=reward+beta*np.einsum("san,n->sa",macro_P,psi)-psi[:,None]
shaped_v,shaped_policy=solve(macro_P,shaped_r)
check(np.max(np.abs(shaped_v-(true_v-psi)))<1e-11, "potential shaping")
check(np.array_equal(shaped_policy,true_policy), "shaping policy equivalence")
g=np.array([1.,0.]); G=np.array([-1.,1.])
threshold=-float(g@G)/(float(g@g)-float(g@G))
check(abs(threshold-.5)<1e-14, "alignment threshold")
check(float(g@((1-.6)*G+.6*g))>0, "alignment improves")
x=0.0
for _ in range(8): x=(x+1)/2
check(0<x<1, "contractive recursion need not converge to zero")
# Degenerate observation loses the orthogonal drift exactly.
sigma=np.array([[1.,0.]])
u=np.array([0.,1.])
check(float(np.linalg.norm(sigma@u))==0, "unobserved Girsanov drift")
results["girsanov"] = dict(time=1,full_noise_relative_entropy=.5,
                           observed_state_relative_entropy=0.0)

# Rank-two per-environment responses do not ensure pooled common rank two.
J1=np.diag([1.,1.,0.]); J2=np.diag([0.,1.,1.])
check(np.linalg.matrix_rank(J1)==2 and np.linalg.matrix_rank(J2)==2,
      "separate response ranks")
check(np.linalg.matrix_rank(np.vstack([J1,J2]))==3, "pooled rank counterexample")
results["summary"] = dict(checks=checks,status="all_passed",
    meaning="Finite instances and counterexamples only; not formal proof verification.",
    numpy_version=np.__version__)
print(json.dumps(results,ensure_ascii=False,indent=2))
\end{verbatim}

\hypertarget{d.3-v3.0-ux539fux6587ux4fddux7559ux7684ux5386ux53f2ux8fd0ux884cux8f93ux51fa}{%
\section{D.3 v3.0 原文保留的历史运行输出}\label{d.3-v3.0-ux539fux6587ux4fddux7559ux7684ux5386ux53f2ux8fd0ux884cux8f93ux51fa}}

以下 JSON 按上传底本文字保留。本次维护没有重新运行该旧 Python 程序；新 Lean 编译的实际证据在第 38.8 节单列。2371 项数值断言不是 2371 个一般定理，也不是本轮新证明的计数。

\begin{verbatim}
{
  "blocker_antichains_checked": 20,
  "resolution": {
    "surjections": 6,
    "fine_antichains": 20,
    "coarse_antichains": 6,
    "assertions": 2316
  },
  "sampling": {
    "seed": 30905,
    "samples": 160000,
    "eta": 0.01200866457583208,
    "observed_max_error": 0.005449999999999955,
    "empirical_response": [
      [
        0.14565,
        0.8509
      ],
      [
        0.15085,
        0.85295
      ],
      [
        0.89955,
        0.84455
      ],
      [
        0.89675,
        0.851
      ]
    ],
    "threshold": 0.2,
    "nonzero_gap": 0.75
  },
  "control": {
    "true_value": [
      7.771727748691103,
      8.944502617801051
    ],
    "estimated_value": [
      7.769815365791166,
      8.941100066407659
    ],
    "true_policy": [
      1,
      0
    ],
    "estimated_policy": [
      1,
      0
    ],
    "epsilon_p_actual": 0.005449999999999955,
    "value_error": 0.0034025513933926277,
    "span_bound_actual": 0.05745151456523852,
    "span_bound_confidence": 0.12659008583256628,
    "policy_regret": 0.0,
    "estimated_action_gaps": [
      0.623443681598804,
      0.17797859268051575
    ],
    "certified_actions_at_confidence_radius": [
      true,
      false
    ],
    "stationary_true": [
      0.10526315789473684,
      0.8947368421052632
    ],
    "stationary_estimated": [
      0.10558679770851953,
      0.8944132022914805
    ],
    "stationary_error": 0.0003236398137826946,
    "mixing_coefficient": 0.050000000000000044,
    "stationary_confidence_bound": 0.012640699553507453
  },
  "stable_partition": [
    [
      0,
      1
    ],
    [
      2,
      3
    ]
  ],
  "path": {
    "horizon": 8,
    "observed_tv": 0.0020629076678654185,
    "actual_defect_bound": 0.042777333714830834
  },
  "xor": {
    "I_Y_X": 0.0,
    "I_Y_E": 0.0,
    "I_Y_XE": 0.6931471805599453
  },
  "qsd": {
    "exit_rate": 1.2679491924311228,
    "q": [
      0.7320508075688773,
      0.2679491924311227
    ],
    "doob_stationary": [
      0.7886751345948128,
      0.2113248654051871
    ]
  },
  "evolution": {
    "next_frequency": [
      0.6315789473684211,
      0.368421052631579
    ],
    "selection": -0.1263157894736841,
    "transmission": -0.10526315789473686,
    "growth_rate": 1.874456264653803
  },
  "girsanov": {
    "time": 1,
    "full_noise_relative_entropy": 0.5,
    "observed_state_relative_entropy": 0.0
  },
  "summary": {
    "checks": 2371,
    "status": "all_passed",
    "meaning": "Finite instances and counterexamples only; not formal proof verification.",
    "numpy_version": "2.3.5"
  }
}
\end{verbatim}

\hypertarget{ux9644ux5f55-e.-ux7ef4ux62a4ux6269ux5c55ux4e0eux590dux73b0ux7ea6ux5b9a}{%
\chapter{附录 E. 维护、扩展与复现约定}\label{ux9644ux5f55-e.-ux7ef4ux62a4ux6269ux5c55ux4e0eux590dux73b0ux7ea6ux5b9a}}

\hypertarget{e.1-ux5355ux4e00ux4e3bux6587ux7684ux7ef4ux62a4ux65b9ux5f0f}{%
\section{E.1 单一主文的维护方式}\label{e.1-ux5355ux4e00ux4e3bux6587ux7684ux7ef4ux62a4ux65b9ux5f0f}}

本文件采用 UTF-8 Markdown，公式采用标准 LaTeX 数学定界符；章节和定理有固定锚点。后续工作可以按编号检索，不要求重新解析 PDF，也不要求把多个历史补丁拼成一个理论。附录 D 的代码可独立复制运行，所需模型和结果已内嵌。

新增实质结论必须同时更新：定义或条件、完整推导、失败例、依赖关系、应用接口与适当验证。修正已有编号时保留编号和变更说明；完全不同的命题使用新编号，避免引用同号不同义。

\hypertarget{e.2-ux8bbaux8bc1ux5347ux7ea7ux95e8ux69db}{%
\section{E.2 论证升级门槛}\label{e.2-ux8bbaux8bc1ux5347ux7ea7ux95e8ux69db}}

\begin{itemize}
\tightlist
\item
  从 D 升为 P：定义不自动生成外部事实，必须有明确前提与证明。
\item
  从 H 升为 E：声明实验设计、样本、对照、误差和竞争模型；模拟内部一致性不算自然证据。
\item
  从 E 升为普适结论：有限实验不完成该升级；需要限定模型类或增加理论机制。
\item
  从局部升为全局：分别补上紧性、全局可测性、支持、可达性或边界条件。
\item
  从平均升为最坏情况：增加覆盖/密度/统一性假设，不能仅换符号。
\item
  从有限时域升为长期：增加混合、谱、复归或非平稳控制条件。
\item
  从领域结果回归核心：移除专属假设后重证，并检查是否仍有非平凡内容。
\end{itemize}

\hypertarget{e.3-ux5efaux8baeux7684ux72ecux7acbux68c0ux9a8cux6b21ux5e8f}{%
\section{E.3 建议的独立检验次序}\label{e.3-ux5efaux8baeux7684ux72ecux7acbux68c0ux9a8cux6b21ux5e8f}}

先检查对象和量词是否定义完整，再查证据是否循环依赖；随后构造退化例、冗余例、隐藏共同原因、支持外例、非唯一最优与接近分岔例。对可以有限化的主张做穷举和数值反证搜索。最后才评价一般性、计算效率及实际解释力。

所有``完整性''声明都应附模型类与协议边界。本版的完整是单一通用数学架构的自足表达，不是免于反驳的终极完成宣言。

\hypertarget{e.4-ux7248ux672cux4ea4ux4ed8ux8bb0ux5f55}{%
\section{E.4 版本交付记录}\label{e.4-ux7248ux672cux4ea4ux4ed8ux8bb0ux5f55}}

v3.0 完整主文包含第 0--31 章、附录 A--E、来源谱系、正文证明、接口契约、应用流程以及完整可运行检验。没有将量子专属结构设为基础或应用前提。

v3.0 历史验证使用合成数据和有限模型，原文记录通过 2371 项断言，其中 2316 项来自分辨率结构的有限枚举；这不表示机器形式化验证了所有定理。维护版 3.1 保留原代码与输出，并另增已编译的形式化连接和见证。若数据、协议、实现或评价改变，应重新生成相关证书，不能仅沿用历史通过标记。

\hypertarget{ux9644ux5f55-f.-ux7ef4ux62a4ux7248-3.1-ux7684ux6e90ux7801ux5bf9ux5e94ux4e0eux8bc1ux636eux5c42ux6b21}{%
\chapter{附录 F. 维护版 3.1 的源码对应与证据层次}\label{ux9644ux5f55-f.-ux7ef4ux62a4ux7248-3.1-ux7684ux6e90ux7801ux5bf9ux5e94ux4e0eux8bc1ux636eux5c42ux6b21}}

\hypertarget{f.1-ux5e95ux672cux7a33ux5b9aux7f16ux53f7ux4e0eux67e5ux9605ux89c4ux5219}{%
\section{F.1 底本、稳定编号与查阅规则}\label{f.1-ux5e95ux672cux7a33ux5b9aux7f16ux53f7ux4e0eux67e5ux9605ux89c4ux5219}}

本书完整底本为 \href{source/core-v3.0/UEOT_Core_Mathematics_v3.0_Complete.md}{v3.0 原始数学主文}，183440 字节，SHA-256：

\begin{verbatim}
ed00dd102157cdafe3a79c45506e86dc574d6cba65feb2df8686e63ce2726303
\end{verbatim}

维护版保留其全部 106 个稳定 P 编号及附录 B 的完整索引，保留原始数值程序和历史输出。新增 M 编号是本书的连接命题编号，不加入冻结 P-ID 集合。

配套 \href{UEOT_MATHEMATICS_SOURCES.md}{完整来源索引} 列出全部 106 条源结果、163 个登记源端点的全限定名、实际 Lean 文件链接以及 11／89／6 的登记处分；本附录给全书阅读路线和新增连接的直接对应。登记来源是 \href{../formalization/ueot-core/docs/compression/COMPRESSION_LEDGER.yaml}{COMPRESSION\_LEDGER} 与 \href{../formalization/ueot-core/docs/CORE_COMPRESSION_THEOREM_INDEX.csv}{原 106 条章节索引}。账本中源 SHA 与上述上传底本完全一致。

旧 PID\_STATUS 的某些内容保留较早检查点，不能与后来的压缩账本计数混用。本文使用 11／89／6 说明已登记的压缩结构；不把一张状态表当作数学证明，也不以模块数量证明理论优越。

下列缩写只为避免在每段重复长前缀；拼接后就是完整 Lean 名称：

\begin{longtable}[]{@{}
  >{\raggedright\arraybackslash}p{(\columnwidth - 2\tabcolsep) * \real{0.5000}}
  >{\raggedright\arraybackslash}p{(\columnwidth - 2\tabcolsep) * \real{0.5000}}@{}}
\toprule\noalign{}
\begin{minipage}[b]{\linewidth}\raggedright
记号
\end{minipage} & \begin{minipage}[b]{\linewidth}\raggedright
完整命名空间
\end{minipage} \\
\midrule\noalign{}
\endhead
\bottomrule\noalign{}
\endlastfoot
V3 & UEOT.V3 \\
CMP & UEOT.V3.Compression \\
TC & \nolinkurl{UEOT.V3.Compression.TheoryCompletion} \\
MC & \nolinkurl{UEOT.V3.Compression.TheoryCompletion.MinCore} \\
UC & \nolinkurl{UEOT.V3.Compression.TheoryCompletion.UnifiedClosure} \\
SC & \nolinkurl{UEOT.V3.Compression.TheoryCompletion.ScientificClosure} \\
SV & \nolinkurl{UEOT.V3.Compression.TheoryCompletion.UnifiedClosure.SupportViabilityClosure} \\
POST & \nolinkurl{UEOT.V3.Compression.TheoryCompletion.UnifiedClosure.PosteriorSupportCoverageCheck} \\
STRICT & \nolinkurl{UEOT.V3.Compression.TheoryCompletion.UnifiedClosure.PosteriorSupportStrictWitness} \\
REG & \nolinkurl{UEOT.V3.Compression.TheoryCompletion.UnifiedClosure.RegisteredGoodJointCheck} \\
SELECT & \nolinkurl{UEOT.V3.Compression.TheoryCompletion.UnifiedClosure.RegisteredSelectionWitness} \\
VIABLE-JOINT & \nolinkurl{UEOT.V3.Compression.TheoryCompletion.UnifiedClosure.ViableRegisteredJointCheck} \\
\end{longtable}

例如，\nolinkurl{SV.initial_viability_iff_exists_zero_uncovered_policy} 的全名为 \nolinkurl{UEOT.V3.Compression.TheoryCompletion.UnifiedClosure.SupportViabilityClosure.initial_viability_iff_exists_zero_uncovered_policy}。文件名可以不同于命名空间：本轮模块拆分保留原有公开数学名称，不通过复制声明来制造新的证明端点。

\hypertarget{f.2-ux5168-core-ux7684ux9605ux8bfbux4e0eux5047ux8bbeux5730ux56fe}{%
\section{F.2 全 Core 的阅读与假设地图}\label{f.2-ux5168-core-ux7684ux9605ux8bfbux4e0eux5047ux8bbeux5730ux56fe}}

\begin{longtable}[]{@{}
  >{\raggedright\arraybackslash}p{(\columnwidth - 4\tabcolsep) * \real{0.3333}}
  >{\raggedright\arraybackslash}p{(\columnwidth - 4\tabcolsep) * \real{0.3333}}
  >{\raggedright\arraybackslash}p{(\columnwidth - 4\tabcolsep) * \real{0.3333}}@{}}
\toprule\noalign{}
\begin{minipage}[b]{\linewidth}\raggedright
本书位置
\end{minipage} & \begin{minipage}[b]{\linewidth}\raggedright
直接源入口
\end{minipage} & \begin{minipage}[b]{\linewidth}\raggedright
需连同定理保留的条件
\end{minipage} \\
\midrule\noalign{}
\endhead
\bottomrule\noalign{}
\endlastfoot
第 1--4、32 章：因子、预测、核降落 & \href{../formalization/ueot-core/UEOT/V3/Compression/QuotientDescent.lean}{QuotientDescent}、\href{../formalization/ueot-core/UEOT/V3/PredictionUpdate.lean}{PredictionUpdate} & 集合降落与可测降落分开；递归 Bayes 保留给定同一联合核和零分母分支 \\
第 5 章：信息与压缩界 & \href{../formalization/ueot-core/UEOT/V3/InformationPacking.lean}{InformationPacking}、\href{../formalization/ueot-core/UEOT/V3/InformationEntropy.lean}{InformationEntropy} & 有限打包、共同解码器及均匀输入熵是明确不同子结论 \\
第 6--7、11、32 章：接口与误差传输 & \href{../formalization/ueot-core/UEOT/V3/Compression/TransportCertificate.lean}{TransportCertificate}、\href{../formalization/ueot-core/UEOT/V3/PathError.lean}{PathError} & 传输链中间对象一致；实际核的共同质量与抽象递推分别证明 \\
第 8、35 章：生存固定点 & \href{../formalization/ueot-core/UEOT/V3/ViabilitySource.lean}{ViabilitySource} & 有限状态／动作、固定同一源、初始分布支撑可行核；核可以为空 \\
第 9、22 章：恢复与路径信息 & \href{../formalization/ueot-core/UEOT/V3/RecoveryDynkin.lean}{RecoveryDynkin}、\href{../formalization/ueot-core/UEOT/V3/GirsanovPathKL.lean}{GirsanovPathKL} & Dynkin 期望证书、Girsanov 密度和可积条件是输入；可观察 KL 的等号需额外条件 \\
第 10、21、25、32 章：不变与正谱结构 & \href{../formalization/ueot-core/UEOT/V3/Compression/OccupationLimitInvariance.lean}{\nolinkurl{OccupationLimitInvariance}}、\href{../formalization/ueot-core/UEOT/V3/Compression/PositiveEigenstructure.lean}{PositiveEigenstructure} & 极限／紧性与正特征数据不能由归一化公式自动产生 \\
第 12--13、34 章：对象、载体与恢复 & \href{../formalization/ueot-core/UEOT/V3/Compression/TheoryCompletion/OperationalObject.lean}{OperationalObject}、\href{../formalization/ueot-core/UEOT/V3/Compression/TheoryCompletion/AutopoiesisClosure.lean}{AutopoiesisClosure} & 访问协议、物理破坏、同一实例、程序来源及预期资源预算各自保留 \\
第 14、32、34 章：尺度与分辨率 & \href{../formalization/ueot-core/docs/compression/hierarchy/H0_HIERARCHY_INVENTORY.md}{H0 目录}、\href{../formalization/ueot-core/UEOT/V3/Compression/TheoryCompletion/ObjectScaleCalculus.lean}{ObjectScaleCalculus} & 八类层级不混型；当前 RG 流记录没有与实际尺度映射的相容方程 \\
第 15、34 章：父形成与 binding & \href{../formalization/ueot-core/UEOT/V3/Compression/TheoryCompletion/ScientificClosure/C4FBT.lean}{C4FBT} & 两类传输及父形成存在分别是输入；不能从形成单独推出 binding \\
第 16、36--38 章：统计、路径与实际覆盖 & \href{../formalization/ueot-core/UEOT/V3/Compression/TheoryCompletion/UnifiedClosure/StatisticalRiskCalibration.lean}{\nolinkurl{StatisticalRiskCalibration}}、\href{../formalization/ueot-core/UEOT/V3/Compression/TheoryCompletion/UnifiedClosure/DeploymentCoverage.lean}{DeploymentCoverage} & 注册采样法则、同源参数、共同失败事件和实际占据不能替换为训练均值 \\
第 17--18、23、26、34 章：目的来源与规范 & \href{../formalization/ueot-core/UEOT/V3/Compression/TheoryCompletion/TeleologicalGauge.lean}{TeleologicalGauge}、\href{../formalization/ueot-core/UEOT/V3/Compression/TheoryCompletion/UnifiedClosure/ProgramPurposeInformationBoundary.lean}{\nolinkurl{ProgramPurposeInformationBoundary}} & 表示存在、数值 gauge、目的排序类和来源识别是不同问题 \\
第 19、22 章：连续时间 HJB & \href{../formalization/ueot-core/UEOT/V3/DiffusionHJBVerification.lean}{DiffusionHJBVerification} & \nolinkurl{V3.DiffusionHJBVerification.Model.ItoRun} 提供实际占据、终端、Itô 期望与极限义务 \\
第 20、28、31、35 章：同源组织控制商 & \href{../formalization/ueot-core/UEOT/V3/Compression/TheoryCompletion/UnifiedClosure/OrganizationalCoreAssembly.lean}{\nolinkurl{OrganizationalCoreAssembly}}、\href{../formalization/ueot-core/UEOT/V3/Compression/TheoryCompletion/UnifiedClosure/UnifiedCoreClosureCertificate.lean}{\nolinkurl{UnifiedCoreClosureCertificate}} & 注册属性与奖励实际进入稳定分割；标签的物理正确性另需证据 \\
第 23、33 章：定量缺陷与双驱动 & \href{../formalization/ueot-core/UEOT/V3/Compression/TheoryCompletion/MinCore/RelationDefectNoGo.lean}{RelationDefectNoGo}、\href{../formalization/ueot-core/UEOT/V3/Compression/TheoryCompletion/MinCore/CalibratedDefectComparison.lean}{\nolinkurl{CalibratedDefectComparison}}、\href{../formalization/ueot-core/UEOT/V3/Compression/TheoryCompletion/MinCore/DefectGaugeCovariance.lean}{DefectGaugeCovariance}、\href{../formalization/ueot-core/UEOT/V3/Compression/TheoryCompletion/UnifiedClosure/DualDriveGaugeCompleteness.lean}{\nolinkurl{DualDriveGaugeCompleteness}} & 相同零关系不确定数值；双边校准带单位；驱动唯一性需要独立锚 \\
第 24--26、34 章：价值、亲子与谱系 & \href{../formalization/ueot-core/UEOT/V3/Compression/TheoryCompletion/LineageEvolution.lean}{LineageEvolution} & SameObject 与 Offspring 分开；Perron 桥消费同一繁殖算子的特征证书 \\
第 27、34 章：参数化信念与科学接口 & \href{../formalization/ueot-core/UEOT/V3/Compression/TheoryCompletion/ScientificClosure/C1ParametricBelief.lean}{C1ParametricBelief}、\href{../formalization/ueot-core/UEOT/V3/Compression/TheoryCompletion/ScientificClosure/C5DualDriveTest.lean}{C5DualDriveTest}、\href{../formalization/ueot-core/UEOT/V3/Compression/TheoryCompletion/ScientificClosure/C6PurposeMechanism.lean}{C6PurposeMechanism} & 联合可测性有正则条件；秩只给必要条件；目的识别需要忠实性 \\
第 29、34 章：现实证据 & \href{../formalization/ueot-core/docs/compression/theory_completion/scientific_closure/SCIENTIFIC_CLOSURE_STATUS.json}{Scientific Closure 状态} & package closure 不等于全部领域端口关闭；现实支持不由 Lean 编译给出 \\
第 35 章：安全最优与修复 & \href{../formalization/ueot-core/UEOT/V3/Compression/TheoryCompletion/UnifiedClosure/CoreSourceAllTimesAssembly.lean}{\nolinkurl{CoreSourceAllTimesAssembly}}、\href{../formalization/ueot-core/UEOT/V3/Compression/TheoryCompletion/UnifiedClosure/SafeOptimalControlFromViability.lean}{\nolinkurl{SafeOptimalControlFromViability}}、\href{../formalization/ueot-core/UEOT/V3/Compression/TheoryCompletion/UnifiedClosure/RepairProgramSafeControlBridge.lean}{\nolinkurl{RepairProgramSafeControlBridge}} & 原始源核不变、初始可行、安全动作限制、可信程序底座；安全不自动等于最优 \\
\end{longtable}

表中的路径是数学源入口，完整 106 条全限定名见配套索引。第 34 章的 P12 状态另参照 \href{../formalization/ueot-core/docs/compression/theory_completion/P12_FINAL_AUTOPOIESIS_AUDIT.md}{P12 条件边界审计}；本书采用其中 PARTIAL 的实质含义，不将工作包 CLOSED 改写成无条件自创生原理。

\hypertarget{f.3-ux65b0ux589eux7ae0ux5185ux547dux9898ux7684ux7cbeux786eux7aefux70b9}{%
\section{F.3 新增章内命题的精确端点}\label{f.3-ux65b0ux589eux7ae0ux5185ux547dux9898ux7684ux7cbeux786eux7aefux70b9}}

下列列表给出主要端点；同一命题正文中列出的辅助定理说明关键证明环节。并非每个 M 编号都等于一个单独 Lean declaration；多个条件结果的文字组合和本书手算例已在正文明确区分。

\begin{longtable}[]{@{}
  >{\raggedright\arraybackslash}p{(\columnwidth - 2\tabcolsep) * \real{0.5000}}
  >{\raggedright\arraybackslash}p{(\columnwidth - 2\tabcolsep) * \real{0.5000}}@{}}
\toprule\noalign{}
\begin{minipage}[b]{\linewidth}\raggedright
章内命题
\end{minipage} & \begin{minipage}[b]{\linewidth}\raggedright
精确命名空间与声明
\end{minipage} \\
\midrule\noalign{}
\endhead
\bottomrule\noalign{}
\endlastfoot
M-32.1 & \nolinkurl{CMP.QuotientDescent.fiberCompatible_iff_existsUnique_descend} \\
M-32.2 & \nolinkurl{CMP.TransportCertificate.weighted_chain_bound} \\
M-32.3 & \nolinkurl{CMP.OccupationLimitInvariance.invariant_of_continuous_observable_residual} \\
M-32.4 & \nolinkurl{CMP.PositiveEigenstructure.qsdDoob_bundle}；\nolinkurl{CMP.PositiveEigenstructure.stochasticRightEigen_bundle} \\
M-33.1 & \nolinkurl{MC.same_relation_arbitrary_positive_scale}；\nolinkurl{MC.no_relation_only_quantitative_recovery}；\nolinkurl{MC.relationDefectPresentation_does_not_force_triangle} \\
M-33.2 & \nolinkurl{MC.DefectCalibration.zero_defect_iff}；\nolinkurl{MC.DefectCalibration.forward_budget}；\nolinkurl{MC.DefectCalibration.reverse_budget}；\nolinkurl{MC.DefectCalibration.tendsto_zero_iff} \\
M-33.3 & \nolinkurl{MC.DefectCalibration.triangle_quasi}；\nolinkurl{MC.calibrated_cross_stage_contraction} \\
M-33.4 & \nolinkurl{MC.rescale_contraction_transformed}；\nolinkurl{MC.rescale_local_budget_transformed}；\nolinkurl{MC.rescaled_weighted_chain_bound} \\
M-33.5 & \nolinkurl{UC.same_objective_iff_gauge_related}；\nolinkurl{UC.pi_identified_of_phi_anchor}；\nolinkurl{UC.phi_identified_of_pi_anchor} \\
M-34.1 & \nolinkurl{TC.incomparableBoolContract_has_no_numerical_representation}；\nolinkurl{SC.mechanismObservation_identifies_teleologicalClass}；\nolinkurl{SC.equalMechanismObservation_does_not_determine_purpose} \\
M-34.2 & \nolinkurl{SC.approximate_fbt_realized_bound}；\nolinkurl{SC.certified_formed_structural_continuation}；\nolinkurl{SC.formationTransport_does_not_imply_bindingTransport} \\
M-34.3 & \nolinkurl{TC.no_physicalOnly_synthesizer_for_samePhysical_distinctSeeds}；\nolinkurl{UC.perfect_program_decoder_iff_injective_information} \\
M-35.1 & \nolinkurl{V3.ViabilitySource.p_per_03}；\nolinkurl{UC.single_source_terminal_object_viability_all_times} \\
M-35.2 & \nolinkurl{UC.viable_restricted_greedy_preserves_source_kernel}；\nolinkurl{UC.viable_restricted_optimal_dominates_all_causal_policies} \\
M-35.3 & \nolinkurl{UC.reconstructed_program_source_kernel_agrees_with_safe_action}；\nolinkurl{UC.reconstructed_program_preserves_source_viability}；\nolinkurl{UC.reconstructed_program_joint_and_source_safety} \\
M-36.1 & \nolinkurl{UC.risk_map_mul_fiberPosterior}；\nolinkurl{UC.risk_conditionedRow_mass}；UC.risk\_map\_extend\_eq \\
M-36.2 & \nolinkurl{UC.coverage_tv_extend_same_initial}；\nolinkurl{UC.coverage_tv_extend_same_kernel} \\
M-36.3 & \nolinkurl{UC.coverage_projected_causalLaw_error}；\nolinkurl{UC.coverage_rowLaw_expectation}；\nolinkurl{UC.coverage_same_policy_path_error} \\
M-36.4 & \nolinkurl{UC.coverage_expected_error_le}；\nolinkurl{UC.coverage_same_policy_microscopic_failure_risk} \\
M-36.5 & \nolinkurl{POST.supported_history_actions_bound_actual_coverage}；\nolinkurl{POST.posterior_supported_actions_bound_actual_coverage} \\
M-37.1 & \nolinkurl{SV.mem_traceSupport_iff_realized}；\nolinkurl{SV.mem_coarseLaw_succ_support_iff}；\nolinkurl{SV.traceSupport_eq_fiberTerminalPosterior_support}；\nolinkurl{SV.successorSupport_eq_fiberTerminalPosterior_support} \\
M-37.2 & \nolinkurl{SV.mem_deploymentRowLaw_support_iff}；\nolinkurl{SV.zero_uncovered_iff_support_certified} \\
M-37.3 & \nolinkurl{SV.viablePolicy_zero_uncovered}；\nolinkurl{SV.supported_viable_of_zero_uncovered}；\nolinkurl{SV.initial_viability_iff_exists_zero_uncovered_policy} \\
M-37.4 & \nolinkurl{SV.finiteViabilityCheck_correct}；\nolinkurl{SV.finiteRelationViabilityCheck_correct} \\
M-37.5 & \nolinkurl{STRICT.two_positive_coarse_histories}；\nolinkurl{STRICT.posterior_supports_are_singletons}；\nolinkurl{STRICT.actual_uncovered_probability_zero}；\nolinkurl{STRICT.whole_current_fiber_has_no_common_action}；\nolinkurl{STRICT.null_history_boundary} \\
M-37.6 & \nolinkurl{SV.SafePrefixDeadend.a_complete_zero_uncovered_policy_exists}；\nolinkurl{SV.SafePrefixDeadend.deadend_policy_second_step_not_covered}；\nolinkurl{SV.SafePrefixDeadend.viable_source_with_a_covered_deadend} \\
M-38.1 & \nolinkurl{UC.selectedTrainingDeployment_failure_tower}；\nolinkurl{UC.selectedTrainingDeployment_failure_integral} \\
M-38.2 & \nolinkurl{UC.selectedTrainingDeployment_data_dependent_bound} \\
M-38.3 & \nolinkurl{UC.selectedTrainingDeployment_failure_le}；\nolinkurl{REG.registered_selected_joint_failure_of_cap_on_good} \\
M-38.4 & \nolinkurl{VIABLE-JOINT.viable_registered_selected_joint_failure_on_good} \\
第 38.6 节完整例 & \nolinkurl{SELECT.registered_selection_joint_witness}；\nolinkurl{SELECT.actual_joint_failure_exact}；\nolinkurl{SELECT.allTrue_positive_probability} \\
第 38.7 节紧例与固定律 & \nolinkurl{UC.TrainingDeploymentWitness.sharp_mixture_is_attained}；\nolinkurl{UC.statistical_fixed_policy_training_deployment_witness} \\
\end{longtable}

第 36--38 章的新源按下列单一归属组织：

\begin{longtable}[]{@{}
  >{\raggedright\arraybackslash}p{(\columnwidth - 2\tabcolsep) * \real{0.5000}}
  >{\raggedright\arraybackslash}p{(\columnwidth - 2\tabcolsep) * \real{0.5000}}@{}}
\toprule\noalign{}
\begin{minipage}[b]{\linewidth}\raggedright
模块
\end{minipage} & \begin{minipage}[b]{\linewidth}\raggedright
主要内容
\end{minipage} \\
\midrule\noalign{}
\endhead
\bottomrule\noalign{}
\endlastfoot
\href{../formalization/ueot-core/UEOT/V3/Compression/TheoryCompletion/UnifiedClosure/ObservableActionCoverage.lean}{ObservableActionCoverage} & 整当前观察纤维认证、动作漏出到实际覆盖、预算代换 \\
\href{../formalization/ueot-core/UEOT/V3/Compression/TheoryCompletion/UnifiedClosure/PosteriorSupportCoverage.lean}{PosteriorSupportCoverage} & 复用真实后验、正历史支持动作的实际覆盖界 \\
\href{../formalization/ueot-core/UEOT/V3/Compression/TheoryCompletion/UnifiedClosure/ReachablePolicyExtension.lean}{ReachablePolicyExtension} & 前缀策略相容与非循环的逐步延长充分条件 \\
\href{../formalization/ueot-core/UEOT/V3/Compression/TheoryCompletion/UnifiedClosure/FiniteSupportViability.lean}{FiniteSupportViability} & 策略无关 trace support、实际概率桥、后向递推与双向策略定理 \\
\href{../formalization/ueot-core/UEOT/V3/Compression/TheoryCompletion/UnifiedClosure/FiniteSupportViabilityCheck.lean}{\nolinkurl{FiniteSupportViabilityCheck}} & 有限关系表执行及其到实际源的正确性 \\
\href{../formalization/ueot-core/UEOT/V3/Compression/TheoryCompletion/UnifiedClosure/ControlCoverageWitness.lean}{ControlCoverageWitness} & 代表不足、后验支持严格放松及可行问题中的死路 \\
\href{../formalization/ueot-core/UEOT/V3/Compression/TheoryCompletion/UnifiedClosure/RegisteredDeploymentRisk.lean}{RegisteredDeploymentRisk} & 规范 Good 与仅好事件需要满足的选后联合风险 \\
\href{../formalization/ueot-core/UEOT/V3/Compression/TheoryCompletion/UnifiedClosure/RegisteredDeploymentRiskWitness.lean}{\nolinkurl{RegisteredDeploymentRiskWitness}} & 两注册行、真实产品采样、非恒定选择、精确联合风险 \\
\href{../formalization/ueot-core/UEOT/V3/Compression/TheoryCompletion/UnifiedClosure/ViableRegisteredJoint.lean}{ViableRegisteredJoint} & 用同一支持生存策略消除另给的实际覆盖输入 \\
\end{longtable}

模块归属是查阅地址，不改变已通过端点的数学假设。书稿阶段的检查事实是下一节所列固定 assembled source；最终正式模块的构建与发布收据应以发行附带的真实记录为准。

\hypertarget{f.3.1-ux672cux6b21ux65b0ux589eux6b63ux6587ux7684ux7cbeux786eux5c42ux6b21}{%
\subsection{F.3.1 本次新增正文的精确层次}\label{f.3.1-ux672cux6b21ux65b0ux589eux6b63ux6587ux7684ux7cbeux786eux5c42ux6b21}}

下表补充新节，不改变前表的既有端点、106 个 P-ID 或历史证据。

\begin{longtable}[]{@{}
  >{\raggedright\arraybackslash}p{(\columnwidth - 4\tabcolsep) * \real{0.3333}}
  >{\raggedright\arraybackslash}p{(\columnwidth - 4\tabcolsep) * \real{0.3333}}
  >{\raggedright\arraybackslash}p{(\columnwidth - 4\tabcolsep) * \real{0.3333}}@{}}
\toprule\noalign{}
\begin{minipage}[b]{\linewidth}\raggedright
本次正文
\end{minipage} & \begin{minipage}[b]{\linewidth}\raggedright
形式对应
\end{minipage} & \begin{minipage}[b]{\linewidth}\raggedright
当前范围
\end{minipage} \\
\midrule\noalign{}
\endhead
\bottomrule\noalign{}
\endlastfoot
M-21.1--M-21.7 & 源 UC.conflictRewardControl；一般 TC.InvariantLawGOA、TC.CesaroLimitGOA、TC.KilledQSDGOA、TC.GreedyInvariantGOASet；\nolinkurl{CMP.PositiveEigenstructure.BiEigenData} & 新的参数化实例、显式数值、R 的完整 greedy 与混合不变律族为书面证明，尚无专门已编译端点 \\
M-36.6 & \nolinkurl{UC.FiniteObservationEscapeCheck.tvDist_atomEncode_eq}；\nolinkurl{tvDist_embedded_eq_encoded}；\nolinkurl{tvDist_eq_of_live_reference_masses} & 本版生产接口，含一侧有限载体与完整逃逸质量，不声称对任意两参考全局等距 \\
M-36.7 & \nolinkurl{UC.PathObservationEscapeCheck.escape_path_tv}；\nolinkurl{UC.ArbitraryObservationRiskCheck.coverage_arbitrary_observation_path_error} 及其事件风险端点 & 本版生产接口，实际源、策略、完整占据和原 Q(F) 均保持；Y 有可测单点，F 可测 \\
§36.8 正逃逸边界 & \nolinkurl{UC.FiniteObservationEscapeBoundaryCheck.infinite_ambient_positive_escape_boundary} & 本版生产接口：对固定有限源的 TV 为 5/12，逃逸质量及两个原参考间 TV 为 1/4，两个编码参考间 TV 为零 \\
M-38.5 & \nolinkurl{UC.VariableRegisteredSamples.registered_row_calibration}；\nolinkurl{product_experiment_calibration} & 本版生产接口，实际逐行 N j；一般行间无需独立 \\
M-38.6 & \nolinkurl{UC.VariableRegisteredJointCheck.registered_selected_joint_failure_of_cap_on_good} & 本版任意 Y 接口，\nolinkurl{hfit/hcap/hF/henvelope/hchecked} 在同一个规范 Good \\
M-38.7 & \nolinkurl{UC.ViableVariableRegisteredJointCheck.viable_registered_selected_joint_failure_on_good} & 本版任意 Y 接口，同一个 Good 内部证明同一 viableFamilyPolicy 的零占据，hF 也在该 Good 上 \\
\end{longtable}

原 \#354 研究原型与首版有限 Y 的记录作为历史保留。本版所述七个模块已纳入生产接口；历史检查记录与本版源码的对应仍须区分。完整签名与验收范围见 \href{LEAN_API_SCOPE.md}{LEAN\_API\_SCOPE} 和 \href{../formalization/ueot-core/docs/compression/theory_completion/unified_closure/REGISTERED_INTERFACES_CLOSURE_V1.md}{接口闭环}。下一节的历史检查记录保持原样。

\hypertarget{f.4-ux672cux8f6e-lean-ux68c0ux67e5ux7684ux53efux5b9aux4f4dux8eabux4efd}{%
\section{F.4 本轮 Lean 检查的可定位身份}\label{f.4-ux672cux8f6e-lean-ux68c0ux67e5ux7684ux53efux5b9aux4f4dux8eabux4efd}}

原始检查在项目实施证据中保留。以下用完整源码哈希标识它们；失败历史仍保留原始状态，没有以第三稿通过覆盖第一稿。

\begin{longtable}[]{@{}
  >{\raggedright\arraybackslash}p{(\columnwidth - 4\tabcolsep) * \real{0.3333}}
  >{\raggedright\arraybackslash}p{(\columnwidth - 4\tabcolsep) * \real{0.3333}}
  >{\raggedright\arraybackslash}p{(\columnwidth - 4\tabcolsep) * \real{0.3333}}@{}}
\toprule\noalign{}
\begin{minipage}[b]{\linewidth}\raggedright
检查及源
\end{minipage} & \begin{minipage}[b]{\linewidth}\raggedright
源码 SHA-256
\end{minipage} & \begin{minipage}[b]{\linewidth}\raggedright
实际结论
\end{minipage} \\
\midrule\noalign{}
\endhead
\bottomrule\noalign{}
\endlastfoot
check-03 / \nolinkurl{SupportViabilityPolicyCheck_v1.lean} & \nolinkurl{ceed43d8aa3e4c6b140c0f998105d53ca2bd770b3ae29519d9f3b80562da581f} & 退出 0；12 项标准公理 \\
check-06 / \nolinkurl{SupportViabilityFiniteCheck_v2.lean} & \nolinkurl{59964729e9b415ebe6f856f41fe0b9359ca1dda4318b184f5f500385610d08fc} & 退出 0；17 项标准公理；实际输出 true、false、true \\
check-06 / \nolinkurl{ViableRegisteredJointCheck_v1.lean} & \nolinkurl{c876dbfaefeac67c1ac01fd19fc9c6d54c818b545f4707b3d93eb4c96e3ae6d5} & 退出 0；所打印端点仅标准三公理 \\
check-08 / \nolinkurl{RegisteredSelectionWitnessCheck_v3.lean} & \nolinkurl{6bbdb1ee39e6d65960123bff1410182b2a1a0096a55d5e74d4f0f7afb0f666bf} & 退出 0；12 项标准公理；无警告 \\
\end{longtable}

这里``标准三公理''指 propext、Classical.choice、Quot.sound 的子集。check-03 和 check-06 的旧拼接源有已观察到的 unused-variable／simp 警告，不因此把它们记为 warning-free。check-04 整批和前一轮严格后验／延长混合批次曾含其他源的失败；第 38.8 节只使用其中实际通过的单项。

历史检查成功、当前正式源码构建成功、真实 Git 候选通过、远端合并以及现实实验支持是不同事实。一本数学书需要准确说明它引用哪一种事实，不必把整个研发过程写成哲学论据。本文的条件性数学结论与明确反例先由正文承担，编译证据为其具体实现提供额外核验。

\hypertarget{f.5-ux540cux6e90ux6267ux884cux4e0eux53efux6d4bux9884ux6d4bux7684ux65b0ux589eux5bf9ux5e94}{%
\section{F.5 同源执行与可测预测的新增对应}\label{f.5-ux540cux6e90ux6267ux884cux4e0eux53efux6d4bux9884ux6d4bux7684ux65b0ux589eux5bf9ux5e94}}

第 39 章的来源分工如下：

\begin{itemize}
\tightlist
\item
  程序执行与历史：\texttt{RegisterRepairSource}、\texttt{RegisterRepairHistory}。
\item
  故障和资源：\texttt{RegisterFaultSource}、\texttt{RegisterResourceSource}。
\item
  实际子代的维持：\texttt{RegisterBirthMaintenance}。
\item
  同源控制与记录：\nolinkurl{RegisterFaultControlSource}、\nolinkurl{RegisterFaultControlModel}、\texttt{RegisterFaultJointRecord}。
\item
  实际无限执行、价值和费用：\nolinkurl{RegisterFaultInfiniteRecord}。
\end{itemize}

\nolinkurl{UC.RegisterFaultPurposeInformation.observed_source_identifies_actual_purpose} 用源均值导出原 C6 忠实性。无故障的 \nolinkurl{UC.RegisterProgramOptimality} 与 UC.RegisterJointRecord 仍保留其不同源的原合同，不能用于替代有故障记录律。

C1 的源及传感器对应 \nolinkurl{UC.MeasurablePredictiveSource}、\nolinkurl{UC.MeasurablePredictiveSensor}；\nolinkurl{UC.MeasurablePredictiveJointPath.actual_timevarying_joint_path}、\nolinkurl{actual_finite_record_event}、\nolinkurl{actual_finite_record_lintegral} 保持同一时变可见 belief 策略的完整有限联合记录。新增 \nolinkurl{UC.MeasurablePredictiveLatentPath.actual_finite_prefix_consistency} 与 \nolinkurl{actual_source_timevarying_joint_path} 从实际保留潜状态的转移证明递归 filter 的条件一致性及同源可数记录律。具体核验、独立审查及源码身份见\href{../formalization/ueot-core/docs/compression/theory_completion/unified_closure/actual_infinite_sources_v1/RESULTS_AND_OPEN_OBLIGATIONS.md}{当前来源差额}、\href{../formalization/ueot-core/docs/compression/theory_completion/unified_closure/register_fault_resources_v1/RESULTS_AND_OPEN_OBLIGATIONS.md}{故障包}和\href{../formalization/ueot-core/docs/compression/theory_completion/unified_closure/c1_joint_path_v1/RESULTS_AND_OPEN_OBLIGATIONS.md}{C1 路径包}。两书编译不改变其数学范围，全理论状态仍为 OPEN / NOT\_ESTABLISHED。

第 39.7---39.11 节补入的实际端点为：\nolinkurl{UC.MeasurablePredictiveResponseRecursion.actual_C1_response_kernel_factor}；\nolinkurl{UC.MeasurablePredictiveCanonicalSource.actual_past_future_factorization}、\nolinkurl{actual_response_ae_minimal} 与 \nolinkurl{actual_canonical_recursion}；\nolinkurl{UC.RegisterTimestampedBirthSource} 的同一记录抹除、时钟与实际出生身份消费者；\nolinkurl{UC.BrownianProjectiveSource} 的实际全时间有限维 Gaussian 来源；\nolinkurl{UC.BrownianContinuousSource.process_isBrownian}、measurable\_joint\_process、actual\_ou\_recovery；以及 \nolinkurl{UC.MeasurablePredictiveHistoryControl.actual_action_observation_history_path}、\nolinkurl{actual_finite_action_observation_history_path} 与 \nolinkurl{actual_positive_memory_histories}。完整前提与来源身份分别见\href{../formalization/ueot-core/docs/compression/theory_completion/unified_closure/predictive_genealogy_source_v1/RESULTS_AND_OPEN_OBLIGATIONS.md}{一般响应与时戳包}、\href{../formalization/ueot-core/docs/compression/theory_completion/unified_closure/canonical_brownian_source_v1/RESULTS_AND_OPEN_OBLIGATIONS.md}{规范预测与全时间高斯包}和\href{../formalization/ueot-core/docs/compression/theory_completion/unified_closure/continuous_history_source_v1/RESULTS_AND_OPEN_OBLIGATIONS.md}{连续源与历史控制包}。原冻结结果索引及其登记处分不因新增来源而重计。第 39.12 节进一步给出声明协议类中任意真实历史过去与全部 fresh adaptive 未来的最小性；更广协议仍按第 4.3 节的条件处理。一般 Itô/Girsanov 与全局总装不由这些局部来源自动取得。

第 39.12 节对应 \nolinkurl{UC.MeasurablePredictiveHistoryCanonical} 的 \nolinkurl{actual_past_future_factorization}、\nolinkurl{actual_response_adaptive_ae_minimal} 与 \nolinkurl{futureKernel_response_memory_factor}。第 39.13 节对应 \nolinkurl{UC.BrownianMartingaleSource}；第 39.14---39.15 节的输入表示、原输出与完整 predictable 域对应 \nolinkurl{UC.BrownianIntegralRepresentation}、\nolinkurl{UC.BrownianIntegralClosure.integral_returns_literal}、\nolinkurl{UC.BrownianPredictableDensity.predictableIntegral_returns_literal}。第 39.16 节对应 \nolinkurl{UC.RegisterFaultFinitePurposeProbe} 及 \nolinkurl{UC.RegisterFaultFinitePurposeInfinite.returned_positive_actual_infinite_record}。完整类型与义务见\href{../formalization/ueot-core/docs/compression/theory_completion/unified_closure/adaptive_stochastic_sources_v1/RESULTS_AND_OPEN_OBLIGATIONS.md}{自适应／条件噪声包}、\href{../formalization/ueot-core/docs/compression/theory_completion/unified_closure/history_integral_source_v1/RESULTS_AND_OPEN_OBLIGATIONS.md}{真实过去／积分表示包}和\href{../formalization/ueot-core/docs/compression/theory_completion/unified_closure/predictable_finite_purpose_v1/RESULTS_AND_OPEN_OBLIGATIONS.md}{真正 predictable／有限目的证书包}。全理论仍 OPEN / NOT\_ESTABLISHED。

第 39.17---39.19 节对应 \nolinkurl{UC.BrownianIntegralContinuousVersion}、\nolinkurl{UC.RegisterFaultFeedbackResource}、\nolinkurl{UC.BrownianQuadraticVariation}、\nolinkurl{UC.BrownianIntegralPastLocalization} 与 \nolinkurl{UC.BrownianIntegralCompensator}；第 39.20---39.21 节对应 \nolinkurl{UC.RegisterPopulationHistorySource} 和 \nolinkurl{UC.RegisterPopulationTypePrice}。直接义务与精确类型分别见\href{../formalization/ueot-core/docs/compression/theory_completion/unified_closure/continuous_resource_source_v1/RESULTS_AND_OPEN_OBLIGATIONS.md}{连续积分／反馈资源}和\href{../formalization/ueot-core/docs/compression/theory_completion/unified_closure/compensator_population_source_v1/RESULTS_AND_OPEN_OBLIGATIONS.md}{真实二阶能量／人口来源}。补偿鞅不等于一般分割 QV／Itô；开放资源端口不是封闭源，子代普查不是整个末表。全理论仍 OPEN / NOT\_ESTABLISHED。

第 39.22---39.25 节采用 UC 中已集成的完整标量分割 QV、交叉／加权积分、欧氏乘积／向量闭包积分、资金训练及共同群体／日志学习来源。实际公共消费者、精确前提及总义务差额见\href{../formalization/ueot-core/docs/compression/theory_completion/unified_closure/full_qv_vector_training_source_v1/RESULTS_AND_OPEN_OBLIGATIONS.md}{完整 QV／训练／发现来源}和\href{../formalization/ueot-core/docs/compression/theory_completion/unified_closure/vector_weighted_population_source_v1/RESULTS_AND_OPEN_OBLIGATIONS.md}{向量表示／加权表／共同人口源}。后继 Taylor／Itô、向量可预测密度与带噪 C5 校准尚未作为本书已公开结果；全理论仍 OPEN / NOT\_ESTABLISHED。

\backmatter
\end{document}
